% concatenated from mainV6.tex
\documentclass[12pt]{amsart}
\usepackage{ amsmath, amsthm, amsfonts, amssymb, color}
\usepackage{mathrsfs}
\usepackage{amsfonts, amsmath}
\usepackage{enumitem}
\usepackage{amsmath,amstext,amsthm,amssymb,amsxtra}
\usepackage{txfonts} %also pxfonts
\usepackage[colorlinks, citecolor=blue,pagebackref,hypertexnames=false]{hyperref}
\allowdisplaybreaks
\usepackage{pgf,tikz}
\usetikzlibrary{calc}
\usepackage{caption} % for \captionof

\newcommand\Dir{\operatorname{Dir}}
\newcommand\Neu{\operatorname{Neu}}
\newcommand\LapoD{\Delta^{(0, \infty)}_{\Dir, d}}
\newcommand\LapoN{\Delta^{(0, \infty)}_{\Neu, d}}
\newcommand\LapiD{\Delta^{(1, \infty)}_{\Dir, d}}
\newcommand\LapiN{\Delta^{(1, \infty)}_{\Neu, d}}
\newcommand\TiD{T^{(1, \infty)}_{\Dir, d}}
\newcommand\TiN{T^{(1, \infty)}_{\Neu, d}}%\textit{}
\newcommand\TrD{T^{\R}_{ d}}
\newcommand\Td{\widetilde{T_d}}


\newcommand\mcS{\mathcal{S}}
\newcommand\pv{\mathcal{V}}
\newcommand\mcB{\mathcal{B}}
\newcommand\mcC{\mathcal{C}}
\newcommand\D{\mathcal{D}}
\newcommand\C{\mathbb{C}}
\newcommand\N{\mathbb{N}}
\newcommand\R{\mathbb{R}}
\newcommand\G{\mathbf{G}}
\newcommand\T{\mathbf{T}}
\newcommand\M{\mathbf{m}}
\newcommand\V{\mathbf{V}}
\newcommand\Z{\mathbf{Z}}
%\newcommand\RR{\widetilde{\R}}
\newcommand\dd{\,{\rm d}}
\newcommand{\cdo}{\, \cdot \,}


\newcommand\LaprD{\Delta^{\R}_{ d}}
\newcommand\Lapd{\widetilde{\Delta}_{d_1, d_2} }

\newcommand\Lapdd{\widetilde{\Delta}_{d} }


\newcommand\Lapda{\widetilde{\Delta_{d,a}}}
% \newcommand\Lapdd{\widetilde{\Delta_{d, -2(d-2)}}}
\newcommand\Lapdaperp{{\widetilde{\Delta_{d,a}}^\perp}}


%\usepackage{showkeys}
\textwidth =166mm
\textheight =236mm
\marginparsep=0cm
\oddsidemargin=-0.0cm
\evensidemargin=0.0cm
\headheight=13pt
\headsep=0.8cm
\parskip=0pt
\hfuzz=6pt
\widowpenalty=10000
\setlength{\topmargin}{-0.6cm}


\newcommand{\com}[1]{\vskip.3cm
	\fbox{%
		\color{red}
		\parbox{0.93\linewidth}{\footnotesize #1}}
	\vskip.3cm}



\widowpenalty=10000

\newtheorem{cl}{Claim}
\newtheorem{theorem}{Theorem}[section]
\newtheorem{proposition}[theorem]{Proposition}
\newtheorem{corollary}[theorem]{Corollary}
\newtheorem{lemma}[theorem]{Lemma}
\newtheorem{definition}[theorem]{Definition}
%\theoremstyle{definition}
\newtheorem{assum}{Assumption}[section]
\newtheorem{example}[theorem]{Example}
\newtheorem{remark}[theorem]{Remark}
%\renewcommand{\theequation}
%{\thesection.\arabic{equation}}






\begin{document}

\title[Riesz transform on Grushin spaces]{Riesz transform and its related inequalities for degenerate elliptic operators of Grushin type}

\author{Dangyang He}  


\begin{abstract}

We study the $L^p$ boundedness of the Riesz transform and the reverse Riesz inequality for degenerate elliptic operators of Grushin type. We prove full-range $L^p$ boundedness of the Riesz transform when the degenerate variable has dimension at least two, and obtain the sharp range in the one-dimensional weakly degenerate case, including the endpoint obstruction. In the strongly degenerate one-dimensional regime, we recover full-range boundedness, revealing a striking transition in the behavior of the singular set.

The proof develops a reverse Hölder theory for Grushin harmonic functions near the singular set. The main ingredients are explicit Poisson and Green kernel constructions adapted to the Friedrichs extension and a harmonic annihilation method which isolates the critical part of the Riesz kernel. These techniques illuminate the mechanism behind both the boundedness and unboundedness phenomena, and yield essentially sharp reverse Riesz inequalities for the same class of operators.

\end{abstract}

\subjclass[2020]{Primary 42B20; Secondary 35J70, 35H20.}

\keywords{Riesz transform, reverse Riesz inequality, reverse H\"older inequality, Grushin operators}



%\date{due by 31 March 2021}
\maketitle


 \tableofcontents






\section{Introduction}\label{section_1}

\noindent
Riesz transforms associated with non-negative self-adjoint operators are central objects in harmonic analysis and geometric analysis. On Euclidean space, the classical Riesz transforms are Calder\'on--Zygmund operators and are bounded on $L^p$ for every $1<p<\infty$; see, for example, \cite{R,CZ}. In dimension one, the Riesz transform reduces, up to normalization, to the Hilbert transform. On general metric measure or geometric spaces, however, the corresponding boundedness may fail outside a restricted range of exponents. The sharp range of $p$ is often sensitive to the heat-kernel geometry, the validity of Poincar\'e inequalities, and the structure of harmonic functions. A large body of work relates Riesz transform estimates to heat-kernel gradient bounds, reverse H\"older inequalities for gradients of harmonic functions, and Hodge projection estimates; see, among many others, \cite{ACDH,Shen,CJKS,Jiang,CD2}.

In this paper, we study these questions for a class of degenerate elliptic operators of Grushin type. Let $n,m\geq1$, $0<\alpha<1$, and $\beta\geq0$. We consider the Friedrichs realization on $L^2(\mathbb R^{n+m})$ of
\begin{equation}\label{Grushin}
   L=-\nabla_x\cdot\bigl(|x|^{2\alpha}\nabla_x\bigr)
     +|x|^{2\beta}\Delta_y,
\end{equation}
where $\Delta_y$ denotes the non-negative Laplacian on $\mathbb R^m$. The corresponding intrinsic gradient is
\begin{equation}\label{Grushin2}
    \nabla_L
    =
    \left(|x|^\alpha\nabla_x,\ |x|^\beta\nabla_y\right).
\end{equation}
The degeneracy occurs on the singular set $\{x=0\}$. The corresponding
intrinsic geometry is anisotropic. Indeed, the exponent in the dilation
\begin{equation*}
    \delta_r(x,y)
    =
    \bigl(rx,r^{\beta+1-\alpha}y\bigr)
\end{equation*}
is dictated by the homogeneity of the operator: for every smooth function
$f$,
\begin{equation*}
    L(f\circ\delta_r)
    =
    r^{2(1-\alpha)}(Lf)\circ\delta_r,
    \qquad
    \nabla_L(f\circ\delta_r)
    =
    r^{1-\alpha}(\nabla_Lf)\circ\delta_r.
\end{equation*}
Consequently, the intrinsic distance satisfies
\begin{equation*}
    d(\delta_r z,\delta_r z')
    =
    r^{1-\alpha}d(z,z').
\end{equation*}
% The degeneracy occurs on the singular set $\{x=0\}$. The intrinsic geometry is anisotropic: the $x$- and $y$-scales are coupled through the dilations
% \begin{equation*}
%     \delta_r(x,y)
%     =
%     \bigl(rx,r^{\beta+1-\alpha}y\bigr),
%     \qquad r>0,
% \end{equation*}
% under which
% \begin{equation*}
%     L(f\circ\delta_r)
%     =
%     r^{2(1-\alpha)}(Lf)\circ\delta_r.
% \end{equation*}
The associated distance, volume growth, subelliptic estimates, Poincar\'e inequalities, and heat-kernel upper bounds were studied in the foundational work of Robinson and Sikora \cite{RS,RS2,RS_1D}. Classical Grushin operators and their Riesz transforms, spectral multipliers, and Bochner--Riesz means have also been studied by Jotsaroop, Sanjay, and Thangavelu \cite{JST}, Martini and Sikora \cite{MS}, Chen and Sikora \cite{ChenSikora}, and others.

We say that the Riesz transform associated with $L$ is bounded on $L^p$ if
\begin{equation}\tag{$\mathrm{R}_p$}\label{R_p}
    \|\nabla_LL^{-1/2}f\|_p
    \leq
    C_p\|f\|_p,
    \qquad f\in C_c^\infty(\mathbb R^{n+m}).
\end{equation}
Equivalently, 
\begin{equation*}
    \|\nabla_L f\|_p
    \leq
    C_p\|L^{1/2}f\|_p.
\end{equation*}
Conversely, we say that the reverse Riesz inequality holds on $L^p$ if
\begin{equation}\tag{$\mathrm{RR}_p$}\label{RR_p}
    \|L^{1/2}f\|_p
    \leq
    C_p\|\nabla_L f\|_p,
    \qquad f\in C_c^\infty(\mathbb R^{n+m}).
\end{equation}

We shall use the following Poincar\'e inequalities. For $1\leq p<\infty$, we say that \eqref{P_p_intro} holds if there exist constants $C>0$ and $\lambda\geq1$ such that, for every ball $B$ and every locally Lipschitz function $f$,
\begin{equation}\tag{$\mathrm{P}_p$}\label{P_p_intro}
    \left(\fint_B|f-f_B|^p\right)^{1/p}
    \leq
    C r(B)
    \left(\fint_{\lambda B}|\nabla_Lf|^p\right)^{1/p}.
\end{equation}
Here $f_B:=\fint_Bf$, $r(B)$ denotes the radius of $B$, and $\lambda B$ denotes the concentric dilate of $B$. 

A key role in the proofs is played by reverse H\"older inequalities for gradients of harmonic functions. For $2<p<\infty$, we say that
\eqref{eq:RH:Jiang} holds if there exist constants $C>0$ and $\lambda\geq1$ such that, whenever $u$ is $L$-harmonic in $\lambda B$,
\begin{equation}\tag{$\mathrm{RH}_p$}\label{eq:RH:Jiang}
     \left(\fint_B|\nabla_Lu|^p\right)^{1/p}
     \leq
     \frac{C}{r(B)}
     \fint_{\lambda B}|u|.
\end{equation}
The endpoint condition $(\mathrm{RH}_\infty)$ is defined by replacing the left-hand side with $\operatorname*{sup}_B|\nabla_Lu|$. The use of an enlarged ball reflects the local nature of the Poincar\'e and Caccioppoli inequalities.

In the setting of doubling Dirichlet spaces supporting a scale-invariant $L^2$-Poincar\'e inequality, Coulhon, Jiang, Koskela, and Sikora \cite{CJKS} proved that, for $2<p<\infty$, the $L^p$ gradient estimate for the heat semigroup, \eqref{eq:RH:Jiang}, and \eqref{R_p} are equivalent. Their endpoint result also relates $(\mathrm{RH}_\infty)$ to pointwise gradient estimates for the heat kernel. Related versions under Poincar\'e inequalities outside a compact set were obtained by Jiang \cite{Jiang}. This abstract principle is one of the main engines of the present paper. In the Grushin setting, the principal difficulty is to establish the appropriate reverse H\"older estimates near the singular set $\{x=0\}$, where the Friedrichs realization selects different endpoint conditions in different parameter regimes.

Grushin-type operators have long served as model examples in the analysis of degenerate elliptic and subelliptic phenomena. The terminology originates in the two-dimensional degenerate operators
\begin{equation*}
    -\partial_x^2-x^{2k}\partial_y^2,
    \qquad k\in\mathbb N,
\end{equation*}
introduced by Grushin \cite{Grushin}. The heat-kernel bounds of Robinson and Sikora \cite{RS}, together with the abstract result of Sikora \cite{Sikora}, imply that, for the general operator \eqref{Grushin}, \eqref{R_p} holds for every $1<p\leq2$. Beyond the threshold $p=2$, however, the previously known full-range results concern special algebraic models.

In the case $\alpha=0$, $\beta=1$, and $m=1$, namely for the classical model
\begin{equation*}
    L=\Delta_x-|x|^2\partial_y^2,
\end{equation*}
Jotsaroop, Sanjay, and Thangavelu \cite{JST} proved $L^p$-boundedness for the full range $1<p<\infty$, using a closely related first-order system based on the creation and annihilation operators for the Hermite operator. Subsequently, Robinson and Sikora \cite{RS3} established full $L^p$ bounds for the integer-power models
\begin{equation*}
    L=\Delta_x+|x|^{2\beta}\Delta_y,
    \qquad
    n,m\geq1,\quad \alpha=0,\quad \beta\in\mathbb N,
\end{equation*}
by lifting the corresponding polynomial vector fields to a nilpotent Lie group. This method depends essentially on the integer-power structure and does not extend to the general power-degenerate Friedrichs model \eqref{Grushin}. In particular, once $\alpha>0$, the degeneracy in the $x$-derivatives introduces an endpoint mechanism at $\{x=0\}$ which is invisible to the nilpotent lifting argument. More generally, degenerate divergence-form operators may exhibit qualitative features absent from the uniformly elliptic theory; see \cite{ARSZ}.

The work of Robinson and Sikora naturally leaves open whether the full
$L^p$ range persists for arbitrary powers in the case $\alpha=0$. Our first main result establishes the corresponding full-range phenomenon for $n\geq2$ in the more singular regime $0<\alpha<1$.

\begin{theorem}\label{thm_main1}
Let $L$ be the Grushin operator defined by \eqref{Grushin}, with parameters
$\alpha\in(0,1)$, $\beta\geq0$, $m\geq1$, and $n\geq2$. Let $\nabla_L$ be the
corresponding gradient operator defined in \eqref{Grushin2}. Then
\eqref{R_p} holds for every $1<p<\infty$.
\end{theorem}

The proof of Theorem~\ref{thm_main1} proceeds through
$(\mathrm{RH}_\infty)$. When $n\geq2$, the Grushin space satisfies volume doubling, Gaussian heat-kernel upper bounds, and the required Poincar\'e inequality. The theorem of Coulhon, Jiang, Koskela, and Sikora therefore reduces the $L^p$-boundedness for $p>2$ to the reverse H\"older estimate \eqref{eq:RH:Jiang}. Away from the singular set, the desired estimate follows from standard elliptic arguments after rescaling. The main analytic difficulty is local and lies in proving the estimate on balls intersecting $\{x=0\}$.

In higher dimensions, one must decompose in spherical harmonics in the
$x$-variable. The order of the resulting modified Bessel functions depends on the angular frequency and tends to infinity with it. Consequently, the proof cannot rely only on the usual fixed-order asymptotic formulae for Bessel functions. Instead, one needs estimates that are uniform jointly in the Bessel order, the Fourier frequency in the $y$-variable, and the distance to the singular set. This uniform-in-order analysis is the principal technical point in the proof of the higher-dimensional theorem; see Part~\ref{part1}.

The one-dimensional case is qualitatively different. Suppose first that $n=1$ and $0<\alpha<1/2$. Then
\begin{equation*}
    \int_{-1}^1|x|^{-2\alpha}\,dx<\infty,
\end{equation*}
and functions in the local Friedrichs form domain admit a common trace across $\{x=0\}$. Moreover, weak solutions satisfy a conormal flux-matching condition; see \cite{RS_1D}. Consequently, after decomposing into even and odd parts, the full-line problem reduces to two different half-line problems: the even component satisfies a zero-flux condition at the origin, whereas the odd component satisfies a zero-trace condition. The leading odd branch behaves like
\begin{equation*}
    \operatorname{sgn}(x)|x|^{1-2\alpha},
\end{equation*}
and hence its intrinsic gradient has the singular behavior
$|x|^{-\alpha}$. This produces the exact obstruction at
$p=\alpha^{-1}$; see Part~\ref{part2}.

\begin{theorem}\label{thm:n=1:alpha<1/2}
Let $L$ be the Grushin operator defined by \eqref{Grushin}, with parameters
$\alpha\in(0,\frac12)$, $\beta\geq0$, $m\geq1$, and $n=1$. Let $\nabla_L$ be
the corresponding gradient operator defined in \eqref{Grushin2}. Then
\eqref{R_p} holds if and only if
\begin{equation*}
    1<p<\alpha^{-1}.
\end{equation*}
\end{theorem}

The range in Theorem~\ref{thm:n=1:alpha<1/2} is sharp. The same singular
branch that appears in the Friedrichs half-line problem yields both the
positive range and the counterexample at and beyond the endpoint. Thus the
obstruction is neither an artifact of the method nor a global geometric
phenomenon; it is caused by the local endpoint behavior of the odd component
at the singular set.

The remaining one-dimensional regime, $n=1$ and $1/2\leq\alpha<1$, requires
a different strategy. In this range,
\begin{equation*}
    \int_{-1}^1|x|^{-2\alpha}\,dx=\infty,
\end{equation*}
so the trace-matching argument breaks down. The set $\{x=0\}$ has zero
capacity with respect to the Friedrichs form, and the two sides
$\{x>0\}$ and $\{x<0\}$ are separated by the Friedrichs realization. In
particular, the global Poincar\'e inequality is no longer available in the
form needed to invoke directly the equivalence between \eqref{R_p} and
\eqref{eq:RH:Jiang}; see \cite{RS2}. We therefore postpone the proof of the
Riesz transform estimate in this strongly degenerate regime and first turn
to the reverse Riesz inequality.

The reverse Riesz inequality is governed by a different mechanism. Although
\eqref{R_p} implies $(\mathrm{RR}_{p'})$ by duality; see \cite{CD},
reverse Riesz estimates may hold in ranges where the Riesz transform itself
fails. This phenomenon occurs, for example, for exterior-domain Laplacians
and manifolds with ends; see \cite{KVZ,JL,H2}. For exponents below $2$, a
classical route to \eqref{RR_p} uses volume doubling together with \eqref{P_p_intro}; see, for instance, \cite{AC}. In the strongly
degenerate one-dimensional regime, however, the relevant Poincar\'e
inequality across the singular set fails \cite{RS2}. We therefore develop a different argument, introduced by the author in \cite{H2}.

The guiding principle is that the obstruction to the Riesz transform is
often carried by a leading harmonic term in the kernel. When the square-root
operator is written in bilinear form and paired with the gradient, this
leading harmonic contribution can be eliminated by an integration-by-parts
argument. We refer to this cancellation mechanism as harmonic annihilation.

The model case to keep in mind is the Riesz transform on manifolds with Euclidean ends. The low-energy analysis of Carron, Coulhon, and Hassell \cite{CCH} identifies a leading contribution to the kernel of
$\Delta^{-1/2}$ whose coefficient is harmonic. When the manifold has more than one end, this coefficient is nonconstant, and its gradient is responsible for the failure of the Riesz transform in the endpoint range. For the reverse Riesz problem, however, the harmonicity of this leading coefficient becomes a source of cancellation: after passing to the bilinear form, the leading profile is annihilated. This idea was developed in \cite{H2} and further refined in \cite{H3}.

In the present Grushin setting, the same principle can be implemented
without relying on a global scattering calculus. Following an idea of
Carron \cite{Gilles}, we decompose the relevant kernel into geometric pieces, isolate the critical region in which the leading harmonic profile appears, and subtract this profile before estimating the remainder. The remaining term is then controlled by a Hardy inequality adapted to the Grushin weights. This yields the following reverse Riesz theorem.

\begin{theorem}\label{thm_main2}
Let $L$ be the Grushin operator defined by \eqref{Grushin}, with parameters $\alpha\in(0,1)$, $\beta\geq0$, and $n,m\geq1$. Let $\nabla_L$ be the corresponding gradient operator defined in \eqref{Grushin2}. Then \eqref{RR_p} holds for every
\begin{equation}\label{p2}
   p\in
   \begin{cases}
        \left(1,\frac{1}{1-\alpha}\right)
        \cup
        \left(\frac{1}{1-\alpha},\infty\right),
        & n=1,\quad 0<\alpha<\frac12,\\[3mm]
        (1,\infty),
        & n\geq2
        \quad\text{or}\quad
        n=1,\quad \frac12\leq\alpha<1.
   \end{cases}
\end{equation}
\end{theorem}

The principal novelty in Theorem~\ref{thm_main2} is again one-dimensional.
When $n\geq2$, the full range follows from Theorem~\ref{thm_main1} and
duality. In the weakly degenerate one-dimensional case
$0<\alpha<1/2$, Theorem~\ref{thm:n=1:alpha<1/2} and duality give
\eqref{RR_p} for $p>(1-\alpha)^{-1}$, whereas harmonic annihilation yields the complementary range
$1<p<(1-\alpha)^{-1}$. The excluded exponent
$p=(1-\alpha)^{-1}$ is the natural Hardy threshold. Thus, after the leading
harmonic profile has been annihilated, the reverse Riesz inequality persists
well beyond the range in which the Riesz transform fails.

In the strongly degenerate case $1/2\leq\alpha<1$, the same method yields the
reverse inequality without a global Poincar\'e inequality across the
singular set. Moreover, the reverse Riesz estimate provides an essential
ingredient in the proof of the remaining Riesz transform theorem.

We now return to the strongly degenerate one-dimensional case. The threshold
$\alpha=1/2$ is consistent with the general one-dimensional extension theory
for degenerate divergence-form operators. For the model
\begin{equation*}
    H\varphi=-(c\varphi')',
    \qquad c(0)=0,
\end{equation*}
Robinson and Sikora \cite{RS_1D} characterized the self-adjoint and
submarkovian extensions in terms of the local behavior of $\int c^{-1}$ near
the degeneracy. In particular, when
$c(x)\simeq|x|^\delta$ with $\delta\geq1$, the Friedrichs semigroup separates
the two half-lines, and elements of the corresponding operator domain satisfy
the zero-flux condition
\begin{equation*}
    (c\varphi')(0\pm)=0.
\end{equation*}
For $c(x)=|x|^{2\alpha}$, this is precisely the regime
$\alpha\geq1/2$.

Thus the admissible half-line branch has the same zero-flux endpoint condition
as the even branch in the weakly degenerate case, although a global solution
need not be even. Using this zero-flux branch, we establish the endpoint
reverse H\"older estimate $(\mathrm{RH}_\infty)$ on each separated component.
The absence of a global Poincar\'e inequality is then bypassed by combining
the reverse Riesz inequality from Theorem~\ref{thm_main2} with a Hodge
projector argument. More precisely, the separated structure yields a
strengthened local reverse H\"older estimate of the type appearing in
Shen's criterion \cite{Shen}; together with the reverse Riesz inequality,
this allows us to apply a boundedness criterion for the Hodge projector and
recover the full range of the Riesz transform.

\begin{theorem}\label{thm:n=1:alpha>1/2}
Let $L$ be the Grushin operator defined by \eqref{Grushin}, with parameters
$\alpha\in[\frac12,1)$, $\beta\geq0$, $m\geq1$, and $n=1$. Let $\nabla_L$ be
the corresponding gradient operator defined in \eqref{Grushin2}. Then
\eqref{R_p} holds for every $1<p<\infty$.
\end{theorem}

Combining Theorems~\ref{thm_main1}, \ref{thm:n=1:alpha<1/2}, and
\ref{thm:n=1:alpha>1/2}, we obtain a complete description of the
$L^p$-boundedness of the Riesz transform associated with the Grushin operator
\eqref{Grushin}.

\medskip

\noindent\textbf{Outline.}
The paper is organized as follows. Section~\ref{sec2} introduces the
Friedrichs realization, recalls the Grushin geometry, and proves the
Poincar\'e inequalities needed later. Part~\ref{part1} proves
Theorem~\ref{thm_main1} through anchored reverse H\"older estimates and uniform-in-order Bessel analysis. Part~\ref{part2} proves
Theorem~\ref{thm:n=1:alpha<1/2} using the even--odd half-line decomposition and the endpoint counterexample. Part~\ref{part3} proves the reverse Riesz inequality, Theorem~\ref{thm_main2}, using a Grushin Hardy inequality and harmonic annihilation. Finally, Part~\ref{part4} treats the strongly degenerate one-dimensional regime and proves Theorem~\ref{thm:n=1:alpha>1/2} using separated half-line kernels and the Hodge projector.







% \noindent Riesz transforms associated with non-negative self-adjoint operators are central objects in harmonic analysis and geometric analysis. On Euclidean space, the classical Riesz transforms are Calder\'on--Zygmund operators and are bounded on $L^p$ for every $1<p<\infty$ (cf.~\cite{R,CZ}); in dimension one this includes the Hilbert transform. On general metric-measure or geometric spaces, however, the corresponding boundedness may fail outside a restricted range of exponents. The sharp range of $p$ is often sensitive to the heat-kernel geometry, the validity of Poincar\'e inequalities, and the structure of harmonic functions. A large body of work relates
% Riesz transform estimates to heat-kernel gradient bounds, reverse H\"older inequalities for gradients of harmonic functions, and Hodge projection estimates; see, among many others, \cite{ACDH,Shen,CJKS,Jiang,CD2}.

% In this paper we study these questions for a class of degenerate elliptic operators of Grushin type. Let $n,m\geq1$, $0<\alpha<1$, and $\beta\geq0$. We consider the Friedrichs realization on $L^2(\mathbb R^{n+m})$ of
% \begin{equation}\label{Grushin}
%    L = -\nabla_x\cdot |x|^{2\alpha}\nabla_x + |x|^{2\beta}\Delta_y,
% \end{equation}
% where $\Delta_y$ denotes the non-negative Laplacian on $\mathbb R^m$. The corresponding intrinsic gradient is
% \begin{equation}\label{Grushin2}
%     \nabla_L = \left(|x|^\alpha\nabla_x,\ |x|^\beta\nabla_y\right).
% \end{equation}
% The degeneracy occurs on the singular set $\{x=0\}$. The intrinsic geometry is anisotropic: the $x$-scale and the $y$-scale are coupled by the dilation
% \begin{equation*}
%     x\mapsto rx, \qquad y\mapsto r^{\beta+1-\alpha}y .
% \end{equation*}
% The associated distance, volume growth, Poincar\'e inequalities and heat-kernel upper bounds were studied in the foundational work of Robinson and Sikora \cite{RS,RS2,RS_1D}. Classical Grushin operators, spectral multipliers and Riesz transforms have also been studied by Jotsaroop, Sanjay and Thangavelu \cite{JST}, Martini--Sikora \cite{MS}, Chen--Sikora \cite{ChenSikora}, and others.

% We say that the Riesz transform associated with $L$ is bounded on $L^p$ if
% \begin{equation}\label{R_p}\tag{$\mathrm{R}_p$}
%     \|\nabla_LL^{-1/2}f\|_p
%     \leq
%     C_p\|f\|_p ,
%     \qquad f\in C_c^\infty(\mathbb R^{n+m}).
% \end{equation}
% Equivalently,
% \begin{equation*}
%     \|\nabla_L f\|_p
%     \leq
%     C_p\|L^{1/2}f\|_p .
% \end{equation*}
% Conversely, the reverse Riesz inequality is said to hold on $L^p$ if
% \begin{equation}\label{RR_p}\tag{$\mathrm{RR}_p$}
%     \|L^{1/2}f\|_p
%     \leq
%     C_p\|\nabla_L f\|_p ,
%     \qquad f\in C_c^\infty(\mathbb R^{n+m}).
% \end{equation}


% We shall use the following Poincar\'e inequalities in our discussion. For
% $1\le p<\infty$, we say that \eqref{P_p_intro} holds if there exist constants $C>0$ and
% $\lambda\ge 1$ such that, for every ball $B$ and every locally Lipschitz function $f$,
% \begin{equation}\label{P_p_intro}\tag{$\mathrm{P}_p$}
% \int_B |f-f_B|^p
% \le
% C r(B)^p
% \int_{\lambda B}|\nabla_L f|^p .
% \end{equation}
% Here $f_B:=\fint_B f$ and $r(B)$ denotes the radius of $B$.

% A key role in the proof is played by reverse H\"older inequalities for gradients of
% harmonic functions. For $p>2$, we say that \eqref{eq:RH:Jiang} holds if there exist constants $C>0$ and $\lambda \ge 1$
% such that, whenever $u$ is $L$-harmonic in $\lambda B$, 
% \begin{equation}\label{eq:RH:Jiang}\tag{$\mathrm{RH}_p$}
%      \left( \fint_B |\nabla_L u|^p  \right)^{\frac{1}{p}} \le \frac{C}{r(B)}  \fint_{\lambda B} |u| 
% \end{equation}
% The use of $\lambda B$ rather than $B$ reflects the local nature of the Poincar\'e inequality
% and the Caccioppoli estimate.


% In the presence of the doubling property, Gaussian heat-kernel upper bounds and suitable Poincar\'e inequalities, the work of Coulhon--Jiang--Koskela--Sikora \cite{CJKS} and Jiang \cite{Jiang} invetigates the equivalence among the gradient estimates of heat kernel, \eqref{eq:RH:Jiang} and \eqref{R_p}. This abstract principle is one of the main engines of the paper. In the setting of Grushin spaces, the main difficulty is however to prove the appropriate reverse H\"older estimates near the singular set $\{x=0\}$, where the Friedrichs realization imposes different endpoint conditions in different parameter regimes.

% Grushin-type operators have long served as model examples in the analysis of
% degenerate elliptic and subelliptic phenomena. The terminology comes from the
% two-dimensional degenerate operators
% \begin{equation*}
%     -\partial_x^2-x^{2k}\partial_y^2,
%     \qquad k\in\mathbb N,
% \end{equation*}
% introduced by Grushin \cite{Grushin}. It is known from the work of Sikora \cite{Sikora} and Robinson--Sikora \cite{RS} that, for the general operator \eqref{Grushin}, the estimate \eqref{R_p} holds for all $1<p\leq2$. However, beyond the threshold $p=2$, only special cases were understood. In the case $\alpha=0$, $\beta=1$, $n\ge 1$ and $m=1$, namely for the classical model
% \begin{equation*}
%     L=\Delta_x - |x|^2\partial_y^2,
% \end{equation*}
% Jotsaroop--Sanjay--Thangavelu \cite{JST} prove $L^p$ boundedness of the Riesz transform for the full range $1<p<\infty$, with a slightly different convention for the associated gradient. Subsequently, Robinson--Sikora \cite{RS3} establish full $L^p$ bounds for the special integer-power case
% \begin{equation*}
%     L=\Delta_x+|x|^{2\beta}\Delta_y,
%     \qquad
%     n,m\geq1,\quad \alpha=0,\quad \beta\in\mathbb N,
% \end{equation*}
% by realizing the corresponding vector fields through nilpotent Lie groups. This Lie-group method does not apply to the general power-degenerate Friedrichs model \eqref{Grushin}. Particularly, once $\alpha>0$, the degeneracy in the $x$-derivatives creates an endpoint mechanism at $\{x=0\}$ which is invisible in the nilpotent lifting approach. More generally, degenerate divergence-form operators may display qualitative features absent in the uniformly elliptic theory; see \cite{ARSZ}. 

% Robinson and Sikora conjecture that full $L^p$-boundedness should persist at least in the case $n,m\geq1$, $\alpha=0$ and $\beta\geq0$. Our first main result confirms the corresponding full-range phenomenon for $n\geq2$ in the more singular setting $0<\alpha<1$.

% \begin{theorem}\label{thm_main1}
% Let $L$ be the Grushin operator defined by \eqref{Grushin}, with parameters
% $\alpha\in(0,1)$, $\beta\geq0$, $m\geq1$ and $n\geq2$. Let $\nabla_L$ be the
% corresponding gradient operator defined in \eqref{Grushin2}. Then \eqref{R_p}
% holds for every $1<p<\infty$.
% \end{theorem}

% The proof of Theorem~\ref{thm_main1} proceeds through $(\mathrm{RH}_\infty)$. Since the Grushin spaces satisfy volume doubling and Gaussian heat-kernel upper bounds, and since the relevant Poincar\'e inequalities hold when $n\geq2$, Coulhon-Jiang-Koskela-Sikora theorem reduces the Riesz transform estimate to the validity of \eqref{eq:RH:Jiang}. The main analytic difficulty is local and occurs near the singular set, showing the reverse H\"older inequalities for balls across the hyperplane $\{x=0\}$. Also, in higher dimension one must decompose in spherical harmonics in the $x$-variable and the corresponding modified Bessel order depends on the angular frequency and tends to infinity with it. Thus the estimates cannot rely only on the usual fixed-order asymptotic formulae for Bessel functions; it requires manipulation which is uniform simultaneously in the Bessel order, the Fourier frequency in the $y$-variable, and the distance to the singular set. This uniform-in-order analysis is the principal technical point in the proof of the higher-dimensional theorem; see Part~\ref{part1}.

% The one-dimensional case is qualitatively different. Suppose first that $n=1$ and $0<\alpha<1/2$. Then $\int_{-1}^1 |x|^{-2\alpha}\,dx<\infty$, and functions in the local Friedrichs form domain have a common trace across $\{x=0\}$. Moreover, weak solutions satisfy a conormal flux matching condition (cf.~\cite{RS_1D}). Consequently, after decomposing into even and odd parts, the full-line problem becomes two different half-line problems: the even component satisfies a zero-flux condition at the origin, while the odd part satisfies a zero-trace condition.
% Surprisingly, the intrinsic gradient of the odd branch has a leading singularity $|x|^{-\alpha}$. This produces the exact obstruction at $p=\alpha^{-1}$; see Part~\ref{part2}.

% \begin{theorem}\label{thm:n=1:alpha<1/2}
% Let $L$ be the Grushin operator defined by \eqref{Grushin}, with parameters
% $\alpha\in(0,\frac12)$, $\beta\geq0$, $m\geq1$ and $n=1$. Let $\nabla_L$ be the
% corresponding gradient operator defined in \eqref{Grushin2}. Then \eqref{R_p}
% holds if and only if $1<p<\alpha^{-1}$.
% \end{theorem}

% This result is sharp in the strongest possible sense: the same singular branch which appears in the Friedrichs half-line problem gives both the positive range and the counterexample beyond it. Thus the obstruction is not an artifact of the method or of a global estimate; it is the local endpoint behavior of the odd component at
% the singular set.

% The remaining one-dimensional regime, $n=1$ and $1/2\leq\alpha<1$, requires a
% different strategy. In this range the integral $\int_{-1}^1|x|^{-2\alpha}\,dx$
% diverges, so the trace-matching argument breaks down. The singular point has zero Friedrichs capacity, and the two half-lines are separated by the Friedrichs realization. In particular, the Poincar\'e inequality is no longer available in the form needed to invoke directly the equivalence between \eqref{R_p} and \eqref{eq:RH:Jiang}; see \cite{RS2}. We therefore postpone the proof of the Riesz transform in this strongly degenerate regime and first turn to the reverse Riesz inequality.

% The reverse Riesz inequality is governed by a different mechanism. Although \eqref{R_p} implies $(\mathrm{RR}_{p'})$ by duality; see \cite{CD}, reverse Riesz estimates may hold in ranges where the Riesz transform itself fails. This phenomenon appears, for example, for exterior-domain Laplacians and for manifolds with ends; see \cite{KVZ,JL,H2}. A classical way to obtain \eqref{RR_p} is through
% volume doubling and an $L^p$-Poincar\'e inequality \eqref{P_p_intro}; see, for instance, \cite{AC}. However, in the one-dimensional strongly degenerate regime $n=1$ and $1/2\leq\alpha<1$, the relevant Poincar\'e inequality across the singular set fails; see \cite{RS2}. We therefore follow a different route, inspired by the method introduced in \cite{H2}.

% The guiding principle is that the obstruction to the Riesz transform is often
% carried by a leading harmonic term in the kernel. When the square-root operator is
% lifted to a bilinear form and paired with the gradient, this leading harmonic
% contribution can be eliminated by an integration-by-parts argument. We refer to
% this cancellation mechanism as harmonic annihilation.

% The model case to keep in mind is the Riesz transform on manifolds with Euclidean
% ends. Carron--Coulhon--Hassell \cite{CCH} showed that the leading term in the expansion of
% $\Delta^{-1/2}$ at infinity has a harmonic coefficient; when the manifold has more
% than one end this coefficient is nonconstant, and its gradient is responsible for
% the failure of the Riesz transform in the endpoint range. In the reverse Riesz
% problem, however, the harmonicity of this leading coefficient becomes harmless: after
% passing to the bilinear form, the harmonic leading profile is annihilated. This
% idea was developed in \cite{H2} and further refined in \cite{H3}.

% In the present Grushin setting, the same philosophy is implemented without relying on a global scattering calculus. Inspored by Carron \cite{Gilles}, we decompose the relevant kernel into geometric pieces, isolate the critical region where the leading harmonic profile appears, and subtract this profile before estimating the remainder. The remaining term is then controlled by a Hardy inequality adapted to the Grushin weight. This gives the
% following reverse Riesz theorem.

% \begin{theorem}\label{thm_main2}
% Let $L$ be the Grushin operator defined by \eqref{Grushin}, with parameters
% $\alpha\in(0,1)$, $\beta\geq0$ and $n,m\geq1$. Let $\nabla_L$ be the corresponding
% gradient operator defined in \eqref{Grushin2}. Then \eqref{RR_p} holds for all
% \begin{equation}\label{p2}
%    p\in
%    \begin{cases}
%         \left(1,\frac{1}{1-\alpha}\right)
%         \cup
%         \left(\frac{1}{1-\alpha},\infty\right),
%         & n=1,\quad 0<\alpha<\frac12,\\[2mm]
%         (1,\infty),
%         & n\geq2
%         \quad\text{or}\quad
%         n=1,\quad \frac12\leq\alpha<1.
%     \end{cases}
% \end{equation}
% \end{theorem}

% The main novelty is again one-dimensional. When
% $n\geq2$, the full range of \eqref{RR_p} follows from Theorem~\ref{thm_main1} and duality. In the weakly degenerate one-dimensional case $0<\alpha<1/2$, the missing exponent
% $p=(1-\alpha)^{-1}$ is the natural Hardy threshold. The point is that:
% \begin{itemize}
%     \item In the weakly degenerate case $0<\alpha<1/2$, the reverse inequality is not simply controlled by the same obstruction as the Riesz transform: after the leading harmonic profile is annihilated, the reverse Riesz inequality persists beyond the range in which the Riesz transform fails.
%     \item In the strongly degenerate case $1/2\leq\alpha<1$, one obtains reverse inequality without the framework of Poincar\'e inequality. Moreover, the reverse inequality also provides the main engine for the final Riesz transform result.
% \end{itemize}
  


% To complete the picture of the Riesz transform, we now return to the strongly degenerate one-dimensional case. The threshold $\alpha=1/2$ is consistent with the general one-dimensional extension theory for degenerate divergence-form operators. In the model $H=-(c\varphi')'$, with $c(0)=0$, Robinson and Sikora \cite{RS_1D} characterized the self-adjoint and submarkovian extensions in terms of the local behavior of $\int c^{-1}$ near the degeneracy. In particular, when $c(x)\simeq |x|^\delta$ with $\delta\geq1$, the Friedrichs semigroup separates the two half-lines, and the endpoint condition is the zero-flux condition $(c\varphi')(0\pm)=0$. For $c(x)=|x|^{2\alpha}$, this is precisely the regime $\alpha\geq1/2$. In this case, the finite-energy Friedrichs condition selects the zero-flux branch on each half-line. Structurally speaking, this is the same half-line kernel that appears as the even zero-flux kernel in the weakly degenerate case, although the global solution need not be even. With such zero-flux branch, one can again build the $L^p$ reverse H\"older theory for all $1<p\le \infty$. The lack of global Poincar\'e inequalities is bypassed by combining the reverse Riesz inequality from Theorem~\ref{thm_main2} with a Hodge projector argument. More precisely, the separated structure allows one to improve \eqref{eq:RH:Jiang} to a form used in Shen \cite{Shen}, and hence a criterion of the boundedness of the Hodge projector can be applied. This allows us to recover the full range of the Riesz transform.





% \begin{theorem}\label{thm:n=1:alpha>1/2}
% Let $L$ be the Grushin operator defined by \eqref{Grushin}, with parameters
% $\alpha\in[\frac12,1)$, $\beta\geq0$, $m\geq1$ and $n=1$. Let $\nabla_L$ be the
% corresponding gradient operator defined in \eqref{Grushin2}. Then \eqref{R_p}
% holds for every $1<p<\infty$.
% \end{theorem}

% Combining Theorems~\ref{thm_main1}, \ref{thm:n=1:alpha<1/2} and
% \ref{thm:n=1:alpha>1/2}, we obtain a complete description of the $L^p$-boundedness of the Riesz transform for the Grushin operator \eqref{Grushin}. 


% \medskip

% \noindent \textbf{Outline.} The paper is organized as follows. Section~\ref{sec2} fixes the Friedrichs realization, recalls the Grushin geometry, and proves the Poincar\'e inequalities needed later. Part~\ref{part1} proves Theorem~\ref{thm_main1} through anchored reverse H\"older estimates and uniform-in-order Bessel analysis. Part~\ref{part2} proves Theorem~\ref{thm:n=1:alpha<1/2} by the even--odd half-line decomposition and the endpoint counterexample. Part~\ref{part3} proves the reverse Riesz inequality, Theorem~\ref{thm_main2}, using a Grushin Hardy inequality and harmonic annihilation. Finally, Part~\ref{part4} treats the strongly degenerate one-dimensional regime and proves Theorem~\ref{thm:n=1:alpha>1/2} using separated half-line kernels and the Hodge projector.









\section{Preliminaries}\label{sec2}





\subsection{Friedrichs realization}\label{ssec:Friedrichs}

The singular set $\{x=0\}$ is not merely a technical nuisance; it determines which endpoint branches are admissible in the kernel construction (see Part~\ref{part1}~\ref{part2},~\ref{part4} below). We therefore begin by fixing the Friedrichs realization of $L$ and by recording the interface conditions it imposes. For more detailed discussion in this direction, we refer reader to \cite{RS_1D,RS_flux} and references therein.

 

Let $n,m\ge 1$, $\alpha\in (0,1)$ and $\beta \ge 0$. On
$C_c^\infty(\mathbb R^{n+m})$, define
\begin{equation*}
    \mathfrak q_0(u,v)
    =
    \int_{\mathbb R^{n+m}}
    \left(
        |x|^{2\alpha}\nabla_xu\cdot \overline{\nabla_xv}
        +
        |x|^{2\beta}\nabla_yu\cdot \overline{\nabla_yv}
    \right)dxdy.
\end{equation*}
This non-negative symmetric form is closable in $L^2(\mathbb R^{n+m})$ (cf. \cite[Lemma~2.1]{RS}). We denote its
closure by $(\mathfrak q,\mathcal F)$, where
\begin{equation*}
    \mathcal F
    =
    \overline{C_c^\infty(\mathbb R^{n+m})}^{\|\cdot\|_{\mathcal F}},
    \qquad
    \|u\|_{\mathcal F}^2
    =
    \|u\|_2^2+\mathfrak q(u,u).
\end{equation*}
The Friedrichs realization of $L$ is the non-negative self-adjoint operator associated
with $(\mathfrak q,\mathcal F)$. Thus $u\in\mathcal D(L)$ and $Lu=f$ if and only if
$u\in\mathcal F$, $f\in L^2$, and
\begin{equation*}
    \mathfrak q (u,\varphi)=\langle f,\varphi\rangle_{L^2}
    \qquad
    \text{for all }\varphi\in\mathcal F.
\end{equation*}
We also use the corresponding local notion. If $\Omega\subset\mathbb R^{n+m}$ is open, then $u\in\mathcal F_{\mathrm{loc}}(\Omega)$ means that for every $\widetilde{\Omega} \subset\Omega$ there exists $v\in\mathcal F$ such that $u=v$ a.e. on $\widetilde{\Omega}$. We say that $u$ has locally finite energy if $u\in\mathcal F_{\mathrm{loc}}$.
Equivalently, locally one has
\begin{equation*}
    \int
    \left(
        |x|^{2\alpha}|\nabla_xu|^2
        +
        |x|^{2\beta}|\nabla_yu|^2
    \right)dxdy
    <\infty.
\end{equation*}
A function $u\in\mathcal F_{\mathrm{loc}}(\Omega)$ is called $L$-harmonic in $\Omega$ if
\begin{equation}
    \mathfrak q(u,\varphi)=0
    \qquad
    \text{for all }\varphi\in\mathcal F
    \text{ compactly supported in }\Omega.
\end{equation}

We now record the only subtle point concerning the singular set when $n=1$. Near $x=0$, the relevant part of the energy is $\int |x|^{2\alpha}|\partial_xu|^2dx$. The difference between $0<\alpha<1/2$ and $1/2 \le \alpha<1$ is essential. In the former regime, the form domain has a common trace across $x=0$, while in the latter regime the singular set has zero Friedrichs capacity and separates the two half-lines.

Assume first that $0<\alpha<1/2$. Then $\int_{-1}^{1}|x|^{-2\alpha}dx<\infty$. Hence, for smooth $u$ and $x_1 < 0 < x_2$, Cauchy-Schwarz inequality gives
\begin{align*}
|u(x_2,y)-u(x_1,y)|^2\le \left(\int_{x_1}^{x_2}|t|^{-2\alpha}dt\right) \left(\int_{x_1}^{x_2}|t|^{2\alpha}|\partial_tu(t,y)|^2 dt\right).
\end{align*}
By density, the same estimate holds for functions in the local Friedrichs form
domain. Hence, for almost every $y$, the one-sided traces $u(0^+,y)$ and
$u(0^-,y)$ exist. Moreover, since both factors on the right-hand side above tend to zero as $x_1\to0^-$ and
$x_2\to0^+$, one obtains
\begin{equation}\label{eq_subcritical_trace_matching}
    u(0^+,y)=u(0^-,y)
    \qquad
    \text{for a.e. }y.
\end{equation}
Thus the common trace condition comes from the Friedrichs form domain.

The flux condition comes from the weak equation. Suppose that $u$ is piecewise smooth away from $x=0$, belongs locally to the Friedrichs form domain, and solves
\begin{equation}\label{equation}
    Lu=f
\end{equation}
weakly across $x=0$, with no singular source supported on $\{x=0\}$. Suppose that the one-sided conormal fluxes
\begin{equation*}
    F_+(y):=\lim_{x\to0^+}x^{2\alpha}\partial_xu(x,y),
    \qquad
    F_-(y):=\lim_{x\to0^-}|x|^{2\alpha}\partial_xu(x,y)
\end{equation*}
exist. Integrating by parts separately on $(-\varepsilon,0)$ and $(0,\varepsilon)$
gives the interface contribution
\begin{equation*}
\int_{\mathbb R^m}
    \left[F_-(y)\varphi(0^-,y)-F_+(y)\varphi(0^+,y)\right] dy.    
\end{equation*}
By \eqref{eq_subcritical_trace_matching}, test functions in the Friedrichs space
have a common trace at $x=0$. Therefore the interface contribution becomes
\begin{equation*}
    \int_{\mathbb R^m}
    \left[F_-(y)-F_+(y)\right]\varphi(0,y) dy.
\end{equation*}
Since the weak equation contains only the volume term $\int f\varphi$ and has no
delta contribution on $\{x=0\}$, this interface term must vanish for all test
functions $\varphi$. Hence
\begin{equation}\label{eq_subcritical_flux_matching}
    F_-(y)=F_+(y)
    \qquad\text{for a.e. }y.
\end{equation}


For later use, now consider even and odd weak solutions of \eqref{equation}. If $u_e(x,y)=v_e(|x|,y)$, then
\begin{equation*}
F_+(y)=\lim_{\rho\to0^+}\rho^{2\alpha}\partial_\rho v_e(\rho,y),
    \qquad
    F_-(y)=-\lim_{\rho\to0^+}\rho^{2\alpha}\partial_\rho v_e(\rho,y). 
\end{equation*}
The matching condition \eqref{eq_subcritical_flux_matching} therefore gives $\lim_{\rho\to0+}\rho^{2\alpha}\partial_\rho v_e(\rho,y)=0$. Thus the even solution satisfies the zero-flux condition. If $u_o(x,y)=\mathrm{sgn}(x)v_o(|x|,y)$, then the trace matching condition \eqref{eq_subcritical_trace_matching} gives $v_o(0+,y)=-v_o(0+,y)$ and hence $v_o(0+,y)=0$. Thus the odd solution satisfies the zero-trace condition.


When $1/2\leq\alpha<1$, the situation is different. The integral $\int_{-1}^{1}|x|^{-2\alpha}dx$ diverges, so the above trace argument no longer applies. In fact, the singular point has zero
Friedrichs capacity. This can be seen from the approximation of the jump function $\mathrm{sgn}(x)$.
For $\alpha>1/2$, take $\eta_\varepsilon\in C^\infty([0,\infty))$ such that
\begin{equation*}
    \eta_\varepsilon(r) = \begin{cases}
        0, & r\le \varepsilon,\\
        1, & r\ge 2\varepsilon,
    \end{cases}
\end{equation*}
and $|\eta_\varepsilon'(r)|\leq C\varepsilon^{-1}$. Then
\begin{equation*}
    \int_{-1}^{1}
    |x|^{2\alpha}
    \left|
        \partial_x\bigl(\mathrm{sgn}(x)\eta_\varepsilon(|x|)\bigr)
    \right|^2 dx
    \le
    C\varepsilon^{2\alpha-1}
    \to 0.
\end{equation*}
At $\alpha=1/2$, using the logarithmic cut-off
\begin{equation*}
\eta_\varepsilon(r)= \begin{cases}
        0, & 0\le r\leq\varepsilon^2,\\
        \frac{\log(r/\varepsilon^2)}{\log(1/\varepsilon)},
        & \varepsilon^2<r<\varepsilon,\\
        1, & r\ge \varepsilon,
    \end{cases}
\end{equation*}
one obtains
\begin{equation*}
\int_{-1}^{1}|x|\left|\partial_x\bigl(\mathrm{sgn}(x)\eta_\varepsilon(|x|)\bigr)\right|^2 dx \le \frac{C}{\log(1/\varepsilon)} \to 0.
\end{equation*}
Thus $\mathrm{sgn}(x)$ belongs to the local Friedrichs form domain when
$1/2\leq\alpha<1$, and
\begin{equation}\label{eq:gradient:sgn}
    \nabla_L \mathrm{sgn}(x)=0
    \qquad
    \text{in the form sense}.
\end{equation}
Consequently, $\mathfrak q(\mathrm{sgn}(x),\varphi)=0$ for every compactly supported test function $\varphi$, so $\mathrm{sgn}(x)$ is
locally $L$-harmonic. Therefore, locally near $x=0$, the two half-lines are separated and there is no trace matching condition across $x=0$. 

We now explain why the separated half-line Friedrichs problem carries the zero-flux endpoint condition. Consider a sufficiently regular solution $u$ on the half-line $(0,2)$ whose conormal flux $w(\rho):=\rho^{2\alpha}\partial_\rho u(\rho)$ has a finite limit as $\rho\to0^+$. If $\lim_{\rho\to0^+}w(\rho)=C\ne 0$, then, for all sufficiently small $\rho$, $|\partial_\rho u(\rho)|= \rho^{-2\alpha}|w(\rho)| \gtrsim \rho^{-2\alpha}$. Consequently,
\begin{equation*}
\int_0^1 \rho^{2\alpha}|\partial_\rho u(\rho)|^2 d\rho
\gtrsim
\int_0^1 \rho^{-2\alpha} d\rho
=\infty,    
\end{equation*}
because $\alpha\ge 1/2$. This contradicts the Friedrichs finite-energy condition. Therefore every such Friedrichs solution must satisfy
\begin{equation}\label{eq_supercritical_zero_flux}
    \lim_{\rho\to0^+}\rho^{2\alpha}\partial_\rho u(\rho)=0.
\end{equation}
Equivalently, the local homogeneous normal equation $-\partial_\rho\left(\rho^{2\alpha}\partial_\rho u\right)=0$ has two endpoint branches: $u(\rho)\sim 1$ and
\begin{equation*}
u(\rho)\sim
\begin{cases}
\rho^{1-2\alpha}, & \alpha>1/2,\\
\log\rho, & \alpha=1/2.
\end{cases}    
\end{equation*}
The second branch has nonzero conormal flux and infinite Friedrichs energy, whereas the first branch has zero conormal flux. Thus the Friedrichs extension selects the first branch on each separated half-line. This separation phenomenon is a manifestation of the general capacity-flux criterion for degenerate divergence-form operators: a hypersurface may become an impenetrable barrier for the Friedrichs semigroup when its normal energy flux, vanishes (see \cite{RS_flux}).

For the explicit Fourier-Bessel modes used in Part~\ref{part2} and ~\ref{part4} below, the conormal flux exists and finite energy excludes the nonzero-flux branch.









 

\subsection{Grushin geometry}

We next recall the intrinsic geometry associated with $L$. The relevant distance is obtained by declaring the vector fields $|x|^{\alpha} \nabla_x$ and $|x|^\beta \nabla_y$ to have unit size (cf.~\cite{JS}). The resulting balls have two different forms: remote balls, which stay away from the singular set and are comparable to anisotropic Euclidean boxes with essentially constant coefficients, and anchored balls, which meet the singular set and reflect the true Grushin scaling. This dichotomy is used throughout the paper.

Recall the homogeneous dimension of Grushin space:
\begin{equation*}
    \mathcal Q:= \frac{n+m(\beta+1-\alpha)}{1-\alpha}.
\end{equation*}
We record the following result from \cite[Proposition~5.1, Theorem~6.4, Corollary~6.6]{RS}.

\begin{theorem}\cite{RS}\label{RS_Grushin}
Let $n,m\ge 1$, $\alpha \in (0,1)$ and $\beta \ge 0$. The following estimates hold:
\begin{enumerate}[label=(\roman*)]
\item the distance estimate
\begin{align*}
d(\xi,\eta) = d((x,y);(x',y')) \sim \frac{|x-x'|}{\left( |x| + |x'| \right)^\alpha} + \frac{|y-y'|}{(|x|+|x'|)^{\beta} + |y-y'|^{\frac{\beta}{\beta+1-\alpha}}} =: D(\xi,\eta),
\end{align*}
\item the volume estimate
\begin{align*}
V(\xi,r) = V((x,y),r) \sim \begin{cases}
r^{\mathcal{Q}}, & r\ge |x|^{1-\alpha},\\
r^{n+m}|x|^{n \alpha + m\beta}, & r\le |x|^{1-\alpha},
\end{cases}
\end{align*}
\item the volume doubling property
\begin{align}\label{D}\tag{$\mathrm{D}$}
V(\xi,sr) \le C s^{\mathcal{Q}} V(\xi,r)
\end{align}
for all $\xi \in \mathbb{R}^{n+m}$ and all $s\ge 1$, $r>0$,
\item the heat kernel estimate
\begin{equation}\label{UE}\tag{$\mathrm{UE}$}
e^{-tL}(\xi,\eta) \le \frac{C}{V(\xi, \sqrt{t})} \textrm{exp} \left( - \frac{d(\xi,\eta)^2}{ct} \right).
\end{equation}
\end{enumerate}
\end{theorem}


With Theorem~\ref{RS_Grushin}, we establish the following estimates.


\begin{lemma}\label{le_volume}
Let $n,m\ge 1$, $\alpha\in (0,1)$ and $\beta\ge 0$. Let $\xi=(x,y)\in \mathbb{R}^{n+m}$. There exist $0<c_1<c_2$ and $\varepsilon_0 \in (0,1)$ such that for all $r\le \varepsilon_0 |x|^{1-\alpha}$,
\begin{equation}\label{le_volume2}
B_n(x, c_1 r |x|^\alpha) \times B_m(y, c_1 r |x|^\beta) \subset B(\xi, r) \subset B_n(x, c_2 r |x|^\alpha) \times B_m(y, c_2 r |x|^\beta).
\end{equation}
Moreover, there exist $0<c_3<c_4$ such that for all $r>0$ and all $y\in \mathbb{R}^m$,
\begin{equation}\label{le_volume3}
B_n(0, c_3 r) \times B_m(y, c_3 r^{\beta+1-\alpha}) \subset B((0,y), r^{1-\alpha}) \subset B_n(0, c_4 r) \times B_m(y, c_4 r^{\beta+1-\alpha}).
\end{equation}
\end{lemma}

\begin{proof}
Fix a constant $C_0>1$ such that $C_0^{-1} D(\xi,\eta) \le d(\xi,\eta) \le C_0 D(\xi,\eta)$. First, we establish that
\begin{equation}\label{le_volume1}
B(\xi, r) \subset B_n(x, c_2 r |x|^\alpha) \times B_m(y, c_2 r |x|^\beta).
\end{equation}
Let $\eta = (x',y')\in B(\xi, r)$. Then $D(\xi,\eta)\le C_0 r$. The following claim holds:
\begin{equation}\label{claim1}
|x-x'| \le \frac{|x|}{2},
\end{equation}
provided $\varepsilon_0$ is sufficiently small. Indeed, if this were false, then $|x|+|x'|\le |x-x'| + 2|x| \le 3|x-x'|$, and hence
\begin{equation*}
\frac{|x-x'|}{\left(|x|+|x'|\right)^\alpha} \ge 3^{-\alpha} |x-x'|^{1-\alpha} \ge 3^{-\alpha} 2^{\alpha-1} |x|^{1-\alpha}.
\end{equation*}
Consequently, $D(\xi,\eta)\ge 3^{-\alpha} 2^{\alpha-1} |x|^{1-\alpha}$. Choose $\varepsilon_0 < C_0^{-1} 3^{-\alpha} 2^{\alpha-1}$. Since $r\le \varepsilon_0 |x|^{1-\alpha}$, the previous lower bound contradicts $D(\xi,\eta)\le C_0 r$. This proves \eqref{claim1}.

Therefore,
\begin{equation}\label{eq_volume2}
\frac{|x|}{2} \le |x'| \le \frac{3}{2}|x|\quad \textrm{and}\quad |x|+|x'| \le \frac{5}{2} |x|.
\end{equation}
Since $D(\xi,\eta)\le C_0 r$, we obtain
\begin{align*}
\frac{|x-x'|}{\left(|x|+|x'|\right)^\alpha} \le C_0 r \implies |x-x'|\le C_0 r \left(|x|+|x'|\right)^\alpha \le C_0 \left(\frac{5}{2}\right)^\alpha r |x|^\alpha.
\end{align*}
For the $y$-part of $D(\xi,\eta)$, if $\left(|x|+|x'|\right)^{\beta+1-\alpha}\ge |y-y'|$, then
\begin{align*}
\frac{|y-y'|}{2\left(|x|+|x'|\right)^{\beta}} \le \frac{|y-y'|}{\left(|x|+|x'|\right)^{\beta} + |y-y'|^{\frac{\beta}{\beta+1-\alpha}}} \le C_0 r.
\end{align*}
Combining this with \eqref{eq_volume2} yields
\begin{equation*}
|y-y'|\le 2C_0 r \left(|x|+|x'|\right)^{\beta} \le 2C_0 \left(\frac{5}{2}\right)^\beta r |x|^\beta.
\end{equation*}
If instead $\left(|x|+|x'|\right)^{\beta+1-\alpha}\le |y-y'|$, then
\begin{align*}
|y-y'|^{\frac{1-\alpha}{\beta+1-\alpha}} \le 2 C_0 r \implies |y-y'|\le \left(2C_0\right)^{\frac{\beta+1-\alpha}{1-\alpha}} r^{\frac{\beta+1-\alpha}{1-\alpha}}\le \left(2C_0\right)^{\frac{\beta+1-\alpha}{1-\alpha}} r |x|^{\beta}.
\end{align*}
Thus, setting
\begin{equation*}
c_2 = \max \left( C_0 \left(\frac{5}{2}\right)^\alpha,  2C_0 \left(\frac{5}{2}\right)^\beta, \left(2C_0\right)^{\frac{\beta+1-\alpha}{1-\alpha}} \right),
\end{equation*}
gives \eqref{le_volume1}.

It remains to prove the converse inclusion
\begin{equation*}
B_n(x, c_1 r |x|^\alpha) \times B_m(y, c_1 r |x|^\beta) \subset B(\xi, r).
\end{equation*}
This is immediate. Let $\eta = (x',y')$ satisfy $|x-x'|\le c_1 r |x|^\alpha$ and $|y-y'|\le c_1 r |x|^\beta$. It is plain that
\begin{equation*}
d(\xi,\eta)\le C_0 D(\xi,\eta)\le C_0 \frac{c_1 r |x|^\alpha}{\left(|x|+|x'|\right)^\alpha} + C_0 \frac{c_1 r |x|^\beta}{\left(|x|+|x'|\right)^\beta+|y-y'|^{\frac{\beta}{\beta+1-\alpha}}} \le 2 C_0 c_1 r.
\end{equation*}
Choosing $c_1 = (4C_0)^{-1}$ implies $\eta \in B(\xi,r)$, and hence \eqref{le_volume2} follows.

Finally, \eqref{le_volume3} follows directly from a similar analysis as above and involves no additional difficulty, so the details are omitted.
\end{proof}


\begin{remark}
Note that the constants $c_1$ and $c_2$ depend only on $C_0,\alpha,\beta$. Therefore, by requiring
\begin{equation*}
\varepsilon_0< \min \left( C_0^{-1} 3^{-\alpha} 2^{\alpha-1}, \frac{c_2}{2} \right),
\end{equation*}
it follows that for all $r\le \varepsilon_0 |x|^{1-\alpha}$ and all $\eta = (x',y') \in B_n(x, c_2 r |x|^\alpha) \times B_m(y, c_2 r |x|^\beta)$,
\begin{equation*}
|x|\le |x-x'|+|x'| \le c_2 r |x|^{\alpha} + |x'| \le \frac{|x|}{2} +|x'|,
\end{equation*}
and hence $|x|\le 2|x'|\le 3|x|$.
\end{remark}

\begin{definition}
A Grushin ball $B=B(\xi,r) = B((x,y), r)$ is called
\begin{itemize}
\item remote if
\begin{equation*}
r \le \varepsilon_0 |x|^{1-\alpha},
\end{equation*}
where $\varepsilon_0$ comes from Lemma~\ref{le_volume}, so that for all $\eta = (x',y')\in B$, one has $|x-x'|\le \frac{|x|}{2}$ and hence $|x|\sim |x'|$;
\item anchored if $B=B(\xi,r) = B((0,y),r)$.
\end{itemize}
\end{definition}




We briefly recall the curvature-dimension notation used below. Let $\mathcal L$ be a non-positive diffusion operator with carré du champ
\begin{equation*}
    \Gamma(f,g)
    :=
    2^{-1} \bigl(\mathcal L(fg)-f\mathcal Lg-g \mathcal Lf\bigr),
    \qquad
    \Gamma(f):=\Gamma(f,f).
\end{equation*}
The iterated carré du champ is defined by
\begin{equation*}
    \Gamma_2(f,g)
    :=
    2^{-1} \bigl(\mathcal L\Gamma(f,g)-\Gamma(f, \mathcal Lg)-\Gamma(g, \mathcal Lf)\bigr),
    \qquad
    \Gamma_2(f):=\Gamma_2(f,f).
\end{equation*}
Given $N_1 \in\mathbb R$ and $N_2 \in(0,\infty]$, one says that $\mathcal L$ satisfies the curvature-dimension inequality $\operatorname{CD}(N_1,N_2)$ if
\begin{equation*}
    \Gamma_2(f)
    \ge
    N_1\Gamma(f)+ N_2^{-1} (\mathcal Lf)^2
\end{equation*}
for all sufficiently smooth functions $f$, with the last term omitted when $N_2=\infty$. 









By setting $\mathcal L=-L$, we have the following gradient estimates for the heat kernel.

\begin{lemma}\cite{H4}\label{Grushin_gradient}
Let $n,m\ge 1$, $\alpha\in (0,1)$, and $\beta\ge 0$. Then 
\begin{enumerate}[label=(\roman*)]
    \item for any $f\in C^\infty(\mathbb{R}^n \setminus \{0\} \times \mathbb{R}^m)$, the following curvature-dimension inequality holds:
    \begin{equation*}
        \Gamma_2(f) \ge - \frac{C_{n,m,\alpha,\beta}}{|x|^{2-2\alpha}} \Gamma(f) + C_{n,m,\alpha,\beta}' (Lf)^2;
    \end{equation*}
    \item the following gradient estimate holds:
    \begin{align*}
        \left|\nabla_{L}e^{-tL}(\xi,\eta)\right|
        \le \left(\frac{1}{\sqrt{t}} + \frac{1}{|x|^{1-\alpha}}\right)
        \frac{C}{V(\xi,\sqrt{t})}e^{-\frac{d(\xi,\eta)^2}{ct}},
    \end{align*}
    where $\xi = (x,y) \in \mathbb{R}^n \setminus \{0\} \times \mathbb{R}^m$.
\end{enumerate}
\end{lemma}







\begin{remark}
On remote balls the coefficients are uniformly elliptic after the intrinsic rescaling. Hence the required gradient bound may also follow from standard interior gradient estimates for uniformly parabolic equations; see \cite{Lieberman}, for instance.
\end{remark}







\subsection{Poincar\'e inequality}

We establish the following variants of the Poincar\'e inequality.


\begin{lemma}\label{PE}
Let $n,m\ge 1$, $\alpha\in (0,1)$ and $\beta\ge 0$. Let $B(\xi_0, r) = B((x_0,y_0), r)$ be a remote ball. Then there exist $\gamma \ge 1$ (depending only on the structural constants) and $C_P>0$ such that for all $1\le p < \infty$, the following Poincar\'e inequality holds:
\begin{equation}
    \int_B |f-f_B|^p d\xi \le C_P r^p \int_{\gamma B} |\nabla_L f|^p d\xi,
\end{equation}
where the constant $C_P$ depends only on $n,m,\alpha,\beta,p$.
\end{lemma}

\begin{proof}
Let $B=B(\xi_0,r) = B((x_0,y_0),r)$ be a remote ball, and let $f\in C^\infty(\gamma B)$. By Lemma~\ref{le_volume}, there exists $c_2>0$ such that
\begin{equation*}
    B(\xi_0, r) \subset B_n(x_0,c_2 r |x_0|^\alpha) \times B_m(y_0,c_2 r |x_0|^\beta)=: Q.
\end{equation*}
Define the rescaling map
\begin{equation*}
    \Lambda: (x', y') \mapsto \left( x_0 + r|x_0|^\alpha x', y_0 + r|x_0|^\beta y' \right),\quad \forall x'\in \mathbb{R}^n,\quad y'\in \mathbb{R}^m,
\end{equation*}
and the function $F(x',y') = f\left( \Lambda(x',y') \right)$. Clearly, $\Lambda$ maps $Q_0:=B_n(0,c_2) \times B_m(0,c_2)$ onto $Q$. By the standard Poincar\'e inequality on $\mathbb{R}^{n+m}$ for the convex set $Q_0$,
\begin{align*}
    \int_{Q_0} |F(x',y') - F_{Q_0}|^p dx'dy' \le C(n,m,c_2,p) \int_{Q_0} |\nabla F(x',y')|^p dx'dy'.
\end{align*}
After a change of variables,
\begin{equation*}
    \int_{Q} |f(\xi) - f_{Q}|^p d\xi \le C(n,m,c_2,p) r^p \int_Q \left[|x_0|^{2\alpha} |\nabla_xf(\xi)|^2 + |x_0|^{2\beta} |\nabla_y f(\xi)|^2 \right]^{\frac{p}{2}} d\xi.
\end{equation*}
Since $B$ is remote, for all $(x,y)\in Q$, $|x_0|\le (1+c_2\varepsilon_0) |x|$, and therefore
\begin{align*}
    \int_{Q} |f(\xi) - f_{Q}|^p d\xi &\le C(n,m,c_2,p) r^p \int_Q \left[|x|^{2\alpha} |\nabla_xf(\xi)|^2 + |x|^{2\beta} |\nabla_y f(\xi)|^2 \right]^{\frac{p}{2}} d\xi \\
    &= C(n,m,c_2,p) r^p \int_Q |\nabla_L f(\xi)|^p d\xi.
\end{align*}
Moreover, Lemma~\ref{le_volume} gives $Q \subset \frac{c_2}{c_1} B=: \gamma B$, where $\gamma \ge 1$. It follows that
\begin{align*}
    \int_B |f-f_B|^p d\xi &\le 2^p \int_B |f-f_Q|^p d\xi \\
    &\le C(n,m,c_2,p) r^p \int_Q |\nabla_L f(\xi)|^p d\xi \\
    &\le C(n,m,c_2,p) r^p \int_{\gamma B} |\nabla_L f(\xi)|^p d\xi,
\end{align*}
which completes the proof.
\end{proof}



Next, we introduce a remote-ball version telescoping lemma. The proof is standard and is based on the classical telescoping argument; see \cite{HK} for example. 



\begin{lemma}\label{le_tele}
Let $n,m\ge 1$, $\alpha\in (0,1)$ and $\beta \ge 0$. Let $B=B(\xi_0,r) = B((x_0,y_0),r)$ be a ball such that $3B$ is remote. Then there exists a constant $C>0$, depending only on the doubling constant and on $C_P$, such that for all $f\in C^{1}(2\gamma B)$ and for all $\xi,\eta\in B$,
\begin{equation}\label{eq:remote-lemma-312}
|f(\xi)-f(\eta)|
\le
C\, d(\xi,\eta)\Big(M_2(|\nabla_L f|\mathbf{1}_{2\gamma B})(\xi)+M_2(|\nabla_L f|\mathbf{1}_{2\gamma B})(\eta)\Big),
\end{equation}
where
\[
M_2 h(\xi):=\sup_{B\ni \xi}\left(\frac{1}{|B|}\int_B |h|^2 \right)^{1/2}.
\]
\end{lemma}

\begin{proof}
See Appendix~\ref{app:tele}.
\end{proof}

% Lemma~\ref{le_tele} is a remote-ball telescoping estimate that will be used later for local oscillation bounds; see Part~\ref{part3}.





% \subsection{Modified Bessel functions}
% Throughout, $I_\nu$ and $K_\nu$ denote the standard modified Bessel functions. We use several classical facts about these functions, including their power-series expansions, recurrence formulae, Wronskian identity, integral representation, and small-argument asymptotics; see Watson~\cite[Chapter~III]{Watson} and Abramowitz--Stegun~\cite[Chapter~9]{AS}. For $\nu>-1$ and $t>0$ we also use the normalized branches
% \begin{equation*}
%     \mathcal I_\nu(t):=t^{-\nu}I_\nu(t), \qquad \mathcal K_\nu(t):=t^\nu K_\nu(t).
% \end{equation*}
% Set
% \begin{equation*}
%     S_\nu(t):=\frac{I_{\nu+1}(t)}{I_\nu(t)}.
% \end{equation*}
% It is well known that they are two linearly independent solutions of the modified Bessel equation
% \begin{equation*}
%     t^2 F''(t) + t F'(t) - \left( t^2 + \nu^2 \right)F(t) = 0.
% \end{equation*}
% For small and large arguments, one has
% \begin{align*}
%     I_\nu(t) \sim 
%     \begin{cases}
%         t^\nu, & t\le 1,\\
%         t^{-\frac{1}{2}} e^t, & t\ge 1.
%     \end{cases} \qquad 
%     K_\nu(t) \sim \begin{cases}
%         t^{-\nu}, & t\le 1,\\
%         t^{-\frac{1}{2}} e^{-t}, & t\ge 1.
%     \end{cases}
% \end{align*}


% For later use, we record the following lemmas whose proofs will be postponed to Appendix~\ref{app:B} below.


% \begin{lemma}\label{lem:Bessel_regular_branch}
% Let $\nu_*>-1$. For every integer $N\geq0$, there exists a constant
% $C_{N,\nu_*}>0$ such that, for all $\nu\ge\nu_*$,
% \begin{equation*}
%     \sup_{0\leq t\leq1}
%     \left|
%         \frac{d^k}{dt^k}\mathcal I_\nu(t)
%     \right|
%     \leq
%     C_{N,\nu_*},
%     \qquad
%     0\leq k\leq N.
% \end{equation*}
% Moreover,
% \begin{equation*}
%     \frac{d}{dt}\mathcal I_\nu(t)=t\mathcal I_{\nu+1}(t).
% \end{equation*}

% \end{lemma}






% \begin{lemma}\label{lem_S_derivative_bound}
% Let $\nu_*>-1$. For every integer $N\ge 0$, there exist constants
% $C_{N,\nu_*},A_N>0$ such that, for all $\nu\ge\nu_*$,
% \begin{equation*}
%     \left|S_\nu^{(k)}(t)\right|
%     \le
%     C_{N,\nu_*}(1+\nu+t)^{A_N},
%     \qquad
%     0\le k\leq N,\quad t>0.
% \end{equation*}
% Consequently, for every multi-index $\gamma$ with $|\gamma|\leq N$,
% \begin{equation*}
% \left| \partial_Y^\gamma\bigl(|Y|S_\nu(|Y|)\bigr) \right| \le C_{N,\nu_*}(1+\nu+|Y|)^{A_N}, \qquad Y\in\mathbb R^m.
% \end{equation*}

% \end{lemma}




% \begin{lemma}\label{lem_bessel_inputs_high_frequency}
% Let $\nu_*>-1$. There exist constants $c,C>0$, depending only on $\nu_*$, such
% that, for all $\nu\geq\nu_*$,
% \begin{equation}\label{eq_S_lower_large}
%     \frac{I_{\nu+1}(t)}{I_\nu(t)}
%     \geq
%     c,
%     \qquad
%     t\geq C(1+\nu).
% \end{equation}
% Moreover, $\mathcal K_\nu(t)=t^\nu K_\nu(t)$ is decreasing, and
% \begin{equation}\label{eq_K_normalized_exp_decay}
%     \frac{\mathcal K_\nu(t_2)}{\mathcal K_\nu(t_1)}
%     \leq
%     C\exp\{-c(t_2-t_1)\},
%     \qquad
%     t_2\geq t_1\geq C(1+\nu).
% \end{equation}
% In addition, if $t>0$ and $\nu\ge 0$. Then
% \begin{equation}\label{eq:product:Bessel}
%     I_\nu(t)K_\nu(t)
%     \le
%     \frac{C}{t+\nu}.
% \end{equation}

% \end{lemma}











% \noindent \textbf{Notation.} Throughout the article, set for $n,m\ge 1$, $\alpha \in (0,1)$ and $\beta \ge 0$,
% \begin{align*}
%     \mathfrak a := \beta+1-\alpha,\qquad \mathfrak b := \frac{n+2\alpha-2}{2}, \qquad \mathcal Q:= \frac{n+m(\beta+1-\alpha)}{1-\alpha} = \frac{n+m \mathfrak a}{1-\alpha}.
% \end{align*}
% The notation $c,C$ denote positive constants which could variant from line to line.


\subsection{Modified Bessel functions}

Throughout, $I_\nu$ and $K_\nu$ denote the standard modified Bessel functions. We shall use several classical properties of these functions, including their power-series expansions, recurrence relations, Wronskian identity, integral representations, and small- and large-argument asymptotics; see Watson~\cite[Chapter~III]{Watson} and Abramowitz--Stegun~\cite[Chapter~9]{AS}. For $\nu>-1$ and $t>0$, we also introduce the normalized functions
\begin{equation*}
    \mathcal I_\nu(t):=t^{-\nu}I_\nu(t),
    \qquad
    \mathcal K_\nu(t):=t^\nu K_\nu(t).
\end{equation*}
We further set
\begin{equation*}
    S_\nu(t):=\frac{I_{\nu+1}(t)}{I_\nu(t)}.
\end{equation*}
The functions $I_\nu$ and $K_\nu$ form a fundamental system of solutions to the modified Bessel equation
\begin{equation*}
    t^2F''(t)+tF'(t)-\bigl(t^2+\nu^2\bigr)F(t)=0.
\end{equation*}
For each fixed $\nu>-1$, one has
\begin{equation*}
    I_\nu(t)\sim_\nu
    \begin{cases}
        t^\nu, & 0<t\leq 1,\\
        t^{-1/2}e^t, & t\geq 1.
    \end{cases}
\end{equation*}
Since $K_{-\nu}=K_\nu$, for each fixed $\nu\in\mathbb R$ one has
\begin{equation*}
    K_\nu(t)\sim_\nu
    \begin{cases}
        t^{-|\nu|}, & 0<t\leq 1,\quad \nu\neq 0,\\
        1+|\log t|, & 0<t\leq 1,\quad \nu=0,\\
        t^{-1/2}e^{-t}, & t\geq 1.
    \end{cases}
\end{equation*}
For later use, we record the following lemmas, whose proofs are postponed to Appendix~\ref{app:B}.

\begin{lemma}\label{lem:Bessel_regular_branch}
Let $\nu_*>-1$. For every integer $N\geq 0$, there exists a constant
$C_{N,\nu_*}>0$ such that, for every $\nu\geq\nu_*$,
\begin{equation*}
    \sup_{0\leq t\leq 1}
    \left|
        \frac{d^k}{dt^k}\mathcal I_\nu(t)
    \right|
    \leq C_{N,\nu_*},
    \qquad 0\leq k\leq N.
\end{equation*}
Here $\mathcal I_\nu$ is understood at $t=0$ through its smooth extension. Moreover,
\begin{equation*}
    \frac{d}{dt}\mathcal I_\nu(t)
    =t\mathcal I_{\nu+1}(t).
\end{equation*}
\end{lemma}

\begin{lemma}\label{lem_S_derivative_bound}
Let $\nu_*>-1$. For every integer $N\geq 0$, there exist constants
$C_{N,\nu_*},A_N>0$ such that, for every $\nu\geq\nu_*$,
\begin{equation*}
    \left|S_\nu^{(k)}(t)\right|
    \leq
    C_{N,\nu_*}(1+\nu+t)^{A_N},
    \qquad
    0\leq k\leq N,\quad t>0.
\end{equation*}
Consequently, for every multi-index $\gamma$ satisfying $|\gamma|\leq N$,
\begin{equation*}
    \left|
        \partial_Y^\gamma\bigl(|Y|S_\nu(|Y|)\bigr)
    \right|
    \leq
    C_{N,\nu_*}(1+\nu+|Y|)^{A_N},
    \qquad
    Y\in\mathbb R^m.
\end{equation*}
\end{lemma}

\begin{lemma}\label{lem_bessel_inputs_high_frequency}
Let $\nu_*>-1$. There exist constants $c,C>0$, depending only on $\nu_*$, such that, for every $\nu\geq\nu_*$,
\begin{equation}\label{eq_S_lower_large}
    \frac{I_{\nu+1}(t)}{I_\nu(t)}
    \geq c,
    \qquad
    t\geq C(1+\nu).
\end{equation}
Moreover, the function $\mathcal K_\nu(t)=t^\nu K_\nu(t)$ is decreasing on $(0,\infty)$, and
\begin{equation}\label{eq_K_normalized_exp_decay}
    \frac{\mathcal K_\nu(t_2)}{\mathcal K_\nu(t_1)}
    \leq
    C\exp\bigl\{-c(t_2-t_1)\bigr\},
    \qquad
    t_2\geq t_1\geq C(1+\nu).
\end{equation}
In addition, for every $\nu\geq 0$ and $t>0$,
\begin{equation}\label{eq:product:Bessel}
    I_\nu(t)K_\nu(t)
    \leq
    \frac{C}{t+\nu}.
\end{equation}
\end{lemma}

\noindent\textbf{Notation.}
Throughout the article, for $n,m\geq 1$, $\alpha\in(0,1)$, and $\beta\geq 0$, we set
\begin{equation*}
    \mathfrak a:=\beta+1-\alpha,
    \qquad
    \mathfrak b:=\frac{n+2\alpha-2}{2},
    \qquad
    \mathcal Q
    :=\frac{n+m(\beta+1-\alpha)}{1-\alpha}
    =\frac{n+m\mathfrak a}{1-\alpha}.
\end{equation*}
The symbols $c$ and $C$ denote positive constants that may vary from line to line.





\part{Riesz transform for \texorpdfstring{$n\ge 2$}{n>=2}}\label{part1}


In this part we prove the full-range Riesz transform estimate when the degenerate variable has dimension $n\geq2$. By the abstract criterion recalled in Section~\ref{sec3}, it is enough to prove reverse H\"older inequalities for gradients of Grushin harmonic functions. Remote balls are controlled by local Li--Yau type estimates, since the coefficients are essentially non-degenerate there. The only new difficulty is the anchored case, where the ball meets the singular set.

After rescaling an anchored ball to a fixed cylinder, we represent a localized harmonic function as the sum of a Poisson term and a Green term. The singularity at $\rho=0$ is encoded in modified Bessel functions after Fourier transform in $y$ and spherical harmonic decomposition in the angular variable. Unlike the one-dimensional analysis of Part~\ref{part2}, the Bessel order now depends on the angular frequency $\ell$ and tends to infinity with it. Thus the proof requires Bessel estimates which are uniform simultaneously in the order, the Fourier frequency, and the distance to the singular set.


 



\section{Reduction}\label{sec3}






From Robinson–Sikora \cite{RS}, it is well known that the Grushin space satisfies the volume doubling condition \eqref{D} and the Gaussian upper bound for the heat kernel \eqref{UE}; see also Theorem~\ref{RS_Grushin}. Moreover, it is known from \cite{RS2} that the $L^2$-Poincar\'e inequality holds when $n\ge 2$, or when $n=1$ and $\alpha \in (0,1/2)$. That is,
\begin{equation}\label{P2}\tag{$\mathrm{P}$}
\int_B |f-f_B|^2 d\xi \le C r(B)^2 \int_{\lambda B} |\nabla_L f|^2 d\xi,\quad f\in C^\infty(\lambda B)
\end{equation}
for some $\lambda \ge 1$ and any Grushin ball $B$ with radius $r(B)$.

By \cite{CJKS}, it is known that, under suitable conditions, \eqref{R_p} and \eqref{eq:RH:Jiang} are equivalent. More precisely,

\begin{theorem}\cite[Theorem~1.9]{CJKS}\label{thm:CJKS}
Let $X$ be a non-compact Dirichlet metric measure space endowed with a ``carr\'e du champ''. Assume that $X$ satisfies \eqref{D} and \eqref{UE}. Let $2<p<\infty$. If \eqref{P2} holds, then \eqref{eq:RH:Jiang} and \eqref{R_p} are equivalent.
\end{theorem}






\begin{remark}
We emphasize the following points.
\begin{enumerate}[label=(\roman*)]
\item In fact, \cite[Theorem~1.9]{CJKS} requires only a \eqref{P_p_intro} rather than \eqref{P2}. However, under \eqref{D}, condition \eqref{P2} implies \eqref{P_p_intro} for all $p>2$; see \cite{HK}. For convenience, we therefore state their result in the above form.

\item By \cite[Theorem~1.1]{RS2}, it is known that \eqref{P2} holds on Grushin spaces if either $n\ge 2$, or $n=1$ and $0<\alpha<1/2$. Thus, in these two cases, the Riesz transform problem is reduced to proving \eqref{eq:RH:Jiang}.
\item In the special case $n=1$ and $1/2\le \alpha<1$, the global condition \eqref{P2} fails, and hence Theorem~\ref{thm:CJKS} cannot be applied. In \cite{Jiang}, Jiang refines the above theorem by relaxing \eqref{P2} to a so-called Poincar\'e inequality on the ends, which requires \eqref{P2} only for balls outside a compact set. This refinement makes the equivalence stable under gluing operations, for example on manifolds or Dirichlet spaces with ends. However, in our strongly degenerate case, \eqref{P2} holds only for balls outside the singular hyperplane $\{x=0\}$, which is not compact. Therefore, Jiang's theorem does not apply either. We will address this issue later in Part~\ref{part4}.

\end{enumerate}
\end{remark}





The following two lemmas further reduce the problem to establishing reverse H\"older estimates on anchored balls.


\begin{lemma}\label{reduce_RH}
Let $n\ge 1$, $m\ge 1$, $\alpha\in (0,1)$ and $\beta \ge 0$. Let $1<p\le \infty$. The $L^p$-reverse H\"older inequality \eqref{eq:RH:Jiang} holds provided that it holds for all anchored balls of the form $B((0,y_0), r)$ and for all remote balls $B((x_0,y_0), r)$ with $r\le \varepsilon_0 |x_0|^{1-\alpha}$.
\end{lemma}

\begin{proof}
Suppose that $B=B(\xi_0,r) = B((x_0,y_0), r)$ is neither anchored nor remote. That is, $x_0 \ne 0$ and $r>\varepsilon_0 |x_0|^{1-\alpha}$. Clearly, $B((x_0,y_0), r) \subset B((0,y_0), C_1 r)\subset B((0,y_0), \lambda C_1 r) \subset B((x_0,y_0), C_2 r)$ for some $C_2 \ge C_1 \ge 1$. Therefore, by the $L^p$-reverse H\"older inequality for anchored balls and \eqref{D},
\begin{align*}
    \left( \fint_B |\nabla_L u|^p d\xi \right)^{\frac{1}{p}}
    &\le C \left( \fint_{B((0,y_0), C_1 r)} |\nabla_L u|^p d\xi \right)^{\frac{1}{p}} \\
    &\le \frac{C}{r} \fint_{B((0,y_0), \lambda C_1 r)} |u| d\xi \\
    &\le \frac{C}{r} \fint_{C_2 B} |u| d\xi
\end{align*}
as required.
\end{proof}

\begin{lemma}\label{RH_remote}
Let $n\ge 1$, $m\ge 1$, $\alpha\in (0,1)$ and $\beta \ge 0$. The $L^p$-reverse H\"older inequality holds for remote balls $B((x_0,y_0), r)$ with $r\le \varepsilon_0 |x_0|^{1-\alpha}$ for all $1<p\le \infty$.
\end{lemma}

\begin{proof}
Let $B=B(\xi_0,r)=B((x_0,y_0),r)$ be a remote ball. Choosing $\varepsilon_0$ sufficiently small if necessary, it may be assumed that $4B$ is also remote. It follows that $|x| \sim |x_0|$ for all $(x,y) \in 4B$, and that $r \le \varepsilon_0 |x_0|^{1-\alpha}$. By Lemma~\ref{Grushin_gradient}, on $4B$ the local curvature-dimension inequality is bounded below by $\mathrm{CD}(-C_1/r^{2}, C_2)$ for some $C_1,C_2>0$. Let $u$ be $L$-harmonic in $4B$. Then $\widetilde{u}:=u+\|u\|_{L^\infty(2B)}+\varepsilon$ with $\varepsilon>0$ is positive and $L$-harmonic in $2B$. Hence the local Li--Yau type inequality (see \cite[Theorem~1.4]{ZZ}) gives
\begin{equation}\label{Li-Yau}
   \sup_B |\nabla_L u| = \sup_B |\nabla_L \widetilde{u}| \le \frac{C}{r} \|\widetilde{u}\|_{L^\infty(2B)}\le \frac{C}{r} \left(\|u\|_{L^\infty(2B)}+\varepsilon\right).
\end{equation}
Next, since the Grushin space satisfies \eqref{D} and \eqref{UE}, \cite[Proposition~2.1]{CJKS} yields
\begin{equation}\label{meanvalue}
    \|u\|_{L^\infty(2B)} \le C \fint_{4B} |u| d\xi.
\end{equation}
Combining \eqref{Li-Yau} and \eqref{meanvalue}, and then letting $\varepsilon \to 0$, yields the $L^\infty$-reverse H\"older inequality for remote balls, and hence the $L^p$-reverse H\"older inequality for all $1<p\le \infty$.
\end{proof}


Combining Lemma~\ref{reduce_RH} and Lemma~\ref{RH_remote}, the $L^p$-boundedness of the Riesz transform is now reduced to proving the following result.



\begin{theorem}\label{thm:RH}
Let $n\ge 2$, $m\ge 1$, $\alpha\in (0,1)$ and $\beta \ge 0$. Let $B=B\big((0,y_0), r\big)$ be an anchored Grushin ball. Suppose that $u$ is $L$-harmonic in $4B$. Then for every $2<p\le \infty$, the reverse H\"older inequality \eqref{eq:RH:Jiang} holds. In particular, there exists a constant $C>0$ such that
\begin{equation*}
    \left( \fint_B |\nabla_L u|^p d\xi \right)^{\frac{1}{p}} \le  \frac{C}{r} \fint_{4B} |u| d\xi.
\end{equation*}
\end{theorem}

\begin{proof}[Reduction of Theorem~\ref{thm:RH}]
For $R>0$, consider the Grushin box
\begin{equation*}
Q_R := B_n(0,R) \times B_m \left(0, R^{\mathfrak a} \right).
\end{equation*}
By Lemma~\ref{le_volume},
\begin{equation*}
B\big((0,y_0), r\big) \sim B_n\left(0, r^{\frac{1}{1-\alpha}}\right) \times B_m\left(y_0, r^{\frac{\mathfrak a}{1-\alpha}}\right).
\end{equation*}
Define
\begin{equation*}
v(\xi) = v(x,y) := u\left( r^{\frac{1}{1-\alpha}} x, y_0 + r^{\frac{\mathfrak a}{1-\alpha}} y \right).
\end{equation*}
Then $v$ is $L$-harmonic in $Q_4$. After the change of variables
\begin{equation*}
x \mapsto r^{\frac{1}{1-\alpha}}x, \qquad y \mapsto y_0+ r^{\frac{\mathfrak a}{1-\alpha}}y,
\end{equation*}
it is enough to prove
\begin{equation}\label{eq_RHv}
\left( \fint_{Q_1}|\nabla_L v|^p d\xi \right)^{\frac{1}{p}} \le C \fint_{Q_4} |v| d\xi.
\end{equation}
Moreover, by \eqref{D} and \eqref{UE}, \cite[Proposition~2.1]{CJKS} gives
\begin{equation*}
\| v \|_{L^\infty(Q_3)} \le C \fint_{Q_4} |v| d\xi.
\end{equation*}
Therefore, \eqref{eq_RHv} follows for all $2<p\le \infty$ once
\begin{equation}\label{eq_RH_show}
\sup_{Q_1} |\nabla_L v| \le C \| v \|_{L^\infty(Q_3)}
\end{equation}
is established.

Passing to polar coordinates,
\begin{equation*}
x = \rho \theta, \quad \rho = |x|,\quad \theta \in \mathbb{S}^{n-1},
\end{equation*}
a straightforward computation shows that
\begin{equation*}
L = - \rho^{2\alpha} \partial_\rho^2 - (n-1+2\alpha) \rho^{2\alpha-1} \partial_\rho + \rho^{2\alpha -2} \Delta_{\mathbb{S}^{n-1}} + \rho^{2\beta} \Delta_y.
\end{equation*}
Choose $\chi \in C_c^\infty\left( B_m\big(0, 3^{\mathfrak a}\big) \right)$ such that $\chi = 1$ on $B_m\big(0, 2^{\mathfrak a}\big)$, and define
\begin{equation*}
\omega(\rho,\theta,y) = \chi(y) v(\rho,\theta,y).
\end{equation*}
Then $\omega$ is a smooth function compactly supported in the $y$-variable. Moreover, $\omega = v$ on $Q_2$, and
\begin{equation*}
\sup_{Q_1}|\nabla_L v| = \sup_{Q_1}|\nabla_L \omega|.
\end{equation*}
Consider the cylinder
\begin{equation*}
\mathcal{C}:= (0,2)_\rho \times \mathbb{S}^{n-1}_\theta \times \mathbb{R}^m_y.
\end{equation*}
The function $\omega$ satisfies
\begin{equation*}
L \omega = f:= [L,\chi]v = |x|^{2\beta} v \Delta_y \chi - 2|x|^{2\beta} \nabla_y \chi \cdot \nabla_y v.
\end{equation*}
Clearly, $f$ is supported in $\{2^{\mathfrak a}\le |y|\le 3^{\mathfrak a}\}$. Let
\begin{equation*}
h(\theta,y) := \omega(2,\theta,y).
\end{equation*}
By the Friedrichs representation,
\begin{equation*}
\omega = \mathcal{T}h + \mathcal{S}f,
\end{equation*}
where $\mathcal{T}h$ is the solution to the boundary value problem
\begin{align}\label{BVP}
\begin{cases}
Lg = 0, & \textrm{on}\quad \mathcal{C},\\
g = h, & \textrm{when}\quad \rho=2,\\
g\quad \textrm{has finite energy at}\quad \rho=0,
\end{cases}
\end{align}
and $\mathcal{S}f$ is the solution to the Dirichlet problem
\begin{align}\label{DP}
\begin{cases}
Lg = f, & \textrm{on}\quad \mathcal{C},\\
g = 0, & \textrm{when}\quad \rho=2,\\
g\quad \textrm{has finite energy at}\quad \rho=0.
\end{cases}
\end{align}
Accordingly, \eqref{eq_RH_show}, and hence Theorem~\ref{thm:RH}, will follow once the two estimates
\begin{align*}
\sup_{Q_1} |\nabla_L \mathcal{T}h| \le C \| v \|_{L^\infty(Q_3)},
\end{align*}
and
\begin{align*}
\sup_{Q_1} |\nabla_L \mathcal{S}f| \le C \| v \|_{L^\infty(Q_3)}
\end{align*}
are proved.
\end{proof}





















\section{Kernel construction}



We now construct the Poisson and Green kernels needed for the anchored estimate. The singularity at $\rho=0$ is encoded in a family of modified Bessel equations after Fourier transform in $y$ and spherical harmonic decomposition in the angular variable. The
finite-energy condition selects the regular $I$-branch and excludes the singular $K$-branch. The remainder of the section records the resulting kernels and the estimates required for the proof of the reverse H\"older inequality.


Throughout Part~\ref{part1}, for each $\ell \ge 0$, we use the notation
\begin{equation*}
\lambda_\ell:= \ell(\ell+n-2),\qquad \nu_\ell:= \frac{\sqrt{\mathfrak b^2 + \lambda_\ell}}{\mathfrak a}, \qquad \gamma_\ell:= - \mathfrak b + \sqrt{\mathfrak b^2 + \lambda_\ell}, \qquad \delta_\ell := \mathfrak b+\sqrt{\mathfrak b^2+\lambda_\ell}.
\end{equation*}


% \noindent \textbf{Notation.} Throughout the article, set for $n,m\ge 1$, $\alpha \in (0,1)$ and $\beta \ge 0$,
% \begin{align*}
%     \mathfrak a := \mathfrak a,\qquad \mathfrak b := \frac{n+2\alpha-2}{2}, \qquad \mathcal Q:= \frac{n+m(\mathfrak a)}{1-\alpha} = \frac{n+m \mathfrak a}{1-\alpha}
% \end{align*}




\subsection{The Poisson multiplier}

Let $\{\Theta_{\ell,\mu}\}_{\substack{\ell\ge 0\\ 1\le \mu\le d_\ell}}$ be the spherical harmonics on $\mathbb{S}^{n-1}$, that is,
\begin{equation*}
    \Delta_{\mathbb{S}^{n-1}} \Theta_{\ell,\mu} = \lambda_{\ell} \Theta_{\ell,\mu},
\end{equation*}
where $d_\ell \sim (1+\ell)^{n-2}$ denotes the dimension of the space of spherical harmonics of degree $\ell$; see \cite[Chapter~IV]{SW2}. Consider the equation
\begin{equation*}
    - \rho^{2\alpha} \partial_\rho^2 g - (n-1+2\alpha) \rho^{2\alpha-1} \partial_\rho g + \rho^{2\alpha -2} \Delta_{\mathbb{S}^{n-1}}g + \rho^{2\beta} \Delta_y g = 0.
\end{equation*}
Taking the Fourier transform in the $y$-variable and then expanding in spherical harmonics gives
\begin{align*}
    0 &= - \rho^{2\alpha} \partial_\rho^2 \hat{g} - (n-1+2\alpha) \rho^{2\alpha-1} \partial_\rho \hat{g} + \rho^{2\alpha -2} \Delta_{\mathbb{S}^{n-1}} \hat{g} + \rho^{2\beta} |Y|^2 \hat{g}\\
    &= \sum_{\ell=0}^\infty \sum_{\mu=1}^{d_\ell} \left(  - \rho^{2\alpha} \partial_\rho^2 a_{\ell,\mu} - (n-1+2\alpha) \rho^{2\alpha-1} \partial_\rho a_{\ell,\mu} + \rho^{2\alpha -2} \lambda_{\ell} a_{\ell,\mu} + \rho^{2\beta} |Y|^2 a_{\ell,\mu}  \right) \Theta_{\ell,\mu}(\theta),
\end{align*}
where
\begin{equation*}
    \hat{g}(\rho,\theta,Y) = \int_{\mathbb{R}^m} e^{-iy \cdot Y} g(\rho,\theta,y)\, dy,
\end{equation*}
\begin{equation*}
    a_{\ell,\mu}(\rho,Y) = \int_{\mathbb{S}^{n-1}} \hat{g}(\rho,\theta,Y)\, \overline{\Theta_{\ell,\mu}(\theta)}\, d\theta,
\end{equation*}
and
\begin{equation}\label{eq_spherical_decom}
    \hat{g}(\rho,\theta,Y) = \sum_{\ell=0}^\infty \sum_{\mu=1}^{d_\ell} a_{\ell,\mu}(\rho,Y) \Theta_{\ell,\mu}(\theta).
\end{equation}
By orthogonality,
\begin{equation}\label{eq_equation1}
    \rho^{2} \partial_\rho^2 a_{\ell,\mu} + (n-1+2\alpha) \rho \partial_\rho a_{\ell,\mu} -  \lambda_{\ell} a_{\ell,\mu} - \rho^{2\beta-2\alpha+2} |Y|^2 a_{\ell,\mu} = 0
\end{equation}
for each $\ell,\mu$. Introduce
\begin{equation*}
    t:= \frac{|Y|}{\mathfrak a} \rho^{\mathfrak a}
    \quad \textrm{and}\quad
    U_{\ell,\mu,Y}(t):= \rho^{\mathfrak b} a_{\ell,\mu}(\rho,Y).
\end{equation*}
A direct computation shows that \eqref{eq_equation1} is equivalent to
\begin{equation}\label{eq_equation2}
    t^2 U_{\ell,\mu,Y}'' + t U_{\ell,\mu,Y}' - \left( t^2 + \nu_\ell^2 \right)U_{\ell,\mu,Y} = 0.
\end{equation}
The solutions of the modified Bessel equation \eqref{eq_equation2} are linear combinations of $I_{\nu_\ell}(t)$ and $K_{\nu_\ell}(t)$, where $I_{\nu_\ell}$ and $K_{\nu_\ell}$ denote the modified Bessel functions. Therefore,
\begin{equation}\label{a_l,mu}
    a_{\ell,\mu}(\rho,Y) = C_1(\ell,\mu,Y) \rho^{-\mathfrak b} I_{\nu_\ell}\left( \frac{|Y| \rho^{\mathfrak a}}{\mathfrak a} \right) + C_2(\ell,\mu,Y) \rho^{-\mathfrak b} K_{\nu_\ell}\left( \frac{|Y| \rho^{\mathfrak a}}{\mathfrak a} \right).
\end{equation}
However, since the Friedrichs representation imposes a finite energy condition at $\rho=0$, the following claim holds.

\textit{Claim.} The finite-energy solution of the singular Sturm--Liouville equation \eqref{eq_equation1} excludes the $K_{\nu_\ell}$-branch, whose energy near $\rho=0$ is infinite.

\begin{proof}[Proof of Claim]
Let $g$ be the corresponding finite-energy solution. Then its weighted radial derivative satisfies, for sufficiently small $\varepsilon>0$,
\begin{equation*}
\int_0^\varepsilon \int_{\mathbb S^{n-1}} \int_{\mathbb R^m} \rho^{2\alpha}|\partial_\rho g(\rho,\theta,y)|^2 \rho^{n-1}\, dy\, d\theta\, d\rho < \infty.
\end{equation*}
Since $\partial_\rho$ commutes with both the Fourier transform in the $y$-variable and the spherical harmonic projection in $\theta$, it follows from \eqref{eq_spherical_decom} that
\begin{equation}\label{eq_finite_energy}
    \partial_\rho \hat{g}(\rho,\theta,Y) =
\sum_{\ell=0}^{\infty}\sum_{\mu=1}^{d_\ell}\partial_\rho a_{\ell,\mu}(\rho,Y) \Theta_{\ell,\mu}(\theta).
\end{equation}
By Plancherel's theorem and the orthogonality of the spherical harmonics,
\begin{align*}
&\int_0^\varepsilon \int_{\mathbb S^{n-1}} \int_{\mathbb R^m} \rho^{2\alpha}|\partial_\rho \hat{g}(\rho,\theta,Y)|^2 \rho^{n-1}\, dY\, d\theta\, d\rho\\
&\qquad = \sum_{\ell=0}^{\infty} \sum_{\mu=1}^{d_\ell} \int_{\mathbb R^m} \int_0^\varepsilon \rho^{2\alpha} |\partial_\rho a_{\ell,\mu}(\rho,Y)|^2 \rho^{n-1}\, d\rho\, dY.
\end{align*}
Since the left-hand side is finite, every term on the right-hand side is finite as well. Hence, for each $(\ell,\mu)$ and for a.e. $Y\in \mathbb{R}^m$,
\begin{equation}\label{eq_radial_gradient}
\int_0^\varepsilon \rho^{2\alpha} |\partial_\rho a_{\ell,\mu} (\rho,Y)|^2 \rho^{n-1}\, d\rho < \infty.
\end{equation}
Next, since $K_{\nu_\ell}(t) \sim t^{-\nu_\ell}$ for $t\le 1$, the $K_{\nu_\ell}$-branch in \eqref{a_l,mu} behaves like $\rho^{\widetilde{\gamma}_\ell}$ as $\rho \to 0$, where
\begin{equation*}
    \widetilde{\gamma}_\ell = -\mathfrak b - \sqrt{\mathfrak b^2 + \lambda_\ell}.
\end{equation*}
Therefore, its contribution to \eqref{eq_radial_gradient} is
\begin{equation*}
    \int_0^\varepsilon \rho^{2(\alpha + \widetilde{\gamma}_\ell -1 ) } \rho^{n-1}\, d\rho = \infty,\quad \forall \ell \ge 0,
\end{equation*}
since $n\ge 2$. Consequently, the entire $K_{\nu_\ell}$-branch must be excluded from the solution \eqref{a_l,mu}, that is,
\begin{equation*}
    C_2(\ell,\mu,Y) = 0,\quad \forall \ell,\mu,Y.
\end{equation*}
This proves the \textit{Claim}.
\end{proof}

Define
\begin{equation*}
    \phi_{\ell, Y}(\rho) = \phi_{\ell}(\rho,Y):= \rho^{-\mathfrak b} I_{\nu_\ell}\left( \frac{|Y| \rho^{\mathfrak a}}{\mathfrak a} \right)
    \qquad \textrm{and}\qquad
    M_{\ell, Y}(\rho) = M_{\ell}(\rho,Y):= \frac{\phi_{\ell}(\rho,Y)}{\phi_{\ell}(2,Y)}.
\end{equation*}
It follows from \eqref{BVP} and the \textit{Claim} that
\begin{align*}
    \widehat{\mathcal{T}h}(\rho,\theta,Y) = \sum_{\ell=0}^\infty \sum_{\mu=1}^{d_\ell} C_1(\ell,\mu,Y) \phi_{\ell}(\rho,Y)\, \Theta_{\ell,\mu}(\theta).
\end{align*}
Moreover, the boundary condition gives
\begin{align*}
    \widehat{\mathcal{T}h}(2,\theta,Y) = \hat{h}(\theta,Y) = \sum_{\ell=0}^\infty \sum_{\mu=1}^{d_\ell} b_{\ell,\mu}(Y) \Theta_{\ell,\mu}(\theta),
\end{align*}
where
\begin{equation*}
    b_{\ell,\mu}(Y) = \int_{\mathbb{S}^{n-1}} \hat{h}(\theta,Y)\, \overline{\Theta_{\ell,\mu}(\theta)}\, d\theta.
\end{equation*}
Orthogonality therefore implies
\begin{equation*}
    C_1(\ell,\mu,Y) = \frac{b_{\ell,\mu}(Y)}{\phi_{\ell}(2,Y)}.
\end{equation*}
Taking the inverse Fourier transform yields
\begin{align*}
    \mathcal{T}h(\rho,\theta,y) = \iint_{\mathbb{S}^{n-1}\times \mathbb{R}^m} \mathcal{K}_{\rho}(\theta,y; \theta', y') h(\theta', y')\, d\theta'\, dy',
\end{align*}
where
\begin{equation}\label{eq_kernel_of_T}
    \mathcal{K}_{\rho}(\theta,y; \theta', y') = C_m \sum_{\ell=0}^\infty \Pi_{\ell}(\theta,\theta') \int_{\mathbb{R}^m} e^{i(y-y')\cdot Y} M_{\ell}(\rho,Y)\, dY,
\end{equation}
and
\begin{equation*}
    \Pi_{\ell}(\theta,\theta') = \sum_{\mu=1}^{d_\ell} \Theta_{\ell,\mu}(\theta) \overline{\Theta_{\ell,\mu}(\theta')}.
\end{equation*}

\subsection{The Green kernel}\label{ssec:Green:n>2}

Next, we construct the Green kernel. Let $\mathcal{G}(Z,Z')$ (where $Z = (\rho,\theta,y)$ and $Z' = (\sigma,\theta',y')$) denote the Green kernel associated with $L$ on $\mathcal{C}$ under the Dirichlet boundary condition at $\rho=2$ and the local finite-energy condition at $\rho=0$. In the same spirit, by taking the Fourier transform in the $y$-$y'$ variables and the spherical harmonic decomposition in the $\theta$-$\theta'$ variables, the Green kernel takes the form
\begin{equation}\label{eq_Green's_kernel}
    \mathcal{G}(Z,Z') = \sum_{\ell=0}^\infty \Pi_{\ell}(\theta, \theta') \int_{\mathbb{R}^m} e^{i(y-y') \cdot Y} G_{\ell, Y}(\rho, \sigma) \, dY,
\end{equation}
where $G_{\ell,Y}$ is the one-dimensional Green kernel associated with the operator
\begin{equation*}
    \mathcal{L}_{\ell,Y} := - \rho^{2\alpha} \partial_{\rho}^2 - (n-1+2\alpha) \rho^{2\alpha-1} \partial_\rho + \lambda_\ell \rho^{2\alpha-2} + |Y|^2 \rho^{2\beta}.
\end{equation*}
Equivalently,
\begin{equation}
    \mathcal{S}f(Z) = \iiint_{\mathcal{C}} \mathcal{G}(Z,Z') f(Z') \, \sigma^{n-1} \, d\sigma \, d\theta' \, dy'.
\end{equation}
Recall the function
\begin{equation*}
    \phi_{\ell}(\rho,Y):= \rho^{-\mathfrak b} I_{\nu_\ell}\!\left( \frac{|Y| \rho^{\mathfrak a}}{\mathfrak a} \right).
\end{equation*}
This is the solution to \eqref{eq_equation1} with finite energy at $\rho=0$. To construct the Green kernel, another independent solution satisfying the Dirichlet boundary condition at $\rho=2$ is required. Define
\begin{equation*}
    \psi_{\ell}(\rho, Y):= \rho^{-\mathfrak b} \left[ K_{\nu_\ell} \!\left( \frac{|Y| \rho^{\mathfrak a}}{\mathfrak a} \right) - \frac{K_{\nu_\ell} \!\left( \frac{|Y| 2^{\mathfrak a}}{\mathfrak a} \right)}{I_{\nu_\ell} \!\left( \frac{|Y| 2^{\mathfrak a}}{\mathfrak a} \right)} \, I_{\nu_\ell} \!\left( \frac{|Y| \rho^{\mathfrak a}}{\mathfrak a} \right) \right].
\end{equation*}
Note that $\psi_{\ell}$ solves \eqref{eq_equation1} and satisfies $\psi_{\ell}(2)=0$. Hence,
\begin{equation*}
    G_{\ell,Y}(\rho,\sigma) = \begin{cases}
        A_\ell(\sigma, Y) \, \phi_{\ell}(\rho, Y), & 0 < \rho < \sigma,\\[4pt]
        B_\ell(\sigma, Y) \, \psi_{\ell}(\rho, Y), & \sigma < \rho < 2.
    \end{cases}
\end{equation*}
Continuity of the Green kernel at $\rho=\sigma$ yields the first equation:
\begin{equation*}
    A_\ell(\rho,Y) \, \phi_{\ell}(\rho, Y) = B_\ell(\rho,Y) \, \psi_{\ell}(\rho, Y).
\end{equation*}
The derivative jump condition, namely $\mathcal{L}_{\ell,Y}G_{\ell,Y}(\rho,\sigma) = \frac{\delta_{\sigma}(\rho)}{\sigma^{n-1}}$, gives the second equation:
\begin{equation*}
    \sigma^{n+2\alpha-1} \bigl[ A_\ell(\sigma,Y) \, \phi_{\ell}'(\sigma,Y) - B_\ell(\sigma,Y) \, \psi_{\ell}'(\sigma,Y) \bigr] = 1.
\end{equation*}
Solving these equations gives
\begin{equation*}
    A_\ell(\sigma,Y) = \frac{\psi_{\ell}(\sigma,Y)}{W_{\ell,Y}}, \qquad B_\ell(\sigma,Y) = \frac{\phi_{\ell}(\sigma,Y)}{W_{\ell,Y}},
\end{equation*}
where
\begin{equation*}
    W_{\ell,Y} = t^{n+2\alpha-1} \bigl[ \phi_{\ell}'(t,Y) \, \psi_\ell(t,Y) - \phi_{\ell}(t,Y) \, \psi'_\ell(t,Y) \bigr]
\end{equation*}
is the weighted Wronskian. Using the identity
\begin{equation*}
    I_{\nu_\ell}'(t) \, K_{\nu_\ell}(t) - I_{\nu_\ell}(t) \, K_{\nu_\ell}'(t) = \frac{1}{t},
\end{equation*}
it is readily verified that $W_{\ell,Y} = \mathfrak a = \beta+1-\alpha$.

Therefore,
\begin{align}\label{eq_kernel_Green}
    G_{\ell,Y}(\rho,\sigma) = \mathfrak a^{-1}
    \begin{cases}
        \phi_{\ell}(\rho,Y) \, \psi_{\ell}(\sigma,Y), & 0 < \rho < \sigma,\\[4pt]
        \phi_{\ell}(\sigma,Y) \, \psi_{\ell}(\rho,Y), & \sigma < \rho < 2.
    \end{cases}
\end{align}





\section{Proof of Theorem~\ref{thm:RH}}


This section is devoted to proving Theorem~\ref{thm:RH}.

\begin{theorem}\label{lem:Poisson_kernel}
Let $h\in L^\infty(\mathbb{S}^{n-1} \times \mathbb{R}^m)$, and let $\mathcal{T}h$ be the solution of the homogeneous problem \eqref{BVP} with finite energy at $\rho=0$. Then
\begin{equation*}
    \sup_{0<\rho\le 1} \bigl\| \nabla_L \mathcal{T}h \bigr\|_{L^\infty(\mathbb{S}^{n-1} \times \mathbb{R}^m)} \le C \|h\|_{L^\infty(\mathbb{S}^{n-1} \times \mathbb{R}^m)}.
\end{equation*}
\end{theorem}

\begin{proof}
The proof is based on a multiplier-to-kernel argument. First, we estimate the three multiplier families corresponding to the radial, angular, and base components. Note that
\begin{equation*}
    \nabla_L = \bigl( \rho^\alpha \partial_\rho, \, \rho^{\alpha-1} \nabla_{\mathbb{S}^{n-1}}, \, \rho^{\beta} \nabla_y \bigr).
\end{equation*}
By the explicit formula \eqref{eq_kernel_of_T}, it is enough to estimate the kernels associated with the following three multipliers:
\begin{align*}
m_{\ell,\rho}^{\mathrm{rad}}(Y)&:= \rho^{\alpha} \partial_\rho M_{\ell}(\rho, Y),\\
m_{\ell,\rho}^{\mathrm{ang}}(Y)&:= \rho^{\alpha-1} \sqrt{\lambda_\ell} \, M_{\ell}(\rho, Y),\\
m_{\ell,\rho}^{\mathrm{bas},j}(Y)&:= i Y_j \rho^{\beta} M_{\ell}(\rho, Y), \quad 1 \le j \le m,
\end{align*}
where, in particular, $m_{0,\rho}^{\mathrm{ang}}$ is understood to be zero.

The following estimate will be used, whose proof is postponed to Proposition~\ref{lem_detailed_poisson_multiplier} below: for every integer $N \ge 0$,
\begin{equation}\label{eq_detailed_poisson_multiplier_goal_copy}
    \sum_{|\gamma|\le N}
    \bigl\|
        \partial_Y^\gamma m_{\ell,\rho}^{(*)}
    \bigr\|_{L^1(\mathbb R^m_Y)}
    \le
    C_N e^{-c_N \ell},
    \quad
    *\in \{ \mathrm{rad}, \mathrm{ang}, \mathrm{bas}, j\},
\end{equation}
uniformly for $0<\rho\le 1$ and $\ell\ge 0$.

Next, we use the identity
\begin{equation*}
    (I+\Delta_Y)^N e^{i Y \cdot (y-y')} = (1 + |y-y'|^2)^N e^{iY\cdot(y-y')}.
\end{equation*}
Together with \eqref{eq_detailed_poisson_multiplier_goal_copy}, integration by parts yields
\begin{align*}
    \biggl| \int_{\mathbb{R}^m} e^{iY\cdot(y-y')} m_{\ell,\rho}^{(*)}(Y) \, dY \biggr|
    &= (1+|y-y'|^2)^{-N} \biggl| \int_{\mathbb{R}^m} e^{iY\cdot(y-y')} (I+\Delta_Y)^N m_{\ell,\rho}^{(*)}(Y) \, dY \biggr|\\[4pt]
    &\le C_N (1+|y-y'|^2)^{-N} \sum_{|\gamma|\le 2N} \bigl\|
        \partial_Y^\gamma m_{\ell,\rho}^{(*)}
    \bigr\|_{L^1(\mathbb R^m_Y)}\\[4pt]
    &\le C_N (1+|y-y'|^2)^{-N} e^{-c_N\ell}.
\end{align*}
Choosing $N>\frac{m}{2}$, we obtain
\begin{equation}\label{eq_multiplier_final}
    \sup_{ \substack{0<\rho\le 1\\ y\in \mathbb{R}^{m}} } \biggl\| \int_{\mathbb{R}^m} e^{iY\cdot(y-y')} m_{\ell,\rho}^{(*)}(Y) \, dY \biggr\|_{L^1_{y'}(\mathbb{R}^m)} \le C_N e^{-c_N \ell}, \quad
    *\in \{ \mathrm{rad}, \mathrm{ang}, \mathrm{bas}, j\}.
\end{equation}
On the other hand, the standard spherical projection estimates \cite[Chapter~IV]{SW2} (see also \cite[Theorem~2.9, Section~2.3]{AHan}) yield
\begin{align}\label{eq_sph_pro1}
    \sup_{\theta \in \mathbb{S}^{n-1}} \int_{\mathbb{S}^{n-1}} |\Pi_{\ell}(\theta, \theta')| \, d\theta' \le C (1+\ell)^{\frac{n-2}{2}},
\end{align}
and
\begin{align}\label{eq_sph_pro2}
    \sup_{\theta \in \mathbb{S}^{n-1}} \int_{\mathbb{S}^{n-1}} \bigl| \lambda_\ell^{-1/2} \nabla_{\theta} \Pi_\ell(\theta, \theta') \bigr| \, d\theta' \le  C (1+\ell)^{\frac{n-2}{2}}.
\end{align}
Indeed, \eqref{eq_sph_pro1} follows from Cauchy--Schwarz and the rotational invariance of the spherical projection:
\begin{align*}
    \int_{\mathbb{S}^{n-1}} |\Pi_{\ell}(\theta, \theta')| \, d\theta'
    &\le C \biggl( \int_{\mathbb{S}^{n-1}} |\Pi_\ell(\theta,\theta')|^2 \, d\theta' \biggr)^{\frac{1}{2}}
    = C \, \Pi_\ell(\theta, \theta)^{\frac{1}{2}}
    = C \, d_\ell^{\frac{1}{2}}
    \le C (1+\ell)^{\frac{n-2}{2}}.
\end{align*}
Similarly, for \eqref{eq_sph_pro2},
\begin{align*}
    \int_{\mathbb{S}^{n-1}} \bigl| \lambda_\ell^{-1/2} \nabla_{\theta} \Pi_\ell(\theta, \theta') \bigr| \, d\theta'
    &\le C \lambda_\ell^{-1/2} \biggl( \int_{\mathbb{S}^{n-1}} |\nabla_\theta \Pi_\ell(\theta,\theta')|^2 \, d\theta' \biggr)^{\frac{1}{2}}\\
    &= C \lambda_\ell^{-1/2} \biggl( \sum_{\mu=1}^{d_\ell} |\nabla_\theta \Theta_{\ell,\mu}(\theta)|^2 \biggr)^{\frac{1}{2}}\\
    &\le C \, d_\ell^{\frac{1}{2}}
    \le C (1+\ell)^{\frac{n-2}{2}},
\end{align*}
where the penultimate bound follows from
\begin{align*}
    \sum_{\mu=1}^{d_\ell} |\nabla_\theta \Theta_{\ell,\mu}(\theta)|^2
    &= \int_{\mathbb{S}^{n-1}} \sum_{\mu=1}^{d_\ell} |\nabla_\theta \Theta_{\ell,\mu}(\theta)|^2 \, d\theta\\
    &= \sum_{\mu=1}^{d_\ell} \int_{\mathbb{S}^{n-1}} \Delta_{\mathbb{S}^{n-1}} \Theta_{\ell,\mu}(\theta) \, \Theta_{\ell,\mu}(\theta) \, d\theta\\
    &= \lambda_\ell \, d_\ell.
\end{align*}
Combining \eqref{eq_sph_pro1}, \eqref{eq_sph_pro2}, and \eqref{eq_multiplier_final}, it follows that
\begin{align*}
    \sup_{\substack{0<\rho\le 1\\ y\in \mathbb{R}^{m}\\ \theta \in \mathbb{S}^{n-1}}} \iint_{\mathbb{S}^{n-1} \times \mathbb{R}^m} \bigl| \nabla_L \mathcal{K}_\rho (\theta,y; \theta', y') \bigr| \, d\theta' \, dy' \le C \sum_{\ell=0}^\infty (1+\ell)^{\frac{n-2}{2}} e^{-c\ell} \le C,
\end{align*}
which completes the proof of Theorem~\ref{lem:Poisson_kernel}.
\end{proof}

Before proving the key estimate \eqref{eq_detailed_poisson_multiplier_goal_copy}, further introduce the notation
\begin{equation*}
\mathfrak A_\rho:=\frac{\rho^{\mathfrak a}}{\mathfrak a},
\qquad
\mathfrak A_2:=\frac{2^{\mathfrak a}}{\mathfrak a},
\qquad
\tau_\rho(Y):= |Y| \mathfrak A_\rho,
\qquad
\tau_2(Y):=|Y|\mathfrak A_2,
\end{equation*}
and define
\begin{equation*}
H_{\ell,\rho}(r)
:=
\frac{\mathcal I_{\nu_\ell}(\mathfrak A_\rho r)}
{\mathcal I_{\nu_\ell}(\mathfrak A_2 r)},
\qquad
\mathcal I_\nu(t):=t^{-\nu}I_\nu(t).
\end{equation*}
For simplicity, the dependence of $\tau_\rho$ and $\tau_2$ on $Y$ will often be suppressed.

All Poisson multiplier estimates reduce to the behavior of a single normalized Bessel ratio. The main difficulty here is that the order $\nu_\ell$ is not fixed, but grows linearly with the spherical harmonic frequency $\ell$. The next proposition records the uniform estimates needed later. Its proof separates the transition region $r\lesssim 1+\ell$ from the large-frequency region $r\gg1+\ell$.

\begin{proposition}\label{lem_transition_ratio_derivatives}
For every integer $N \ge 0$ and every $0\le k\le N$, there exist constants $C_N,A_N,c_N>0$ such that
\begin{align}\label{eq_transition_H_derivatives}
    |H_{\ell,\rho}^{(k)}(r)|
    \le \begin{cases}
        C_N(1+\ell)^{A_N}, & r \le C^*(1+\ell),\\[4pt]
        C_N(1+r)^{A_N}e^{-c_N r}, & r > C^*(1+\ell),
    \end{cases}
\end{align}
where $C^*>1$ is a sufficiently large constant.

Consequently, for every multi-index $|\gamma|\le N$,
\begin{align}\label{eq_transition_Y_derivatives}
\bigl|\partial_Y^\gamma H_{\ell,\rho}(|Y|)\bigr| \le \begin{cases}
    C_N (1+\ell)^{A_N}, & |Y|\le C^* (1+\ell),\\[4pt]
    C_N (1+|Y|)^{A_N}e^{-c_N|Y|}, & |Y|>C^*(1+\ell).
\end{cases}
\end{align}
\end{proposition}

\begin{proof}
Using the identity
\begin{equation*}
    tI_\nu'(t)-\nu I_\nu(t)=tI_{\nu+1}(t),
\end{equation*}
it follows that $\mathcal I_\nu'(t)=t^{-\nu}I_{\nu+1}(t)$. Hence,
\begin{equation*}
    \frac{d}{dt}\log\mathcal I_\nu(t) = S_\nu(t)
    \qquad\text{and}\qquad
    \log H_{\ell,\rho}(r)
    =
    -\int_{\mathfrak A_\rho r}^{\mathfrak A_2 r}S_{\nu_\ell}(t) \, dt.
\end{equation*}
On the other hand, note that $\log H_{\ell,\rho}(r) = \log\mathcal I_{\nu_\ell}(\mathfrak A_\rho r) - \log\mathcal I_{\nu_\ell}(\mathfrak A_2 r)$.
% \begin{equation*}
%     \log H_{\ell,\rho}(r)
%     =
%     \log\mathcal I_{\nu_\ell}(\mathfrak A_\rho r)
%     -
%     \log\mathcal I_{\nu_\ell}(\mathfrak A_2 r),
% \end{equation*}
This implies that for $j\ge 1$,
\begin{equation}\label{eq_log_H_derivative_formula}
    \frac{d^j}{dr^j}\log H_{\ell,\rho}(r)
    =
    \mathfrak A_\rho^j S_{\nu_\ell}^{(j-1)}(\mathfrak A_\rho r)
    -
    \mathfrak A_2^j S_{\nu_\ell}^{(j-1)}(\mathfrak A_2 r).
\end{equation}

\textit{Case 1: $r\le C^*(1+\ell)$.} Since $0<\mathfrak A_\rho\le \mathfrak A_2$ and $\nu_\ell\sim 1+\ell$, Lemma~\ref{lem_S_derivative_bound} yields
\begin{equation}\label{eq_log_H}
    \Bigl|
        \frac{d^j}{dr^j}\log H_{\ell,\rho}(r)
    \Bigr|
    \le
    C_j(1+\ell)^{A_j}.
\end{equation}
Moreover, as $t^{-\nu}I_\nu(t)$ is increasing on $(0,\infty)$, $H_{\ell,\rho}(r)$ is positive and uniformly bounded by $1$, i.e., $0<H_{\ell,\rho}(r)\le 1$.
% \begin{equation*}
%     0<H_{\ell,\rho}(r)\le 1.
% \end{equation*}
Applying Faà di Bruno's formula to $H_{\ell,\rho}=\exp\bigl(\log H_{\ell,\rho}\bigr)$, gives, for each integer $k\ge 1$,
\begin{equation}\label{eq_faa_di_bruno_H}
    H_{\ell,\rho}^{(k)}(r)
    =
    H_{\ell,\rho}(r)
    \, P_k\Bigl(
        \frac{d}{dr}\log H_{\ell,\rho}(r),
        \dots,
        \frac{d^k}{dr^k}\log H_{\ell,\rho}(r)
    \Bigr),
\end{equation}
where $P_k$ is a polynomial depending only on $k$. Combine \eqref{eq_faa_di_bruno_H} with \eqref{eq_log_H} to get
\begin{equation}\label{eq_H_k}
    |H_{\ell,\rho}^{(k)}(r)|
    \le
    C_N (1+\ell)^{A_N} H_{\ell,\rho}(r)
    \le
    C_N (1+\ell)^{A_N}.
\end{equation}
This proves \eqref{eq_transition_H_derivatives} in the range $r\le C^*(1+\ell)$. 

It remains to pass from one-dimensional derivatives in $r$ to derivatives in $Y$. If $\tau_2 \ge 1$, that is, $|Y|\ge \mathfrak A_2^{-1}$, then the chain rule yields
\begin{align*}
\bigl|\partial_Y^\gamma H_{\ell,\rho}(|Y|)\bigr|
&\le \sum_{k=0}^{|\gamma|} |Y|^{k-|\gamma|} \bigl|H_{\ell,\rho}^{(k)}(|Y|)\bigr|
\le C_N (1+\ell)^{A_N}, \qquad |Y|\le C^* (1+\ell),
\end{align*}
for all $|\gamma|\le N$. Next, if $\tau_2 \le 1$, that is, $|Y|\le \mathfrak A_2^{-1}$, first observe that $H_{\ell,\rho}(|Y|)$ is in fact a smooth function of $|Y|^2$. Indeed, recalling the power series expansion
\begin{equation*}
    I_\nu(t)
    =
    \frac{1}{\Gamma(\nu+1)}
    \Bigl(\frac{t}{2}\Bigr)^{\!\nu}
    \sum_{q=0}^{\infty}
    \frac{1}{q!(\nu+1)_q}
    \Bigl(\frac{t^2}{4}\Bigr)^{\!q},
\end{equation*}
we find
\begin{equation*}
    \mathcal I_\nu(t)
    :=
    t^{-\nu}I_\nu(t)
    =
    \frac{1}{2^\nu\Gamma(\nu+1)}
     E_\nu(t^2), \qquad  E_\nu(s)
    =
    \sum_{q=0}^{\infty}
    \frac{1}{q!(\nu+1)_q}
    \Bigl(\frac{s}{4}\Bigr)^{\!q}.
\end{equation*}
Note that $E_\nu(s) \ge 1$ for all $s>0$. Also,
\begin{equation*}
    H_{\ell,\rho}(r)
    =
    \frac{
        E_{\nu_\ell}(\mathfrak A_\rho^2 r^2)
    }{
        E_{\nu_\ell}(\mathfrak A_2^2 r^2)
    }.
\end{equation*}
Equivalently,
\begin{equation*}
    H_{\ell,\rho}(|Y|)
    =
    h_{\ell,\rho}(|Y|^2), \qquad h_{\ell,\rho}(s) :=  \frac{
        E_{\nu_\ell}(\mathfrak A_\rho^2 s)
    }{
        E_{\nu_\ell}(\mathfrak A_2^2 s)
    }.
\end{equation*}
% We show next that for every $k\ge 0$, $|h_{\ell,\rho}^{(k)}(s)|\le C_k$ for all $\mathfrak A_2^2 s\le 1$.
% \begin{equation*}
%     |h_{\ell,\rho}^{(k)}(s)|\le C_k, \qquad \forall\, \mathfrak A_2^2 s\le 1.
% \end{equation*}
By \eqref{eq_A_derivatives_small}, for every fixed $N$, $\sup_{0\le s\le 1}\bigl| E_\nu^{(k)}(s) \bigr| \le C_N$ ($0\le k\le N$),
% \begin{equation*}
%     \sup_{0\le s\le 1}
%     \bigl|
%         E_\nu^{(k)}(s)
%     \bigr|
%     \le C_N,
%     \qquad
%     0\le k\le N,
% \end{equation*}
uniformly in $\nu> -1$. The quotient rule therefore gives
\begin{equation}\label{eq_h_deri}
    |h_{\ell,\rho}^{(k)}(s)|\le C_k, \quad \forall\, s\le \mathfrak A_2^{-2}.
\end{equation}
Hence, for every multi-index $\gamma$,
\begin{equation*}
    \bigl| \partial_Y^{\gamma} H_{\ell,\rho}(|Y|)\bigr|
    = \bigl|\partial_Y^\gamma
    \bigl[
        h_{\ell,\rho}(|Y|^2)
    \bigr]\bigr|
    \le \sum_{k=0}^{|\gamma|} C_{\gamma, k} \bigl|P_{\gamma,k}(|Y|)\bigr| \, \bigl|h_{\ell,\rho}^{(k)}\bigl(|Y|^2\bigr)\bigr|,
\end{equation*}
where $P_{\gamma,k}$ is a polynomial depending only on $\gamma$ and $k$. Since $|Y|\le C^*(1+\ell)$, and $|Y|\le \mathfrak A_2^{-1}$, \eqref{eq_h_deri} implies
\begin{equation*}
    \bigl|
        \partial_Y^\gamma H_{\ell,\rho}(|Y|)
    \bigr|
    \le
    C_N (1+\ell)^{A_N},\quad \forall\, |\gamma|\le N,
\end{equation*}
as required.

\textit{Case 2: $r>C^*(1+\ell)$.} By a standard Amos-type estimate \cite[Corollary~1]{KB} (see Lemma~\ref{lem_bessel_inputs_high_frequency}), there exist $c_0, C_1>0$ such that
\begin{equation}\label{eq_amos_lower_bound}
    S_\nu(t)=\frac{I_{\nu+1}(t)}{I_\nu(t)}\geq c_0,\quad \text{if}\quad t\ge C_1(1+\nu).
\end{equation}
Also, since $0<\rho\le 1$, $\frac{\mathfrak A_\rho}{\mathfrak A_2} = \frac{\rho^{\mathfrak a}}{2^{\mathfrak a}} \le 2^{-\mathfrak a}$.
% \begin{equation*}
%     \frac{\mathfrak A_\rho}{\mathfrak A_2} = \frac{\rho^{\mathfrak a}}{2^{\mathfrak a}} \le 2^{-\mathfrak a} <1.
% \end{equation*}
Now, if $C^*$ is chosen sufficiently large, then for every
$t\in(2^{- \mathfrak a} \mathfrak A_2 r,\mathfrak A_2 r]$,
\begin{equation*}
    t\ge 2^{- \mathfrak a} \mathfrak A_2 r\ge \mathfrak a^{-1} C^*(1+\ell) \ge C_1(1+\nu_\ell).
\end{equation*}
Using \eqref{eq_amos_lower_bound}, we deduce, for some $c>0$,
\begin{align*}
\log H_{\ell,\rho}(r) &= -\int_{\mathfrak A_\rho r}^{\mathfrak A_2 r}S_{\nu_\ell}(t) \, dt
\le -\int_{2^{- \mathfrak a} \mathfrak A_2 r}^{\mathfrak A_2 r}S_{\nu_\ell}(t) \, dt \\
&\le -c_0(1-2^{- \mathfrak a})\mathfrak A_2 r \le -cr.
\end{align*}
It follows that
\begin{equation}\label{eq_H_large_decay}
    0<H_{\ell,\rho}(r)\le e^{-cr},
    \qquad
    r>C^* (1+\ell).
\end{equation}
Next, we estimate the derivatives of $H_{\ell,\rho}$. Combining \eqref{eq_log_H_derivative_formula} with Lemma~\ref{lem_S_derivative_bound}, 
% and noting that $0< \mathfrak A_\rho \le \mathfrak A_2$, 
yields, for $j\ge 1$,
\begin{equation}\label{eq_log_H_derivative_polynomial}
    \Bigl|
        \frac{d^j}{dr^j}\log H_{\ell,\rho}(r)
    \Bigr|
    \le
    C_j(1+r)^{A_j},
    \qquad
    r>C^*(1+\ell).
\end{equation}
Therefore, by \eqref{eq_H_large_decay}, \eqref{eq_log_H_derivative_polynomial}, and \eqref{eq_faa_di_bruno_H}, Faà di Bruno's formula gives
\begin{equation*}
    |H_{\ell,\rho}^{(k)}(r)| \le C_N (1+r)^{A_N}e^{-c_N r},\qquad 0\le k \le N,
\end{equation*}
which proves \eqref{eq_transition_H_derivatives} in the range $r>C^*(1+\ell)$. Estimate \eqref{eq_transition_Y_derivatives} then follows immediately from the chain rule for derivatives of radial functions, since $|Y|>C^*(1+\ell)\ge C^*$. This completes the proof of Proposition~\ref{lem_transition_ratio_derivatives}.
\end{proof}




\begin{remark}\label{re:H}
The proof of Proposition~\ref{lem_transition_ratio_derivatives} uses the growth of the order only to keep the constants uniform in $\ell$. If the order is fixed, the same argument applies to every $\mu>-1$. More precisely, for
\begin{equation*}
    H_{\mu,\rho}(r):=
\frac{\mathcal I_\mu(\mathfrak A_\rho r)}{\mathcal I_\mu(\mathfrak A_2 r)},
\end{equation*}
one has, for every integer $N\ge 0$,
\begin{equation*}
    |H_{\mu,\rho}^{(k)}(r)|
\le
\begin{cases}
C_{N,\mu}, & r\le C^*,\\[4pt]
C_{N,\mu}(1+r)^{A_N}e^{-c_N r}, & r>C^*,
\end{cases}
\end{equation*}
for $0\le k\le N$, uniformly in $0<\rho\leq1$. The corresponding
$Y$-derivative estimates for $H_{\mu,\rho}(|Y|)$ follow in the same way. In
Part~\ref{part2} we apply this fixed-order version with $\mu=\nu$ and with
$\mu=-\nu\in(-1,0)$. 
\end{remark}

We now translate the normalized Bessel-ratio estimates into estimates for the
actual Poisson multipliers. The three multiplier families correspond respectively
to the radial derivative, the angular derivative, and the $y$-derivatives in
$\nabla_L$. The factor $(\rho/2)^{\gamma_\ell}$ is crucial: for $\rho\leq1$ it
converts polynomial growth in $\ell$ into exponential summability in $\ell$.





\begin{proposition}\label{lem_detailed_poisson_multiplier}
For every integer $N \ge 0$, there exist constants $C_N, c_N > 0$ such that
\begin{equation}\label{eq_detailed_poisson_multiplier_goal}
    \sum_{|\gamma|\le N}
    \bigl\|
        \partial_Y^\gamma m_{\ell,\rho}^{(*)}
    \bigr\|_{L^1(\mathbb R^m_Y)}
    \le
    C_N e^{-c_N\ell},
    \quad
    *\in \{\mathrm{rad}, \mathrm{ang}, \mathrm{bas}, j\},
\end{equation}
uniformly for $0<\rho\le 1$ and $\ell\ge 0$.
\end{proposition}

\begin{proof}
It suffices to prove that for every multi-index $|\gamma|\le N$, there exist constants $A_N, C_N, c_N>0$ such that
\begin{equation}\label{eq_pointwise_split_bound_Y}
    \bigl|
        \partial_Y^\gamma m_{\ell,\rho}^{(*)}(Y)
    \bigr|
    \le
    C_N
    (1+\ell+|Y|)^{A_N}
    \bigl[
        e^{-c_N \ell}\mathbf 1_{|Y|\le C^* (1+\ell)}
        +
        e^{-c_N |Y|}\mathbf 1_{|Y|>C^* (1+\ell)}
    \bigr],
\end{equation}
uniformly for $0<\rho\le 1$. Once \eqref{eq_pointwise_split_bound_Y} is established, \eqref{eq_detailed_poisson_multiplier_goal} follows immediately. Indeed, \eqref{eq_pointwise_split_bound_Y} suggests
\begin{align*}
    \int_{|Y|\le C^* (1+\ell)}
    (1+\ell+|Y|)^{A_N}e^{-c_N \ell} \, dY
    &\le
    C_N (1+\ell)^{A_N + m}e^{-c_N \ell}
    \le
    C_N e^{-c_N'\ell},
\end{align*}
and
\begin{align*}
    \int_{|Y|> C^* (1+\ell)}
    (1+\ell+|Y|)^{A_N}e^{-c_N |Y|} \, dY
    &\le
    C_N e^{-c_N''\ell}.
\end{align*}
Therefore, $\bigl\| \partial_Y^\gamma m_{\ell,\rho}^{(*)} \bigr\|_{L^1_Y} \le C_N e^{-c_N \ell}$.
% \begin{equation*}
%     \bigl\|
%         \partial_Y^\gamma m_{\ell,\rho}^{(*)}
%     \bigr\|_{L^1_Y}
%     \le
%     C_N e^{-c_N \ell}.
% \end{equation*}
Summing over $|\gamma|\le N$ proves \eqref{eq_detailed_poisson_multiplier_goal}.

It remains to prove \eqref{eq_pointwise_split_bound_Y}. Note that for $\ell\ge 0$,
\begin{equation}\label{eq_M_H}
    M_{\ell,\rho}(Y) = \Bigl( \frac{\rho}{2} \Bigr)^{\!-\mathfrak b} \frac{I_{\nu_\ell}(\tau_\rho)}{I_{\nu_\ell}(\tau_2)} = \Bigl( \frac{\rho}{2} \Bigr)^{\!\gamma_\ell} H_{\ell,\rho}(|Y|).
\end{equation}
Moreover, for $\ell\ge 1$,
\begin{align*}
    \rho^\alpha \partial_\rho M_{\ell, Y}(\rho)
    = \gamma_\ell \rho^{\alpha-1} \Bigl(\frac{\rho}{2}\Bigr)^{\!\gamma_\ell} H_{\ell,\rho}(|Y|) + \rho^{\alpha} \Bigl(\frac{\rho}{2}\Bigr)^{\!\gamma_\ell} \partial_\rho H_{\ell,|Y|}(\rho).
\end{align*}
The identity
\begin{align*}
    \frac{\partial_\rho H_{\ell, |Y|}(\rho)}{H_{\ell, |Y|}(\rho)}
    = \partial_\rho \log H_{\ell, |Y|}(\rho)
    = |Y| \rho^{\mathfrak a-1} \frac{\mathcal{I}'_{\nu_\ell}(\mathfrak A_\rho |Y|)}{\mathcal{I}_{\nu_\ell}(\mathfrak A_\rho |Y|)}
    = |Y| \rho^{\mathfrak a-1} S_{\nu_\ell}(\mathfrak A_\rho |Y|),
\end{align*}
gives
\begin{align}\label{eq_rad_M_H}
    \rho^\alpha \partial_\rho M_{\ell, Y}(\rho)
    = \gamma_\ell 2^{\alpha-1} \Bigl(\frac{\rho}{2}\Bigr)^{\!\gamma_\ell+\alpha-1} H_{\ell, \rho}(|Y|)
    + \mathfrak a 2^{\alpha -1} \Bigl(\frac{\rho}{2}\Bigr)^{\!\gamma_\ell+\alpha - 1} |\mathfrak A_\rho Y| \, S_{\nu_\ell}(\mathfrak A_\rho |Y|) H_{\ell, \rho}(|Y|).
\end{align}
We split the argument into two cases.

\textit{Case 1: $|Y|\le C^* (1+\ell)$.} 
%Note that for $\ell\ge 0$,
% \begin{equation}\label{eq_M_H}
%     M_{\ell,\rho}(Y) = \Bigl( \frac{\rho}{2} \Bigr)^{\!-\mathfrak b} \frac{I_{\nu_\ell}(\tau_\rho)}{I_{\nu_\ell}(\tau_2)} = \Bigl( \frac{\rho}{2} \Bigr)^{\!\gamma_\ell} H_{\ell,\rho}(|Y|).
% \end{equation}
By Proposition~\ref{lem_transition_ratio_derivatives} and \eqref{eq_M_H},
\begin{equation}\label{eq_M_deri}
    \bigl| \partial_Y^\gamma M_{\ell,\rho}(Y) \bigr| \le C_N \Bigl( \frac{\rho}{2} \Bigr)^{\!\gamma_\ell} (1+\ell)^{A_N} \le C_N (1+\ell)^{A_N} e^{-c \ell},\quad \forall\, |\gamma|\le N,
\end{equation}
where the last inequality follows from
\begin{equation*}
     \Bigl(\frac{\rho}{2}\Bigr)^{\!\gamma_\ell} \le e^{- \gamma_\ell \log 2} \le C e^{-c \ell}
\end{equation*}
since $\rho \le 1$ and $\gamma_\ell \sim 1+\ell$. Therefore, since $\alpha-1+\gamma_\ell>0$ for $\ell \ge 1$, $m_{0,\rho}^{(\mathrm{ang})}=0$, and $\sqrt{\lambda_\ell}$ grows at most polynomially in $\ell$, we obtain
\begin{equation*}
    \bigl|\partial_Y^\gamma m_{\ell,\rho}^{(\mathrm{ang})}(Y)\bigr|
    \le
    C_N
    (1+\ell+|Y|)^{A_N}
        e^{-c_N \ell}\mathbf 1_{|Y|\le C^* (1+\ell)},\quad |\gamma|\le N.
\end{equation*}
For the base multiplier $m_{\ell,\rho}^{(\mathrm{bas}, j)}(Y) = i Y_j \rho^{\beta} M_{\ell,\rho}(Y)$, with $|Y|\le C^*(1+\ell)$, $\gamma_\ell + \beta >0$, and \eqref{eq_M_deri}, Leibniz's rule gives
\begin{equation*}
    \bigl|\partial_Y^\gamma m_{\ell,\rho}^{(\mathrm{bas},j)}(Y)\bigr|
    \le
    C_N
    (1+\ell+|Y|)^{A_N}
        e^{-c_N \ell}\mathbf 1_{|Y|\le C^* (1+\ell)},\quad |\gamma|\le N,\quad 1\le j\le m.
\end{equation*}
Next, for the radial multiplier, that is, for $*=\mathrm{rad}$, by \eqref{eq_rad_M_H},
% \begin{align*}
%     \rho^\alpha \partial_\rho M_{\ell, Y}(\rho)
%     = \gamma_\ell \rho^{\alpha-1} \Bigl(\frac{\rho}{2}\Bigr)^{\!\gamma_\ell} H_{\ell,\rho}(|Y|) + \rho^{\alpha} \Bigl(\frac{\rho}{2}\Bigr)^{\!\gamma_\ell} \partial_\rho H_{\ell,|Y|}(\rho),\qquad \ell\ge 1.
% \end{align*}
% The identity
% \begin{align*}
%     \frac{\partial_\rho H_{\ell, |Y|}(\rho)}{H_{\ell, |Y|}(\rho)}
%     = \partial_\rho \log H_{\ell, |Y|}(\rho)
%     = |Y| \rho^{\mathfrak a-1} \frac{\mathcal{I}'_{\nu_\ell}(\mathfrak A_\rho |Y|)}{\mathcal{I}_{\nu_\ell}(\mathfrak A_\rho |Y|)}
%     = |Y| \rho^{\mathfrak a-1} S_{\nu_\ell}(\mathfrak A_\rho |Y|),
% \end{align*}
% gives
% \begin{align}\label{eq_rad_M_H}
%     \rho^\alpha \partial_\rho M_{\ell, Y}(\rho)
%     = \gamma_\ell 2^{\alpha-1} \Bigl(\frac{\rho}{2}\Bigr)^{\!\gamma_\ell+\alpha-1} H_{\ell, \rho}(|Y|)
%     + \mathfrak a 2^{\alpha -1} \Bigl(\frac{\rho}{2}\Bigr)^{\!\gamma_\ell+\alpha - 1} |\mathfrak A_\rho Y| \, S_{\nu_\ell}(\mathfrak A_\rho |Y|) H_{\ell, \rho}(|Y|).
% \end{align}
and $\gamma_\ell +\alpha-1>0$, Proposition~\ref{lem_transition_ratio_derivatives} together with Lemma~\ref{lem_S_derivative_bound} yields
\begin{equation*}
    \bigl| \partial_Y^{\gamma} m_{\ell,\rho}^{(\mathrm{rad})}(Y) \bigr|
    \le C_{\alpha,N} (1+\ell + |Y|)^{A_N} e^{-c_N \ell} \mathbf 1_{|Y|\le C^* (1+\ell)},\quad |\gamma|\le N,\qquad \ell \ge 1.
\end{equation*}
The case $\ell = 0$ is handled similarly. Indeed, since $\gamma_0=0$, $M_{0,\rho}(Y) = H_{0,\rho}(|Y|)$ and hence
\begin{equation*}
    m_{0,\rho}^{(\mathrm{rad})}(Y) = \rho^\alpha \partial_\rho M_{0, Y}(\rho)
    = 2^\beta \Bigl(\frac{\rho}{2}\Bigr)^{\!\beta}|Y| S_{\nu_0}(\mathfrak A_\rho |Y|) H_{0, \rho}(|Y|).
\end{equation*}
% Since $\mathfrak A_\rho\le C$, 
Clearly,
\begin{equation*}
    |Y|S_{\nu_0}(\mathfrak A_\rho |Y|)\le C(1+\nu_0+\mathfrak A_\rho |Y|)\le C(1+|Y|).
\end{equation*}
Moreover, by Lemma~\ref{lem_S_derivative_bound} and the chain rule, for $|\gamma|\ge 1$,
\begin{equation*}
    |\partial_Y^{\gamma}(|Y|S_{\nu_0}(\mathfrak A_\rho |Y|))|
    \le C_N \mathfrak A_\rho^{|\gamma|-1} (1+\nu_0+\mathfrak A_\rho |Y|)^{A_N}
    \le C_N (1+|Y|)^{A_N}.
\end{equation*}
The required $Y$-derivative estimate for $m_{0,\rho}^{(\mathrm{rad})}$ follows at once. 

\textit{Case 2: $|Y|>C^* (1+\ell)$.} With \eqref{eq_M_H}, \eqref{eq_rad_M_H} in mind, using Proposition~\ref{lem_transition_ratio_derivatives}, Lemma~\ref{lem_S_derivative_bound}, the same argument as in \textit{Case 1} implies 
\begin{equation*}
    \bigl|
        \partial_Y^\gamma m_{\ell,\rho}^{(*)}(Y)
    \bigr|
    \le
    C_N
    (1+\ell+|Y|)^{A_N}
        e^{-c_N |Y|}\mathbf 1_{|Y|> C^* (1+\ell)},\quad |\gamma|\le N,\qquad * \in \{\mathrm{rad}, \mathrm{ang}, \mathrm{bas}, j\}.
\end{equation*}
for all $\ell \ge 0$. This proves \eqref{eq_pointwise_split_bound_Y}, and hence Proposition~\ref{lem_detailed_poisson_multiplier}.
\end{proof}







\begin{remark}\label{re:multiplier}
The proof of Proposition~\ref{lem_detailed_poisson_multiplier} has two independent
components: the fixed-order Bessel estimates for the multiplier in the Fourier
variable $Y$, and the summation over spherical harmonics. In the one-dimensional
case of Part~\ref{part2} the second component disappears. Therefore, once Proposition~\ref{lem_transition_ratio_derivatives} is replaced by its fixed-order version from the preceding Remark~\ref{re:H}, the same multiplier-to-kernel argument applies
verbatim. The only new feature in the odd half-line branch is the explicit factor
$\rho^{1-2\alpha}$, whose intrinsic radial derivative contributes the loss $\rho^{-\alpha}$.
\end{remark}

We next turn from the Poisson operator to the Green operator. The Poisson term is
controlled by differentiating normalized boundary multipliers. For the Green term,
the main issue is row integrability of the one-dimensional Green kernel and of its
intrinsic derivatives. 


We first record the following Bessel estimates.

\begin{lemma}\label{lem:Bessel:integral}
Let $\nu>-1$ and $s>|\nu|$. Let $\varepsilon \in \{0,1\}$. The following two estimates hold:
\begin{enumerate}[label=(\roman*)]
    \item if $q\ge -\nu - \varepsilon$ and $q+s\le 2$, then
    \begin{equation*}
        \sup_{t>0}I(q,s,\nu,\varepsilon,t):=\sup_{t>0} t^q I_{\nu+\varepsilon}(t) \int_t^\infty u^{s-1} K_{|\nu|}(u) du \le C,
    \end{equation*}
    \item if $q+s+\nu-|\nu|\ge 0$ and $q+s\le 2$, then
    \begin{equation*}
        \sup_{t>0}II(q,s,\nu,t):= \sup_{t>0} t^q K_{|\nu|}(t) \int_0^t u^{s-1} I_{\nu}(u) du \le C.
    \end{equation*}
\end{enumerate}
\end{lemma}


\begin{proof}
We first show $(i)$. Assume first $t\le 1$. Then by small argument asymptotics of Bessel functions, $I$ is bounded by 
\begin{align*}
    t^{q+\nu+\varepsilon} \left( \int_t^1 u^{s-1-|\nu|} du + \int_1^\infty u^{s-\frac{3}{2}} e^{-u} du   \right) \le C t^{q+\nu+\varepsilon} \le C.
\end{align*}
While if $t>1$, use large argument asymptotic formula instead. $I$ attains upper bound
\begin{align*}
    t^{q-\frac{1}{2}} e^t \int_t^\infty u^{s-\frac{3}{2}} e^{-u} du \le C t^{q+s-2}\le C, 
\end{align*}
where the second to last inequality follows from the elementary estimate $\int_t^\infty u^{\delta} e^{-u}du\le C_\delta t^\delta e^{-t}$ for $\delta\in \mathbb{R}$.

Next, we treat $(ii)$. If $t\le 1$, a parallel analysis gives
\begin{align*}
    II(q,s,\nu,t) \le C t^{q-|\nu|} \int_0^t u^{s-1+\nu} du \le C t^{q+s+\nu-|\nu|}\le C. 
\end{align*}
On the other hand, if $t>1$, one deduces 
\begin{align*}
   II(q,s,\nu,t)\le C t^{q-\frac{1}{2}} e^{-t} \int_0^t u^{s-\frac{3}{2}} e^u du \le C t^{q+s-2} \le C,
\end{align*}
where the second to last inequality is a consequence of estimate $\int_0^t u^{\delta} e^{u}du\le C_\delta t^\delta e^{t}$ for $\delta\in \mathbb{R}$.



\end{proof}




The next estimate is purely mode-wise and will later be
combined with integration by parts in the Fourier variable and summation in $\ell$.





\begin{proposition}\label{prop:Green:Shur}
For $\ell\ge 1$ and $0<\rho\le 2$, we have
\begin{equation}\label{eq_basic_R_schur_improved}
    \int_0^2 G_{\ell,Y}(\rho,\sigma) \, \sigma^{n-1} \, d\sigma \le \begin{cases}
        C \rho^{1-\alpha} (1+\ell)^{-2}, & |Y|\le C^*(1+\ell),\\[4pt]
        C \rho^{1-\alpha} (1+\ell+\tau_\rho)^{-2}, & |Y|> C^*(1+\ell).
    \end{cases}
\end{equation}
In particular,
\begin{equation}\label{eq_basic_R_schur_ell=0}
    \int_0^2 G_{0,Y}(\rho,\sigma) \, \sigma^{n-1} \, d\sigma \le \begin{cases}
        C, & |Y|\le C^*,\\[4pt]
        C (1+\tau_\rho)^{-2}, & |Y|> C^*.
    \end{cases}
\end{equation}
Moreover, for $\ell \ge 0$,
\begin{align}\label{eq_basic_R_schur}
\sup_{0<\rho<2} \rho^{2\beta} \int_0^2 G_{\ell,Y}(\rho,\sigma) \, \sigma^{n-1} \, d\sigma \le \begin{cases}
        C (1+\ell)^{-2}, & |Y|\le C^*(1+\ell),\\[4pt]
        C |Y|^{-2}, & |Y|> C^*(1+\ell).
    \end{cases}
\end{align}
\end{proposition}

\begin{proof}
\textit{Case 1: $|Y|\le C^*(1+\ell)$.} 
%The finite-energy solution at $\rho=0$ is
% \begin{equation*}
%     \phi_{\ell,Y}(\rho) = \rho^{-\mathfrak b}I_{\nu_\ell}(\tau_\rho).
% \end{equation*}
% The solution satisfying the Dirichlet condition at $\rho=2$ is
% \begin{equation*}
%     \psi_{\ell,Y}(\rho) = \rho^{-\mathfrak b} \bigl[ K_{\nu_\ell}(\tau_\rho) - \eta_{\ell,Y} I_{\nu_\ell}(\tau_\rho)\bigr],\qquad \eta_{\ell,Y}
%     :=
%     \frac{K_{\nu_\ell}(\tau_2)}{I_{\nu_\ell}(\tau_2)}.
% \end{equation*}
The one-dimensional Green kernel takes the form (cf.~\eqref{eq_kernel_Green}):
\begin{align}\label{eq_kernel_Green_copy}
G_{\ell,Y}(\rho,\sigma) = \mathfrak a^{-1}
\begin{cases}
\phi_{\ell}(\rho,Y)\,\psi_{\ell}(\sigma,Y), &0<\rho<\sigma,\\[4pt]
\phi_{\ell}(\sigma,Y)\,\psi_{\ell}(\rho,Y), &\sigma < \rho<2,
\end{cases}
\end{align}
where $\phi_{\ell,Y}(\rho) = \rho^{-\mathfrak b}I_{\nu_\ell}(\tau_\rho)$ and
\begin{equation*}
\psi_{\ell,Y}(\rho) = \rho^{-\mathfrak b} \bigl[ K_{\nu_\ell}(\tau_\rho) - \eta_{\ell,Y} I_{\nu_\ell}(\tau_\rho)\bigr],\qquad \eta_{\ell,Y} := \frac{K_{\nu_\ell}(\tau_2)}{I_{\nu_\ell}(\tau_2)}.
\end{equation*}


Suppose first $\ell\ge 1$. With the notation
\begin{equation*}
\gamma_\ell:= -\mathfrak b+\mathfrak a\nu_\ell = -\mathfrak b+\sqrt{\mathfrak b^2+\lambda_\ell},\quad \delta_\ell := \mathfrak b+\mathfrak a\nu_\ell = \mathfrak b+\sqrt{\mathfrak b^2+\lambda_\ell},
\end{equation*}
both $\gamma_\ell$ and $\delta_\ell$ are of size $1+\ell$. Since $t^{-\nu}I_\nu(t)$ is increasing on $(0,\infty)$, for
$0<\sigma\leq\rho$,
\begin{equation}\label{eq_phi_ratio}
    \frac{\phi_{\ell,Y}(\sigma)}{\phi_{\ell,Y}(\rho)}
    =
    \Bigl(\frac{\sigma}{\rho}\Bigr)^{\!-\mathfrak b}
    \frac{I_{\nu_\ell}(\tau_\sigma)}{I_{\nu_\ell}(\tau_\rho)}
    \le
    \Bigl(\frac{\sigma}{\rho}\Bigr)^{\!\gamma_\ell}.
\end{equation}
On the other hand, for $0<\rho\le 2$, there exists $0<c<1$ such that
\begin{equation}\label{eq_psi_comparable_K}
    c\rho^{-\mathfrak b}K_{\nu_\ell}(\tau_\rho)
    \le
    \psi_{\ell,Y}(\rho)
    \le
    \rho^{-\mathfrak b}K_{\nu_\ell}(\tau_\rho),\qquad \ell \ge 0.
\end{equation}
Indeed, as $t^\nu K_\nu(t)$ is decreasing on $(0,\infty)$, the definition of $\eta_{\ell,Y}$ gives
\begin{equation*}
\frac{\eta_{\ell,Y}I_{\nu_\ell}(\tau_\rho)}{K_{\nu_\ell}(\tau_\rho)} =
\frac{I_{\nu_\ell}(\tau_\rho)}{I_{\nu_\ell}(\tau_2)}
\frac{K_{\nu_\ell}(\tau_2)}{K_{\nu_\ell}(\tau_\rho)}\le
\Bigl(\frac{\rho}{2}\Bigr)^{\!2\mathfrak a\nu_\ell}\le
4^{-\mathfrak a\nu_0} < 1.
\end{equation*}
% \begin{equation*}
% \frac{\eta_{\ell,Y}I_{\nu_\ell}(\tau_\rho)}{K_{\nu_\ell}(\tau_\rho)}\le
% \Bigl(\frac{\rho}{2}\Bigr)^{\!2\mathfrak a\nu_\ell}\le
% 4^{-\mathfrak a\nu_0} < 1.
% \end{equation*}
This establishes \eqref{eq_psi_comparable_K}. Consequently, if $\rho\le \sigma \le 2$, then
\begin{equation}\label{eq_psi_ratio}
    \frac{\psi_{\ell,Y}(\sigma)}{\psi_{\ell,Y}(\rho)}
    \leq
    C
    \Bigl(\frac{\rho}{\sigma}\Bigr)^{\!\delta_\ell}.
\end{equation}
By Lemma~\ref{lem_bessel_inputs_high_frequency} and $\nu_\ell \ge \nu_0 >0$,
\begin{equation}\label{eq_IK_product_bound}
I_{\nu_\ell} (t) K_{\nu_\ell} (t) \le \frac{C}{1+\nu_\ell}, \qquad \ell \ge 0, \qquad t>0.
\end{equation}
Combining \eqref{eq_psi_comparable_K} and \eqref{eq_IK_product_bound},
\begin{equation}\label{eq_phi_psi_product_bound}
    \phi_{\ell,Y}(\rho)\psi_{\ell,Y}(\rho)
    \le
    C\rho^{-2\mathfrak b}(1+\ell)^{-1},
    \qquad
    0<\rho\leq1,
    \qquad
    \ell\ge 1.
\end{equation}
Fix $0<\rho\le 1$. Split
\begin{equation*}
\begin{aligned}
    \int_0^2 G_{\ell,Y}(\rho,\sigma) \, \sigma^{n-1} \, d\sigma
    &=
    \mathfrak a^{-1} \psi_{\ell,Y}(\rho)
    \int_0^\rho
        \phi_{\ell,Y}(\sigma)\,\sigma^{n-1}\,d\sigma          \\
    &\quad+
    \mathfrak a^{-1} \phi_{\ell,Y}(\rho)
    \int_\rho^2
        \psi_{\ell,Y}(\sigma)\,\sigma^{n-1}\,d\sigma          \\
    &=: I_1+I_2.
\end{aligned}
\end{equation*}
For $I_1$, \eqref{eq_phi_ratio} together with \eqref{eq_phi_psi_product_bound} implies
\begin{align}\label{eq_I1_bound_schur}
    I_1 &\le \mathfrak a^{-1} \phi_{\ell,Y}(\rho)\psi_{\ell,Y}(\rho) \int_0^\rho \Bigl(\frac{\sigma}{\rho}\Bigr)^{\!\gamma_\ell} \sigma^{n-1} \, d\sigma\\ \nonumber
    &= \frac{\mathfrak a^{-1}}{n+\gamma_\ell} \, \rho^n \phi_{\ell,Y}(\rho)\psi_{\ell,Y}(\rho)\\ \nonumber
    &\le C (1+\ell)^{-2}\rho^{n-2\mathfrak b}\\ \nonumber
    &\le C (1+\ell)^{-2}\rho^{1-\alpha}.
\end{align}
For $I_2$, using \eqref{eq_psi_ratio} and \eqref{eq_phi_psi_product_bound},
\begin{align*}
    I_2 &\le C\phi_{\ell,Y}(\rho)\psi_{\ell,Y}(\rho)
    \int_\rho^2 \Bigl(\frac{\rho}{\sigma}\Bigr)^{\!\delta_\ell}\sigma^{n-1} \, d\sigma \\
    &\le C (1+\ell)^{-1} \rho^{\delta_\ell - 2 \mathfrak b} \int_\rho^2 \sigma^{n-\delta_\ell - 1} \, d\sigma \\
    &\le C \begin{cases}
        (1+\ell)^{-2} \rho^{n-2\mathfrak b}, & \delta_\ell >n,\\
        (1+\ell)^{-1} \rho^{n-2\mathfrak b} |\log{\rho}|, & \delta_\ell = n,\\
        (1+\ell)^{-1} |n-\delta_\ell|^{-1} \rho^{\delta_\ell-n}, &\delta_\ell<n.
    \end{cases}
\end{align*}
For $\delta_\ell \ge n$, it is clear that $I_2 \le C (1+\ell)^{-2} \rho^{1-\alpha}$ since $n-2\mathfrak b=2-2\alpha$ and $\rho\le 1$ (the logarithmic blow-up $|\log{\rho}|$ can be bounded by $\rho^{-\varepsilon}$, with $0<\varepsilon<1-\alpha$). While if $\delta_\ell <n$, first note that for $n\ge 2$, 
\begin{equation}\label{eq_delta_lower}
    \delta_\ell \ge \delta_1\geq n+\alpha-1.
\end{equation}
It follows that $\rho^{\delta_\ell-n} \le \rho^{\delta_1-n} \le \rho^{n+\alpha-1-2\mathfrak b} = \rho^{1-\alpha}$. The coefficients in $\ell$ are harmless as there are only finite many $\ell\ge 1$ such that $\delta_\ell < n$. In summary, 
\begin{equation*}
    I_2 \le C (1+\ell)^{-2} \rho^{1-\alpha},\qquad 0<\rho \le 1.
\end{equation*}
Combining this with \eqref{eq_I1_bound_schur} proves \eqref{eq_basic_R_schur_improved} and \eqref{eq_basic_R_schur} with $\ell \ge 1$.


% To continue, note first that, since $n\ge 2$,
% \begin{equation}\label{eq_delta_lower}
%     \delta_1\geq n+\alpha-1.
% \end{equation}
% A similar computation shows that $\delta_\ell \ge \delta_2 > n$ for $\ell \ge 2$. Hence
% \begin{align}\label{eq_sigma_l>=2}
%     \int_\rho^2 \sigma^{n-\delta_\ell - 1} \, d\sigma \le \frac{\rho^{n-\delta_\ell}}{\delta_\ell - n} \le C(1+\ell)^{-1} \rho^{n-\delta_\ell},\qquad \ell \ge 2.
% \end{align}
% On the other hand, if $\ell = 1$, then
% \begin{align}\label{eq_l=1}
%     \int_\rho^2 \sigma^{n-\delta_1 - 1} \, d\sigma \le \begin{cases}
%         C, & \delta_1\le n,\\[4pt]
%         C \rho^{n-\delta_1}, & \delta_1>n.
%     \end{cases}
% \end{align}
% Therefore, if $\ell \ge 2$ or $\ell=1$ with $\delta_1 >n$, then
% \begin{equation*}
%     I_2 \le C (1+\ell)^{-2} \rho^{n-2\mathfrak b} \le C (1+\ell)^{-2} \rho^{1-\alpha},
% \end{equation*}
% whereas if $\ell =1$ and $\delta_1 \le n$, then
% \begin{equation*}
%     I_2\le C (1+\ell)^{-1} \rho^{\delta_1-2\mathfrak b} \le C \rho^{n+\alpha-1 - 2\mathfrak b} = C \rho^{1-\alpha}.
% \end{equation*}
% Together with \eqref{eq_I1_bound_schur}, this proves \eqref{eq_basic_R_schur_improved}, and \eqref{eq_basic_R_schur} for $\ell \ge 1$ follows at once.

It remains to treat the case $\ell=0$. Observe that \eqref{eq_psi_comparable_K} and \eqref{eq_IK_product_bound} remain valid. From the explicit formula \eqref{eq_kernel_Green_copy} and the monotonicity of $t^{-\nu}I_{\nu}(t)$ and $t^\nu K_\nu(t)$,
\begin{equation*}
    G_{0,Y}(\rho,\sigma) \le C \begin{cases}
        \rho^{\gamma_0} \sigma^{-\delta_0}, & 0<\rho \le \sigma <2,\\[4pt]
        \rho^{-\delta_0} \sigma^{\gamma_0}, & 0<\sigma <\rho<2.
    \end{cases}
\end{equation*}
Therefore, since $\gamma_0 = 0$ and $\delta_0 = n-2(1-\alpha) <n$, for all $0<\rho \le 1$,
\begin{align*}
    \int_0^2 G_{0,Y}(\rho,\sigma)\,\sigma^{n-1}\,d\sigma
    &\le C \rho^{-\delta_0} \int_0^\rho \sigma^{n-1}\,d\sigma\\
    &\quad+ C \int_\rho^2 \sigma^{n-\delta_0 -1}\,d\sigma \\
    &\le C.
\end{align*}
This proves \eqref{eq_basic_R_schur_ell=0} and \eqref{eq_basic_R_schur} with $\ell=0$.




\textit{Case 2: $|Y|> C^*(1+\ell)$.} Set for simplicity $k_{\ell,Y}(\rho):=\rho^{-\mathfrak b}K_{\nu_\ell}(\tau_\rho)$. 
% \begin{equation*}
%     k_{\ell,Y}(\rho):=\rho^{-\mathfrak b}K_{\nu_\ell}(\tau_\rho),\quad 0<\rho <2.
% \end{equation*}
By \eqref{eq_psi_comparable_K}, it is enough to estimate $G_{\ell,Y}(\rho,\sigma)$ with $\psi_{\ell,Y}$ replaced by $k_{\ell,Y}$. Fix $0<\rho<2$. From \eqref{eq_kernel_Green_copy},
\begin{align}\label{eq_weighted_split}
    \rho^{2\beta}
    \int_0^2
        G_{\ell,Y}(\rho,\sigma)
    \, \sigma^{n-1} \, d\sigma                                      \nonumber
    \le
    C\rho^{2\beta}
    \phi_{\ell,Y}(\rho)k_{\ell,Y}(\rho)
    \biggl[\int_0^\rho \frac{\phi_{\ell,Y}(\sigma)}{\phi_{\ell,Y}(\rho)} \, \sigma^{n-1} \, d\sigma + \int_\rho^2 \frac{k_{\ell,Y}(\sigma)}{k_{\ell,Y}(\rho)} \, \sigma^{n-1} \, d\sigma \biggr].
\end{align}
By Lemma~\ref{lem_bessel_inputs_high_frequency}, for $\nu_\ell\ge \nu_0>0$,
\begin{equation}\label{eq_phi_k_product}
    \phi_{\ell,Y}(\rho)k_{\ell,Y}(\rho)
    =
    \rho^{-2\mathfrak b}
    I_{\nu_\ell}(\tau_\rho)K_{\nu_\ell}(\tau_\rho)
    \le
    C\rho^{-2\mathfrak b}(1 + \ell + \tau_\rho)^{-1}.
\end{equation}
Put $r=\sigma/\rho\in(0,1]$. Then
\begin{equation*}
\frac{\phi_{\ell,Y}(\sigma)}{\phi_{\ell,Y}(\rho)} = r^{\gamma_\ell} \frac{\mathcal I_{\nu_\ell}(\tau_\rho r^{\mathfrak a})}{\mathcal I_{\nu_\ell}(\tau_\rho)}.
\end{equation*}
Assume first that $\tau_\rho \le  C_0(1+\ell)$ with $C_0>0$ to be determined later. Since $\mathcal I_\nu$ is increasing,
\begin{align*}
    \int_0^\rho
        \frac{\phi_{\ell,Y}(\sigma)}
             {\phi_{\ell,Y}(\rho)}
    \, \sigma^{n-1} \, d\sigma
    &\le C\rho^n \int_0^1 r^{n-1+\gamma_\ell}\,dr
    \le C\frac{\rho^n}{1+\ell}
    \le C\frac{\rho^n}{1+ \ell + \tau_\rho}.
\end{align*}
Next assume that $\tau_\rho>C_0(1+\ell)$. Choose $C_0>1$ sufficiently large so that $2^{- \mathfrak a} C_0 (1+\ell) \ge C(1+\ell)$, where $C>0$ is the constant from \eqref{eq_S_lower_large}. Lemma~\ref{lem_bessel_inputs_high_frequency} \eqref{eq_S_lower_large} then gives for $0<r<1/2$ (thus $r^{\mathfrak a}\le 2^{- \mathfrak a}$),
\begin{align*}
\frac{\mathcal I_{\nu_\ell}(\tau_\rho r^{\mathfrak a})}{\mathcal I_{\nu_\ell}(\tau_\rho)}
&=
\exp\Bigl\{-\int_{\tau_\rho r^{\mathfrak a}}^{\tau_\rho}\frac{I_{\nu_\ell+1}(t)}{I_{\nu_\ell}(t)}\,dt\Bigr\} \\
&\le
\exp\Bigl\{-\int_{\tau_\rho 2^{-\mathfrak a}}^{\tau_\rho}\frac{I_{\nu_\ell+1}(t)}{I_{\nu_\ell}(t)}\,dt\Bigr\}
\le C e^{-c \tau_\rho},
\end{align*}
and for $1/2\le r \le 1$ (thus $\tau_\rho r^{\mathfrak a} \ge 2^{- \mathfrak a} C_0 (1+\ell) \ge C(1+\ell)$),
\begin{equation*}
\frac{\mathcal I_{\nu_\ell}(\tau_\rho r^{\mathfrak a})}{\mathcal I_{\nu_\ell}(\tau_\rho)}
=\exp\Bigl\{-\int_{\tau_\rho r^{\mathfrak a}}^{\tau_\rho}\frac{I_{\nu_\ell+1}(t)}{I_{\nu_\ell}(t)}\,dt\Bigr\}
\le C e^{-c \tau_\rho(1-r^{\mathfrak a})}
\le C e^{-c \tau_\rho(1-r)}.
\end{equation*}
Hence
\begin{align*}
   \int_0^\rho \frac{\phi_{\ell,Y}(\sigma)}{\phi_{\ell,Y}(\rho)}\,\sigma^{n-1}\,d\sigma
   &\le
    C\rho^n \int_0^{1/2} r^{n-1+\gamma_\ell} e^{-c\tau_\rho}\,dr\\
    &\quad+ C\rho^n \int_{1/2}^1 r^{n-1+\gamma_\ell}e^{-c\tau_\rho(1-r)}\,dr\\
    &\le C\frac{\rho^n}{1+\ell+\tau_\rho}.
\end{align*}
Thus, in all cases,
\begin{equation}\label{eq_I_ratio_integral_final}
    \int_0^\rho
        \frac{\phi_{\ell,Y}(\sigma)}
             {\phi_{\ell,Y}(\rho)}
    \, \sigma^{n-1} \, d\sigma
    \le
    C\frac{\rho^n}{1+\ell+\tau_\rho}.
\end{equation}

For the $K$-branch, note that now $r=\sigma/\rho \ge 1$. Then
\begin{equation}\label{eq_K_ratio}
    \frac{k_{\ell,Y}(\sigma)}
         {k_{\ell,Y}(\rho)}
    =
    r^{-\delta_\ell}
    \frac{\mathcal K_{\nu_\ell}(\tau_\rho r^{\mathfrak a})}
         {\mathcal K_{\nu_\ell}(\tau_\rho)},
    \qquad
    \mathcal K_\nu(t):=t^\nu K_\nu(t),
\end{equation}
where
\begin{equation*}
    \delta_\ell=\mathfrak b+\mathfrak a\nu_\ell =
    \mathfrak b+\sqrt{\mathfrak b^2+\lambda_\ell}
    \sim 1+\ell.
\end{equation*}
Moreover, for $\ell\ge 1$, \eqref{eq_delta_lower} implies $\delta_\ell\ge n+\alpha-1$.

\textit{Case 2a: $\delta_\ell > n$.} If $\tau_\rho\le C_0(1+\ell)$, then $\mathcal K_\nu$ is decreasing and hence
\begin{align}\label{eq_formula3}
\int_\rho^2 \frac{k_{\ell,Y}(\sigma)}{k_{\ell,Y}(\rho)}\,\sigma^{n-1}\,d\sigma
&\le \rho^n
    \int_1^{2/\rho} r^{n-1-\delta_\ell}\,dr
    \le C\frac{\rho^n}{1+\ell}
    \le C\frac{\rho^n}{1+ \ell + \tau_\rho}.
\end{align}
If instead $\tau_\rho>C_0 (1+\ell)$, where $C_0$ is again chosen large enough so that Lemma~\ref{lem_bessel_inputs_high_frequency} applies, then \eqref{eq_K_normalized_exp_decay} yields
\begin{equation*}
    \frac{\mathcal K_{\nu_\ell}(\tau_\rho r^{\mathfrak a})}
         {\mathcal K_{\nu_\ell}(\tau_\rho)}
    \le
    C e^{-c\tau_\rho(r^{\mathfrak a}-1)}
    \le C e^{-c\tau_\rho(r-1)}
\end{equation*}
since $r\ge 1$. Therefore,
\begin{equation}\label{eq_K_ratio_large_ell}
    \int_\rho^2\frac{k_{\ell,Y}(\sigma)}{k_{\ell,Y}(\rho)}\,\sigma^{n-1}\,d\sigma
    \le C\rho^n  \int_1^\infty e^{-c\tau_\rho(r-1)}\,dr
    \le C\frac{\rho^n}{1+\ell+\tau_\rho}.
\end{equation}
Combining this with \eqref{eq_phi_k_product},
\begin{align*}
    \rho^{2\beta}
    \int_0^2
        G_{\ell,Y}(\rho,\sigma)
    \, \sigma^{n-1} \, d\sigma
    \le \frac{C \rho^{n-2 \mathfrak b+2\beta} }{(1+\ell + \tau_\rho)^2}
    = \frac{C \rho^{2 \mathfrak a} }{(1+\ell + \tau_\rho)^2}
    \le C |Y|^{-2},
\end{align*}
as $\rho^{2 \mathfrak a} = \mathfrak a^2 |Y|^{-2} \tau_\rho^2$. This proves \eqref{eq_basic_R_schur} with $\delta_\ell >n$.

\textit{Case 2b: $\delta_\ell \le n$.} On the region $\sigma<\rho$, the change of variables $u=\tau_\sigma$ gives
\begin{align}\label{eq_finite_left_scaling}
\rho^{2\beta}
    k_{\ell,Y}(\rho)
    \int_0^\rho
        \phi_{\ell,Y}(\sigma) \, \sigma^{n-1} \, d\sigma
    =
    C |Y|^{-2}
    \tau_\rho^{\frac{2\beta-\mathfrak b}{\mathfrak a}}
    K_{\nu_\ell}(\tau_\rho)
    \int_0^{\tau_\rho}
        u^{\frac{n-\mathfrak b}{\mathfrak a}-1}
        I_{\nu_\ell}(u) \, du.
\end{align}
Since only finitely many values of $\ell$ are involved here, the standard small and large argument asymptotic formulas for the modified Bessel functions suffice, and there is no need to keep track of the $\ell$-dependence. We show that the function of $\tau_\rho$ on the right-hand side of \eqref{eq_finite_left_scaling} is uniformly bounded on $(0,\infty)$. By Lemma~\ref{lem:Bessel:integral} $(ii)$, this function is nothing but $II(q,s,\nu_\ell,\tau_\rho)$ with $q=\frac{2\beta-\mathfrak b}{\mathfrak a}$, $s=\frac{n-\mathfrak b}{\mathfrak a}$ and $\nu_\ell \ge \nu_0>0$. It is easy to check 
\begin{align}\label{check}
    q+s+\nu-|\nu| = q+s = \frac{2\beta - 2\mathfrak b+n}{\mathfrak a} = 2 \le 2,
\end{align}
inferring the uniform boundedness of $II(q,s,\nu_\ell,\tau_\rho)$ for all $\tau_\rho >0$, and hence \eqref{eq_finite_left_scaling} is $O(|Y|^{-2})$ as required.

On the other hand,
% Near $\tau_\rho=0$,
% \begin{equation*}
%     \tau_\rho^{\frac{2\beta-\mathfrak b}{\mathfrak a}}
%     K_{\nu_\ell}(\tau_\rho)
%     \int_0^{\tau_\rho}
%         u^{\frac{n-\mathfrak b}{\mathfrak a}-1}
%         I_{\nu_\ell}(u) \, du
%     \sim \tau_\rho^{\frac{2\beta-\mathfrak b}{\mathfrak a} - \nu_\ell} \int_0^{\tau_\rho} u^{\frac{n-\mathfrak b}{\mathfrak a}-1+\nu_\ell}\,du
%     \sim \tau_\rho^2 \le 1.
% \end{equation*}
% As $\tau_\rho \to \infty$, the contribution coming from $\int_0^1 u^{\frac{n-\mathfrak b}{\mathfrak a}-1} I_{\nu_\ell}(u) \, du$
% % \begin{equation*}
% %     \int_0^1 u^{\frac{n-\mathfrak b}{\mathfrak a}-1} I_{\nu_\ell}(u) \, du
% % \end{equation*}
% is uniformly bounded because of the exponential decay of $K_{\nu_\ell}(\tau_\rho)$. For the remaining part,
% \begin{align}\label{eq_formula1}
%     \tau_\rho^{\frac{2\beta-\mathfrak b}{\mathfrak a}}
%     K_{\nu_\ell}(\tau_\rho)
%     \int_1^{\tau_\rho}
%         u^{\frac{n-\mathfrak b}{\mathfrak a}-1}
%         I_{\nu_\ell}(u) \, du
%     \sim \tau_\rho^{\frac{2\beta-\mathfrak b}{\mathfrak a}-\frac{1}{2}} e^{-\tau_\rho} \int_1^{\tau_\rho} u^{\frac{n-\mathfrak b}{\mathfrak a}-\frac{3}{2}} e^{u} \, du.
% \end{align}
% If $\frac{n-\mathfrak b}{\mathfrak a}\ge \frac{3}{2}$, then \eqref{eq_formula1} is $O\bigl( \tau_\rho^{\frac{2\beta-\mathfrak b}{\mathfrak a} -\frac{1}{2}+\frac{n-\mathfrak b}{\mathfrak a} - \frac{3}{2}} \bigr) = O(1)$.
% % \begin{equation*}
% %     O\bigl( \tau_\rho^{\frac{2\beta-\mathfrak b}{\mathfrak a} -\frac{1}{2}+\frac{n-\mathfrak b}{\mathfrak a} - \frac{3}{2}} \bigr) = O(1).
% % \end{equation*}
% While if $\frac{n-\mathfrak b}{\mathfrak a} < \frac{3}{2}$, then
% \begin{align*}
% \int_1^{\tau_\rho} u^{\frac{n-\mathfrak b}{\mathfrak a}-\frac{3}{2}} e^{-(\tau_\rho-u)}\,du
% &\le e^{-\tau_\rho/2} \int_1^{\tau_\rho/2} u^{\frac{n-\mathfrak b}{\mathfrak a}-\frac{3}{2}}\,du\\
% &\quad+ C \tau_\rho^{\frac{n-\mathfrak b}{\mathfrak a}-\frac{3}{2}} \int_{\tau_\rho/2}^{\tau_\rho} e^{-(\tau_\rho - u)}\,du\\
% &\le C \tau_\rho^{\frac{n-\mathfrak b}{\mathfrak a}-\frac{3}{2}},
% \end{align*}
% and \eqref{eq_formula1} is again $O(1)$. Hence, for $\delta_\ell \le n$, \eqref{eq_finite_left_scaling} is bounded by $C |Y|^{-2}$.
consider the region $\sigma \ge \rho$. After the same change of variables,
\begin{align*}
\rho^{2\beta}
    \phi_{\ell,Y}(\rho)
    \int_\rho^2
        k_{\ell,Y}(\sigma)\,\sigma^{n-1}\,d\sigma
    =
    C |Y|^{-2}
    \tau_\rho^{\frac{2\beta-\mathfrak b}{\mathfrak a}}
    I_{\nu_\ell}(\tau_\rho)
    \int_{\tau_\rho}^\infty
        u^{\frac{n-\mathfrak b}{\mathfrak a}-1}
        K_{\nu_\ell}(u) \, du = I(q,s,\nu_\ell
        ,0,\tau_\rho),
\end{align*}
with the same $q,s$ as in the case $\sigma>\rho$. With \eqref{check} and 
\begin{align*}
q+\nu_\ell = \frac{2\beta-\mathfrak b + \sqrt{\mathfrak b^2 + \lambda_\ell}}{\mathfrak a} \ge \frac{2\beta}{\mathfrak a}\ge 0,
\end{align*}
Lemma~\ref{lem:Bessel:integral} gives $I(\tau_\rho) \le C$ for all $\tau_\rho >0$. This closes the proof of \eqref{eq_basic_R_schur} for $\delta_\ell\le n$.

% bounded on $(0,\infty)$. Near $\tau_\rho=0$,
% \begin{align*}
%     \tau_\rho^{\frac{2\beta-\mathfrak b}{\mathfrak a}}
%     I_{\nu_\ell}(\tau_\rho)
%     \int_{\tau_\rho}^\infty
%         u^{\frac{n-\mathfrak b}{\mathfrak a}-1}
%         K_{\nu_\ell}(u) \, du
%     &=
%     \tau_\rho^{\frac{2\beta-\mathfrak b}{\mathfrak a}}
%     I_{\nu_\ell}(\tau_\rho)
%     \int_{\tau_\rho}^1
%         u^{\frac{n-\mathfrak b}{\mathfrak a}-1}
%         K_{\nu_\ell}(u) \, du\\
%     &\quad+ \tau_\rho^{\frac{2\beta-\mathfrak b}{\mathfrak a}}
%     I_{\nu_\ell}(\tau_\rho)
%     \int_{1}^\infty
%         u^{\frac{n-\mathfrak b}{\mathfrak a}-1}
%         K_{\nu_\ell}(u) \, du\\
%     &\le C \tau_\rho^{\frac{2\beta-\mathfrak b}{\mathfrak a}+\nu_\ell} \int_{\tau_\rho}^1 u^{\frac{n-\mathfrak b}{\mathfrak a}-\nu_\ell-1}\,du \\
%     &\quad+ C \tau_\rho^{\frac{2\beta-\mathfrak b}{\mathfrak a}+\nu_\ell} \int_1^\infty u^{\frac{n-\mathfrak b}{\mathfrak a}-\frac{3}{2}} e^{-u}\,du.
% \end{align*}
% The last term is plainly $O\bigl( \tau_\rho^{\frac{2\beta-\mathfrak b}{\mathfrak a}+\nu_\ell} \bigr) = O(1)$. For the preceding term, if $\frac{n-\mathfrak b}{\mathfrak a}-\nu_\ell <0$, then it is $O\bigl( \tau_\rho^{\frac{2\beta-\mathfrak b}{\mathfrak a}+\nu_\ell + \frac{n-\mathfrak b}{\mathfrak a}-\nu_\ell} \bigr) = O(1)$. If instead $\frac{n-\mathfrak b}{\mathfrak a}-\nu_\ell \ge 0$, then it is bounded by $O\bigl( \tau_\rho^{\frac{2\beta-\mathfrak b}{\mathfrak a}+\nu_\ell} |\log \tau_\rho| \bigr) = O(1)$, as $\tau_\rho \to 0$. Now suppose $\tau_\rho \to \infty$. Then
% \begin{align}\label{eq_formula2}
%     \tau_\rho^{\frac{2\beta-\mathfrak b}{\mathfrak a}}
%     I_{\nu_\ell}(\tau_\rho)
%     \int_{\tau_\rho}^\infty
%         u^{\frac{n-\mathfrak b}{\mathfrak a}-1}
%         K_{\nu_\ell}(u) \, du
%     \le C \tau_\rho^{\frac{2\beta-\mathfrak b}{\mathfrak a}-\frac{1}{2}} e^{\tau_\rho} \int_{\tau_\rho}^\infty u^{\frac{n-\mathfrak b}{\mathfrak a} - \frac{3}{2}} e^{-u}\,du.
% \end{align}
% If $\frac{n-\mathfrak b}{\mathfrak a} - \frac{3}{2} \le 0$, then the right-hand side is bounded by
% \begin{equation*}
%     C \tau_\rho^{\frac{2\beta-\mathfrak b}{\mathfrak a}-\frac{1}{2} + \frac{n-\mathfrak b}{\mathfrak a} - \frac{3}{2} } e^{\tau_\rho} \int_{\tau_\rho}^\infty e^{-u}\,du \le C.
% \end{equation*}
% If $\frac{n-\mathfrak b}{\mathfrak a} - \frac{3}{2} > 0$, then using $\tau_\rho + u \le \tau_\rho(1+u)$, \eqref{eq_formula2} is bounded by
% \begin{align*}
%     C \tau_\rho^{\frac{2\beta-\mathfrak b}{\mathfrak a}-\frac{1}{2}}\int_{\tau_\rho}^\infty u^{\frac{n-\mathfrak b}{\mathfrak a} - \frac{3}{2}} e^{-(u-\tau_\rho)}\,du
%     &= C \tau_\rho^{\frac{2\beta-\mathfrak b}{\mathfrak a}-\frac{1}{2}}\int_{0}^\infty (\tau_\rho + u)^{\frac{n-\mathfrak b}{\mathfrak a} - \frac{3}{2}} e^{-u}\,du \\
%     &\le C \int_0^\infty (1+u)^{\frac{n-\mathfrak b}{\mathfrak a} - \frac{3}{2}} e^{-u}\,du \le C.
% \end{align*}
% Thus \eqref{eq_basic_R_schur} also holds for $\delta_\ell \le n$.

To complete the argument, it remains to derive a pointwise estimate for the $K$-branch when $\delta_\ell \le n$. Clearly, \eqref{eq_K_ratio_large_ell} remains valid when $\tau_\rho > C_0(1+\ell)$. For the case $\tau_\rho \le C_0(1+\ell)$, however, \eqref{eq_delta_lower} only gives
\begin{align*}
    \int_\rho^2 \frac{k_{\ell,Y}(\sigma)}{k_{\ell,Y}(\rho)} \, \sigma^{n-1} \, d\sigma
    &\le C \rho^n \int_1^{2/\rho} r^{n-\delta_\ell-1}\,dr\\
    &\le C \begin{cases}
        \rho^n |\log{\rho}|, & \delta_\ell = n,\\
        |n-\delta_\ell|^{-1} \rho^{\delta_\ell}, & \delta_\ell < n.
    \end{cases}
\end{align*}
For $\delta_\ell=n$, it is plain that the integral is bounded by $C(1+\ell+\tau_\rho)^{-1} \rho^{n+\alpha-1}$, where one can bound the logarithmic term by $C_\alpha \rho^{-(1-\alpha)}$. As for $\delta_\ell<n$, it satisfies upper bound
\begin{align}\label{equation2}
    \begin{cases}
        (1+\ell+\tau_\rho)^{-1} \rho^{n+\alpha-1}, & \ell \ge 1,\\
        (1+\tau_\rho)^{-1} \rho^{2\mathfrak b}, & \ell=0.
    \end{cases}
\end{align}
% \begin{align*}
%     \int_\rho^2 \frac{k_{\ell,Y}(\sigma)}{k_{\ell,Y}(\rho)} \, \sigma^{n-1} \, d\sigma
%     &\le C \rho^n \int_1^{2/\rho} r^{n-\delta_\ell-1}\,dr\\
%     &\le C \rho^{n} \int_1^{2/\rho} r^{1-\alpha-1}\,dr\\
%     &\le C \rho^{n+\alpha-1}
%     \le C \frac{\rho^{n+\alpha-1}}{1+\ell+\tau_\rho}
% \end{align*}
% for $\ell = 1$, and it is $O((1+\tau_\rho)^{-1}\rho^{\delta_0}) = O((1+\tau_\rho)^{-1} \rho^{2\mathfrak b})$ for $\ell = 0$, since $\delta_0=2\mathfrak b = n+2\alpha-2<n$.
Concluded from \eqref{eq_phi_k_product}, \eqref{eq_I_ratio_integral_final}, \eqref{eq_K_ratio_large_ell} and \eqref{equation2}, for all $0<\rho \le 1$,
\begin{align*}
    \int_0^2
        G_{\ell,Y}(\rho,\sigma)
    \, \sigma^{n-1} \, d\sigma
    &\le C \rho^{-2 \mathfrak b} (1+\ell+\tau_\rho)^{-1} \begin{cases}
        (1+\ell+\tau_\rho)^{-1} \rho^{n+\alpha-1}, & \ell \ge 1,\\
        (1+\tau_\rho)^{-1} \rho^{2\mathfrak b}, & \ell=0.
    \end{cases} \\
    &\le C \begin{cases}
        (1+\ell+\tau_\rho)^{-2} \rho^{1-\alpha}, & \ell \ge 1,\\
        (1+\tau_\rho)^{-2}, & \ell=0.
    \end{cases} 
\end{align*}
% for $\ell = 1$, while for $\ell=0$ it is $(1+\tau_\rho)^{-1} O(\rho^{-2\mathfrak b} \rho^{2 \mathfrak b}) = (1+\tau_\rho)^{-1} O(1)$. 
This completes the proof of \eqref{eq_basic_R_schur_improved}, \eqref{eq_basic_R_schur_ell=0}, and hence Proposition~\ref{prop:Green:Shur}.

\end{proof}








Define operators
\begin{equation*}
    D_r:=\rho^\alpha\partial_\rho,
    \qquad
    D_\theta:=\rho^{\alpha-1}\sqrt{\lambda_\ell},
    \qquad
    D_{y,j}:=\rho^\beta Y_j.
\end{equation*}
We also have the following first order derivative estimates for the Green kernel.








\begin{proposition}\label{lem_basic_D_R_schur_explicit}
For $\ell \ge 0$,
\begin{align}\label{eq_basic_D_R_schur}
    \sup_{0<\rho\leq1} \int_0^2 |D_* G_{\ell,Y}(\rho,\sigma)|\,\sigma^{n-1}\,d\sigma \le \begin{cases}
        C, & |Y|\le C^*(1+\ell),\\[4pt]
        C|Y|, & |Y|>C^*(1+\ell),
    \end{cases}
\end{align}
where $* \in \{r, \theta, y,j\}$.
\end{proposition}

\begin{proof}
Suppose first $\ell\ge 1$. We claim that for $0<\rho\leq1$,
\begin{equation}\label{eq_log_derivative_phi_psi}
    |\rho^\alpha\phi_{\ell,Y}'(\rho)|
    \leq
    C(1+\ell+\tau_\rho) \rho^{\alpha-1}\phi_{\ell,Y}(\rho),
\end{equation}
and
\begin{equation}\label{eq_log_derivative_psi}
    |\rho^\alpha\psi_{\ell,Y}'(\rho)|
    \leq
    C(1+\ell+\tau_\rho) \rho^{\alpha-1}\psi_{\ell,Y}(\rho).
\end{equation}
To verify \eqref{eq_log_derivative_phi_psi}, use the identity $tI_\nu'(t)-\nu I_\nu(t)=tI_{\nu+1}(t)$,
% \begin{equation*}
%     tI_\nu'(t)-\nu I_\nu(t)=tI_{\nu+1}(t),
% \end{equation*}
which gives
\begin{align*}
\rho^\alpha\frac{\phi_{\ell,Y}'(\rho)}{\phi_{\ell,Y}(\rho)}
&=\rho^{\alpha-1}\Bigl[-\mathfrak b+\mathfrak a \tau_\rho \frac{I_{\nu_\ell}'(\tau_\rho)}{I_{\nu_\ell}(\tau_\rho)}\Bigr] =\rho^{\alpha-1}\Bigl[-\mathfrak b+\mathfrak a\nu_\ell+\mathfrak a \tau_\rho \frac{I_{\nu_\ell+1}(\tau_\rho)}{I_{\nu_\ell}(\tau_\rho)}\Bigr].
\end{align*}
Since $-\mathfrak b+\mathfrak a\nu_\ell =\gamma_\ell \sim 1+\ell$, and $0<\frac{I_{\nu+1}(t)}{I_\nu(t)}< 1$,
\begin{equation*}
    \Bigl|\rho^\alpha\frac{\phi_{\ell,Y}'(\rho)}{\phi_{\ell,Y}(\rho)}\Bigr|
    \le
    C\rho^{\alpha-1}\bigl(1+\ell + \tau_\rho\bigr).
\end{equation*}
This proves \eqref{eq_log_derivative_phi_psi}.

To prove \eqref{eq_log_derivative_psi}, write $F_{\ell,Y}(t):=K_{\nu_\ell}(t)-\eta_{\ell,Y}I_{\nu_\ell}(t)$ so that $\psi_{\ell,Y}(\rho)=\rho^{-\mathfrak b}F_{\ell,Y}(\tau_\rho)$.
% \begin{equation*}
%     F_{\ell,Y}(t):=K_{\nu_\ell}(t)-\eta_{\ell,Y}I_{\nu_\ell}(t), \qquad
%     \psi_{\ell,Y}(\rho)=\rho^{-\mathfrak b}F_{\ell,Y}(\tau_\rho).
% \end{equation*}
\eqref{eq_psi_comparable_K} implies that $K_{\nu_\ell}(\tau_\rho)$ and $F_{\ell,Y}(\tau_\rho)$ are comparable, i.e.,
\begin{equation}\label{eq_F_comparable_K}
    cK_{\nu_\ell}(\tau_\rho)
    \le
    F_{\ell,Y}(\tau_\rho)
    \leq
    K_{\nu_\ell}(\tau_\rho),
\end{equation}
with $0<c<1$ independent of $\ell,Y,\rho$. Moreover, by Lemma~\ref{lem_S_derivative_bound}, for all $t>0$ and $\nu \ge \nu_0>0$,
\begin{equation}\label{eq_I_log_derivative_standard}
    \bigl|t\frac{I_\nu'(t)}{I_\nu(t)}\bigr|
    =
    \bigl| \nu + t S_{\nu}(t) \bigr|
    \le
    C(1+\nu+t).
\end{equation}
We also claim that
\begin{equation}\label{eq_K_log_derivative_standard}
    \Bigl|t\frac{K_\nu'(t)}{K_\nu(t)}\Bigr|
    \leq
    C(1+\nu+t),
    \quad
    t>0,\quad \nu\ge \nu_0>0.
\end{equation}
Indeed, since $(t^\nu K_\nu(t))'= -t^\nu K_{\nu-1}(t)$,
% \begin{equation*}
%     (t^\nu K_\nu(t))'= -t^\nu K_{\nu-1}(t),
% \end{equation*}
and $K_{\nu}(t)$ is increasing as a function of $\nu$ on $(0,\infty)$, for each $t>0$ and $\nu> 1$,
\begin{equation*}
\Bigl| t\frac{K_\nu'(t)}{K_\nu(t)} \Bigr| \le \nu + t \frac{K_{\nu-1}(t)}{K_\nu(t)} \le 1+\nu+t.
\end{equation*}
If instead $\nu_0 \le \nu \le 1$, then $1-\nu \le 1-\nu_0$, and therefore
\begin{align*}
    t\frac{K_{\nu-1}(t)}{K_{\nu}(t)}
    &=t \frac{K_{1-\nu}(t)}{K_{\nu}(t)}
    = t\frac{K_{1-\nu}(t)}{K_{1-\nu_0}(t)} \frac{K_{1-\nu_0}(t)}{K_{\nu_0}(t)} \frac{K_{\nu_0}(t)}{K_{\nu}(t)} \\
    &\le t\frac{K_{1-\nu_0}(t)}{K_{\nu_0}(t)}
    \le C_{\nu_0} \begin{cases}
        t^{2\nu_0}, & t\le 1,\\
        t, & t\ge 1.
    \end{cases}
\end{align*}
This proves \eqref{eq_K_log_derivative_standard}.

Combining \eqref{eq_I_log_derivative_standard}, \eqref{eq_K_log_derivative_standard}, and \eqref{eq_F_comparable_K},
\begin{equation*}
    |\tau_\rho F_{\ell,Y}'(\tau_\rho)|
    \le
    C(1+\nu_\ell+\tau_\rho) |F_{\ell,Y}(\tau_\rho)|.
\end{equation*}
Consequently,
\begin{equation}\label{eq:normal2}
\begin{aligned}
    \rho^\alpha
    \Bigl|
        \frac{\psi_{\ell,Y}'(\rho)}{\psi_{\ell,Y}(\rho)}
    \Bigr|
    &\le
    \rho^{\alpha-1}
    \bigl[
        |\mathfrak b|
        +
        \mathfrak a
        \bigl|
            \tau_\rho\frac{F_{\ell,Y}'(\tau_\rho)}{F_{\ell,Y}(\tau_\rho)}
        \bigr|
    \bigr]                                                \\
    &\le
    C\rho^{\alpha-1}(1+\nu_\ell+\tau_\rho)                    \\
    &\le C(1+\ell+\tau_\rho) \rho^{\alpha-1}.
\end{aligned}
\end{equation}
This establishes \eqref{eq_log_derivative_psi}.

We now estimate \eqref{eq_basic_D_R_schur}. By \eqref{eq_kernel_Green}, for
$0<\rho\le 1$,
\begin{align}\label{D:bdd:G}
    \int_0^2
        |D_r G_{\ell,Y}(\rho,\sigma)|
    \, \sigma^{n-1} \, d\sigma
    &\le
    C (1+\ell+\tau_\rho) \rho^{\alpha-1}
    \int_0^2
        G_{\ell,Y}(\rho,\sigma)
    \, \sigma^{n-1} \, d\sigma.
\end{align}
Indeed, when $\sigma\ge \rho$, the $\rho$-derivative falls on $\phi_{\ell,Y}(\rho)$, so \eqref{eq_log_derivative_phi_psi} applies; when $\sigma\le \rho$, the $\rho$-derivative falls on $\psi_{\ell,Y}(\rho)$, so \eqref{eq_log_derivative_psi} applies. It follows by Proposition~\ref{prop:Green:Shur} that if $|Y|\le C^*(1+\ell)$ (thus $\tau_\rho = \mathfrak a^{-1} |Y| \rho^{\mathfrak a} \le C(1+\ell)$), then for $\ell\ge 1$,
% \begin{equation*}
%     \tau_\rho = \mathfrak a^{-1} |Y| \rho^{\mathfrak a} \le C(1+\ell),
%     \qquad \rho\le 1,
% \end{equation*}
\begin{align*}
    \int_0^2 &|D_* G_{\ell,Y}(\rho,\sigma)| \, \sigma^{n-1} \, d\sigma \\
    &\le  \begin{cases}
        C (1+\ell)\rho^{\alpha-1} (1+\ell)^{-2}\rho^{1-\alpha}, & *=r,\\
    C\sqrt{\lambda_\ell} \rho^{\alpha-1} (1+\ell)^{-2}\rho^{1-\alpha}, & *=\theta,\\
    C|Y_j|\rho^\beta (1+\ell)^{-2}\rho^{1-\alpha}, & *=y,j.
    \end{cases}\\
    &\le C.
\end{align*}
% \begin{align*}
%     \int_0^2 |D_r G_{\ell,Y}(\rho,\sigma)|
%     \, \sigma^{n-1} \, d\sigma
%     &\le
%     C(1+\ell)\rho^{\alpha-1}
%     (1+\ell)^{-2}\rho^{1-\alpha}
%     \le C.
% \end{align*}
% Likewise,
% \begin{align*}
%     \int_0^2
%         |D_\theta
%         G_{\ell,Y}(\rho,\sigma)|
%     \, \sigma^{n-1} \, d\sigma
%     &\le
%     C(1+\ell) \rho^{\alpha-1}
%     \int_0^2
%         G_{\ell,Y}(\rho,\sigma)
%     \, \sigma^{n-1} \, d\sigma \\
%     &\le
%     C(1+\ell) \rho^{\alpha-1}
%     (1+\ell)^{-2}\rho^{1-\alpha}
%     \le C.
% \end{align*}
% Finally,
% \begin{align*}
%     \int_0^2
%         |D_{y,j} G_{\ell,Y}(\rho,\sigma)|
%     \, \sigma^{n-1} \, d\sigma
%     &\le
%     C|Y|\rho^\beta
%     \int_0^2
%         G_{\ell,Y}(\rho,\sigma)
%     \, \sigma^{n-1} \, d\sigma \\
%     &\le
%     C (1+\ell)^{-1} \rho^{\mathfrak a}
%     \le C.
% \end{align*}
This proves \eqref{eq_basic_D_R_schur} for $\ell\ge 1$ and $|Y|\le C^*(1+\ell)$.

Suppose next $|Y|>C^*(1+\ell)$. Again by Proposition~\ref{prop:Green:Shur}, for $\ell \ge 1$,
\begin{align*}
    \int_0^2 &|D_* G_{\ell,Y}(\rho,\sigma)| \, \sigma^{n-1} \, d\sigma \\
    &\le  \begin{cases}
        C (1+\ell+\tau_\rho) \rho^{\alpha-1} \frac{\rho^{n+\alpha-1}}{(1+\ell+\tau_\rho)^2}, & *=r,\\
    C (1+\ell) \rho^{\alpha-1} \frac{\rho^{n+\alpha-1}}{(1+\ell+\tau_\rho)^2}, & *=\theta,\\
    C|Y|\rho^\beta
    \frac{\rho^{n+\alpha-1}}{(1+\ell+\tau_\rho)^2}, & *=y,j.
    \end{cases}\\
    &\le C|Y|.
\end{align*}
% \begin{align*}
%     \int_0^2 |D_r G_{\ell,Y}(\rho,\sigma)| \, \sigma^{n-1} \, d\sigma
%     &\le C (1+\ell+\tau_\rho) \rho^{\alpha-1} \frac{\rho^{n+\alpha-1}}{(1+\ell+\tau_\rho)^2}
%     \le C.
% \end{align*}
% Similarly,
% \begin{align*}
%     \int_0^2 |D_\theta G_{\ell,Y}(\rho,\sigma)| \, \sigma^{n-1} \, d\sigma
%     &\le C (1+\ell) \rho^{\alpha-1} \frac{\rho^{n+\alpha-1}}{(1+\ell+\tau_\rho)^2}
%     \le C.
% \end{align*}
% Finally,
% \begin{align*}
%     \int_0^2
%         |D_{y,j} G_{\ell,Y}(\rho,\sigma)|
%     \, \sigma^{n-1} \, d\sigma
%     &\le
%     C|Y|\rho^\beta
%     \frac{\rho^{n+\alpha-1}}{(1+\ell+\tau_\rho)^2}
%     \le C |Y|.
% \end{align*}

To complete, assume $\ell=0$. The angular component, namely the $D_\theta$ part, is absent. The base component $D_{y,j}$ can be treated exactly as above, using \eqref{eq_basic_R_schur_ell=0}. For the radial part $D_r$, note that
\begin{align*}
\Bigl|\rho^\alpha\frac{\phi_{0,Y}'(\rho)}{\phi_{0,Y}(\rho)}\Bigr|
=\rho^{\alpha-1}\Bigl|-\mathfrak b+\mathfrak a\nu_0 + \mathfrak a \tau_\rho \frac{I_{\nu_0+1}(\tau_\rho)}{I_{\nu_0}(\tau_\rho)}\Bigr|= \mathfrak a \rho^{\alpha-1} \tau_\rho \frac{I_{\nu_0+1}(\tau_\rho)}{I_{\nu_0}(\tau_\rho)}.
\end{align*}
This together with \eqref{eq:normal2} gives
\begin{align*}
    \int_0^2
        |D_r G_{0,Y}(\rho,\sigma)|
    \, \sigma^{n-1} \, d\sigma &\le C (1+\tau_\rho) \rho^{\alpha-1} \int_0^\rho \psi_{0,Y}(\rho) \phi_{0,Y}(\sigma) \, \sigma^{n-1} \, d\sigma\\
    &\quad + C \tau_\rho \rho^{\alpha-1} \frac{I_{\nu_0+1}(\tau_\rho)}{I_{\nu_0}(\tau_\rho)}  \int_\rho^2 \phi_{0,Y}(\rho) \psi_{0,Y}(\sigma) \, \sigma^{n-1} \, d\sigma\\
    &\le C (1+\tau_\rho) \rho^{\alpha-1-\mathfrak b} K_{\nu_0}(\tau_\rho) \int_0^\rho \sigma^{n-\mathfrak b-1} I_{\nu_0}(\tau_\sigma) \, d\sigma\\
    &\quad + C \tau_\rho \rho^{\alpha-1-\mathfrak b} I_{\nu_0+1}(\tau_\rho)  \int_\rho^2 \sigma^{n-\mathfrak b-1} K_{\nu_0}(\tau_\sigma) \, d\sigma\\
    &\le  C \rho^{1-\alpha} (1+\tau_\rho) \tau_\rho^{\frac{2\alpha-2-\mathfrak b}{\mathfrak a}} K_{\nu_0}(\tau_\rho) \int_0^{\tau_\rho} u^{\frac{n-\mathfrak b}{\mathfrak a}-1} I_{\nu_0}(u) du\\
    &\quad + C \rho^{1-\alpha} \tau_\rho^{\frac{2\alpha-2-\mathfrak b}{\mathfrak a}+1} I_{\nu_0+1}(\tau_\rho) \int_{\tau_\rho}^\infty u^{\frac{n-\mathfrak b}{\mathfrak a}-1} K_{\nu_0}(u) du.
\end{align*}
By Lemma~\ref{lem:Bessel:integral}, the above is bounded by 
\begin{equation*}
    C \rho^{1-\alpha} II \left(\frac{2\alpha-2-\mathfrak b}{\mathfrak a},\frac{n-\mathfrak b}{\mathfrak a},\nu_0,\tau_\rho \right) +  C \rho^{1-\alpha} II \left(\frac{2\alpha-2-\mathfrak b}{\mathfrak a}+1,\frac{n-\mathfrak b}{\mathfrak a},\nu_0,\tau_\rho \right)\le C,
\end{equation*}
and
\begin{equation*}
   C \rho^{1-\alpha} I\left(\frac{2\alpha-2-\mathfrak b}{\mathfrak a}+1, \frac{n-\mathfrak b}{\mathfrak a},\nu_0, 1, \tau_\rho\right) \le C,
\end{equation*}
as a consequence of $\frac{2\alpha-2-\mathfrak b}{\mathfrak a}+ \frac{n-\mathfrak b}{\mathfrak a}=0$ and $\rho\le 1$. This completes the proof of Proposition~\ref{lem_basic_D_R_schur_explicit}.

% Changing variables gives
% \begin{equation*}
%     I = C \rho^{1-\alpha} (1+\tau_\rho) \tau_\rho^{\frac{2\alpha-2-\mathfrak b}{\mathfrak a}} K_{\nu_0}(\tau_\rho) \int_0^{\tau_\rho} u^{\frac{n-\mathfrak b}{\mathfrak a}-1} I_{\nu_0}(u) du,
% \end{equation*}
% and
% \begin{equation*}
%     II\le  C \rho^{1-\alpha} \tau_\rho^{\frac{2\alpha-2-\mathfrak b}{\mathfrak a}+1} I_{\nu_0+1}(\tau_\rho) \int_{\tau_\rho}^\infty u^{\frac{n-\mathfrak b}{\mathfrak a}-1} K_{\nu_0}(u) du.
% \end{equation*}
% Suppose first $\tau_\rho \le 1$. 
% Using small argument asymptotics of Bessel functions yields 
% \begin{equation*}
%     I\le C \rho^{1-\alpha} \tau_\rho^{\frac{2\alpha-2-\mathfrak b}{\mathfrak a}-\nu_0} \tau_\rho^{\frac{n-\mathfrak b}{\mathfrak a}+\nu_0}= C \rho^{1-\alpha}\le C.
% \end{equation*}
% Likewise,
% \begin{equation*}
%     II\le C \rho^{1-\alpha} \tau_\rho^{\frac{2\alpha-2-\mathfrak b}{\mathfrak a}+\nu_0+2} = C \rho^{1-\alpha} \tau_\rho^{\frac{2\beta}{\mathfrak a}} \le C.
% \end{equation*}
% Now, if $\tau_\rho >1$, then the large argument asymptotics of Bessel functions implies
% \begin{align*}
%     I&\le C  \rho^{1-\alpha} \tau_\rho^{\frac{2\alpha-2-\mathfrak b}{\mathfrak a}+1} \tau_\rho^{-\frac{1}{2}} e^{-\tau_\rho}   \left( 1 + \int_1^{\tau_\rho} u^{\frac{n-\mathfrak b}{\mathfrak a}-1} u^{-\frac{1}{2}} e^{u}  du \right)\\
%     &\le C + C  \rho^{1-\alpha} \tau_\rho^{\frac{2\alpha-2-\mathfrak b}{\mathfrak a}+1} \tau_\rho^{-\frac{1}{2}} e^{-\tau_\rho} \tau_\rho^{\frac{n-\mathfrak b}{\mathfrak a}-1} \tau_\rho^{-\frac{1}{2}} e^{\tau_\rho}\\
%     &= C + C \rho^{1-\alpha} \tau_\rho^{-1} \le C.
% \end{align*}
% Meanwhile,
% \begin{align*}
%     II&\le C \rho^{1-\alpha} \tau_\rho^{\frac{2\alpha-2-\mathfrak b}{\mathfrak a}+\frac{1}{2}} e^{\tau_\rho} \int_{\tau_\rho}^\infty u^{\frac{n-\mathfrak b}{\mathfrak a}-\frac{3}{2}} e^{-u} du\\
%     &= C \rho^{1-\alpha} \tau_\rho^{\frac{2\alpha-2-\mathfrak b}{\mathfrak a}+\frac{1}{2}} \int_0^\infty (\tau_\rho+u)^{\frac{n-\mathfrak b}{\mathfrak a}-\frac{3}{2}} e^{-u} du.
% \end{align*}
% If $\frac{n-\mathfrak b}{\mathfrak a}-\frac{3}{2}\le 0$, then $(\tau+u)^{\frac{n-\mathfrak b}{\mathfrak a}-\frac{3}{2}}\le \tau^{\frac{n-\mathfrak b}{\mathfrak a}-\frac{3}{2}}$ and $II$ is bounded by $C\rho^{1-\alpha} \tau_\rho^{-1}\le C$. While if $\frac{n-\mathfrak b}{\mathfrak a}-\frac{3}{2}> 0$, then $(\tau_\rho+u)^{\frac{n-\mathfrak b}{\mathfrak a}-\frac{3}{2}} \le \tau_\rho^{\frac{n-\mathfrak b}{\mathfrak a}-\frac{3}{2}} (1+u)^{\frac{n-\mathfrak b}{\mathfrak a}-\frac{3}{2}}$ and thus $II$ is again $O(1)$.






% near $\rho=0$, the finite-energy branch admits the expansion (if $|Y|\le C^*$.)
% \begin{equation*}
%     \phi_{0,Y}(\rho)
%     =
%     C(Y)
%     \bigl[
%         1+O\bigl(|Y|^2\rho^{2\mathfrak a}\bigr)
%     \bigr],
% \end{equation*}
% and hence
% \begin{equation*}
%     \rho^\alpha\partial_\rho\phi_{0,Y}(\rho)
%     =
%     O\bigl(|Y|^2\rho^{2\mathfrak a+\alpha-1}\bigr)
%     =
%     O\bigl(|Y|^2\rho^{2\mathfrak a}\bigr),
% \end{equation*}
% which is bounded because $2\mathfrak a>0$. The boundary branch $\psi_{0,Y}(\rho)$ is smooth and positive on $0<\rho\le 1$, and its weighted derivative is uniformly bounded for $|Y|\le C^*$. Therefore,
% \begin{equation*}
%     \sup_{0<\rho\le1}
%     \int_0^2
%         |D_r G_{0,Y}(\rho,\sigma)|
%     \, \sigma^{n-1} \, d\sigma
%     \le C.
% \end{equation*}




% This completes the proof of Proposition~\ref{lem_basic_D_R_schur_explicit}.
\end{proof}







The following estimate is the technical heart of the Green-kernel analysis. Its
proof avoids differentiating the explicit Bessel formula repeatedly. Instead, it
uses the resolvent identity. The key observation is that $\mathcal L_{\ell,Y}$ depends on
$Y$ only through the potential $|Y|^2\rho^{2\beta}$. Hence every $Y$-derivative
of the Green operator can be written as a finite sum of products of resolvents and
multiplication operators. This turns $Y$-differentiation into row-norm estimates
for the already controlled Green kernel.









\begin{proposition}\label{lem:second_to_last}
For every $N> m+2$, we have
\begin{equation}
    \sup_{0<\rho\le 1} \int_{\mathbb{R}^m} \int_0^2 \sum_{|\gamma|=N} \bigl| \partial_Y^\gamma D_* G_{\ell,Y}(\rho,\sigma) \bigr| \, \sigma^{n-1} \, d\sigma \, dY \le C_{N} (1+\ell)^{-N+m+2}
\end{equation}
for all $\ell \ge 0$.
\end{proposition}

\begin{proof}
Define the Hilbert space $\mathcal{H}:= L^2((0,2), \rho^{n-1}d\rho)$. Recall the operator on $\mathcal{H}$ given by
\begin{equation*}
    \mathcal{L}_{\ell,Y} = - \rho^{2\alpha} \partial_\rho^2 - (n-1+2\alpha) \rho^{2\alpha-1} \partial_\rho + \lambda_\ell \rho^{2\alpha -2} + |Y|^2 \rho^{2\beta}.
\end{equation*}
The crucial point is that $\mathcal{L}_{\ell,Y}$ depends on $Y$ only through the potential $|Y|^2 \rho^{2\beta}$. Therefore
\begin{equation}
    \partial_{Y_j} \mathcal{L}_{\ell,Y} = 2Y_j \rho^{2\beta},\quad \text{and}\quad \partial_{Y_i} \partial_{Y_j} \mathcal{L}_{\ell,Y} = 2 \delta_{ij} \rho^{2\beta},
\end{equation}
where $\delta_{ij} = 1$ if $i=j$ and vanishes otherwise. Moreover, all higher $Y$-derivatives of $\mathcal{L}_{\ell,Y}$ vanish. Since, for every $Y\in \mathbb{R}^m$, the function $|Y|^2 \rho^{2\beta}$ defines a bounded multiplication operator on $\mathcal{H}$, the Friedrichs domains agree for all $Y_1, Y_2 \in \mathbb{R}^m$, namely
\begin{equation*}
    \mathcal{D}(\mathcal{L}_{\ell, Y_1}) = \mathcal{D}(\mathcal{L}_{\ell, Y_2}).
\end{equation*}

Define $R_{\ell,Y}:= \mathcal{L}_{\ell,Y}^{-1}$. Then, for any $u\in \mathcal{H}$ and any $t>0$,
\begin{align}\label{eq:operator:calculus}
    R_{\ell,Y+t e_j}
    \bigl( \mathcal{L}_{\ell, Y+t e_j} - \mathcal{L}_{\ell, Y} \bigr) R_{\ell,Y}u
    &= R_{\ell,Y+t e_j} \mathcal{L}_{\ell, Y+t e_j} R_{\ell,Y}u - R_{\ell,Y+t e_j} \mathcal{L}_{\ell, Y} R_{\ell,Y}u \nonumber\\
    &= R_{\ell,Y+t e_j} \mathcal{L}_{\ell, Y+t e_j} R_{\ell,Y}u - R_{\ell,Y+t e_j}u \nonumber\\
    &= R_{\ell,Y}u - R_{\ell,Y+t e_j}u,
\end{align}
where the last identity uses the fact that
\begin{equation*}
    R_{\ell,Y}u \in \mathcal{D}(\mathcal{L}_{\ell,Y}) = \mathcal{D}(\mathcal{L}_{\ell,Y+te_j}).
\end{equation*}
In addition,
\begin{equation*}
    \mathcal{L}_{\ell, Y+t e_j} - \mathcal{L}_{\ell, Y} = (2t Y_j + t^2) \rho^{2\beta}.
\end{equation*}
Dividing both sides of \eqref{eq:operator:calculus} by $t$ and letting $t \to 0$ gives
\begin{equation*}
    \partial_{Y_j}R_{\ell,Y} = - R_{\ell,Y} \bigl( 2Y_j \rho^{2\beta} \bigr) R_{\ell,Y}.
\end{equation*}
Iterating this identity, for every multi-index $\gamma$ with $|\gamma|=N$, the operator $\partial_{Y}^\gamma R_{\ell,Y}$ is a finite sum of terms of the form
\begin{equation}
    R_{\ell,Y} B_1 R_{\ell,Y} B_2 \dots B_k R_{\ell,Y},
\end{equation}
where each $B_i$ is one of the multiplication operators
\begin{equation*}
    2Y_j \rho^{2\beta} \quad \text{or}\quad 2 \delta_{ij}\rho^{2\beta}.
\end{equation*}
Suppose that $q_1$ of the $B_i$'s are of the first type and $q_2$ of the second type. Then
\begin{equation*}
    q_1+2q_2 = N\quad \text{and}\quad k=q_1+q_2.
\end{equation*}
For an integral operator $T$ with kernel $K_T$ acting on $\mathcal H$, define its row norm by
\begin{equation*}
    \| T \|_{\mathrm{row}}:= \sup_{0<\rho <2} \int_0^2 K_T(\rho,\sigma) \, \sigma^{n-1} \, d\sigma.
\end{equation*}
In particular, a multiplication operator $M_b$ satisfies $\|M_b\|_{\mathrm{row}}=\|b\|_\infty$. Hence
\begin{equation*}
    \| 2Y_j \rho^{2\beta}\|_{\mathrm{row}}\le C|Y|
    \qquad\text{and}\qquad
    \| 2 \delta_{ij}\rho^{2\beta}\|_{\mathrm{row}}\le C.
\end{equation*}
Therefore, by Proposition~\ref{prop:Green:Shur}, for $|Y|\le C^*(1+\ell)$,
\begin{align*}
    \bigl\| B_1 R_{\ell,Y} B_2 \dots B_k R_{\ell,Y} \bigr\|_{\mathrm{row}}
    &\le C |Y|^{q_1} (1+\ell)^{-2k}\\
    &\le C (1+\ell)^{q_1-2q_1-2q_2}
    = C (1+\ell)^{-N}.
\end{align*}
Consequently, for $|Y|\le C^*(1+\ell)$,
\begin{align}\label{eq_G_low}
    \sup_{0<\rho\le 1}\int_0^2 &\sum_{|\gamma|=N} \bigl| \partial_Y^\gamma D_* G_{\ell,Y}(\rho,\sigma) \bigr| \, \sigma^{n-1} \, d\sigma \nonumber\\
    &\le \sum_{|\gamma|=N} C_\gamma \sup_{0<\rho\le 1} \int_0^2 \bigl| D_* G_{\ell,Y}(\rho,\sigma) \bigr| \, \sigma^{n-1} \, d\sigma \; (1+\ell)^{-N} \nonumber\\
    &\le C_N (1+\ell)^{-N},\quad *\in \{ r, \theta\}.
\end{align}
For $D_{y,j} = Y_j \rho^\beta$, Leibniz's rule yields
\begin{equation*}
    \sum_{|\gamma|=N}\partial_{Y}^\gamma \bigl( \rho^\beta Y_j R_{\ell,Y} \bigr)
    = \rho^\beta \sum_{|\gamma|=N} \sum_{\gamma_1+\gamma_2 = \gamma} \bigl(\partial_{Y}^{\gamma_1} Y_j\bigr) \partial_Y^{\gamma_2}R_{\ell,Y}
    = D_{y,j} \sum_{|\gamma|=N}\partial_Y^\gamma R_{\ell,Y} + \rho^\beta \sum_{|\gamma|=N-1} \partial_Y^\gamma R_{\ell,Y}.
\end{equation*}
Accordingly,
\begin{align}\label{eq_low_worst}
    \sup_{0<\rho\le 1}\int_0^2 &\sum_{|\gamma|=N} \bigl| \partial_Y^\gamma D_{y,j} G_{\ell,Y}(\rho,\sigma) \bigr| \, \sigma^{n-1} \, d\sigma \nonumber\\
    &\le (1+\ell)^{-N}\sum_{|\gamma|=N} C_\gamma \sup_{0<\rho\le 1} \int_0^2 \bigl| D_{y,j} G_{\ell,Y}(\rho,\sigma) \bigr| \, \sigma^{n-1} \, d\sigma  \nonumber\\
    &\quad+  (1+\ell)^{-N+1}\sum_{|\gamma|=N} C_\gamma \sup_{0<\rho\le 1} \int_0^2 \bigl| \rho^\beta G_{\ell,Y}(\rho,\sigma) \bigr| \, \sigma^{n-1} \, d\sigma \nonumber\\
    &\le C_N (1+\ell)^{-N+1}.
\end{align}

Now consider the region $|Y|\ge C^*(1+\ell)$. Proposition~\ref{prop:Green:Shur} gives
\begin{align*}
    \bigl\| B_1 R_{\ell,Y} B_2 \dots B_k R_{\ell,Y} \bigr\|_{\mathrm{row}}
    &\le \|B_1 R_{\ell,Y}\|_{\mathrm{row}} \dots \|B_k R_{\ell,Y}\|_{\mathrm{row}}
    \le C|Y|^{q_1-2k}.
\end{align*}
Hence, for $|Y|\ge C^*(1+\ell)$,
\begin{align}\label{eq_G_hihg}
    \sup_{0<\rho\le 1}\int_0^2 &\sum_{|\gamma|=N} \bigl| \partial_Y^\gamma D_* G_{\ell,Y}(\rho,\sigma) \bigr| \, \sigma^{n-1} \, d\sigma \nonumber\\
    &\le \sum_{|\gamma|=N} C_\gamma \sup_{0<\rho\le 1} \int_0^2 \bigl| D_* G_{\ell,Y}(\rho,\sigma) \bigr| \, \sigma^{n-1} \, d\sigma \; |Y|^{q_1-2k} \nonumber\\
    &\le C_N |Y|^{1-N},\quad * \in \{ r, \theta\}.
\end{align}
The case $D_* = D_{y,j}$ is analogous, and the corresponding estimate is bounded by $C_N |Y|^{2-N}$.

Finally, integrate \eqref{eq_G_low} and \eqref{eq_G_hihg} in the $Y$-variable:
\begin{align*}
\int_{\mathbb{R}^m} \int_0^2 \sum_{|\gamma|=N} \bigl| \partial_Y^\gamma D_* G_{\ell,Y}(\rho,\sigma) \bigr| \, \sigma^{n-1} \, d\sigma \, dY
&= \int_{|Y|\le C^*(1+\ell)} + \int_{|Y|>C^*(1+\ell)}\\
&\le C(1+\ell)^{-N+1+m} + C (1+\ell)^{-N+2+m}\\
&\le C (1+\ell)^{-N+m+2},
\end{align*}
for all $0<\rho\le 1$, as required.
\end{proof}

\begin{corollary}\label{cor_weighted_Y_moment_Green}
For every integer $N>m+3$, we have
\begin{equation}\label{eq_weighted_Y_moment_DG}
    \sup_{0<\rho\le 1}
    \int_{\mathbb R^m}
    \int_0^2
    |Y| \sum_{|\gamma|=N}\bigl|\partial_Y^\gamma D_*G_{\ell,Y}(\rho,\sigma)\bigr|\,\sigma^{n-1}\,d\sigma\,dY\le C_N(1+\ell)^{-N+m+3}
\end{equation}
for all $\ell \ge 0$.
\end{corollary}

\begin{proof}
On the low-frequency region $|Y|\le C^*(1+\ell)$, Proposition~\ref{lem:second_to_last}, in particular \eqref{eq_low_worst}, gives, for $|\gamma|=N$,
\begin{equation*}
    \sup_{0<\rho\leq1}
    \int_0^2
    \bigl|
        \partial_Y^\gamma D_*G_{\ell,Y}(\rho,\sigma)
    \bigr|
    \,\sigma^{n-1}\,d\sigma
    \leq
    C_N(1+\ell)^{-N+1},
\end{equation*}
where the exponent $-N+1$ is the worst case, corresponding to the base derivative $D_{y,j}$. Multiplying by $|Y|\le C^*(1+\ell)$ and integrating over the ball $B_{\mathbb{R}^m}(0, C^*(1+\ell))$ gives
\begin{align*}
    \int_{|Y|\le C^*(1+\ell)} |Y| \sup_{0<\rho\leq1} \int_0^2 \sum_{|\gamma|=N} \bigl|\partial_Y^\gamma D_*G_{\ell,Y}(\rho,\sigma)\bigr|\,\sigma^{n-1}\,d\sigma\,dY
    \le C_N(1+\ell)^{-N+m+2}.
\end{align*}
On the high-frequency region $|Y|>C^*(1+\ell)$, Proposition~\ref{lem:second_to_last}, in particular \eqref{eq_G_hihg}, implies
\begin{equation*}
    \sup_{0<\rho\le 1}
    \int_0^2
    \sum_{|\gamma|=N}
    \bigl|
        \partial_Y^\gamma D_*G_{\ell,Y}(\rho,\sigma)
    \bigr|
    \,\sigma^{n-1}\,d\sigma
    \le
    C_N |Y|^{2-N},
\end{equation*}
again taking the worst case $D_*=D_{y,j}$. Therefore,
\begin{align*}
    \int_{|Y|>C^*(1+\ell)} |Y| \sup_{0<\rho\leq1} \int_0^2 \sum_{|\gamma|=N}\bigl|\partial_Y^\gamma D_*G_{\ell,Y}(\rho,\sigma)\bigr|\,\sigma^{n-1}\,d\sigma\,dY
    &\le C_N  \int_{|Y|>C^*(1+\ell)} |Y|^{3-N}\,dY\\
    &\le C_N(1+\ell)^{-N+m+3},
\end{align*}
provided $N>m+3$. This completes the proof of Corollary~\ref{cor_weighted_Y_moment_Green}.
\end{proof}







The preceding estimates are mode-wise. We now return to physical space by summing the spherical harmonics and inverting the Fourier transform in $y$. The sets $E_1$ and $E_2$ are separated in the $y$-variable. This separation allows us to integrate by parts in $Y$ and use the high-order estimates from Proposition~\ref{lem:second_to_last}. The resulting kernel bounds are the precise estimates needed to control the localized Green term in the proof of the anchored reverse H\"older inequality.





\begin{theorem}\label{lem:Green_kernel}
Let $\mathcal{G}(Z,Z')$ be the Green kernel for $L$ on $\mathcal{C}$ satisfying the finite energy condition at $\rho=0$ and the Dirichlet boundary condition at $\rho=2$, where $Z=(\rho,\theta,y)$ and $Z' = (\sigma,\theta',y')$. Let
\begin{equation*}
    E_1 = \{|y|\le 1\},\quad E_2 = \bigl\{2^{\mathfrak a}\le |y| \le 3^{\mathfrak a} \bigr\}.
\end{equation*}
Then
\begin{equation}\label{1}
    \sup_{ \substack{0<\rho \le 1\\ \theta \in \mathbb{S}^{n-1}\\ y\in E_1}}   \int_0^2 \int_{\mathbb{S}^{n-1}} \int_{E_2} \bigl| \nabla_L \mathcal{G}(Z,Z') \bigr| \, \sigma^{n-1} \, dy' \, d\theta' \, d\sigma \le C.
\end{equation}
Moreover, for each $1\le k \le m$,
\begin{equation}\label{2}
    \sup_{ \substack{0<\rho \le 1\\ \theta \in \mathbb{S}^{n-1}\\ y\in E_1}}   \int_0^2 \int_{\mathbb{S}^{n-1}} \int_{E_2} \bigl| \nabla_L \partial_{y_k'} \mathcal{G}(Z,Z') \bigr| \, \sigma^{n-1} \, dy' \, d\theta' \, d\sigma \le C.
\end{equation}
\end{theorem}

\begin{proof}
We first prove \eqref{1}. By \eqref{eq_Green's_kernel}, it suffices to estimate
\begin{align*}
    I_1&:=\int_0^2 \int_{\mathbb{S}^{n-1}} \int_{E_2} \biggl|\sum_{\ell=0}^\infty \Pi_{\ell}(\theta, \theta') \int_{\mathbb{R}^m} e^{i(y-y') \cdot Y} D_* G_{\ell, Y}(\rho, \sigma) \, dY\biggr| \, \sigma^{n-1} \, dy' \, d\theta' \, d\sigma,\\
    &\qquad *\in\{r, y,j\},
\end{align*}
and
\begin{align*}
    I_2&:=\int_0^2 \int_{\mathbb{S}^{n-1}} \int_{E_2} \biggl|\sum_{\ell=0}^\infty \lambda_\ell^{-1/2}\nabla_\theta\Pi_{\ell}(\theta, \theta') \int_{\mathbb{R}^m} e^{i(y-y') \cdot Y} D_\theta G_{\ell, Y}(\rho, \sigma) \, dY\biggr| \, \sigma^{n-1} \, dy' \, d\theta' \, d\sigma.
\end{align*}
Since $y'\in E_2$ and $y\in E_1$, we have $|y-y'|\ge c$. Observe that for any integer $N\ge 1$,
\begin{equation*}
    e^{i(y-y')\cdot Y} = \Bigl(\frac{(y-y')}{i|y-y'|^2} \cdot \nabla_Y \Bigr)^{\!N} e^{i(y-y')\cdot Y}.
\end{equation*}
Choose $N$ sufficiently large. Integrating by parts $N$ times and applying Lemma~\ref{lem:second_to_last}, we obtain
\begin{align*}
    I_1 &\le  \int_{\mathbb{S}^{n-1}} \int_{E_2} \sum_{\ell=0}^\infty |\Pi_{\ell}(\theta, \theta')| \, |y-y'|^{-N}\int_{\mathbb{R}^m} \int_0^2 \sum_{|\gamma| =N} \bigl|\partial_Y^{\gamma} D_* G_{\ell, Y}(\rho, \sigma)\bigr| \, \sigma^{n-1} \, d\sigma \, dY \, dy' \, d\theta' \\
    &\le C \sum_{\ell=0}^\infty (1+\ell)^{\frac{n-2}{2}} (1+\ell)^{-N+m+2}\le C,
\end{align*}
provided $N>m+2+\frac{n}{2}$. Similarly,
\begin{align*}
    I_2 &\le  \int_{\mathbb{S}^{n-1}} \int_{E_2} \sum_{\ell=0}^\infty \lambda_\ell^{-1/2}|\nabla_\theta\Pi_{\ell}(\theta, \theta')| \, |y-y'|^{-N} \int_{\mathbb{R}^m} \int_0^2 \sum_{|\gamma|=N}\bigl|\partial_Y^\gamma D_\theta G_{\ell, Y}(\rho, \sigma)\bigr| \, \sigma^{n-1} \, d\sigma \, dY \, dy' \, d\theta' \\
    &\le C \sum_{\ell=0}^\infty (1+\ell)^{\frac{n-2}{2}-N+m+2}\le C.
\end{align*}
This proves \eqref{1}.

It remains to prove \eqref{2}. Fix $1\le k \le m$. Since
\begin{equation*}
    \partial_{y_k'}e^{i(y-y')\cdot Y}
    =
    -iY_k e^{i(y-y')\cdot Y},
\end{equation*}
it is enough to estimate oscillatory integrals of the form
\begin{equation*}
    \int_{\mathbb R^m}
        e^{i(y-y')\cdot Y}
        Y_k D_*G_{\ell,Y}(\rho,\sigma) \, dY,
\end{equation*}
where $D_*\in\{D_r,D_\theta,D_{y,j}:1\le j\le m\}$. Again, integrate by parts $N$ times. Since
$|y-y'|\ge c$, it follows that for all $0<\rho \le 1$,
\begin{align*}
    \biggl|\int_{\mathbb R^m} e^{i(y-y')\cdot Y} Y_k D_*G_{\ell,Y}(\rho,\sigma) \, dY \biggr|
    &\le C_N |y-y'|^{-N} \int_{\mathbb R^m} \sum_{|\gamma|=N} \bigl| \partial_Y^\gamma \bigl[ Y_k D_*G_{\ell,Y}(\rho,\sigma)\bigr]\bigr| \, dY\\
    &\le C_N \sum_{|\gamma|=N} \int_{\mathbb R^m}|Y|\,\bigl|\partial_Y^\gamma D_*G_{\ell,Y}(\rho,\sigma)\bigr| \, dY\\
    &\quad+ C_N\sum_{|\gamma|=N-1} \int_{\mathbb R^m}\bigl|\partial_Y^\gamma D_*G_{\ell,Y}(\rho,\sigma)\bigr| \, dY.
\end{align*}
After integrating in $\sigma$ and $y'$, Corollary~\ref{cor_weighted_Y_moment_Green} gives
\begin{align*}
    \sup_{0<\rho\leq1}
    \int_0^2
    \int_{E_2}
    \biggl|
    \int_{\mathbb R^m}
        e^{i(y-y')\cdot Y}
        Y_k D_*G_{\ell,Y}(\rho,\sigma) \, dY
    \biggr|
    \,\sigma^{n-1}\,dy'\,d\sigma \le C_N(1+\ell)^{-N+m+3}.
\end{align*}
Using the spherical projection estimates once more, the series in $\ell$ is summable provided
\begin{equation*}
    N>m+3+\frac{n}{2}.
\end{equation*}
Choosing such an $N$ proves \eqref{2}, and hence Theorem~\ref{lem:Green_kernel}.
\end{proof}

We are now in a position to complete the proof of Theorem~\ref{thm:RH}.

\begin{proof}[Proof of Theorem~\ref{thm:RH}]
Recall from the reduction of Theorem~\ref{thm:RH} that
\begin{equation*}
    \omega(\rho,\theta,y)=\chi(y)v(\rho,\theta,y),
\end{equation*}
where $\chi\in C_c^\infty(B_m(0,3^{\mathfrak a}))$ and $\chi=1$ on $B_m(0,2^{\mathfrak a})$. Since $v$ is $L$-harmonic in $Q_4$, we have on the cylinder $\mathcal{C} = (0,2)_\rho \times \mathbb S^{n-1}_\theta \times \mathbb R^m_y$ that $L\omega=f:=[L,\chi]v$. More explicitly,
\begin{equation*}
    f(\sigma,\theta',y')
    =
    \sigma^{2\beta}v(\sigma,\theta',y')\Delta_{y'}\chi(y')
    -
    2\sigma^{2\beta}\nabla_{y'}\chi(y')\cdot
    \nabla_{y'}v(\sigma,\theta',y').
\end{equation*}
Clearly, $\operatorname{supp}(f) \subset E_2$. By the Friedrichs representation, we decompose (see \eqref{DP} and \eqref{BVP}) $\omega=\mathcal T h+\mathcal S f$, where $h(\theta,y)=\omega(2,\theta,y)$. Since $\omega=v$ on $Q_2$, it is enough to prove
\begin{equation*}
    \sup_{Q_1}|\nabla_L\omega|
    \le
    C\|v\|_{L^\infty(Q_3)}.
\end{equation*}
By Theorem~\ref{lem:Poisson_kernel},
\begin{equation*}
    \sup_{Q_1}|\nabla_L\mathcal T h|
    \leq
    C\|h\|_{L^\infty(\mathbb S^{n-1}\times\mathbb R^m)}
    \leq
    C\|v\|_{L^\infty(Q_3)}.
\end{equation*}
It remains to estimate $\nabla_L\mathcal S f$. Write $f=f_1+f_2$, where
\begin{equation*}
    f_1(\sigma,\theta',y')
    :=
    \sigma^{2\beta}v(\sigma,\theta',y')\Delta_{y'}\chi(y')
\end{equation*}
and
\begin{equation*}
    f_2(\sigma,\theta',y')
    :=
    -
    2\sigma^{2\beta}
    \sum_{j=1}^m
        \partial_{y_j'}\chi(y')
        \partial_{y_j'}v(\sigma,\theta',y').
\end{equation*}
We first consider $f_1$. By Theorem~\ref{lem:Green_kernel}, estimate
\eqref{1} yields
\begin{align*}
    \sup_{Q_1}
    |\nabla_L\mathcal S f_1|
    &\le C\|v\|_{L^\infty(Q_3)}
    \sup_{\substack{0<\rho\le 1\\ \theta\in\mathbb S^{n-1}\\ y\in E_1}}
    \int_0^2
    \int_{\mathbb S^{n-1}}
    \int_{E_2} |\nabla_L \mathcal{G}(Z,Z')| \, \sigma^{n-1} \, dy' \, d\theta' \, d\sigma \\
    &\le C\|v\|_{L^\infty(Q_3)}.
\end{align*}
Next, consider $f_2$. For $Z=(\rho,\theta,y)\in Q_1$,
\begin{align*}
    \nabla_L\mathcal S f_2(Z) =
    -2\sum_{j=1}^m
    \int_0^2
    \int_{\mathbb S^{n-1}}
    \int_{E_2}
        \nabla_L\mathcal{G}(Z,Z')
        \sigma^{2\beta}
        \partial_{y_j'}\chi(y')
        \partial_{y_j'}v(Z')
    \,\sigma^{n-1}\, dy' \, d\theta' \, d\sigma .
\end{align*}
Integrating by parts in the $y_j'$-variable,
\begin{align*}
    \nabla_L\mathcal S f_2(Z)
    &= 2\sum_{j=1}^m
    \int_0^2
    \int_{\mathbb S^{n-1}}
    \int_{E_2}
        \partial_{y_j'}
        \bigl[
            \nabla_L\mathcal{G}(Z,Z')
            \sigma^{2\beta}
            \partial_{y_j'}\chi(y')
        \bigr]
        v(Z')
    \,\sigma^{n-1}\, dy' \, d\theta' \, d\sigma\\
    &= 2\sum_{j=1}^m
    \int_0^2
    \int_{\mathbb S^{n-1}}
    \int_{E_2}
        \nabla_L\partial_{y_j'}\mathcal{G}(Z,Z')
        \sigma^{2\beta}
        \partial_{y_j'}\chi(y')
        v(Z')
    \,\sigma^{n-1}\, dy' \, d\theta' \, d\sigma\\
    &\quad -2 \int_0^2 \int_{\mathbb S^{n-1}}
    \int_{E_2}
        \nabla_L\mathcal{G}(Z,Z')
        \sigma^{2\beta}
        \Delta_{y'}\chi(y')
        v(Z')
    \,\sigma^{n-1}\, dy' \, d\theta' \, d\sigma.
\end{align*}
By Theorem~\ref{lem:Green_kernel}, using \eqref{2} for the first term and \eqref{1} for the second, we obtain
\begin{align*}
    \sup_{Q_1}|\nabla_L \mathcal S f_2|
    &\le
    C\|v\|_{L^\infty(Q_3)}
    \sup_{\substack{0<\rho\leq1\\ \theta\in\mathbb S^{n-1}\\ y\in E_1}}
    \int_0^2
    \int_{\mathbb S^{n-1}}
    \int_{E_2}
        |\nabla_L \partial_{y_j'} \mathcal{G}(Z,Z')|
    \,\sigma^{n-1}\, dy' \, d\theta' \, d\sigma\\
    &\quad +
    C\|v\|_{L^\infty(Q_3)}
    \sup_{\substack{0<\rho\leq1\\ \theta\in\mathbb S^{n-1}\\ y\in E_1}}
    \int_0^2
    \int_{\mathbb S^{n-1}}
    \int_{E_2}
        |\nabla_L \mathcal{G}(Z,Z')|
    \,\sigma^{n-1}\, dy' \, d\theta' \, d\sigma\\
    &\le C\|v\|_{L^\infty(Q_3)}.
\end{align*}
Combining the estimates for $\mathcal T h$ and $\mathcal S f$ gives
\begin{equation*}
    \sup_{Q_1}|\nabla_L v| = \sup_{Q_1}|\nabla_L\omega|
    \le
    C\|v\|_{L^\infty(Q_3)},
\end{equation*}
as required. The proof of Theorem~\ref{thm:RH} is complete.
\end{proof}











\part{Riesz transform for \texorpdfstring{$n=1$}{n=1} and \texorpdfstring{$0<\alpha<\frac{1}{2}$}{0<alpha<1/2}}\label{part2}

We now turn to the one-dimensional weakly degenerate regime $n=1$ and
$0<\alpha<1/2$. In this range the Friedrichs form domain still imposes a common
trace across $x=0$, while the weak equation imposes matching of the conormal
flux. Consequently, the even and odd components of a solution satisfy different
half-line endpoint conditions: the even component satisfies a zero-flux condition,
whereas the odd component satisfies a zero-trace condition. This distinction is
responsible for the sharp threshold $p<\alpha^{-1}$: the even zero-flux branch
remains harmless, while the odd zero-trace branch produces exactly the singular
factor $|x|^{-\alpha}$ in the intrinsic gradient.





Throughout Part~\ref{part2}, we assume $n=1$, $m\ge 1$, $\beta\ge 0$, and $\alpha \in \big(0,\frac{1}{2}\big)$. and we use the notation
\begin{align*}
\nu:= \frac{-\mathfrak b}{\mathfrak a} = \frac{1-2\alpha}{2(\beta+1-\alpha)},\qquad \rho = |x|, \qquad \mathfrak A_t:= \frac{t^{\mathfrak a}}{\mathfrak a},\qquad \tau_t := |Y| \mathfrak A_t.
\end{align*}
Note that for $n=1$, $\mathfrak b = \alpha-\frac{1}{2} <0$ (unlike the case $n\ge 2$) and $\nu\in (0,1)$.



\section{Kernel construction for \texorpdfstring{$n=1$}{n=1}}\label{sec:kernel:n=1:alpha<1/2}

By \cite[Theorem~1.1]{RS2}, the Poincar\'e inequality \eqref{P2} still holds in the case $n=1$ and $0<\alpha<1/2$. It follows from Theorem~\ref{thm:CJKS}, Lemma~\ref{reduce_RH}, and Lemma~\ref{RH_remote} that showing the $L^p$-boundedness of the Riesz transform is equivalent to proving the following theorem.





\begin{theorem}\label{thm:RH2}
Let $n = 1$, $m \ge 1$, $\alpha\in (0,\frac{1}{2})$, and $\beta \ge 0$. Let $B=B\big((0,y_0), r\big)$ be an anchored Grushin ball. Suppose that $u$ is $L$-harmonic in $4B$. Then for every $2<p < \alpha^{-1}$, there exists a constant $C>0$ such that
\begin{equation*}
    \biggl( \fint_B |\nabla_L u|^p \, d\xi \biggr)^{\!\frac{1}{p}} \le  \frac{C}{r} \fint_{4B} |u| \, d\xi.
\end{equation*}
\end{theorem}

\begin{proof}[Reduction of Theorem~\ref{thm:RH2}]
By an argument analogous to the \textit{Reduction of Theorem~\ref{thm:RH}}, it suffices to prove
\begin{equation*}
    \biggl( \fint_{Q_1} |\nabla_L v|^p \, d\xi \biggr)^{\!\frac{1}{p}} \le C \fint_{Q_4} |v| \, d\xi
\end{equation*}
for $1<p<\alpha^{-1}$, where $v$ is $L$-harmonic in $Q_4$.

Let $\chi \in C_c^\infty\bigl( B_m( 0, 3^{\mathfrak a} ) \bigr)$ with $\chi = 1$ on $B_m( 0, 2^{\mathfrak a} )$, and set $\omega(x,y) = \chi(y) v(x,y)$. Then
\begin{equation*}
    |\nabla_L v| = |\nabla_L \omega| \qquad \text{on } Q_1.
\end{equation*}
Consider the cylinder
\begin{equation*}
    \widetilde{\mathcal{C}}:= (-2,2)_x \times \mathbb{R}^m_y.
\end{equation*}
A straightforward computation gives
\begin{equation*}
    L\omega = f:= |x|^{2\beta} v \, \Delta_y \chi - 2 |x|^{2\beta} \nabla_y \chi \cdot \nabla_y v,
\end{equation*}
and $f$ is compactly supported in $\{2^{\mathfrak a}\le |y|\le 3^{\mathfrak a}\}$. Let
\begin{equation*}
    h_-(y):= \omega(-2,y), \qquad h_+(y):= \omega(2,y).
\end{equation*}
By the Friedrichs representation,
\begin{equation*}
    \omega = \mathcal{T}h + \mathcal{S}f,
\end{equation*}
where $\mathcal{T}h$ and $\mathcal{S}f$ are the unique solutions of the following problems, respectively:
\begin{equation}\label{eq:pde:n=1}
    \begin{cases}
        Lg = 0, & \text{in } \widetilde{\mathcal{C}},\\
        g(\pm 2,y) = h_\pm(y), \\
        \lim_{x\to 0^-} |x|^{2\alpha}\partial_x g = \lim_{x\to 0^+} |x|^{2\alpha}\partial_x g,\\
        g(0^-,y) = g(0^+,y),
    \end{cases} \qquad
    \begin{cases}
        Lg = f, & \text{in } \widetilde{\mathcal{C}},\\
        g(-2,y)=g(2,y)=0,\\
        \lim_{x\to 0^-} |x|^{2\alpha}\partial_x g = \lim_{x\to 0^+} |x|^{2\alpha}\partial_x g,\\
        g(0^-,y) = g(0^+,y),
    \end{cases}
\end{equation}
where the trace-matching and flux-matching Friedrichs conditions were discussed in Subsection~\ref{ssec:Friedrichs}; see \eqref{eq_subcritical_trace_matching} and \eqref{eq_subcritical_flux_matching}.

Define the reflection operator
\begin{equation*}
    \mathfrak Ru(x,y):= u(-x,y).
\end{equation*}
Since $L$ is invariant under reflection, the Friedrichs form is also reflection-invariant, namely
\begin{equation*}
    \mathfrak q(\mathfrak Ru,\mathfrak Rv) = \mathfrak q(u,v).
\end{equation*}
Equivalently,
\begin{equation*}
    \mathfrak q(\mathfrak Ru,v) = \mathfrak q(u,\mathfrak Rv)
\end{equation*}
for all admissible test functions $v$. Suppose that $g\in \mathcal{F}_{\mathrm{loc}}(\widetilde{\mathcal{C}})$ solves the equation $Lg=F$. Set
\begin{equation*}
    g_e:= \frac{g+\mathfrak Rg}{2},\qquad g_o:=\frac{g-\mathfrak Rg}{2},\qquad F_e:=\frac{F+\mathfrak RF}{2},\qquad F_o:=\frac{F-\mathfrak RF}{2}.
\end{equation*}
Then
\begin{equation*}
    Lg_e=F_e,\qquad L g_o = F_o.
\end{equation*}
Indeed, for every compactly supported Friedrichs test function $v$,
\begin{align*}
    \mathfrak q(g_e,v)
    &= \frac{1}{2} \mathfrak q(g,v) + \frac{1}{2} \mathfrak q(\mathfrak R g,v)\\
    &= \frac{1}{2} \mathfrak q(g,v) + \frac{1}{2} \mathfrak q(g,\mathfrak Rv)\\
    &= \frac{1}{2} \int F v \, dx\,dy  + \frac{1}{2} \int F (\mathfrak Rv) \, dx\,dy\\
    &= \int F_e v \, dx\,dy.
\end{align*}
The proof for $g_o$ is identical. Thus the even and odd components are not merely formal algebraic pieces, but genuine Friedrichs weak solutions of the corresponding even and odd equations.

Now set
\begin{equation*}
    h_{e}:= \frac{h_+ + h_-}{2},\qquad h_{o}:= \frac{h_+ - h_-}{2},\qquad f_e:= \frac{f + \mathfrak R f}{2},\qquad f_o:= \frac{f - \mathfrak R f}{2},
\end{equation*}
where `e' and `o' denote the even and odd parts, respectively. Recalling the endpoint Friedrichs conditions for even and odd solutions discussed in Subsection~\ref{ssec:Friedrichs}, decompose
\begin{align*}
    &\mathcal{T}h(x,y) = \mathcal{T}_e h_e(|x|, y) + \operatorname{sgn}(x)\, \mathcal{T}_o h_o(|x|,y),\\
    &\mathcal{S}f(x,y) = \mathcal{S}_e f_e(|x|,y) + \operatorname{sgn}(x)\, \mathcal{S}_o f_o(|x|,y),
\end{align*}
where $\mathcal{T}_{e}h_{e}$ and $\mathcal{T}_o h_o$ are the unique solutions of the half-line problems
\begin{align}\label{eq:evpdd:Poisson}
\begin{cases}
        Lg = 0, & \text{in } (0,2)\times \mathbb{R}^m,\\
        g(2,y)= h_e(y),\\
        \lim_{x\to 0^+} x^{2\alpha} \partial_x g = 0,
    \end{cases} \qquad
    \begin{cases}
        Lg = 0, & \text{in } (0,2)\times \mathbb{R}^m,\\
        g(2,y)= h_o(y),\\
        g(0^+,y) = 0.
    \end{cases}
\end{align}
Similarly, $\mathcal{S}_e f_e$ and $\mathcal{S}_o f_o$ solve
\begin{align}\label{eq:evpdd:Green}
\begin{cases}
        Lg = f_e, & \text{in } (0,2)\times \mathbb{R}^m,\\
        g(2,y)= 0,\\
        \lim_{x\to 0^+} x^{2\alpha} \partial_x g = 0,
    \end{cases} \qquad
    \begin{cases}
        Lg = f_o, & \text{in } (0,2)\times \mathbb{R}^m,\\
        g(2,y)= 0,\\
        g(0^+,y) = 0.
    \end{cases}
\end{align}
The uniqueness of these Friedrichs problems follows directly from the energy identity.

Accordingly, the problem has been reduced to solving \eqref{eq:evpdd:Poisson} and \eqref{eq:evpdd:Green}, and to estimating
\begin{equation*}
    |\nabla_L \mathcal{T}_e h_e|,\qquad |\nabla_L \mathcal{T}_o h_o|, \qquad |\nabla_L \mathcal{S}_o f_o|, \qquad |\nabla_L \mathcal{S}_e f_e|.
\end{equation*}
\end{proof}







Thus the decomposition is not merely formal: it is the Friedrichs decomposition of
the full-line problem into the zero-flux even half-line problem and the zero-trace
odd half-line problem. The uniqueness of the two half-line Friedrichs problems
follows from the energy identity. Therefore the estimates below may be proved
separately for the even and odd kernels and then recombined on the full line.

In the rest of this section, we construct the kernels of $\mathcal{T}$ and $\mathcal{S}$ explicitly, as we did in Part~\ref{part2}. The main statements and proofs will be postponed to Section~\ref{sec:pf:n=1:alpha<1/2}.






\subsection{The Poisson multiplier for \texorpdfstring{$n=1$}{n=1}}

We first treat equation \eqref{eq:evpdd:Poisson}. Consider the partial differential equation
\begin{equation*}
    Lg = -\partial_x (|x|^{2\alpha}\partial_x g) + |x|^{2\beta} \Delta_y g = 0.
\end{equation*}
Taking the Fourier transform in the $y$-variable, it is enough to consider, for $\rho>0$,
\begin{align*}
    - \rho^{2\alpha} \partial_\rho^2 \hat{g}(\rho,Y) - 2\alpha \rho^{2\alpha-1} \partial_\rho \hat{g}(\rho,Y) + \rho^{2\beta} |Y|^2 \hat{g}(\rho,Y) = 0,
\end{align*}
where
\begin{equation*}
    \hat{g}(\rho,Y) = \int_{\mathbb{R}^m} g(\rho,y) e^{-i y \cdot Y} \, dy.
\end{equation*}
Introduce
\begin{equation*}
    t:= \frac{|Y|}{\mathfrak a} \rho^{\mathfrak a}, \qquad U_Y(t) := \rho^{\mathfrak b} \hat{g}(\rho, Y).
\end{equation*}
Then the equation becomes
\begin{equation*}
    t^2 U_Y''(t) + t U_Y'(t) - (t^2+\nu^2)U_Y(t) = 0.
\end{equation*}
The modified Bessel equation above has two basic solutions:
\begin{equation*}
    \phi_Y(\rho):=\rho^{-\mathfrak b} I_\nu\!\left( \frac{|Y|}{\mathfrak a} \rho^{\mathfrak a} \right), \qquad \varphi_Y(\rho):=\rho^{-\mathfrak b} K_\nu\!\left( \frac{|Y|}{\mathfrak a} \rho^{\mathfrak a} \right).
\end{equation*}
The solution is required to satisfy either the zero-flux or the zero-trace condition. Note that for $0<\nu<1$, as $\rho \to 0$,
\begin{align*}
\phi_Y(\rho) &= \frac{1}{2^\nu \Gamma(\nu+1)} \Bigl(\frac{|Y|}{\mathfrak a}\Bigr)^{\!\nu} \rho^{-2\mathfrak b} + O_{\nu,Y}(\rho^{2\mathfrak a - 2\mathfrak b}), \\
\varphi_Y(\rho) &= 2^{\nu-1} \Gamma(\nu) \Bigl(\frac{|Y|}{\mathfrak a}\Bigr)^{\!-\nu} + 2^{-\nu-1} \Gamma(-\nu) \Bigl(\frac{|Y|}{\mathfrak a}\Bigr)^{\!\nu} \rho^{-2\mathfrak b} + O_{\nu,Y}(\rho^{2\mathfrak a}).
\end{align*}
Set
\begin{equation*}
    c_\nu := 2^{-1} \Gamma(-\nu) \Gamma(\nu+1) = -\frac{\pi}{2\sin{\pi \nu}}.
\end{equation*}
To match the zero-flux condition, define
\begin{align*}
    \Phi_Y(\rho)
    :&= \varphi_Y(\rho) - c_\nu \phi_Y(\rho)\\
    &= \rho^{-\mathfrak b} \Bigl( K_\nu\Bigl(\frac{|Y|}{\mathfrak a}\rho^\mathfrak a \Bigr) + \frac{\pi}{2\sin{\pi \nu}} I_\nu \Bigl(\frac{|Y|}{\mathfrak a}\rho^\mathfrak a \Bigr) \Bigr)\\
    &= \frac{\pi}{2\sin{\pi \nu}} \rho^{-\mathfrak b} I_{-\nu}\Bigl(\frac{|Y|}{\mathfrak a}\rho^\mathfrak a\Bigr)\\
    &= O_{\nu,Y}(1) + O_{\nu,Y}(\rho^{2\mathfrak a}),
\end{align*}
where the second-to-last equality follows from the identity
\begin{equation*}
    K_\nu(t) = \frac{\pi}{2\sin{\pi \nu}} \bigl( I_{-\nu}(t) - I_\nu(t) \bigr), \qquad 0<\nu<1, \quad t>0.
\end{equation*}
In particular,
\begin{equation*}
    \lim_{\rho \to 0} \rho^{2\alpha} \partial_\rho \Phi_Y(\rho) = \lim_{\rho \to 0} \rho^{2\beta+1} = 0.
\end{equation*}
To match the zero-trace condition, we simply take the solution $\phi_Y(\rho)$.

The boundary condition at $\rho =2$ must also be imposed. Define the multipliers
\begin{equation*}
    M_e(\rho,Y):= \frac{\Phi_Y(\rho)}{\Phi_Y(2)} = \frac{\rho^{-\mathfrak b}}{2^{-\mathfrak b}} \frac{I_{-\nu}\!\left(\frac{|Y|}{\mathfrak a}\rho^\mathfrak a\right)}{I_{-\nu}\!\left(\frac{|Y|}{\mathfrak a}2^\mathfrak a\right)}, \qquad M_o(\rho,Y):= \frac{\phi_Y(\rho)} {\phi_Y(2)} = \frac{\rho^{-\mathfrak b}}{2^{-\mathfrak b}} \frac{I_{\nu}\!\left(\frac{|Y|}{\mathfrak a}\rho^\mathfrak a\right)}{I_{\nu}\!\left(\frac{|Y|}{\mathfrak a}2^\mathfrak a\right)}.
\end{equation*}
Their values at $|Y|=0$ are understood by continuity, namely
\begin{equation*}
    M_e(\rho,0)=1 \qquad \text{and} \qquad M_o(\rho,0)= (\rho/2)^{1-2\alpha}.
\end{equation*}
Then $M_e(x,Y) \hat{h}_e(Y)$ and $M_o(x,Y) \hat{h}_o(Y)$ solve the Fourier-transformed equations corresponding to \eqref{eq:evpdd:Poisson}. Therefore, for $x\in (-2,2)$,
\begin{align*}
    \widehat{\mathcal{T}h}(x,Y) = M_e(|x|,Y) \hat{h}_e(Y) + \operatorname{sgn}(x) M_o(|x|,Y) \hat{h}_o(Y),
\end{align*}
and hence
\begin{align}
    \mathcal{T}h(x,y) &= \mathcal{T}_e h_e(x,y) + \mathcal{T}_o h_o(x,y) \nonumber\\
    &=\int_{\mathbb{R}^m} \biggl( \int_{\mathbb{R}^m} M_e(|x|, Y) e^{i (y-y')\cdot Y} \, dY \biggr) h_e(y') \, dy' \nonumber\\
    &\quad+ \operatorname{sgn}(x) \int_{\mathbb{R}^m} \biggl( \int_{\mathbb{R}^m} M_o(|x|, Y) e^{i (y-y')\cdot Y} \, dY \biggr) h_o(y') \, dy'.
\end{align}

\subsection{The Green kernel for \texorpdfstring{$n=1$}{n=1}}

We next turn to equation \eqref{eq:evpdd:Green}. The construction of the Green kernel is similar to that in the case $n\ge 2$, although here the equation is split into its even and odd parts.

Let $\mathcal{G}(Z,Z')$ (with $Z = (x,y)$ and $Z' = (x',y')$) denote the Green kernel associated with $L$ on $\widetilde{\mathcal{C}}$ under the Dirichlet boundary condition at $|x|=2$ and the local finite-energy condition at $x=0$. As before, taking the Fourier transform in the $y$-$y'$ variables shows that the Green kernel takes the form
\begin{equation}\label{eq_Green_kernel:n=1}
    \mathcal{G}(Z,Z') =  \int_{\mathbb{R}^m} e^{i(y-y') \cdot Y} G_{Y}(x, x') \, dY,
\end{equation}
where $G_{Y}$ is the one-dimensional Green kernel associated with the operator
\begin{equation*}
    \mathcal{L}_{Y} := - |x|^{2\alpha} \partial_{x}^2 - 2\alpha |x|^{2\alpha-1} \partial_x + |Y|^2 |x|^{2\beta}.
\end{equation*}
Equivalently,
\begin{equation*}
    \mathcal{S}f(Z) = \iint_{\widetilde{\mathcal{C}}} \mathcal{G}(Z,Z') f(Z') \, dx' \, dy'.
\end{equation*}
Similarly, decompose the Green kernel $G_Y(x,x')$ into even and odd parts corresponding to the Fourier-transformed half-line problems in \eqref{eq:evpdd:Green}. First, to impose the vanishing condition at $x=2$, recall the function
\begin{equation}\label{eq:psi}
    \psi_{Y}(\rho):= \rho^{-\mathfrak b} \Bigl[ K_{\nu}\Bigl( \frac{|Y| \rho^{\mathfrak a}}{\mathfrak a} \Bigr) - \eta_Y I_{\nu} \Bigl( \frac{|Y| \rho^{\mathfrak a}}{\mathfrak a} \Bigr) \Bigr], \qquad \eta_Y:= \frac{K_{\nu} \bigl( \frac{|Y| 2^{\mathfrak a}}{\mathfrak a} \bigr)}{I_{\nu} \bigl( \frac{|Y| 2^{\mathfrak a}}{\mathfrak a} \bigr)}.
\end{equation}
Note that $\psi_{Y}(2)=0$. Next, for the zero-flux condition at $x=0$, use the function $\Phi_Y(x)$. Finally, for the zero-trace condition at $x=0$, use the function $\phi_Y(x)$.

Denote by $G_{e,Y}$ and $G_{o,Y}$ the Green operators corresponding, respectively, to the Fourier-transformed equations in \eqref{eq:evpdd:Green}. A calculation parallel to that in Subsection~\ref{ssec:Green:n>2} yields
\begin{align}\label{eq_kernel_Green:n=1:even}
G_{e,Y}(\rho,\sigma) &= \mathfrak a^{-1} (\eta_Y - c_\nu)^{-1}
\begin{cases}
\Phi_{Y}(\rho) \, \psi_{Y}(\sigma), &0<\rho<\sigma,\\[4pt]
\Phi_{Y}(\sigma) \, \psi_{Y}(\rho), &\sigma < \rho < 2,
\end{cases}\nonumber\\
&=\mathfrak a^{-1} \frac{2\sin{\pi \nu}}{\pi} \frac{I_{\nu}\bigl( \frac{|Y| 2^{\mathfrak a}}{\mathfrak a} \bigr)}{I_{-\nu}\bigl( \frac{|Y| 2^{\mathfrak a}}{\mathfrak a} \bigr)}
\begin{cases}
\Phi_{Y}(\rho) \, \psi_{Y}(\sigma), &0<\rho<\sigma,\\[4pt]
\Phi_{Y}(\sigma) \, \psi_{Y}(\rho), &\sigma < \rho < 2,
\end{cases}
\end{align}
and
\begin{align}\label{eq_kernel_Green:n=1:odd}
G_{o,Y}(\rho,\sigma) = \mathfrak a^{-1}
\begin{cases}
\phi_{Y}(\rho) \, \psi_{Y}(\sigma), &0<\rho<\sigma,\\[4pt]
\phi_{Y}(\sigma) \, \psi_{Y}(\rho), &\sigma < \rho < 2.
\end{cases}
\end{align}
Consequently,
\begin{align}\label{eq:Green:n=1:full}
    \mathcal{S}f(x,y) &= \int_0^2 \int_{\mathbb{R}^m} \biggl( \int_{\mathbb{R}^m} G_{e,Y}(|x|,x') e^{i (y-y')\cdot Y} \, dY \biggr) f_e(x',y') \, dy' \, dx' \nonumber\\
    &\quad + \operatorname{sgn}(x) \int_0^2 \int_{\mathbb{R}^m} \biggl( \int_{\mathbb{R}^m} G_{o,Y}(|x|,x') e^{i (y-y')\cdot Y} \, dY \biggr) f_o(x',y') \, dy' \, dx'.
\end{align}













\section{Proof of Theorem~\ref{thm:RH2}}\label{sec:pf:n=1:alpha<1/2}




This section is devoted to proving Theorem~\ref{thm:RH2} and hence completing the positive result of Theorem~\ref{thm:n=1:alpha<1/2}.

Recall the notation
\begin{equation*}
    \mathcal{I}_\nu (t) := t^{-\nu} I_\nu(t), \qquad S_\nu(t):= \frac{I_{\nu+1}(t)}{I_\nu(t)},\qquad H_{\nu,\rho}(r):= \frac{\mathcal{I}_\nu(\mathfrak A_\rho r)}{\mathcal{I}_\nu(\mathfrak A_2 r)}, \qquad \nu>-1.
\end{equation*}






\begin{lemma}\label{lem:Poisson:n=1}
Let $N\ge 0$, $\mu\in \{-\nu, \nu \}$, and $0<\rho \le 1$. Then for every multi-index $|\gamma|\le N$,
\begin{equation*}
        |\partial_Y^{\gamma} H_{\mu,\rho}(|Y|)| \le \begin{cases}
            C_N, & |Y|\le C^*,\\[4pt]
            C_N (1+|Y|)^{A_N} e^{-c_N |Y|}, & |Y|\ge C^*,
        \end{cases}
\end{equation*}
where $C^*>1$ is a sufficiently large constant.
\end{lemma}

\begin{proof}
By Lemma~\ref{lem_S_derivative_bound}, the estimates for $S_\mu$ imply the corresponding estimates for $H_{\mu,\rho}$. The proof is the same as that of Lemma~\ref{lem_transition_ratio_derivatives} and Proposition~\ref{lem_transition_ratio_derivatives}; see also Remark~\ref{re:H}.
\end{proof}

\begin{theorem}\label{prop:Poisson:n=1}
Let $n= 1$, $m\ge 1$, $\alpha\in \big(0,\frac{1}{2} \big)$, and $\beta \ge 0$. Let $0<\rho\le 1$. Then for every $N\ge 0$ and all $|\gamma|\le N$,
\begin{align}\label{eq:prop:Poisson:n=1:even}
    \bigl| \partial_Y^{\gamma} m_{e,\rho}^{(*)}(Y) \bigr| \le \begin{cases}
        C_N, & |Y|\le C^*,\\[4pt]
        C_N (1+|Y|)^{A_N} e^{-c_N |Y|}, & |Y|\ge C^*,
    \end{cases}
\end{align}
while
\begin{align}\label{eq:prop:Poisson:n=1:odd}
    \bigl| \partial_Y^{\gamma} m_{o,\rho}^{(*)}(Y) \bigr| \le \rho^{-\alpha} \begin{cases}
        C_N, & |Y|\le C^*,\\[4pt]
        C_N (1+|Y|)^{A_N} e^{-c_N |Y|}, & |Y|\ge C^*,
    \end{cases}
\end{align}
where $*\in\{\mathrm{rad}, \mathrm{bas},j\}$.
\end{theorem}

\begin{proof}
We first consider the even multiplier $M_{e,\rho}(Y) = M_{e}(\rho,Y)$. Clearly,
\begin{equation*}
    M_e(\rho,Y) = H_{-\nu,\rho}(|Y|).
\end{equation*}
Thus an argument parallel to that in Lemma~\ref{lem_detailed_poisson_multiplier}, with the order $\nu_\ell$ replaced by the fixed order $-\nu$, applies here as well. In particular, note that
\begin{equation*}
    m_{e,\rho}^{\mathrm{rad}}(Y) = \rho^\alpha \partial_\rho H_{-\nu,\rho}(|Y|) = \rho^\beta |Y| S_{-\nu}\bigl(\mathfrak A_\rho |Y| \bigr) H_{-\nu,\rho}(|Y|).
\end{equation*}
Combining this with Lemma~\ref{lem:Poisson:n=1} gives \eqref{eq:prop:Poisson:n=1:even} for $*=\mathrm{rad}$. The case $*=\mathrm{bas},j$ is entirely similar.

Next, consider $M_{o,\rho}(Y)$. Observe that
\begin{equation*}
    M_{o}(\rho,Y) = \Bigl(\frac{\rho}{2}\Bigr)^{\!1-2\alpha} H_{\nu,\rho}(|Y|).
\end{equation*}
Consequently,
\begin{align*}
    m_{o,\rho}^{\mathrm{rad}}(Y) &= 2^{2\alpha-1}(1-2\alpha) \rho^{-\alpha} H_{\nu,\rho}(|Y|) + 2^{2\alpha-1} \rho^{1-\alpha} \partial_\rho H_{\nu,\rho}(|Y|)\\
    &= 2^{2\alpha-1}(1-2\alpha) \rho^{-\alpha} H_{\nu,\rho}(|Y|) + 2^{2\alpha-1} \rho^{\beta+1-2\alpha} |Y| S_{\nu}(\mathfrak A_\rho |Y|) H_{\nu,\rho}(|Y|).
\end{align*}
Since $\rho^{\beta+1-2\alpha}\le \rho^{-\alpha}$, as $\beta+1-\alpha >0$, \eqref{eq:prop:Poisson:n=1:odd} for $*=\mathrm{rad}$ follows from Lemma~\ref{lem:Poisson:n=1}. The case $*=\mathrm{bas},j$ is analogous.
\end{proof}

\begin{remark}\label{re:even:Poisson}
The proof of the even multiplier estimate in Theorem~\ref{prop:Poisson:n=1} does not use the restriction $0<\alpha<1/2$, except through the notation associated with the weakly degenerate regime. The relevant branch is the zero-flux Friedrichs branch, which is expressed by the normalized Bessel ratio with order $-\nu>-1$. The same argument therefore applies to the zero-flux branch in the strongly degenerate case $1/2\leq\alpha<1$, after replacing $-\nu$ by the corresponding non-negative order. This observation will be used in Part~\ref{part4}.
\end{remark}

\begin{proposition}\label{prop:Green:n=1}
Let $n=1$, $m\ge 1$, $\alpha\in \big(0,\frac{1}{2}\big)$, and $\beta\ge 0$. Then for $0<\rho\le 2$,
\begin{align}\label{eq:prop:Green:1}
    \int_0^2 |G_{o,Y}(\rho,\sigma)| \, d\sigma \le C\,\frac{\rho^{2-2\alpha} + \rho^{1-2\alpha}}{1+\tau_\rho}\le C.
\end{align}
Moreover, for $|Y|>C^*$,
\begin{align}\label{eq:prop:Green:2}
    \sup_{0<\rho\le 2} \rho^{2\beta}\int_0^2 |G_{o,Y}(\rho,\sigma)| \, d\sigma \le C|Y|^{-2}.
\end{align}
In addition, for $0<\rho\le 1$,
\begin{equation}\label{eq:prop:Green:3}
     \int_0^2 |D_* G_{o,Y}(\rho,\sigma)| \, d\sigma \le \rho^{-\alpha}\begin{cases}
         C, & |Y|\le C^*,\\[4pt]
         C|Y|, & |Y|>C^*,
     \end{cases}
\end{equation}
where $*\in\{r,y,j\}$.
\end{proposition}

\begin{proof}
The proof parallels those of Proposition~\ref{prop:Green:Shur}, Lemma~\ref{lem_basic_D_R_schur_explicit}, Proposition~\ref{lem:second_to_last}, and Corollary~\ref{cor_weighted_Y_moment_Green}. Recall that
\begin{align*}
G_{o,Y}(\rho,\sigma) = \mathfrak a^{-1}
\begin{cases}
\phi_{Y}(\rho) \, \psi_{Y}(\sigma), &0<\rho<\sigma,\\[4pt]
\phi_{Y}(\sigma) \, \psi_{Y}(\rho), &\sigma < \rho < 2.
\end{cases}
\end{align*}
It follows from the argument in Proposition~\ref{prop:Green:Shur} that
\begin{equation*}
    \frac{\phi_Y(\sigma)}{\phi_Y(\rho)}\le \Bigl(\frac{\sigma}{\rho}\Bigr)^{\!-2\mathfrak b}\quad \text{if}\quad \sigma\le \rho,\qquad \frac{\psi_Y(\sigma)}{\psi_Y(\rho)}\le C \quad \text{if}\quad \rho \le \sigma\le 2.
\end{equation*}
Moreover, Lemma~\ref{lem_bessel_inputs_high_frequency} and \eqref{eq:product:Bessel} give $\phi_Y(\rho) \psi_Y(\rho) \le C\rho^{-2\mathfrak b}(1+\tau_\rho)^{-1}$.
% \begin{equation*}
%     \phi_Y(\rho) \psi_Y(\rho) \le \frac{C\rho^{-2\mathfrak b}}{1+\tau_\rho}.
% \end{equation*}
Hence, for $0<\rho\le 2$, a straightforward computation shows
\begin{align*}
    \int_0^2 G_{o,Y}(\rho,\sigma) \, d\sigma
    &\le C \phi_Y(\rho) \psi_Y(\rho)  \left(\int_0^\rho \frac{\phi_Y(\sigma)}{\phi_Y(\rho)} \, d\sigma  +  \int_\rho^2 \frac{\psi_Y(\sigma)}{\psi_Y(\rho)} \, d\sigma \right)\le C\,\frac{\rho^{2-2\alpha} + \rho^{1-2\alpha}}{1+\tau_\rho}.
\end{align*}
Since $1-2\alpha >0$, this proves \eqref{eq:prop:Green:1}. In particular, if $\rho \le 1$, then $\int_0^2 G_{o,Y}(\rho,\sigma) \, d\sigma \le C\rho^{1-2\alpha}(1+\tau_\rho)^{-1}$.
% \begin{equation*}
%     \int_0^2 G_{o,Y}(\rho,\sigma) \, d\sigma \le C\rho^{1-2\alpha}(1+\tau_\rho)^{-1}.
% \end{equation*}

Estimate \eqref{eq:prop:Green:2} follows from the same argument as in \textit{Case 2} of the proof of Proposition~\ref{prop:Green:Shur}, using the asymptotic expansions of the Bessel functions. Indeed, for $0<\rho\le 2$,
\begin{align*}
     \rho^{2\beta}\int_0^2 |G_{o,Y}(\rho,\sigma)| \, d\sigma &\le \rho^{2\beta} \int_0^\rho \phi_Y(\sigma) \psi_Y(\rho) d\sigma + \rho^{2\beta} \int_\rho^2 \phi_Y(\rho) \psi_Y(\sigma) d\sigma\\
     &\le C \rho^{2\beta -\mathfrak b} K_\nu(\tau_\rho) \int_0^\rho \sigma^{-\mathfrak b} I_\nu(\tau_\sigma) d\sigma + C \rho^{2\beta -\mathfrak b} I_\nu(\tau_\rho) \int_\rho^\infty \sigma^{-\mathfrak b} K_\nu(\tau_\sigma) d\sigma\\
     &=C|Y|^{-2} \tau_\rho^{\frac{2\beta-\mathfrak b}{\mathfrak a}} K_\nu(\tau_\rho) \int_0^{\tau_\rho} u^{\frac{1-\mathfrak b}{\mathfrak a}-1} I_\nu(u) du + C|Y|^{-2} \tau_\rho^{\frac{2\beta-\mathfrak b}{\mathfrak a}} I_\nu(\tau_\rho) \int_{\tau_\rho}^\infty u^{\frac{1-\mathfrak b}{\mathfrak a}-1} K_\nu(u) du\\
     &= C|Y|^{-2} \left( II\left(\frac{2\beta-\mathfrak b}{\mathfrak a}, \frac{1-\mathfrak b}{\mathfrak a},\nu, \tau_\rho \right) + I\left(\frac{2\beta-\mathfrak b}{\mathfrak a}, \frac{1-\mathfrak b}{\mathfrak a},\nu,0, \tau_\rho \right) \right) \le C |Y|^{-2},
\end{align*}
where the last inequality follows by Lemma~\ref{lem:Bessel:integral}.

We next prove the gradient estimate \eqref{eq:prop:Green:3}. By Proposition~\ref{lem_basic_D_R_schur_explicit} (see~\eqref{D:bdd:G}), for all $0<\rho\le 1$,
\begin{align*}
    \int_0^2 |D_r G_{o,Y}(\rho,\sigma)| \, d\sigma
    &\le C(1+\tau_\rho) \rho^{\alpha-1} \int_0^2 G_{o,Y}(\rho,\sigma) \, d\sigma \le C \rho^{-\alpha}.
\end{align*}
This together with \eqref{eq:prop:Green:1} yield
\begin{align*}
    \int_0^2 |D_r G_{o,Y}(\rho,\sigma)| \, d\sigma \le C \rho^{-\alpha},\qquad 0<\rho\le 1.
\end{align*}
Finally, consider the base derivative $D_{y,j}$. Clearly,
\begin{align*}
    \int_0^2 |D_{y,j}G_{o,Y}(\rho,\sigma)| \, d\sigma
    &\le C|Y| \rho^{\beta} \int_0^2 G_{o,Y}(\rho,\sigma) \, d\sigma \le C|Y| \rho^{\beta+1-2\alpha}\le C |Y| \rho^{-\alpha}.
\end{align*}
This proves \eqref{eq:prop:Green:3}.
\end{proof}

\begin{proposition}\label{prop:Green:n=1:even}
Let $n=1$, $m\ge 1$, $\alpha\in (0,1)$, and $\beta\ge 0$. Then for $|Y|\le C^*$,
\begin{align}\label{eq:prop:Green:even:1}
\sup_{0<\rho\le 2} \int_0^2 |G_{e,Y}(\rho,\sigma)| \, d\sigma \le C.
\end{align}
Moreover, for $|Y|>C^*$,
\begin{align}\label{eq:prop:Green:even:2}
    \sup_{0<\rho\le 2} \rho^{2\beta}\int_0^2 |G_{e,Y}(\rho,\sigma)| \, d\sigma \le C|Y|^{-2}.
\end{align}
In addition, for $0<\rho\le 1$,
\begin{equation}\label{eq:prop:Green:even:3}
     \int_0^2 |D_* G_{e,Y}(\rho,\sigma)| \, d\sigma \le \begin{cases}
         C, & |Y|\le C^*,\\[4pt]
         C|Y|, & |Y|>C^*,
     \end{cases}
\end{equation}
where $*\in\{r,y,j\}$.
\end{proposition}

\begin{proof}
Recall that
\begin{align}\label{eq:normal}
G_{e,Y}(\rho,\sigma) = \mathfrak a^{-1} \frac{2\sin{\pi \nu}}{\pi} \frac{I_{\nu}\!\bigl( \frac{|Y| 2^{\mathfrak a}}{\mathfrak a} \bigr)}{I_{-\nu}\!\bigl( \frac{|Y| 2^{\mathfrak a}}{\mathfrak a} \bigr)}
\begin{cases}
\Phi_{Y}(\rho) \, \psi_{Y}(\sigma), &0<\rho<\sigma,\\[4pt]
\Phi_{Y}(\sigma) \, \psi_{Y}(\rho), &\sigma < \rho < 2.
\end{cases}
\end{align}
To prove \eqref{eq:prop:Green:even:1}, define
\begin{equation*}
    \mathcal{U}_Y(\rho):= \frac{\rho^{- \mathfrak b} I_{-\nu}(\tau_\rho)}{2^{-\mathfrak b} I_{-\nu}(\tau_2)},\qquad \mathcal{V}_Y(\rho):= 2^{-\mathfrak b} \rho^{-\mathfrak b} \bigl[ I_\nu(\tau_2) K_\nu(\tau_\rho) - K_\nu(\tau_2) I_\nu(\tau_\rho) \bigr],
\end{equation*}
so that $G_{e,Y}(\rho,\sigma)$ can be written in the normalized form
\begin{equation*}
G_{e,Y}(\rho,\sigma) = \mathfrak a^{-1}
\begin{cases}
\mathcal{U}_{Y}(\rho) \, \mathcal{V}_{Y}(\sigma), &0<\rho<\sigma,\\[4pt]
\mathcal{U}_{Y}(\sigma) \, \mathcal{V}_{Y}(\rho), &\sigma < \rho < 2.
\end{cases}
\end{equation*}
Since $\mathfrak b = - \mathfrak a \nu$, it is clear that $\mathcal{U}_Y(\rho) = H_{-\nu,\rho}(|Y|)$. Moreover, using the identity $K_\nu(t) = \frac{\pi}{2 \sin{\pi \nu}} \bigl( I_{-\nu}(t) - I_\nu(t) \bigr)$,
% \begin{equation*}
%     K_\nu(t) = \frac{\pi}{2 \sin{\pi \nu}} \bigl( I_{-\nu}(t) - I_\nu(t) \bigr),
% \end{equation*}
one obtains
\begin{equation*}
    \mathcal{V}_Y(\rho) = C_{\alpha,\beta,\nu} \, \mathcal{I}_\nu(\tau_2) \, \mathcal{I}_{-\nu}(\tau_2) \Bigl[ H_{-\nu,\rho}(|Y|) - \Bigl(\frac{\rho}{2}\Bigr)^{\!1-2\alpha} H_{\nu,\rho}(|Y|) \Bigr].
\end{equation*}
If $|Y|\le C^*$, then $\tau_2\le C$, and therefore $\mathcal{I}_\nu(\tau_2) \mathcal{I}_{-\nu}(\tau_2)$ is bounded by a constant depending only on $\alpha,\beta,\nu$. Hence Lemma~\ref{lem:Poisson:n=1} implies
\begin{equation*}
    |\mathcal{U}_Y(\rho)|+|\mathcal{V}_Y(\rho)|\le C.
\end{equation*}
Furthermore, the identity $\rho^{\alpha} \partial_\rho H_{\pm \nu, \rho}(|Y|) = \rho^{\beta} |Y| S_{\pm \nu}(\mathfrak A_\rho |Y|) H_{\pm \nu, \rho}(|Y|)$, together with
% \begin{equation*}
%     \rho^{\alpha} \partial_\rho H_{\pm \nu, \rho}(|Y|) = \rho^{\beta} |Y| S_{\pm \nu}(\mathfrak A_\rho |Y|) H_{\pm \nu, \rho}(|Y|),
% \end{equation*}
Lemma~\ref{lem_S_derivative_bound} gives
\begin{equation*}
    |D_r \mathcal{U}_Y(\rho)|\le C
    \qquad\text{and}\qquad
    |D_r \mathcal{V}_Y(\rho)|\le C \rho^{-\alpha}.
\end{equation*}
Immediately, for all $0<\rho\le 1$,
\begin{align*}
    \int_0^2 |D_r G_{e,Y}(\rho,\sigma)| \, d\sigma \le C \int_0^\rho \rho^{-\alpha}\, d\sigma + C\int_\rho^2 \, d\sigma \le C \rho^{1-\alpha} + C\le C.
\end{align*}
The estimate for the base derivative $D_{y,j}$ is straightforward. Therefore \eqref{eq:prop:Green:even:1} and \eqref{eq:prop:Green:even:3} for $|Y|\le C^*$ follow directly.

To prove \eqref{eq:prop:Green:even:2}, first note that for $|Y|\ge C^*$, the factor $\frac{I_\nu(\tau_2)}{I_{-\nu}(\tau_2)}\le C$ since $\tau_2 \ge \mathfrak a^{-1} 2^{\mathfrak a}$. Next, write
\begin{align*}
    \rho^{2\beta}\int_0^2 |G_{e,Y}(\rho,\sigma)| \, d\sigma &\le C \rho^{2\beta} \psi_Y(\rho) \int_0^\rho \Phi_Y(\sigma) \,d\sigma + C \rho^{2\beta} \Phi_Y(\rho) \int_\rho^2 \psi_Y(\sigma)\, d\sigma\\
    &\le C \rho^{2\beta - \mathfrak b} K_\nu(\tau_\rho) \int_0^\rho \sigma^{-\mathfrak b} I_{-\nu}(\tau_\sigma) \, d\sigma + C \rho^{2\beta - \mathfrak b} I_{-\nu}(\tau_\rho) \int_\rho^\infty \sigma^{-\mathfrak b} K_{\nu}(\tau_\sigma) \, d\sigma.
\end{align*}
By change of variables, it equals
\begin{align*}
    &C |Y|^{-2} \tau_\rho^{\frac{2\beta-\mathfrak b}{\mathfrak a}} K_\nu(\tau_\rho) \int_0^{\tau_\rho} u^{\frac{1-\mathfrak b}{\mathfrak a}-1}  I_{-\nu}(u) \, du + C |Y|^{-2} \tau_\rho^{\frac{2\beta-\mathfrak b}{\mathfrak a}} I_{-\nu}(\tau_\rho) \int_{\tau_\rho}^\infty u^{\frac{1-\mathfrak b}{\mathfrak a}-1}  K_{\nu}(u) \,  du \\
    &= C |Y|^{-2} II\left(\frac{2\beta-\mathfrak b}{\mathfrak a}, \frac{1-\mathfrak b}{\mathfrak a}, -\nu, \tau_\rho \right) + C |Y|^{-2} I\left(\frac{2\beta-\mathfrak b}{\mathfrak a}, \frac{1-\mathfrak b}{\mathfrak a}, -\nu, 0, \tau_\rho \right)\le C|Y|^{-2},
\end{align*}
where the last inequality follows by Lemma~\ref{lem:Bessel:integral} and $\frac{2\beta-\mathfrak b}{\mathfrak a}+ \frac{1-\mathfrak b}{\mathfrak a}=2$ and $\nu\in (0,1)$. This proves \eqref{eq:prop:Green:even:2}.

% Estimate \eqref{eq:prop:Green:even:2} follows by the same argument as in \textit{Case 2} of the proof of Proposition~\ref{prop:Green:Shur}, again using the asymptotic expansions of the Bessel functions.

It remains to prove \eqref{eq:prop:Green:even:3} for $|Y|\ge C^*$. Use the expression \eqref{eq:normal}. The prefactor $I_{\nu}( \tau_2 )/I_{-\nu}( \tau_2 )$ is again harmless. The key observation is that
\begin{align*}
\rho^\alpha\frac{\Phi_{Y}'(\rho)}{\Phi_{Y}(\rho)}
&= \mathfrak a \rho^{\alpha-1}\tau_\rho \frac{I_{1-\nu}(\tau_\rho)}{I_{-\nu}(\tau_\rho)}.
\end{align*}
% Together with $0<\frac{I_{1-\nu}(t)}{I_{-\nu}(t)}< 1$, this yields
% \begin{equation*}
%     \Bigl|\rho^\alpha\frac{\Phi_{Y}'(\rho)}{\Phi_{Y}(\rho)}\Bigr|
%     \le
%     C\rho^{\alpha-1}\tau_\rho.
% \end{equation*}
Also, the same argument used for \eqref{eq:normal2} shows that
\begin{equation*}
    \Bigl|\rho^\alpha \frac{\psi_Y'(\rho)}{\psi_Y(\rho)} \Bigr| \le C (1+\tau_\rho) \rho^{\alpha-1}.
\end{equation*}
Therefore, for $0<\rho\le 1$, Lemma~\ref{lem:Bessel:integral} implies
\begin{align*}
    &\int_0^2 |D_r G_{e,Y}(\rho,\sigma)| \, d\sigma
    \le C(1+\tau_\rho) \rho^{\alpha-1} \int_0^\rho \psi_Y(\rho) \Phi_Y(\sigma) \, d\sigma + C\tau_\rho \rho^{\alpha-1} \frac{I_{1-\nu}(\tau_\rho)}{I_{-\nu}(\tau_\rho)} \int_\rho^2 \psi_Y(\sigma) \Phi_Y(\rho) \, d\sigma \\
    &\le C (1+\tau_\rho) \rho^{\alpha-1-\mathfrak b} K_\nu(\tau_\rho) \int_0^\rho \sigma^{-\mathfrak b} I_{-\nu}(\tau_\sigma) \, d\sigma + C\tau_\rho \rho^{\alpha-1-\mathfrak b} I_{1-\nu}(\tau_\rho) \int_\rho^\infty \sigma^{-\mathfrak b} K_\nu(\tau_\sigma) \, d\sigma\\
    &= C |Y|^{-\theta} (1+\tau_\rho) \tau_\rho^{\theta-\nu-\frac{1}{\mathfrak a}} K_\nu(t) \int_0^{\tau_\rho} u^{\nu+\frac{1}{\mathfrak a}-1}I_{-\nu}(u) \, du + C |Y|^{-\theta} t^{\theta+1-\nu-\frac{1}{\mathfrak a}} I_{1-\nu}(t) \int_t^\infty u^{\nu+\frac{1}{\mathfrak a}-1} K_\nu(u) \, du\\
    &= C |Y|^{-\theta} \left[ II\left( \frac{\mathfrak b - \alpha}{\mathfrak a}, \frac{1-\mathfrak b}{\mathfrak a},-\nu,\tau_\rho \right)  +  II\left( \frac{\mathfrak b - \alpha}{\mathfrak a}+1, \frac{1-\mathfrak b}{\mathfrak a},-\nu,\tau_\rho \right)    + I\left( \frac{\mathfrak b - \alpha}{\mathfrak a}+1, \frac{1-\mathfrak b}{\mathfrak a}, -\nu,1,\tau_\rho \right) \right]\\
    &\le C |Y|^{-\theta}\le C,
\end{align*}
where $\theta:= \frac{1-\alpha}{\mathfrak a}\in (0,1]$.




% We claim that both $I$ and $II$ are bounded by $C|Y|^{-\theta}$, where
% \begin{equation*}
%     \theta:= \frac{1-\alpha}{\mathfrak a} \in (0,1].
% \end{equation*}
% Indeed, a change of variables gives
% \begin{align*}
%     I
%     &= \mathfrak a^{-1} \rho^{\alpha+2 \mathfrak a \nu} (1+\tau_\rho) \tau_\rho^{-\nu-\frac{1}{\mathfrak a}} K_\nu(\tau_\rho) \int_0^{\tau_\rho} u^{\nu+\frac{1}{\mathfrak a}-1} I_{-\nu}(u) \, du \\
%     &= C_{\alpha,\beta} |Y|^{-\theta} I_{1}(\tau_\rho),
% \end{align*}
% where
% \begin{equation*}
%     I_{1}(t) := (1+t) t^{\theta-\nu-\frac{1}{\mathfrak a}} K_\nu(t) \int_0^t u^{\nu+\frac{1}{\mathfrak a}-1}I_{-\nu}(u) \, du.
% \end{equation*}
% It remains to show that $\sup_{t>0}I_{11}(t)<\infty$. If $0<t\le 1$, then the small-argument asymptotics of the modified Bessel functions gives
% \begin{align*}
%     I_{1}(t)
%     \le C t^{\theta-2\nu-\frac{1}{\mathfrak a}} \int_0^t u^{\nu+\frac{1}{\mathfrak a}-1-\nu} \, du
%     \le C t^{\theta-2\nu}
%     \le C,
% \end{align*}
% since $\theta-2\nu=\frac{\alpha}{\mathfrak a}>0$. If $t>1$, then splitting the integral at $1$ yields
% \begin{align*}
%     I_{1}(t)
%     &\le C t^{\theta+1-\nu-\frac{1}{\mathfrak a}-\frac{1}{2}} e^{-t} \biggl(1 + \int_1^t u^{\nu+\frac{1}{\mathfrak a}-\frac{3}{2}} e^u \, du\biggr)\\
%     &\le C + C t^{\theta-1}
%     \le C,
% \end{align*}
% where we used the elementary estimate $\int_1^t u^{q} e^u \, du\le C_q e^t t^q$ and the fact that $\theta \le 1$.

% % The term $I_2$ is treated in the same way. Indeed,
% % \begin{equation*}
% %     I_2 = C_{\alpha,\beta} |Y|^{-\theta} \biggl( \tau_\rho^{\theta-\nu-\frac{1}{\mathfrak a}} K_\nu(\tau_\rho) \int_0^{\tau_\rho} u^{\nu+\frac{1}{\mathfrak a}-1} I_{-\nu}(u) \, du \biggr),
% % \end{equation*}
% % and the same argument as for $I_1$ shows that the expression in parentheses is $O(1)$.

% Likewise, $II$ is bounded by $C_{\alpha,\beta} |Y|^{-\theta} II_1(\tau_\rho)$, where
% \begin{equation*}
%     II_1(t):= t^{\theta+1-\nu-\frac{1}{\mathfrak a}} I_{-\nu}(t) \int_t^\infty u^{\nu+\frac{1}{\mathfrak a}-1} K_\nu(u) \, du.
% \end{equation*}
% If $t\le 1$, then
% \begin{equation*}
%     II_1(t)
%     \le C t^{\theta+2-2\nu-\frac{1}{\mathfrak a}} \biggl( \int_t^1 u^{\frac{1}{\mathfrak a}-1} \, du + \int_1^\infty e^{-u} u^{\nu+\frac{1}{\mathfrak a}-\frac{3}{2}} \, du \biggr)
%     = O(1),
% \end{equation*}
% since $\theta+2-2\nu-\frac{1}{\mathfrak a} = 2-\theta \ge 0$. If $t>1$, then
% \begin{equation*}
%     II_1(t)
%     \le C t^{\theta+1-\nu-\frac{1}{\mathfrak a}-\frac{1}{2}} e^{t} \int_t^\infty e^{-u} u^{\nu+\frac{1}{\mathfrak a}-\frac{3}{2}} \, du \le C t^{\theta-1},
% \end{equation*}
% which is again $O(1)$. Hence, for $|Y|\ge C^*$,
% \begin{equation*}
%     \int_0^2 |D_r G_{e,Y}(\rho,\sigma)| \, d\sigma \le C|Y|^{-\theta}\le C.
% \end{equation*}


To close, consider the base derivative $D_{y,j}=i\rho^{\beta}Y_j$. By Lemma~\ref{lem:Bessel:integral} and a parallel argument of \eqref{eq:prop:Green:even:2}, for $|Y|\ge C^*$,
\begin{align*}
    \int_0^2 |D_{y,j}G_{e,Y}(\rho,\sigma)| \, d\sigma \le C|Y|^{-\theta} \left[II\left(\frac{\beta-\mathfrak b}{\mathfrak a}, \frac{1-\mathfrak b}{\mathfrak a}, -\nu, \tau_\rho \right) + I\left(\frac{\beta-\mathfrak b}{\mathfrak a}, \frac{1-\mathfrak b}{\mathfrak a}, -\nu, 0, \tau_\rho \right) \right]\le C.
\end{align*}
This proves \eqref{eq:prop:Green:even:3} and completes the proof of Proposition~\ref{prop:Green:n=1:even}.
\end{proof}








It follows that

\begin{theorem}\label{thm:Green:n=1:secondtolast}
For every $N>m+3$, the following estimates hold for all $0<\rho\le 1$:
\begin{equation}
    \int_{\mathbb{R}^m} \int_0^2 \sum_{|\gamma|=N} \bigl| \partial_Y^{\gamma} D_* G_{\#,Y}(\rho,\sigma) \bigr| \, d\sigma \, dY \le C_N \kappa_{\#}(\rho),
\end{equation}
and
\begin{equation}
    \int_{\mathbb{R}^m} \int_0^2 |Y| \sum_{|\gamma|=N} \bigl| \partial_Y^{\gamma} D_* G_{\#,Y}(\rho,\sigma) \bigr| \, d\sigma \, dY \le C_N \kappa_{\#}(\rho),
\end{equation}
where $*\in\{r,y,j\}$, $\# \in \{ e,o\}$, and
\begin{equation*}
    \kappa_{\#}(\rho) = \begin{cases}
        1, & \# = e,\\[4pt]
        \rho^{-\alpha}, & \# = o.
    \end{cases}
\end{equation*}
\end{theorem}

\begin{proof}
The proof is the fixed-order analogue of Proposition~\ref{lem:second_to_last}
and Corollary~\ref{cor_weighted_Y_moment_Green}. We record the details needed to make the
reduction explicit.

For each fixed $Y\in\mathbb R^m$, the half-line operators defining
$G_{e,Y}$ and $G_{o,Y}$ depend on $Y$ only through the same potential
$|Y|^2\rho^{2\beta}$. Hence the resolvent identity used in the proof of
Proposition~\ref{lem:second_to_last} remains valid without modification:
\begin{equation*}
    \partial_{Y_j}R_Y = -R_Y\bigl(2Y_j\rho^{2\beta}\bigr)R_Y,
\end{equation*}
and the higher $Y$-derivatives are finite sums of products of resolvents and
multiplication operators of the form $2Y_j\rho^{2\beta}$ or $2\delta_{ij}\rho^{2\beta}$.

For the even kernel, Proposition~\ref{prop:Green:n=1:even} provides the same row bounds as in the higher-dimensional case, uniformly for $0<\rho\leq1$. Therefore, the argument of Proposition~\ref{lem:second_to_last} gives
\begin{equation*}
    \int_{\mathbb R^m}\int_0^2
\sum_{|\gamma|=N}\bigl|\partial_Y^\gamma D_*G_{e,Y}(\rho,\sigma)\bigr|
\, d\sigma \, dY
\leq C_N.
\end{equation*}
The corresponding weighted estimate, with one additional factor $|Y|$, follows by the same argument.

For the odd kernel, Proposition~\ref{prop:Green:n=1} has the same structure, except for the additional factor $\rho^{-\alpha}$ in the intrinsic derivative bounds. Since this factor is independent of $Y$, it passes unchanged through the
resolvent-differentiation argument. Consequently,
\begin{equation*}
    \int_{\mathbb R^m}\int_0^2
\sum_{|\gamma|=N}\bigl|\partial_Y^\gamma D_*G_{o,Y}(\rho,\sigma)\bigr|
\, d\sigma \, dY
\leq C_N\rho^{-\alpha},
\end{equation*}
and the same bound remains valid with one extra factor $|Y|$, provided
$N>m+3$. Combining the even and odd estimates proves both assertions.
\end{proof}

We conclude this section by proving Theorem~\ref{thm:RH2}.

\begin{proof}[Proof of Theorem~\ref{thm:RH2}]
By repeating the arguments used in Theorem~\ref{lem:Poisson_kernel}, Theorem~\ref{lem:Green_kernel}, and the proof of Theorem~\ref{thm:RH}, together with Theorem~\ref{thm:Green:n=1:secondtolast}, the following estimates are obtained: let $n=1$, $m\ge 1$, $\alpha\in (0,\frac{1}{2})$, and $\beta \ge 0$. Then for $0<|x|\le 1$ and $y\in \mathbb{R}^m$,
\begin{itemize}
    \item $|\nabla_L \mathcal{T}_e h_e(x,y)| + |x|^{\alpha} |\nabla_L \mathcal{T}_o h_o(x,y)| \le C \|h\|_\infty \le C \|v\|_{L^\infty(Q_3)}$,
    \item $|\nabla_L \mathcal{S}_e f_e(x,y)| + |x|^{\alpha} |\nabla_L \mathcal{S}_o f_o(x,y)| \le C \|v\|_{L^\infty(Q_3)}$.
\end{itemize}
Therefore,
\begin{align*}
    \biggl( \fint_{Q_1} |\nabla_L v|^p \, d\xi \biggr)^{\!\frac{1}{p}}
    &= \biggl( \fint_{Q_1} |\nabla_L \omega|^p \, d\xi \biggr)^{\!\frac{1}{p}}\\
    &\le C\biggl( \fint_{Q_1} |\nabla_L \mathcal{T}_e h_e|^p + |\nabla_L \mathcal{T}_o h_o|^p + |\nabla_L \mathcal{S}_e f_e|^p + |\nabla_L \mathcal{S}_o f_o|^p \biggr)^{\!\frac{1}{p}}\\
    &\le C \|v\|_{L^\infty(Q_3)} \biggl( \fint_{Q_1} |x|^{-\alpha p} \, d\xi \biggr)^{\!\frac{1}{p}}\\
    &\le C \biggl(\fint_{Q_4} |v| \, d\xi \biggr) \biggl( \int_0^1 s^{-\alpha p} \, ds \biggr)^{\!\frac{1}{p}}\\
    &\le C \fint_{Q_4} |v| \, d\xi,
\end{align*}
since $p <\alpha^{-1}$. This completes the proof of Theorem~\ref{thm:RH2}.
\end{proof}










\section{The unboundedness of the Riesz transform}



In this section we complete the proof of Theorem~\ref{thm:n=1:alpha<1/2} by showing the unboundedness of the Riesz transform. By Theorem~\ref{thm:CJKS}, the negative result is equivalent to the failure of \eqref{eq:RH:Jiang}. 

\begin{proposition}\label{prop_RH_failure_n1}
Let $n=1$, $m\ge 1$, $\alpha\in (0,\frac{1}{2})$, and $\beta\ge 0$. Then \eqref{eq:RH:Jiang} fails for every $p\ge \alpha^{-1}$.
\end{proposition}

\begin{proof}
Consider
\begin{equation}\label{eq_counterexample_function}
    u(x,y):=\operatorname{sgn}(x)|x|^{1-2\alpha}.
\end{equation}
Since $\alpha \in (0,\frac{1}{2})$, the function $u$ is bounded near $x=0$ and has locally finite energy. Indeed,
\begin{equation*}
\partial_xu(x,y) = (1-2\alpha)|x|^{-2\alpha}, \qquad x\neq0,
\end{equation*}
and therefore
\begin{equation*}
|\nabla_Lu(x,y)| = |x|^\alpha |\partial_x u(x,y)| = (1-2\alpha)|x|^{-\alpha} \in L^2_{\mathrm{loc}}.
\end{equation*}
We first verify that $u$ is weakly $L$-harmonic. For $x\neq0$,
\begin{equation*}
    |x|^{2\alpha}\partial_xu(x,y)
    =
    1-2\alpha.
\end{equation*}
Thus the weighted flux is the same constant on both sides of the singular line. Consequently, for every test function $\varphi\in C_c^\infty$,
\begin{align*}
\int L u \, \varphi \, dx\,dy &= \int |x|^{2\alpha} \partial_x u \, \partial_x \varphi \, dx\,dy\\
&= (1-2\alpha)\int \partial_x\varphi \, dx\,dy = 0.
\end{align*}
Hence $Lu=0$ in the weak sense.

Now let $B$ be any anchored ball centered on the singular line $\{x=0\}$. Such a ball contains a product sub-box of the form $\{|x|<\varepsilon\}\times B_y$ with $|B_y|>0$. It follows that
\begin{align*}
\int_B|\nabla_Lu|^p \, dx\,dy \ge c\int_{-\varepsilon}^{\varepsilon}|x|^{-\alpha p} \, dx.
\end{align*}
The latter integral diverges whenever $p\ge \alpha^{-1}$. On the other hand, $u$ is bounded on $\lambda B$, so
\begin{equation*}
\fint_{\lambda B}|u|<\infty, \qquad \lambda \ge 1.
\end{equation*}
This proves Proposition~\ref{prop_RH_failure_n1}.
\end{proof}





\begin{remark}\label{rem_counterexample_other_regimes}
The counterexample \eqref{eq_counterexample_function} is specific to the case $n=1$ and $0<\alpha<1/2$. Indeed, this construction is no longer available when $n=1$ and $\alpha\ge 1/2$. If
$\alpha>1/2$, the function $|x|^{1-2\alpha}$ is unbounded and
\begin{equation*}
\int_0^1
x^{2\alpha}
\bigl|
\partial_x x^{1-2\alpha}
\bigr|^2 \, dx
=
C\int_0^1x^{-2\alpha}\,dx
=
\infty.
\end{equation*}
At $\alpha=1/2$, the second independent solution is $\log|x|$, which again has
infinite energy. Although $\operatorname{sgn}(x)$ is then locally Friedrichs harmonic,
its form gradient vanishes, so it produces no failure of a reverse H\"older inequality. 
\end{remark}












\part{Reverse Riesz inequality}\label{part3}




In this part, we study the reverse Riesz inequality. Inspired by \cite{Gilles}, the decomposition below separates the region where the two $x$-variables are of comparable scale from the region where one point lies much farther from the singular set than the other. The most delicate contribution comes from the critical remote interaction. Its leading term is harmonic and therefore cannot be controlled by size estimates alone; the harmonic annihilation argument removes this obstruction and reduces the problem to Hardy-type inequalities.

Although our main interest is in the case $n=1$, the argument below applies verbatim to the general case $n\ge 2$. We therefore do not restrict ourselves to $n=1$. Throughout Part~\ref{part3}, we assume that $n,m\ge 1$, $\alpha\in(0,1)$, and $\beta\ge 0$.

\section{Analysis of the Riesz kernel}

Let $\kappa > 2^{10}$ be large. We split the double space $\mathbb{R}^{n+m}\times \mathbb{R}^{n+m}$ into
\begin{align*} 
    \mathcal{D}_1 &= \{(\xi,\eta)\in \mathbb{R}^{n+m}\times \mathbb{R}^{n+m}: |x-x'| \ge \kappa^{-1} |x| \ \text{and}\ |x'| \le \kappa |x| \},\\
    \mathcal{D}_2 &= \{(\xi,\eta)\in \mathbb{R}^{n+m}\times \mathbb{R}^{n+m}: |x-x'| \le \kappa^{-1} |x| \ \text{and}\ |x'| \le \kappa |x| \},\\
    \mathcal{D}_3 &= \{(\xi,\eta)\in \mathbb{R}^{n+m}\times \mathbb{R}^{n+m}: |x'| \ge \kappa |x| \}.
\end{align*}
Note that in $\mathcal{D}_1$, we have $|x-x'| \sim |x|$, and in $\mathcal{D}_3$, $|x|+|x'|\sim |x'|$. Set $\mathcal R:= \nabla_L L^{-1/2}$ and
\begin{align*}
    \mathcal{R}_i(\xi,\eta) := \nabla_L L^{-1/2}(\xi,\eta) \mathbf{1}_{\mathcal{D}_i} = c \int_0^\infty \nabla_L e^{-tL}(\xi,\eta) \mathbf{1}_{\mathcal{D}_i} \, \frac{dt}{\sqrt{t}}.
\end{align*}





\subsection{Estimates of the good part}

In this subsection, we estimate the good part of the Riesz transform, i.e., $\mathcal{R}_1$.






\begin{proposition}\label{R1}
The good part of the Riesz transform is bounded on $L^p$ for all $\frac{n}{n-1+\alpha}<p< \infty$, i.e.,
\begin{equation*}
    \| \mathcal{R}_1f\|_{p}\le C \|f\|_p,\qquad \forall\,
     \frac{n}{n-1+\alpha}<p< \infty,
\end{equation*}
for all $f\in C_c^\infty(\mathbb{R}^{n+m})$.
\end{proposition}

\begin{proof}
See Proposition~\ref{T1,T2} below.
\end{proof}

\begin{lemma}\label{kernel_good}
Let $n,m\ge 1$, $\alpha\in (0,1)$, and $\beta \ge 0$. Then the good part of the Riesz kernel satisfies the estimate
\begin{align}\label{eq_kernel_good}
    |\mathcal{R}_1(\xi,\eta)| \le C\mathbf{1}_{|x'|\le \kappa|x|} \begin{cases}
        |x|^{-\mathcal{Q}(1-\alpha)}, & |x|^{\mathfrak a} \ge |y-y'|,\\[4pt]
        |y-y'|^{-\frac{\mathcal{Q}(1-\alpha)}{\mathfrak a}} + |x|^{-(1-\alpha)} |y-y'|^{-\frac{(\mathcal{Q}-1)(1-\alpha)}{\mathfrak a}}, & |x|^{\mathfrak a} \le |y-y'|,
    \end{cases}
\end{align}
where $\xi=(x,y)$ and $\eta= (x',y')$.
\end{lemma}

\begin{proof}
For $(\xi,\eta)\in \mathcal{D}_1$, one has $\kappa^{-1}|x| \le |x-x'| \le (1+\kappa)|x|$, so $|x|\sim |x-x'|$ and $|x|+|x'|\sim |x|$. Split the integral at the threshold $t = |x|^{2-2\alpha}$. Then
\begin{align*}
    |\mathcal{R}_1(\xi,\eta)|
    &\le \int_0^\infty \bigl| \nabla_L e^{-tL}(\xi,\eta) \bigr| \mathbf{1}_{\mathcal{D}_1} \, \frac{dt}{\sqrt{t}}\\
    &\le C\int_0^{|x|^{2-2\alpha}} \frac{\mathbf{1}_{\mathcal{D}_1}}{V(\xi,\sqrt{t})} \exp\!\Bigl( - \frac{d(\xi,\eta)^2}{c t} \Bigr) \frac{dt}{t}\\
    &\quad + C\int_{|x|^{2-2\alpha}}^\infty \frac{\mathbf{1}_{\mathcal{D}_1}}{|x|^{1-\alpha} V(\xi,\sqrt{t})} \exp\!\Bigl( - \frac{d(\xi,\eta)^2}{c t} \Bigr) \frac{dt}{\sqrt{t}} \\
    &:= I_1 + I_2.
\end{align*}
For $I_1$, since $\sqrt{t}\le |x|^{1-\alpha}$, Lemma~\ref{le_volume} implies that
\begin{equation*}
    V(\xi,\sqrt{t}) \sim t^{\frac{n+m}{2}} |x|^{n\alpha+m\beta}.
\end{equation*}
Hence,
\begin{align*}
    I_1
    &\le C |x|^{-n\alpha - m\beta} d(\xi,\eta)^{-n-m}\mathbf{1}_{\mathcal{D}_1} \int_{d(\xi,\eta)^2/|x|^{2-2\alpha}}^\infty e^{-s} s^{\frac{n+m}{2}-1} \, ds\\
    &\le C |x|^{-n\alpha - m\beta} d(\xi,\eta)^{-n-m}\mathbf{1}_{\mathcal{D}_1} \exp\!\Bigl( - \frac{d(\xi,\eta)^2}{2|x|^{2-2\alpha}} \Bigr)\\
    &\le C |x|^{-n\alpha-m\beta+\varepsilon(1-\alpha)} d(\xi,\eta)^{-n-m-\varepsilon} \mathbf{1}_{\mathcal{D}_1},\quad \forall\, \varepsilon\ge 0\\
    &= C d(\xi,\eta)^{-\mathcal{Q}} \mathbf{1}_{\mathcal{D}_1},
\end{align*}
where in the last step we choose $\varepsilon = \frac{n\alpha+m\beta}{1-\alpha}$.

Next, consider $I_2$. Since $\sqrt{t}\ge |x|^{1-\alpha}$, Lemma~\ref{le_volume} gives
\begin{equation*}
    V(\xi,\sqrt{t}) \sim t^{\frac{\mathcal{Q}}{2}}.
\end{equation*}
Therefore,
\begin{align*}
    I_2
    &\le C |x|^{\alpha-1} d(\xi,\eta)^{-\mathcal{Q}+1} \mathbf{1}_{\mathcal{D}_1} \int_0^{d(\xi,\eta)^2/|x|^{2-2\alpha}} e^{-s} s^{\frac{\mathcal{Q}-1}{2}-1} \, ds\\
    &\le C |x|^{\alpha-1} d(\xi,\eta)^{-\mathcal{Q}+1} \mathbf{1}_{\mathcal{D}_1}.
\end{align*}
Now note that for $(\xi,\eta)\in \mathcal{D}_1$,
\begin{align*}
    d(\xi,\eta)
    &\sim \frac{|x-x'|}{(|x|+|x'|)^\alpha} + \frac{|y-y'|}{(|x|+|x'|)^\beta + |y-y'|^{\frac{\beta}{\mathfrak a}}}\\
    &\sim \begin{cases}
        |x|^{1-\alpha} + |x|^{-\beta} |y-y'|, & |x|^{\mathfrak a} \ge |y-y'|,\\[4pt]
        |x|^{1-\alpha} + |y-y'|^{\frac{1-\alpha}{\mathfrak a}}, & |x|^{\mathfrak a} \le |y-y'|.
    \end{cases}
\end{align*}
Substituting these expressions into the bounds for $I_1$ and $I_2$ yields the stated estimate.
\end{proof}

According to Lemma~\ref{kernel_good}, define the following operators. Note that
\begin{equation*}
    V\bigl((0,y), |x|^{1-\alpha}\bigr) \sim |x|^{\mathcal{Q}(1-\alpha)} \sim |B_n(0,|x|)| \times |B_m(y,|x|^{\mathfrak a})|.
\end{equation*}
Set
\begin{align}\label{def_T1}
    T_1&: u \mapsto \frac{1}{V\bigl( (0,y), |x|^{1-\alpha} \bigr)} \int_{B_n(0,\kappa|x|)} \int_{B_m(y,|x|^{\mathfrak a})} |u(\eta)| \, d\eta,\\[4pt]
    \label{def_T2}
    T_2&: u \mapsto |x|^{-\varepsilon} \int_{B_n(0,\kappa|x|)} \int_{B_m(y, |x|^{\mathfrak a})^c} \frac{|u(\eta)|}{|y-y'|^{\frac{\mathcal{Q}(1-\alpha)-\varepsilon}{\mathfrak a}}} \, d\eta, \qquad \varepsilon = 0,\, 1-\alpha.
\end{align}

It is immediate from Lemma~\ref{kernel_good} that Proposition~\ref{R1} follows once the $L^p$-boundedness of $T_1$ and $T_2$ is established.

\begin{proposition}\label{T1,T2}
Let $n,m\ge 1$, $\alpha\in (0,1)$, and $\beta \ge 0$. Then
\begin{enumerate}[label=(\roman*)]
    \item $T_1$ is bounded on $L^p$ for all $1<p\le \infty$,
    \item $T_2$ is bounded on $L^p$ for all $1<p\le \infty$ if $\varepsilon=0$,
    \item $T_2$ is bounded on $L^p$ for all $\frac{n}{n-1+\alpha}<p\le  \infty$ if $\varepsilon=1-\alpha$.
\end{enumerate}
\end{proposition}

\begin{proof}
For $(i)$, Lemma~\ref{le_volume} implies that
\begin{equation*}
    B_n(0,\kappa|x|) \times B_m\bigl(y,|x|^{\mathfrak a}\bigr) \subset  B\bigl( (0,y), \kappa c_3^{-1}|x|^{1-\alpha} \bigr).
\end{equation*}
The conclusion then follows from \eqref{D} and the maximal theorem.

We next prove $(ii)$ and $(iii)$. Applying a standard dyadic decomposition to the $y'$-space gives
\begin{align*}
    |T_2f(\xi)|
    &\le |x|^{-\varepsilon} \sum_{j=0}^\infty \int_{B_n(0,\kappa |x|)} \int_{B_m( y, 2^{j+1}|x|^{\mathfrak a} )\setminus B_m ( y, 2^{j}|x|^{\mathfrak a} )} |y-y'|^{-\frac{\mathcal{Q}(1-\alpha)-\varepsilon}{\mathfrak a}} |f(\eta)| \, dy' \, dx'\\
    &\le C |x|^{-\varepsilon} \sum_{j=0}^\infty 2^{-j\frac{\mathcal{Q}(1-\alpha)-\rho}{\mathfrak a}} |x|^{-\mathcal{Q}(1-\alpha)+\varepsilon} \int_{B_n(0,\kappa|x|)} \int_{B_m( y, 2^{j+1}|x|^{\mathfrak a} )} |f(\eta)| \, d\eta.
\end{align*}
Let $q>1$. Applying H\"older's inequality twice yields
\begin{equation}
    |T_2f(\xi)|\le C |x|^{-\mathcal{Q}(1-\alpha)} \sum_{j=0}^\infty 2^{-j\frac{\mathcal{Q}(1-\alpha)-\varepsilon}{\mathfrak a} + j\frac{m}{q'}} |x|^{\frac{\mathcal{Q}(1-\alpha)}{q'}} \biggl( \int_{B_n(0,\kappa|x|) \times B_m(y, 2^{j+1}|x|^{\mathfrak a})} |f(\eta)|^q \, d\eta \biggr)^{\!\frac{1}{q}}.
\end{equation}
Furthermore, Lemma~\ref{le_volume} gives
\begin{align*}
    B_n(0,\kappa|x|) \times B_m\bigl(y, 2^{j+1}|x|^{\mathfrak a}\bigr)
    &\subset B_n\bigl(0,\kappa 2^{\frac{j+1}{\mathfrak a}}|x|\bigr) \times B_m\bigl(y, \kappa^{\mathfrak a} 2^{j+1}|x|^{\mathfrak a}\bigr)\\
    &\subset B\bigl((0,y), c_3^{-1} \kappa^{1-\alpha} 2^{\frac{(j+1)(1-\alpha)}{\mathfrak a}} |x|^{1-\alpha}\bigr).
\end{align*}
Hence,
\begin{equation*}
    |T_2f(\xi)|\le C M_q f(\xi) \sum_{j=0}^\infty 2^{-j\bigl( \frac{n}{q'(\mathfrak a)} - \frac{\varepsilon}{\mathfrak a} \bigr)} \le C M_q f(\xi),
\end{equation*}
provided that
\begin{equation*}
    q > \frac{n}{n-\varepsilon} = \begin{cases}
        1, &\varepsilon=0,\\[4pt]
        \frac{n}{n-1+\alpha}, &\varepsilon=1-\alpha.
    \end{cases}
\end{equation*}
The result now follows from the maximal theorem.
\end{proof}














\subsection{Estimates of the diagonal part}




In this subsection, our aim is to prove the following.

\begin{proposition}\label{R2}
The diagonal part of the Riesz transform is bounded on $L^p$ for all $p>2$, in the sense that
\begin{equation*}
    \|\mathcal{R}_2f\|_p \le C \|f\|_p,\qquad p>2,
\end{equation*}
for all $f\in C_c^\infty(\mathbb{R}^n \setminus \{0\} \times \mathbb{R}^m)$.
\end{proposition}

\begin{proof}
The proof follows from Proposition~\ref{Rj,2}, Lemma~\ref{Rj,1,2}, and Proposition~\ref{Rj,1,1} below.
\end{proof}

We introduce the following covering lemma; see also \cite{Gilles}.

\begin{lemma}\label{lemma_remotecover}
Let $\varepsilon>0$ and $\delta \ge 1$. There exists a sequence of balls $\{B_n^j = B_n(x_j, r_j)\}_{j \ge 0}$ such that
\begin{enumerate}[label=(\roman*)]
    \item\label{(1)} $\mathbb{R}^n = \bigcup_{j\ge 0} B_n^j$, where $B_n^0 = B_n(0,\varepsilon)$ and, for $j\ge 1$, $B_n^j = B_n(x_j,r_j)$ is remote in the sense that $r_j \le \frac{|x_j|}{2}$,
    \item\label{(2)} for all $j \ge 1$, $2^{-10-\delta}|x_j|\le r_j \le 2^{-9-\delta} |x_j|$,
    \item\label{(3)} there exists a finite overlap constant $C_L>0$ such that $\sum_j \mathbf{1}_{B_n^j}(x) \le C_L$ for all $x\in \mathbb{R}^n$, and $C_L$ does not depend on $\varepsilon$.
\end{enumerate}
\end{lemma}


\begin{proof}
See Appendix~\ref{app:tele}.
\end{proof}



\begin{remark}\label{remark_delta}
As is clear from the proof, the finite overlap constant depends only on $n$ and $\delta$. However, we may choose $\delta$ sufficiently large, depending only on the geometric structure of the Grushin spaces, for instance on $c_1,c_2,c_3,c_4,\varepsilon_0,\gamma,n,m,\alpha,\beta$. In fact, for later use, we may choose $\delta \ge 1$ such that
\begin{equation}\label{delta}
    \delta > \max \Bigl( \frac{2}{\mathfrak a} - 9,\; - \log_2\!\Bigl({\frac{2^7 c_1 \varepsilon_0}{9}}\Bigr),\;    - \log_2\!\Bigl({\frac{2^6 c_1 \varepsilon_0}{3\gamma}}\Bigr),\; \frac{-\log_2\!\bigl({\frac{c_1 \varepsilon_0}{36}}\bigr)}{\mathfrak a}-9,\;    \frac{-\log_2\!\bigl({\frac{c_1 \varepsilon_0}{24 \gamma}}\bigr)}{\mathfrak a}-9    \Bigr).
\end{equation}
\end{remark}

Let $\varepsilon>0$ and let $\delta \ge 1$ be chosen later. By Lemma~\ref{lemma_remotecover}, there exists a sequence of balls $\{B_n^j=B_n(x_j,r_j)\}_j$ satisfying properties \eqref{(1)}, \eqref{(2)}, and \eqref{(3)} of Lemma~\ref{lemma_remotecover}. Let $\{\mathcal{X}_j\}_j$ be a smooth partition of unity subordinate to $\{B_n^j\}_j$.

By choosing $\kappa$ sufficiently large, say $\kappa = 2^{10}$, the support property of $\mathcal{D}_2$ ensures that
\begin{equation*}
    |\mathcal{R}_2f(\xi)| \le \sum_{j=0}^\infty \bigl| \mathbf{1}_{4B_n^j}(x) \mathcal{R} \bigl( f \mathcal{X}_j \bigr)(\xi) \bigr| =: \sum_{j=0}^\infty R_j f(\xi).
\end{equation*}
Suppose that $f\in C_c^\infty(B_n(0,\varepsilon)^c \times \mathbb{R}^m)$. Since $\mathcal{X}_0$ is supported in $B_n(0,\varepsilon)$, the term $R_0f$ vanishes.

Next, by the finite overlap property, Proposition~\ref{R2} follows once the following statements are established for $p>2$:
\begin{itemize}
    \item for all $f\in C_c^\infty(B_n(0,\varepsilon)^c \times \mathbb{R}^m)$, one has $\sup_{j\ge 1} \| R_jf\|_{p}\le C\|f\|_p$,
    \item the constant $C$ above does not depend on $\varepsilon$.
\end{itemize}
Then Proposition~\ref{R2} follows by letting $\varepsilon \to 0$.

To this end, we further decompose $R_j$ ($j\ge 1$) as $R_j = R_{j,1}+R_{j,2}$, where
\begin{align}
    R_{j,1}f &= \mathbf{1}_{4B_n^j}(x) \int_0^{r_j^{2-2\alpha}} \nabla_L e^{-tL}\bigl(f \mathcal{X}_j\bigr) \, \frac{dt}{\sqrt{t}},\\[4pt]
    R_{j,2}f &= \mathbf{1}_{4B_n^j}(x) \int_{r_j^{2-2\alpha}}^\infty \nabla_L e^{-tL}\bigl(f \mathcal{X}_j\bigr) \, \frac{dt}{\sqrt{t}}.
\end{align}

\begin{proposition}\label{Rj,2}
We have
    \begin{equation*}
        \sup_{j\ge 1} \|R_{j,2}\|_{p\to p} <\infty
    \end{equation*}
    for all $\frac{n}{n-1+\alpha}<p\le \infty$.
\end{proposition}

\begin{proof}
To simplify the notation, set
\begin{equation}\label{sigma}
    \sigma = \frac{|y-y'|}{r_j^\beta + |y-y'|^{\frac{\beta}{\mathfrak a}}}.
\end{equation}
For $x,x'\in 4 B_n^j$, Lemma~\ref{lemma_remotecover} implies that $|x|,|x'|\sim r_j$, and hence
\begin{equation*}
    \bigl|\mathbf{1}_{4B_n^j}(x) \nabla_L e^{-tL}(\xi,\eta) \mathbf{1}_{B_n^j}(x')\bigr|\le \Bigl( \frac{1}{\sqrt{t}} + \frac{1}{r_j^{1-\alpha}} \Bigr) \frac{C}{V(\xi,\sqrt{t})} \exp\!\Bigl( - \frac{\sigma^2}{c t} \Bigr).
\end{equation*}
A straightforward computation gives
\begin{align*}
    |R_{j,2}(\xi,\eta)|
    &\le C \mathbf{1}_{4B_n^j}(x) \mathbf{1}_{B_n^j}(x') r_j^{\alpha-1} \sigma^{-\mathcal{Q}+1} \int_0^{\sigma^2/r_j^{2-2\alpha}} e^{-s} s^{\frac{\mathcal{Q}-1}{2}-1} \, ds\\
    &\le C \mathbf{1}_{4B_n^j}(x) \mathbf{1}_{B_n^j}(x') r_j^{(\alpha-1)(1+\varepsilon)} \sigma^{-\mathcal{Q}+1+\varepsilon},\quad \forall\, 0\le \varepsilon < \mathcal{Q}-1.
\end{align*}
Observe that
\begin{equation*}
    \sigma \sim \begin{cases}
        r_j^{-\beta} |y-y'|, & |y-y'|\le r_j^{\mathfrak a},\\[4pt]
        |y-y'|^{\frac{1-\alpha}{\mathfrak a}}, & |y-y'|\ge r_j^{\mathfrak a}.
    \end{cases}
\end{equation*}
It follows that
\begin{align*}
    |R_{j,2}f(\xi)|
    &\le C \mathbf{1}_{4B_n^j}(x) r_j^{\beta(\mathcal{Q}-1-\varepsilon) - (1-\alpha)(1+\varepsilon)} \int_{B_n^j \times B_m( y, r_j^{\mathfrak a} )} \frac{|f(\eta)|}{|y-y'|^{\mathcal{Q}-1-\varepsilon}} \, d\eta\\
    &\quad + C \mathbf{1}_{4B_n^j}(x) r_j^{\alpha-1} \int_{B_n^j \times B_m( y, r_j^{\mathfrak a} )^c} \frac{|f(\eta)|}{|y-y'|^{\frac{\mathcal{Q}(1-\alpha)-(1-\alpha)}{\mathfrak a}}} \, d\eta\\
    &:= I_1f + I_2f.
\end{align*}
Exactly the same argument as in Proposition~\ref{T1,T2}~(iii), with $|x|$ replaced by $r_j$, shows that $I_2$ is bounded on $L^p$ for $p>\frac{n}{n-1+\alpha}$. It remains to estimate $I_1$. Fix $q>1$. Using Lemma~\ref{lemma_remotecover} and H\"older's inequality, one obtains
\begin{align}\label{I1}
    |I_1f(\xi)|
    &\le C  \mathbf{1}_{4B_n^j}(x) r_j^{\beta(\mathcal{Q}-1-\varepsilon) - (1-\alpha)(1+\varepsilon)} \biggl( \int_{B_n(0, 2^{11+\delta} r_j) \times B_m (y, r_j^{\mathfrak a} )} |f(\eta)|^q \, d\eta \biggr)^{\!\frac{1}{q}} \nonumber\\
    &\qquad\times \biggl( \int_{B_n^j} \int_{B_m(y,r_j^{\mathfrak a})} \frac{dy'\,dx'}{|y-y'|^{q'(\mathcal{Q}-1-\varepsilon)}} \biggr)^{\!\frac{1}{q'}}.
\end{align}
First, Lemma~\ref{le_volume} gives
\begin{equation*}
    B_n\bigl(0, 2^{11+\delta} r_j\bigr) \times B_m \bigl(y, r_j^{\mathfrak a} \bigr) \subset B\bigl( (0,y), 2^{11+\delta} c_3^{-1} r_j^{1-\alpha} \bigr).
\end{equation*}
Since $x\in 4B_n^j \subset B_n(0, 2^{11+\delta} r_j)$, the first factor in \eqref{I1} is bounded by
\begin{equation}\label{I1(1)}
    C r_j^{\frac{\mathcal{Q}(1-\alpha)}{q}} M_q f(\xi).
\end{equation}
Second, if we choose $\mathcal{Q}-1-\frac{m}{q'}<\varepsilon<\mathcal{Q}-1$, then a direct computation shows that the second factor in \eqref{I1} equals
\begin{equation}\label{I1(2)}
    C |B_n^j|^{\frac{1}{q'}} \biggl(\int_0^{r_j^{\mathfrak a}} s^{-q'(\mathcal{Q}-1-\varepsilon) +m-1} \, ds \biggr)^{\!\frac{1}{q'}} = C r_j^{\frac{n + (\mathfrak a)( m-q'(\mathcal{Q}-1-\varepsilon) )}{q'}}.
\end{equation}
Combining \eqref{I1(1)} and \eqref{I1(2)}, we conclude that
\begin{equation*}
    |I_1f(\xi)| \le C M_q f(\xi),
\end{equation*}
and the $L^p$-boundedness of $I_1$ for $1<p\le \infty$ follows from the maximal theorem. This completes the proof of Proposition~\ref{Rj,2}.
\end{proof}

Our next goal is to analyze $R_{j,1}$. Decompose
\begin{equation}\label{furthersplit}
    R_{j,1}(\xi,\eta) = R_{j,1}(\xi,\eta) \mathbf{1}_{|y-y'|\le r_{j}^{\mathfrak a}} + R_{j,1}(\xi,\eta) \mathbf{1}_{|y-y'|\ge r_{j}^{\mathfrak a}}:= R_{j,1}^1(\xi,\eta) + R_{j,1}^2(\xi,\eta).
\end{equation}

\begin{lemma}\label{Rj,1,2}
We have
    \begin{equation*}
        \sup_{j\ge 1} \|R_{j,1}^2\|_{p\to p} <\infty
    \end{equation*}
    for all $1<p\le \infty$.
\end{lemma}

\begin{proof}
Recall \eqref{sigma}. By Lemma~\ref{Grushin_gradient} and Lemma~\ref{RS_Grushin}, the kernel satisfies the estimate
\begin{align*}
    |R_{j,1}^2(\xi,\eta)|
    &\le C \mathbf{1}_{4B_n^j}(x) \mathbf{1}_{B_n^j}(x') \mathbf{1}_{|y-y'|\ge r_{j}^{\mathfrak a}} r_j^{-n\alpha-m\beta} \sigma^{-n-m} \int_{\sigma^2/r_j^{2-2\alpha}}^\infty e^{-s} s^{\frac{n+m}{2}-1} \, ds\\
    &\le C \mathbf{1}_{4B_n^j}(x) \mathbf{1}_{B_n^j}(x') \mathbf{1}_{|y-y'|\ge r_{j}^{\mathfrak a}} r_j^{-n\alpha-m\beta+\rho(1-\alpha)} \sigma^{-n-m-\varepsilon}\\
    &\le C \mathbf{1}_{4B_n^j}(x) \mathbf{1}_{B_n^j}(x') \mathbf{1}_{|y-y'|\ge r_{j}^{\mathfrak a}} |y-y'|^{-\frac{\mathcal{Q}(1-\alpha)}{\mathfrak a}},
\end{align*}
where in the last step we choose $\varepsilon = \frac{n\alpha+m\beta}{1-\alpha}$. The conclusion then follows from Lemma~\ref{T1,T2}~(ii), with $|x|$ replaced by $r_j$.
\end{proof}

It remains to prove the $L^p$-boundedness of $R_{j,1}^1$. Recall the following theorem from \cite{ACDH}.

\begin{theorem}[{\cite[Theorem~2.4]{ACDH}}]\label{thm_ACDH}
Let $(M,d,\mu)$ be a metric space satisfying the doubling condition. Suppose that $T$ is a bounded sublinear operator which is bounded on $L^2(M,\mu)$, and let $A_r$, $r>0$, be a family of linear operators acting on $L^2(M,\mu)$. Let $E_1,E_2$ be two subsets of $M$ such that $E_1\subset E_2$ and $\mu(E_2)<\infty$. Assume that
\begin{align}\label{ACDH1}
    M^{\#}_{E_2}f(\xi) := \sup_{B\ni \xi} \biggl(\frac{1}{\mu(B\cap E_2)} \int_{B\cap E_2} |T(I-A_{r(B)})f|^2 \, d\mu\biggr)^{\!\frac{1}{2}},
\end{align}
is bounded from $L^p(E_1)$ to $L^p(E_2)$ for all $p\in (2,\infty)$, and that there exists a sublinear operator $S$ bounded from $L^p(E_1)$ to $L^p(E_2)$ for all $p\in (2,\infty)$ such that
\begin{align}\label{ACDH2}
    \sup_{\eta\in B\cap E_2} |TA_{r(B)}f(\eta)| \le C\bigl( M_{E_2}(|Tf|^2)^{\frac{1}{2}}(\xi) + Sf(\xi) \bigr),
\end{align}
for all $f\in L^2$ supported in $E_1$, all balls $B\subset M$, and all $\xi\in B\cap E_2$, where $r(B)$ denotes the radius of $B$. If $2<p<\infty$ and $Tf\in L^p(E_2)$ whenever $f\in L^p(E_1)$, then $T$ is bounded from $L^p(E_1)$ to $L^p(E_2)$, and its operator norm is bounded by a constant depending only on the operator norm of $T$ on $L^2$, the $L^p(E_1)\to L^p(E_2)$ operator norms of $M_{E_2}^{\#}$ and $S$, the doubling constant $C_D$, $p$, and the constants $C$ from \eqref{ACDH1} and \eqref{ACDH2}.
\end{theorem}

\begin{remark}
In \eqref{ACDH2}, $M_E$ denotes the Hardy--Littlewood maximal operator relative to a measurable subset $E$:
\begin{equation*}
    M_E f(\xi) := \sup_{B\ni \xi} \frac{1}{\mu(B \cap E)} \int_{B\cap E} |f| \, d\mu, \qquad f\in L_{\mathrm{loc}}^1(M).
\end{equation*}
We emphasize that the assumption $\mu(E_2)<\infty$ is crucial in the proof of Theorem~\ref{thm_ACDH}; see \cite[Lemma~2.5, 2.6]{ACDH}. However, the operator norm $\|T\|_{L^p(E_1)\to L^p(E_2)}$ does not depend on the size of $E_2$, that is, on $\mu(E_2)$.
\end{remark}

\begin{proposition}\label{Rj,1,1}
We have
    \begin{equation*}
        \sup_{j\ge 1} \|R_{j,1}^1\|_{p\to p} <\infty
    \end{equation*}
    for all $2<p<\infty$.
\end{proposition}




\begin{proof}
For each $j\ge 1$, let
\begin{equation*}
    \bigl\{B_m^k = B_m \bigl(y_k, r_j^{\mathfrak a}\bigr)\bigr\}_k
\end{equation*}
be a maximal $r_j^{\mathfrak a}$-separated subset of $\mathbb{R}^m$ such that
\begin{enumerate}[label=(\roman*)]
    \item $\mathbb{R}^m = \bigcup_k B_m^k$,
    \item for $k\ne l$, $\frac{B_m^k}{2} \cap \frac{B_m^l}{2} = \emptyset$,
    \item $\sum_k \mathbf{1}_{B_m^k}(y)\le C$ for all $y\in \mathbb{R}^m$,
\end{enumerate}
where the finite overlap constant $C>0$ is independent of $r_j$.

Let $\{\mathcal{Y}_k\}_k$ be a smooth partition of unity subordinate to $\{B_m^k\}_k$. Since $R_{j,1}^1(\xi,\eta)$ is supported in the region $|y-y'|\le r_j^{\mathfrak a}$, we infer that
\begin{equation*}
    |R_{j,1}^1f| \le \sum_k \bigl| \mathbf{1}_{4B_n^j \times 4B_m^k} R_{j,1} \bigl( f \mathcal{Z}_{j,k} \bigr) \bigr|,
\end{equation*}
where
\begin{equation*}
    \mathcal{Z}_{j,k}(\eta) = \mathcal{X}_j(x') \mathcal{Y}_k(y').
\end{equation*}
Therefore, by the finite overlap property, it is enough to prove that
\begin{equation*}
    \sup_{j,k} \| R_{j,1} \|_{L^p( B_n^j \times B_m^k ) \to L^p( 4B_n^j \times 4B_m^k )} < \infty.
\end{equation*}
To this end, we apply Theorem~\ref{thm_ACDH} with $A_r = e^{-r^2 L}$, $E_1 = B_n^j \times B_m^k$, $E_2 = 4(B_n^j \times B_m^k)$, and
\begin{equation*}
    T = \int_0^{r_j^{2-2\alpha}} \nabla_L e^{-tL} \, \frac{dt}{\sqrt{t}}.
\end{equation*}
The $L^2$-boundedness of $T$ is immediate. Indeed,
\begin{align*}
    \|T\|_{2\to 2}
    &\le \bigl\| \nabla_L L^{-1/2} \bigr\|_{2\to 2} \; \biggl\| \int_0^{r_j^{2-2\alpha}} L^{1/2} e^{-tL} \, \frac{dt}{\sqrt{t}} \biggr\|_{2\to 2}\\
    &\le \sup_{\lambda >0} \int_0^{r_j^{2-2\alpha}} \lambda^{1/2} e^{-t \lambda} \, \frac{dt}{\sqrt{t}}
    \le \Gamma\Bigl(\frac{1}{2}\Bigr).
\end{align*}
It remains to verify \eqref{ACDH1} and \eqref{ACDH2}. Throughout the proof, let $B\subset \mathbb{R}^{n+m}$ be a ball of radius $r(B)=r$.

\textit{Claim 1:} \eqref{ACDH1} holds, in the sense that
\begin{equation*}
    M_{E_2}^{\#}f(\xi)^2 \le C \bigl(M_{E_2}(|f|^2)(\xi) + M f(\xi)\bigr)
\end{equation*}
for all $f\in L^2$ supported in $E_1$ and all $\xi \in E_2$.

The argument follows that of \cite[Subsection~3.2, Section~4]{ACDH}. We include it here for completeness. Let $f\in L^2$ be supported in $E_1$, and let $\xi \in B$. Write
\begin{equation*}
    f = f \mathbf{1}_{2B} + f ( 1- \mathbf{1}_{2B} ) =: f_1+f_2.
\end{equation*}
For $f_1$, the $L^2$-boundedness of $T$ and $A_r$ gives
\begin{align}\label{ACDHclaim1}
    \int_{B \cap E_2} |T(I-A_r)f_1|^2
    \le C \int_{2B \cap E_1} |f|^2
    \le C|B \cap E_2| \, M_{E_2}(|f|^2)(\xi),
\end{align}
where the last inequality follows from \eqref{D}.

Next, for $z\in B \cap E_2$, one has
\begin{align}\label{f2}
    |T(I-A_r)f_2(z)|\le \int_{\mathbb{R}^{n+m}} |K_r(z,\eta)| \, |f_2(\eta)| \, d\eta,
\end{align}
where
\begin{align*}
    K_r(z,\eta) = \int_0^\infty g_r(t) \mathbf{1}_{E_2}(z) \nabla_L e^{-tL}(z,\eta) \mathbf{1}_{E_1}(\eta) \, dt,
\end{align*}
and
\begin{equation*}
    g_r(t) = \frac{\mathbf{1}_{[0,r_j^{2-2\alpha}]}(t)}{\sqrt{t}} - \frac{\mathbf{1}_{[r^2,\, r_j^{2-2\alpha}+r^2]}(t)}{\sqrt{t-r^2}}.
\end{equation*}
Since $z\in E_2$, Lemma~\ref{Grushin_gradient} implies
\begin{align}\label{Kr}
    |K_r(z,\eta)|
    &\le C \mathbf{1}_{E_2}(z) \mathbf{1}_{E_1}(\eta) \int_0^{r_j^{2-2\alpha}} |g_r(t)| \exp\!\Bigl( - \frac{d(z,\eta)^2}{c t} \Bigr) \frac{dt}{\sqrt{t} \, V(z,\sqrt{t})}\nonumber\\
    &\quad + C \mathbf{1}_{E_2}(z) \mathbf{1}_{E_1}(\eta) \int_{r_j^{2-2\alpha}}^{r_j^{2-2\alpha}+r^2} r_j^{\alpha-1} |g_r(t)| \exp\!\Bigl(-\frac{d(z,\eta)^2}{c t}\Bigr) \frac{dt}{V(z,\sqrt{t})}.
\end{align}
Now $f_2$ is supported in $(2B)^c$, and therefore $d(z,\eta)\ge r$. By \cite[Lemma~3.1]{CD},
\begin{equation}\label{CD}
    \int_{\mathbb{R}^{n+m}} \exp\!\Bigl( - \frac{d(z,\eta)^2}{c t} \Bigr) |v(\eta)| \, d\eta \le C V(z,\sqrt{t}) \exp\!\Bigl( -\frac{r^2}{c t} \Bigr) \inf_{a\in B(z,r)} Mv(a).
\end{equation}
Combining \eqref{CD}, \eqref{Kr}, and \eqref{f2}, we obtain
\begin{align*}
    |T(I-A_r)f_2(z)|
    &\le C Mf(\xi) \biggl( \int_0^{r_j^{2-2\alpha}} |g_r(t)| \exp\!\Bigl(-\frac{r^2}{c t}\Bigr) \frac{dt}{\sqrt{t}} \\
    &\qquad\qquad\qquad\qquad + \int_{r_j^{2-2\alpha}}^{r_j^{2-2\alpha}+r^2} |g_r(t)| \exp\!\Bigl(-\frac{r^2}{c t}\Bigr) \frac{dt}{r_j^{1-\alpha}} \biggr)\\
    &:= C Mf(\xi) (I_1+I_2).
\end{align*}
It remains to verify that both integrals are finite.

\textit{Case 1: $r\le r_j^{1-\alpha}$.} In this case, split $I_1 = \int_0^{r^2}+\int_{r^2}^{r_j^{2-2\alpha}}$. Then
\begin{equation*}
    \int_0^{r^2} \dots = \int_0^{r^2} \exp\!\Bigl( - \frac{r^2}{c t} \Bigr) \frac{dt}{t} = C \int_1^\infty e^{-s} \, \frac{ds}{s}<\infty.
\end{equation*}
Moreover,
\begin{align*}
    \int_{r^{2}}^{r_j^{2-2\alpha}} \dots
    &= \int_{r^{2}}^{r_j^{2-2\alpha}} \biggl( \frac{1}{\sqrt{t}} - \frac{1}{\sqrt{t-r^2}} \biggr) \exp\!\Bigl( - \frac{r^2}{c t} \Bigr) \frac{dt}{\sqrt{t}}\\
    &\le C\int_0^1 \Bigl| 1- \frac{1}{\sqrt{1-s}} \Bigr| e^{-s} \, \frac{ds}{s} < \infty.
\end{align*}
For $I_2$,
\begin{align*}
    I_2
    = r_j^{\alpha-1}\int_{r_j^{2-2\alpha}}^{r_j^{2-2\alpha}+r^2} \frac{1}{\sqrt{t-r^2}} \exp\!\Bigl( - \frac{r^2}{c t} \Bigr) \, dt
    \le r_j^{\alpha-1} \int_{0}^{r_j^{2-2\alpha}} \frac{ds}{\sqrt{s}}<\infty.
\end{align*}

\textit{Case 2: $r\ge r_j^{1-\alpha}$.} In this case,
\begin{align*}
    I_1 \le \int_0^{r_j^{2-2\alpha}} \exp\!\Bigl( - \frac{r^2}{c t} \Bigr) \frac{dt}{t}
    \le \int_0^{r^2} \exp\!\Bigl( - \frac{r^2}{c t} \Bigr) \frac{dt}{t}<\infty,
\end{align*}
and
\begin{align*}
    I_2
    = r_j^{\alpha-1} \int_{r^2}^{r^2+r_j^{2-2\alpha}} \exp\!\Bigl( - \frac{r^2}{c t} \Bigr)\frac{dt}{\sqrt{t-r^2}}
    \le r_j^{\alpha-1} \int_{0}^{r_j^{2-2\alpha}} \frac{ds}{\sqrt{s}}<\infty.
\end{align*}
Hence, for all $z\in B\cap E_2$,
\begin{equation*}
    |T(I-A_r)f_2(z)| \le C Mf(\xi).
\end{equation*}
Combining this with \eqref{ACDHclaim1}, \textit{Claim}~1 follows.

\textit{Claim 2:} \eqref{ACDH2} holds, in the sense that
\begin{equation*}
    \sup_{z\in B\cap E_2} |TA_{r(B)}f(z)| \le C\Bigl( M_{E_2}(|Tf|^2)^{\frac{1}{2}}(\xi) + Sf(\xi) \Bigr),
\end{equation*}
where
\begin{equation*}
    Sf = M_2( Mf ) + r_j^{\alpha-1} Mh + r_j^{\alpha-1} h
\end{equation*}
and
\begin{equation*}
    h = \int_0^{r_j^{2-2\alpha}}e^{-tL}f \, \frac{dt}{\sqrt{t}}.
\end{equation*}
Note that $TA_r f = |\nabla_L e^{-r^2L}h|$ and that $|\nabla_L h| = Tf$. Let $\xi,z \in B \cap E_2$. Set
\begin{equation*}
    C_\delta = 2^{-7-\delta} \vee 2^{2-(\delta+9)(\mathfrak a)}.
\end{equation*}
By Lemma~\ref{lemma_remotecover} and Lemma~\ref{le_volume},
\begin{align}\label{B}
    E_2
    &\subset B_n\bigl( x_j, 2^{-7-\delta}|x_j| \bigr) \times B_m\bigl( y_k, 2^{2-(\delta+9)(\mathfrak a)} |x_j|^{\mathfrak a} \bigr)\nonumber\\
    &\subset B\bigl( (x_j,y_k), C_\delta c_1^{-1} |x_j|^{1-\alpha} \bigr):= \widetilde{B}.
\end{align}
By the choice of $\delta$ in Remark~\ref{remark_delta}, it is easy to verify that both $9\widetilde{B}$ and $6\gamma \widetilde{B}$ are remote.

\textit{Case 1: $r\le r_j^{1-\alpha}$.} By stochastic completeness (\cite[Lemma]{RS}),
\begin{align}
    |\nabla_Le^{-r^2L}h(z)|
    &= |\nabla_Le^{-r^2L}( h(z) - h(\xi) )| \nonumber\\
    &\le \int_{\mathbb{R}^{n+m}} |\nabla_L e^{-r^2L}(z,\eta)| \, |h(\eta)-h(\xi)| \, d\eta\nonumber\\
    &\le \int_{3\widetilde{B}} |\nabla_L e^{-r^2L}(z,\eta)| \, |h(\eta)-h(\xi)| \, d\eta
      + \int_{(3\widetilde{B})^c} |\nabla_L e^{-r^2L}(z,\eta)| \, |h(\eta)-h(\xi)| \, d\eta \nonumber\\
    &:= I_1 + I_2.
\end{align}
By Lemma~\ref{PE} and Lemma~\ref{le_tele}, for $\eta \in 3\widetilde{B}$,
\begin{equation*}
    |h(\eta)-h(\xi)| \le C d(\xi,\eta)\bigl(M_2(|\nabla_L h|\mathbf{1}_{6\gamma \widetilde{B}})(\xi)+M_2(|\nabla_L h|\mathbf{1}_{6\gamma \widetilde{B}})(\eta)\bigr).
\end{equation*}
Hence, by Lemma~\ref{Grushin_gradient},
\begin{align*}
    I_1
    &\le \frac{C}{V(z,r)} \int_{2\widetilde{B}} \frac{d(\xi,\eta)}{r} \exp\!\Bigl( - \frac{d(z,\eta)^2}{c r^2} \Bigr) \bigl( M_2(|\nabla_L h|\mathbf{1}_{6\gamma \widetilde{B}})(\xi)+M_2(|\nabla_L h|\mathbf{1}_{6\gamma \widetilde{B}})(\eta) \bigr) \, d\eta.
\end{align*}
Since $d(\xi,\eta)\le d(z,\eta) + 2r$, the factor $\frac{d(\xi,\eta)}{r}$ can be absorbed into the exponential. Thus $I_1$ is bounded by a constant multiple of
\begin{align*}
    \frac{M_2(|\nabla_L h|\mathbf{1}_{6\gamma \widetilde{B}})(\xi)}{V(z,r)}
    \int_{\mathbb{R}^{n+m}} \exp\!\Bigl( - \frac{d(z,\eta)^2}{c' r^2} \Bigr) \, d\eta
    + \frac{1}{V(z,r)} \int_{\mathbb{R}^{n+m}}  \exp\!\Bigl( - \frac{d(z,\eta)^2}{c'r^2} \Bigr) M_2(|\nabla_L h|\mathbf{1}_{6\gamma \widetilde{B}})(\eta) \, d\eta.
\end{align*}
By \cite[Lemma~3.1]{CD}, the first term is bounded by $M_2(|\nabla_L h|\mathbf{1}_{6\gamma \widetilde{B}})(\xi)$. For the second term, if $\frac{d(\xi,\eta)}{r}\le 4$, then
\begin{equation*}
    \exp\!\biggl(-\frac{d(z,\eta)^2}{c'r^2} \biggr) \le 1 \le e^{16} \exp\!\biggl(-\frac{d(\xi,\eta)^2}{r^2} \biggr).
\end{equation*}
On the other hand, if $\frac{d(\xi,\eta)}{r} > 4$, then $\frac{d(z,\eta)}{r}\ge \frac{d(\xi,\eta)}{2r}$, and therefore
\begin{equation*}
    \exp\!\biggl(-\frac{d(z,\eta)^2}{c'r^2} \biggr) \le  \exp\!\biggl(-\frac{d(\xi,\eta)^2}{4c' r^2} \biggr).
\end{equation*}
Consequently,
\begin{equation}
    I_1 \le C \bigl(M_2(|\nabla_L h|\mathbf{1}_{6\gamma \widetilde{B}})(\xi) + M M_2( |\nabla_L h| \mathbf{1}_{6\gamma \widetilde{B}} )(\xi) \bigr)\le C M_2(|\nabla_L h|\mathbf{1}_{6\gamma \widetilde{B}})(\xi),
\end{equation}
where the last inequality follows from \cite{CR}.

Recall that
\begin{equation*}
    M_2(|\nabla_L h|\mathbf{1}_{6\gamma \widetilde{B}})(\xi) = \sup_{\Bar{B} \ni \xi} \biggl( \frac{1}{|\Bar{B}|} \int_{\Bar{B} \cap 6\gamma \widetilde{B}}  |\nabla_L h(\tau)|^2 \, d\tau \biggr)^{\!\frac{1}{2}}.
\end{equation*}
Let $\Bar{B}\ni \xi$. First,
\begin{align}\label{Rkey11}
    \frac{1}{|\Bar{B}|}\int_{\Bar{B}\cap E_2} |\nabla_L h(\tau)|^2 \mathbf{1}_{6\gamma \widetilde{B}} \, d\tau
    &\le  \frac{|\Bar{B}\cap E_2|}{|\Bar{B}|} M_{E_2}\bigl(|\nabla_Lh|^2\bigr)(\xi)\nonumber\\
    &\le M_{E_2}\bigl(|\nabla_Lh|^2\bigr)(\xi).
\end{align}
Next,
\begin{align}\label{Rkey}
    \int_{\Bar{B} \setminus E_2} |\nabla_L h(\tau)|^2 \mathbf{1}_{6\gamma \widetilde{B}} \, d\tau
    \le \int_{(\Bar{B} \cap 6\gamma \widetilde{B}) \setminus E_2} \biggl| \int_0^{r_j^{2-2\alpha}} \int_{E_1} |\nabla_L e^{-tL}(\tau, \eta)| \, |f(\eta)| \, d\eta \, \frac{dt}{\sqrt{t}} \biggr|^2 d\tau.
\end{align}
Since $\tau \in 6 \gamma \widetilde{B} \setminus E_2$ and $\eta \in E_1$, one has $d(\tau,\eta)\ge c r_j^{1-\alpha}$. Therefore Lemma~\ref{Grushin_gradient} yields
\begin{align*}
    \eqref{Rkey}
    &\le C \int_{(\Bar{B} \cap 6\gamma \widetilde{B}) \setminus E_2} \biggl| \int_0^{r_j^{2-2\alpha}} \int_{d(\tau,\eta)\ge c r_j^{1-\alpha}} \frac{1}{V(\tau,\sqrt{t})} \exp\!\Bigl(-\frac{d(\tau,\eta)^2}{c t}\Bigr) |f(\eta)| \, d\eta \, \frac{dt}{t} \biggr|^2 d\tau\\
    &\le C \int_{(\Bar{B} \cap 6\gamma \widetilde{B}) \setminus E_2} \biggl| \int_0^{r_j^{2-2\alpha}} \exp\!\Bigl(-\frac{r_j^{2-2\alpha}}{c t}\Bigr) Mf(\tau) \, \frac{dt}{t} \biggr|^2 d\tau\\
    &\le C \int_{\Bar{B}} Mf(\tau)^2 \, d\tau.
\end{align*}
Combining this with \eqref{Rkey11}, we conclude that
\begin{equation}\label{est_I1}
    I_1\le C M_2(|\nabla_L h|\mathbf{1}_{6\gamma \widetilde{B}})(\xi) \le C \bigl( M_{E_2}(|\nabla_L h|^2)^{\frac{1}{2}}(\xi) +  M_2( Mf )(\xi) \bigr).
\end{equation}

For $I_2$, since $\eta \in (3 \widetilde{B})^c$, one has $d(z,\eta)\ge c r_j^{1-\alpha}$. Hence
\begin{align*}
    I_2
    &\le \frac{C}{r V(z,r)} \int_{d(z,\eta)\ge 2r(\widetilde{B})} \exp\!\Bigl(-\frac{d(z,\eta)^2}{c r^2}\Bigr) |h(\eta)| \, d\eta\\
    &\quad + \frac{C|h(\xi)|}{r V(z,r)} \int_{d(z,\eta)\ge 2r(\widetilde{B})} \exp\!\Bigl(-\frac{d(z,\eta)^2}{c r^2}\Bigr) \, d\eta\\
    &\le C\frac{r_j^{1-\alpha}}{r} \exp\!\Bigl(-\frac{r_j^{2-2\alpha}}{c r^2}\Bigr) \frac{Mh(\xi)}{r_j^{1-\alpha}} + C\frac{r_j^{1-\alpha}}{r} \exp\!\Bigl(-\frac{r_j^{2-2\alpha}}{c r^2}\Bigr) \frac{|h(\xi)|}{r_j^{1-\alpha}}\\
    &\le C r_j^{\alpha-1}\bigl( Mh(\xi) + |h(\xi)| \bigr),
\end{align*}
where the second inequality follows from \cite[Lemma~3.1]{CD} and the fact that $r(\widetilde{B}) \sim r_j^{1-\alpha}$.

Summarizing, in the case $r\le r_j^{1-\alpha}$, we have for all $z,\xi \in B \cap E_2$,
\begin{equation}\label{rsmall}
    |TA_rf(z)| \le C \bigl( M_{E_2}(|Tf|^2)^{\frac{1}{2}}(\xi) +  M_2( Mf )(\xi) + r_j^{\alpha-1} Mh(\xi) + r_j^{\alpha-1} |h(\xi)| \bigr).
\end{equation}

\textit{Case 2: $r\ge r_j^{1-\alpha}$.} In this case, Lemma~\ref{Grushin_gradient} directly gives, for all $z,\xi\in B \cap E_2$,
\begin{align*}
    |\nabla_L e^{-r^2L}h(z)|
    &\le \frac{C r_j^{\alpha-1}}{V(z,r)}  \int_{\mathbb{R}^{n+m}} \exp\!\Bigl( - \frac{d(z,\eta)^2}{c r^2} \Bigr) |h(\eta)| \, d\eta\\
    &\le \frac{C r_j^{\alpha-1}}{V(z,r)}  \int_{\mathbb{R}^{n+m}} \exp\!\Bigl( - \frac{d(\xi,\eta)^2}{c' r^2} \Bigr) |h(\eta)| \, d\eta\\
    &\le C r_j^{\alpha-1} \frac{V(\xi,r)}{V(z,r)} Mh(\xi)\\
    &\le C r_j^{\alpha-1} Mh(\xi),
\end{align*}
where the last inequality follows from the doubling property. Together with \eqref{rsmall}, this proves \textit{Claim}~2.

Finally, by \textit{Claim}~1, the operator $M_{E_2}^{\#}$ is bounded from $L^p(E_1)$ to $L^p(E_2)$ for all $2<p<\infty$. Moreover, the sublinear operator $S$ satisfies
\begin{equation*}
    \|Sf\|_p \le \| M_2( Mf ) \|_p + r_j^{\alpha-1} \| Mh \|_p + r_j^{\alpha-1} \|h\|_p \le C \|f\|_p
\end{equation*}
for all $2<p<\infty$, where the last inequality follows from the maximal theorem and the contractivity of the heat semigroup, namely
\begin{equation*}
    \|h\|_p \le C r_j^{1-\alpha}\|f\|_p.
\end{equation*}
Therefore, Theorem~\ref{thm_ACDH} applies, and the proof of Proposition~\ref{Rj,1,1} is complete.
\end{proof}






\section{Proof of Theorem~\ref{thm_main2}}

\subsection{Hardy inequality}

We first establish the following weighted Hardy inequality.


\begin{theorem}\label{hardy_x}
Let $L$ be the Grushin-type operator with parameters $n,m\ge 1$, $\alpha \in (0,1)$,
and $\beta \ge 0$, as defined in \eqref{Grushin}, and let $\nabla_L$ be the associated
gradient operator given in \eqref{Grushin2}. Then the following Hardy-type inequality
holds:
\begin{equation}\label{eq_hardy_x}
    \int_{\mathbb R^{n+m}}
    \frac{|u(\xi)|^p}{|x|^{p(1-\alpha)}}\,d\xi
    \leq
    C
    \int_{\mathbb R^{n+m}}
    |\nabla_L u(\xi)|^p \, d\xi,
\end{equation}
for all
\[
    1<p<\frac{n}{1-\alpha}
\]
and all $u\in C_c^\infty(\mathbb R^{n+m})$.
\end{theorem}

\begin{proof}
The proof follows the idea in \cite{ambrosio}. Because of the singular set
$\{x=0\}$, we use an $\varepsilon$-regularization argument.

For $\varepsilon>0$, define
\begin{itemize}
    \item $r_\varepsilon=\bigl(\varepsilon^2+|x|^2\bigr)^{\frac12}$,
    \item $\nabla_L^\varepsilon=\bigl(r_\varepsilon^\alpha\nabla_x,\,
    r_\varepsilon^\beta\nabla_y\bigr)$,
    \item $\operatorname{div}_L^\varepsilon(h)
    =
    \operatorname{div}_x(r_\varepsilon^\alpha h_x)
    +
    \operatorname{div}_y(r_\varepsilon^\beta h_y)$, where $h=(h_x,h_y)$ is a sufficiently smooth vector field.
\end{itemize}
If $\operatorname{div}_L^\varepsilon(h)>0$, then Hölder's inequality
(see \cite[Theorem~3.5]{ambrosio}) yields, for any
$u\in C_c^\infty(\mathbb R^{n+m})$,
\begin{equation}\label{divergence_x}
    \int_{\mathbb R^{n+m}}
    |u|^p\operatorname{div}_L^\varepsilon(h)\,d\xi
    \leq
    C
    \int_{\mathbb R^{n+m}}
    |h|^p
    \bigl(
        \operatorname{div}_L^\varepsilon(h)
    \bigr)^{1-p}
    |\nabla_L^\varepsilon u|^p \, d\xi.
\end{equation}
Choose
\begin{equation}\label{choice_h_x}
    h_\varepsilon(\xi)
    =
    \bigl(
        r_\varepsilon^{\alpha(p-1)-p}x,\,
        0
    \bigr).
\end{equation}
We claim that $\operatorname{div}_L^\varepsilon(h_\varepsilon)>0$. Indeed, a direct computation gives
\begin{align*}
    \operatorname{div}_L^\varepsilon(h_\varepsilon)
    &=
    \operatorname{div}_x
    \bigl(
        r_\varepsilon^\alpha r_\varepsilon^{\alpha(p-1)-p}x
    \bigr)
    =r_\varepsilon^{-p(1-\alpha)}
    \Bigl[
        n
        -
        p(1-\alpha)\frac{|x|^2}{r_\varepsilon^2}
    \Bigr].
\end{align*}
Since $0\leq \frac{|x|^2}{r_\varepsilon^2}\leq 1$, it follows that
\begin{equation}\label{>0_x}
    n
    -
    p(1-\alpha)\frac{|x|^2}{r_\varepsilon^2}
    \geq
    n-p(1-\alpha)>0,
\end{equation}
where the last inequality follows from the assumption $p<\frac{n}{1-\alpha}$. Consequently,
\begin{equation}\label{div_lower_x}
    \operatorname{div}_L^\varepsilon(h_\varepsilon)
    \geq
    C r_\varepsilon^{-p(1-\alpha)}
    >0,
\end{equation}
where $C>0$ is independent of $\varepsilon$.

We next estimate the right-hand side of \eqref{divergence_x}. First,
\[
    |h_\varepsilon|^p
    =
    r_\varepsilon^{p(\alpha(p-1)-p)}|x|^p.
\]
Second, using the lower bound for $\operatorname{div}_L^\varepsilon(h_\varepsilon)$, we obtain
\begin{align*}
    |h_\varepsilon|^p
    \bigl(
        \operatorname{div}_L^\varepsilon(h_\varepsilon)
    \bigr)^{1-p}
    &\leq
    C
    r_\varepsilon^{p(\alpha(p-1)-p)}
    |x|^p
    \bigl(
        r_\varepsilon^{-p(1-\alpha)}
    \bigr)^{1-p}                                      \\
    &=
    C
    |x|^p
    r_\varepsilon^{p(\alpha(p-1)-p)+p(1-\alpha)(p-1)}.
\end{align*}
A simple calculation shows that
\[
    p(\alpha(p-1)-p)+p(1-\alpha)(p-1)= -p.
\]
Therefore,
\begin{equation}\label{RHS_upper_x}
    |h_\varepsilon|^p
    \bigl(
        \operatorname{div}_L^\varepsilon(h_\varepsilon)
    \bigr)^{1-p}
    \leq
    C\frac{|x|^p}{r_\varepsilon^p}
    \leq C,
\end{equation}
where $C>0$ is independent of $\varepsilon$. Combining \eqref{divergence_x}, \eqref{div_lower_x}, and \eqref{RHS_upper_x}, we arrive at
\begin{equation}\label{hardy_x_final}
    \int_{\mathbb R^{n+m}}
    |u|^p r_\varepsilon^{-p(1-\alpha)}\,d\xi
    \leq
    C
    \int_{\mathbb R^{n+m}}
    |\nabla_L^\varepsilon u|^p \, d\xi,
\end{equation}
where $C>0$ is independent of $\varepsilon$.

Letting $\varepsilon\to0$, one has $r_\varepsilon\to |x|$ pointwise away from $\{x=0\}$. Applying Fatou's lemma to the left-hand side of
\eqref{hardy_x_final}, and the dominated convergence theorem to the right-hand side of
\eqref{hardy_x_final}, noting that $u$ has compact support, we obtain
\[
    \int_{\mathbb R^{n+m}}
    \frac{|u(\xi)|^p}{|x|^{p(1-\alpha)}}\,d\xi
    \leq
    C
    \int_{\mathbb R^{n+m}}
    |\nabla_L u(\xi)|^p \, d\xi.
\]
This proves \eqref{eq_hardy_x}.
\end{proof}

\begin{remark}\label{re:Hardy}
For later use, we record the following extension of the Hardy inequality. Assume $n=1$ and $1/2\leq\alpha<1$, and let $u$ be a distribution such that 
\begin{enumerate}[label=(\roman*)] 
    \item $|u(\xi)| \to 0$ as $d(\xi,0) \to \infty$,
    \item $|\nabla_L u| \in L^p(\mathbb{R}^{1+m})$ for some $1<p<\frac{1}{1-\alpha}$.
\end{enumerate}
Then we claim that the Hardy inequality \eqref{eq_hardy_x} also holds for $u$. Indeed, we prove the estimate on $\Omega_+:=(0,\infty)\times\mathbb R^m$; the
argument on $\Omega_-:=(-\infty,0)\times\mathbb R^m$ is identical. Change variables via $r:=x^{1-\alpha}$ and set $U(r,y):= u(r^{(1-\alpha)^{-1}},y)$. For a.e.\ $y$, the function $U(\cdot,y)$ is locally absolutely continuous on
$(0,\infty)$. Moreover, condition (i) implies $U(r,y)\to 0$ as $r\to \infty$.
Since $1<p<(1-\alpha)^{-1}$, the one-dimensional Hardy inequality (see, for example, \cite[Example~1.2.8]{G}) gives
\begin{equation*}
\int_0^\infty |U(r,y)|^p r^{\frac{1}{1-\alpha}-1-p} \, dr \le C_{p,\alpha} \int_0^\infty |\partial_r U(r,y)|^p r^{\frac{1}{1-\alpha}-1} \, dr.
\end{equation*}
Equivalently,
\begin{equation*}
\int_0^\infty \frac{|u(x,y)|^p}{x^{p(1-\alpha)}} \, dx \le C_{p,\alpha} \int_0^\infty x^{\alpha p}|\partial_x u(x,y)|^p \, dx. 
\end{equation*}
Integrating in $y$ yields
\begin{equation*}
\Bigl\| \frac{u}{x^{1-\alpha}} \Bigr\|_{L^p(\Omega_+)} \le C_{p,\alpha} \|x^\alpha\partial_xu\|_{L^p(\Omega_+)} \le C_{p,\alpha} \| \nabla_L u\|_{L^p(\Omega_+)}.
\end{equation*}
The same estimate on $\Omega_-$ gives \eqref{eq_hardy_x}.
\end{remark}



\begin{remark}
It is also worth mentioning that the $L^2$ Hardy and Rellich inequalities for local Dirichlet forms and weighted second-order operators, including certain Grushin-type examples, were studied by Robinson in the quadratic form setting; see \cite{Robinson_Hardy}. However, the form we needed later is an $L^p$ inequality for the intrinsic Grushin gradient, valid in the sharp range. 
\end{remark}








\subsection{Harmonic annihilation}
We next prove Theorem~\ref{thm_main2} using the so-called harmonic annihilation method, developed from \cite{H2}, to eliminate the dangerous harmonic leading term arising in the regime $\mathcal{D}_3$. In fact, by \cite{RS}, \eqref{R_p} holds for all $1<p\le 2$, and therefore, by duality, \eqref{RR_p} holds for all $p\ge 2$. 







\begin{theorem}\label{prop_RR}
Let $\alpha\in (0,1)$, $\beta \ge 0$, and $n,m\ge 1$. The reverse Riesz transform associated to $L$ is bounded on $L^p$ for all 
\begin{equation}\label{p3}
   p\in \begin{cases}
        \bigl(1,\frac{1}{1-\alpha}\bigr) \cup \bigl(\frac{1}{1-\alpha},\infty \bigr), & n=1 \text{ and } \alpha \in \bigl(0,\frac{1}{2}\bigr),\\[4pt]
        (1,\infty), & n\ge 2 \text{ or } n=1 \text{ and } \alpha\in \bigl[\frac{1}{2},1\bigr),
    \end{cases}
\end{equation}
and all $f\in C_c^\infty(\mathbb{R}^{n+m})$.
\end{theorem}

\begin{proof}
By Theorem~\ref{thm_main1} and duality, \eqref{RR_p} holds for all $1<p<\infty$ when $n\ge 2$. Also, Theorem~\ref{thm:n=1:alpha<1/2} gives \eqref{R_p} for $1<p<\alpha^{-1}$. Hence, by duality, \eqref{RR_p} holds for $p>(1-\alpha)^{-1}$ when $n=1$ and $\alpha\in (0,\frac{1}{2})$. Moreover, \cite[Theorem~8.1]{RS} yields \eqref{R_p} for $1<p\le 2$, which in turn implies \eqref{RR_p} for $p\ge 2$. Therefore, it remains to prove the reverse inequality for $1<p<\frac{n}{1-\alpha}$, namely,
\begin{equation*}
    \| L^{1/2}f\|_p \le C \| \nabla_L f\|_p,\qquad 1<p<\frac{n}{1-\alpha}.
\end{equation*}

Let $\vartheta_\varepsilon$ ($\varepsilon>0$) be a smooth function on $\mathbb{R}^n$ such that
\begin{itemize}
    \item $\operatorname{supp}(\vartheta_\varepsilon) \subset B_n(0,4 \varepsilon/3)$,
    \item $\vartheta_\varepsilon=1$ on $B_n(0,\varepsilon)$,
    \item $\|\vartheta_\varepsilon\|_\infty + \varepsilon \|\nabla \vartheta_\varepsilon\|_\infty \lesssim 1$.
\end{itemize}
Then $\vartheta_{\kappa^{-1} |x'|}(x)$ is a smooth function supported in the region
\begin{equation*}
    \{(x,x')\in \mathbb{R}^n \times \mathbb{R}^n : 4|x'|\ge 3\kappa |x|\}.
\end{equation*}
We decompose the Riesz kernel as
\begin{align}\label{eq_Grushin_1}
    \mathcal{R}(\xi,\eta) = \int_0^\infty \vartheta_{\kappa^{-1} |x'|}(x)& \nabla_L e^{-tL}(\xi,\eta) \, \frac{dt}{\sqrt{\pi t}} \nonumber\\
    &+ \int_0^\infty \bigl[1 - \vartheta_{\kappa^{-1} |x'|}(x) \bigr] \nabla_L e^{-tL}(\xi,\eta) \, \frac{dt}{\sqrt{\pi t}}
    := \mathcal{S}_1(\xi,\eta) + \mathcal{S}_2(\xi,\eta).
\end{align}
Note that $\mathcal{S}_2$ is supported in the region
\begin{align*}
    \{(\xi,\eta)\in \mathbb{R}^{n+m}\times \mathbb{R}^{n+m} \setminus \{\xi=\eta\}: |x'|\le \kappa |x|\} \subset \mathcal{D}_1 \cup \mathcal{D}_2.
\end{align*}
Hence $\mathcal{S}_2(\xi,\eta)$ is bounded by $|\mathcal{R}_1(\xi,\eta)| + |\mathcal{R}_2(\xi,\eta)|$, up to a constant.

Let $1<p<\frac{n}{1-\alpha}$, let $f\in C_c^\infty(\mathbb{R}^{n+m})$, and let $g\in C_c^\infty(\mathbb{R}^{n}\setminus \{0\} \times \mathbb{R}^m)$ satisfy $\|g\|_{p'}=1$. Since $C_c^\infty(\mathbb{R}^{n}\setminus \{0\} \times \mathbb{R}^m)$ is dense in $L^{p'}(\mathbb{R}^{n+m})$, duality reduces the problem to proving
\begin{align*}
    \bigl| \langle L^{1/2}f, g \rangle \bigr| \le C \|\nabla_L f\|_p \|g\|_{p'},\qquad \forall\, 1<p<\frac{n}{1-\alpha},
\end{align*}
for some constant $C>0$ independent of $g$.

We use the identity
\begin{equation*}
    L^{1/2}f
    =
    \frac{1}{\sqrt{\pi}}
    \int_0^\infty
    Le^{-tL}f\,\frac{dt}{\sqrt t}.
\end{equation*}
Strictly speaking, this formula is understood in the truncated sense: for
$0<t_1<t_2<\infty$, the integral on the right-hand side is interpreted as
$\lim_{\substack{t_1\to 0\\ t_2 \to \infty}} \int_{t_1}^{t_2}$; see \cite{AC} for details. However, all estimates below are uniform in $t_1$ and $t_2$. For simplicity of notation, we continue to write the integral over $(0,\infty)$.

Next, by positivity and self-adjointness of $L$,
\begin{align*}
    \langle L^{1/2}f, g \rangle
    = \biggl\langle \int_0^\infty L e^{-tL}f \, \frac{dt}{\sqrt{\pi t}}, g \biggr\rangle
    &= \bigl\langle \nabla_L f, \int_0^\infty \nabla_L e^{-tL}g \, \frac{dt}{\sqrt{\pi t}} \bigr\rangle\\
    &= \langle \nabla_L f, \mathcal{S}_1g \rangle + \langle \nabla_L f, \mathcal{S}_2 g \rangle.
\end{align*}
It follows from Proposition~\ref{R1} and Proposition~\ref{R2} that for all $p\in \bigl(1, \frac{n}{1-\alpha}\bigr)$,
\begin{align}\label{S2}
|\langle \nabla_L f, \mathcal{S}_2 g \rangle|
\le \|\nabla_L f\|_p \|\mathcal{S}_2g\|_{p'}
\le C \|\nabla_L f\|_p \|\mathcal{R}_1g+\mathcal{R}_2g\|_{p'}
\le C \|\nabla_L f\|_p \|g\|_{p'}.
\end{align}
Consider the bilinear form
\begin{equation*}
    \mathcal{B}(f,g) := \int_{\mathbb{R}^{n+m}} \nabla_L f(\xi) \cdot \int_{\mathbb{R}^{n+m}} \int_0^\infty \vartheta_{\kappa^{-1} |x'|}(x) \nabla_L e^{-tL}(\xi,\eta) g(\eta)\,\frac{dt}{\sqrt{\pi t}}\,d\eta\,d\xi.
\end{equation*}
Integration by parts gives
\begin{align}\label{eq:HA}
    \mathcal{B}(f,g)
    &= \int_{\mathbb{R}^{n+m}} f(\xi) \int_{\mathbb{R}^{n+m}} \int_0^\infty \vartheta_{\kappa^{-1} |x'|}(x) L_{\xi} e^{-tL}(\xi,\eta) g(\eta) \, \frac{dt}{\sqrt{\pi t}} \, d\eta \, d\xi \nonumber\\
    &\quad - \int_{\mathbb{R}^{n+m}} f(\xi) \int_{\mathbb{R}^{n+m}} \int_0^\infty (\nabla_L)_x [\vartheta_{\kappa^{-1} |x'|}(x)] \cdot \nabla_L e^{-tL}(\xi,\eta) g(\eta) \, \frac{dt}{\sqrt{\pi t}} \, d\eta \, d\xi.
\end{align}
In view of \eqref{S2}, it is enough to prove
\begin{equation*}
    |\mathcal{B}(f,g)| \le C \|\nabla_L f\|_p \|g\|_{p'},\qquad 1<p<\frac{n}{1-\alpha}.
\end{equation*}

\begin{lemma}\label{lemma_key2}
Under the assumptions of Proposition~\ref{prop_RR}, the following estimates hold:
\begin{align*}
    &\int_0^\infty  \bigl|(\nabla_L)_x [\vartheta_{\kappa^{-1} |x'|}(x)] \cdot \nabla_L e^{-tL}(\xi,\eta)\bigr| + \bigl|\vartheta_{\kappa^{-1} |x'|}(x) \partial_t e^{-tL}(\xi,\eta)\bigr| \, \frac{dt}{\sqrt{t}} \\
    &\qquad \lesssim 
    \begin{cases}
        |x'|^{-(1-\alpha)(\mathcal{Q}+1)} \mathbf{1}_{|x|\le c_1|x'|}, & |y-y'|\le |x'|^{\mathfrak a},\\[4pt]
        |x'|^{-2+2\alpha} |y-y'|^{-\frac{(1-\alpha)(\mathcal{Q}-1)}{\mathfrak a}} \mathbf{1}_{c_2|x|\le |x'| \le c_3|x|} + |y-y'|^{-\frac{(1-\alpha)(\mathcal{Q}+1)}{\mathfrak a}}\mathbf{1}_{|x|\le c_4|x'|}, & |y-y'|\ge |x'|^{\mathfrak a},
    \end{cases}
\end{align*}
for some constants $c_1,c_2,c_3,c_4>0$.
\end{lemma}

\begin{proof}
Note that if $|x'|\ge \frac{3\kappa}{4} |x|$, then $|x-x'|\sim |x'|$ provided $\kappa$ is sufficiently large. Hence, by Theorem~\ref{RS_Grushin},
\begin{align*}
    d(\xi,\eta) \sim |x'|^{1-\alpha} + \frac{|y-y'|}{|x'|^\beta + |y-y'|^{\frac{\beta}{\mathfrak a}}} := \sigma.
\end{align*}
Clearly,
\begin{align}\label{eq_delta}
    \sigma \sim \begin{cases}
        |x'|^{1-\alpha}, & |y-y'|\le |x'|^{\mathfrak a},\\[4pt]
        |y-y'|^{\frac{1-\alpha}{\mathfrak a}}, & |y-y'|\ge |x'|^{\mathfrak a}.
    \end{cases}
\end{align}
Moreover, for $x$ in the support of $\nabla_L \vartheta_{\kappa^{-1} |x'|}$, one has $|x|\sim |x'|$. In addition,
\begin{equation*}
    |\nabla_L \vartheta_{\kappa^{-1} |x'|}|\lesssim |x|^\alpha |x'|^{-1} \mathbf{1}_{|x|\sim |x'|}.
\end{equation*}
By Theorem~\ref{RS_Grushin} and Lemma~\ref{Grushin_gradient},
\begin{align*}
    \int_0^{|x|^{2-2\alpha}} \bigl|\nabla_L [\vartheta_{\kappa^{-1} |x'|}(x)]\bigr| \, \bigl|\nabla_L e^{-tL}(\xi,\eta)\bigr| \, \frac{dt}{\sqrt{t}}
    &\lesssim \int_0^{|x|^{2-2\alpha}} \frac{\mathbf{1}_{|x|\sim |x'|}}{|x'|^{1-\alpha}} \frac{e^{-\frac{\sigma^2}{c t}}}{|x|^{n\alpha+m\beta}t^{\frac{n+m}{2}}} \frac{dt}{t}\\
    &\lesssim |x'|^{\alpha-1-n\alpha-m\beta} \mathbf{1}_{|x|\sim |x'|} \int_0^{|x'|^{2-2\alpha}} t^{-\frac{n+m}{2}-1} e^{-\frac{\sigma^2}{c t}} \, dt\\
    &\lesssim |x'|^{\alpha-1-n\alpha-m\beta} \sigma^{-n-m} e^{-\frac{\sigma^2}{c|x'|^{2-2\alpha}}} \mathbf{1}_{|x|\sim |x'|}\\
    &\lesssim_\varepsilon |x'|^{\alpha-1-n\alpha -m\beta+(1-\alpha)\varepsilon} \sigma^{-n-m-\varepsilon} \mathbf{1}_{|x|\sim |x'|},\quad \forall\, \varepsilon \ge 0\\
    &\lesssim \sigma^{-\mathcal{Q}-1} \mathbf{1}_{|x|\sim|x'|},
\end{align*}
where we choose $\varepsilon = \mathcal{Q}+1-n-m>0$.

On the other hand,
\begin{align*}
    \int_{|x|^{2-2\alpha}}^\infty \bigl|\nabla_L [\vartheta_{\kappa^{-1} |x'|}(x)]\bigr| \, \bigl|\nabla_L e^{-tL}(\xi,\eta)\bigr| \, \frac{dt}{\sqrt{t}}
    &\lesssim \int_{|x|^{2-2\alpha}}^\infty \frac{\mathbf{1}_{|x|\sim |x'|}}{|x'|^{1-\alpha}} \frac{e^{-\frac{\sigma^2}{c t}}}{t^{\frac{\mathcal{Q}}{2}}} \frac{dt}{|x|^{1-\alpha}\sqrt{t}}\\
    &\lesssim |x'|^{-2+2\alpha} \mathbf{1}_{|x|\sim |x'|} \int_{|x'|^{2-2\alpha}}^\infty e^{-\frac{\sigma^2}{c t}} t^{-\frac{\mathcal{Q}+1}{2}} \, dt\\
    &\lesssim |x'|^{-2+2\alpha} \sigma^{-\mathcal{Q}+1} \mathbf{1}_{|x|\sim|x'|}.
\end{align*}
Similarly, using the time-derivative estimate for the heat kernel, namely
\begin{equation*}
    \bigl|\partial_t e^{-tL}(\xi,\eta) \bigr| \le \frac{C}{t V(\eta,\sqrt{t})} e^{-\frac{d(\xi,\eta)^2}{c t}}, \qquad \forall\, \xi,\eta\in \mathbb{R}^{n+m},\quad \forall\, t>0,
\end{equation*}
which is a standard consequence of the analyticity of the heat semigroup; see \cite{D1989} or \cite{G1}, we obtain
\begin{align*}
    \int_0^\infty \bigl|\vartheta_{\kappa^{-1} |x'|}(x)\bigr| \, &\bigl|L_{\xi} e^{-tL}(\xi,\eta) \bigr| \, \frac{dt}{\sqrt{t}} \\
    &\lesssim  \int_0^\infty \frac{t^{-\frac{3}{2}} \mathbf{1}_{|x|\le c|x'|}}{V(\eta,\sqrt{t})} e^{-\frac{d(\xi,\eta)^2}{c t}} \, dt\\
    &\lesssim \int_0^{|x'|^{2-2\alpha}} \frac{ \mathbf{1}_{|x|\le c|x'|}}{|x'|^{n\alpha+m\beta}(\sqrt{t})^{n+m}} e^{-\frac{\sigma^2}{c t}} \frac{dt}{t^{3/2}}\\
    &\quad + \int_{|x'|^{2-2\alpha}}^\infty \frac{ \mathbf{1}_{|x|\le c|x'|}}{(\sqrt{t})^{\mathcal{Q}}} e^{-\frac{\sigma^2}{c t}} \frac{dt}{t^{3/2}}\\
    &\lesssim |x'|^{-n\alpha-m\beta+(1-\alpha)\varepsilon} \sigma^{-n-m-1-\varepsilon} \mathbf{1}_{|x|\le c|x'|} + \sigma^{-\mathcal{Q}-1} \mathbf{1}_{|x|\le c|x'|},\quad \forall\, \varepsilon\ge 0\\
    &\lesssim \sigma^{-\mathcal{Q}-1} \mathbf{1}_{|x|\le c|x'|},
\end{align*}
where we choose $\varepsilon = \mathcal{Q}-n-m \ge 0$. The result follows from \eqref{eq_delta}.
\end{proof}

With Lemma~\ref{lemma_key2} in mind, define the operators
\begin{gather*}
    T_3: u \mapsto |x|^{1-\alpha} \int_{B_n(0,c|x|)^c} |x'|^{-(1-\alpha)(\mathcal{Q}+1)} \int_{B_m(y,|x'|^{\mathfrak a})} |u(x',y')| \, dy' \, dx',\\
    T_4: u \mapsto |x|^{1-\alpha} \int_{B_n(0,c|x|)^c} \int_{B_m(y,|x'|^{\mathfrak a})^c} |y-y'|^{-\frac{(1-\alpha)(\mathcal{Q}+1)}{\mathfrak a}}|u(x',y')| \, dy' \, dx',\\
    T_5: u \mapsto |x|^{1-\alpha} \int_{c_1|x|\le |x'|\le c_2|x|} |x'|^{-2+2\alpha} \int_{B_m(y,|x'|^{\mathfrak a})^c} |y-y'|^{-\frac{(1-\alpha)(\mathcal{Q}-1)}{\mathfrak a}}|u(x',y')| \, dy' \, dx'.
\end{gather*}

\begin{lemma}\label{T3T4}
$T_3$ and $T_4$ are bounded on $L^p$ for all $1<p\le \infty$.
\end{lemma}

\begin{proof}
The proof is similar to that of Lemma~\ref{T1,T2}. By dyadic decomposition,
\begin{equation*}
    |T_3 f(\xi)|\le |x|^{1-\alpha} \sum_{j\ge 0} \int_{B_n(0,2^{j+1}\kappa |x|) \setminus B_n(0, 2^j \kappa |x|)} \bigl(2^j |x| \bigr)^{-(1-\alpha)(\mathcal{Q}+1)} \int_{B_m(y,|x'|^{\mathfrak a})} |f(\eta)| \, d\eta.
\end{equation*}
The $j$th term is bounded by
\begin{align*}
    2^{-j(1-\alpha)(\mathcal{Q}+1)} |x|^{-(1-\alpha)(\mathcal{Q}+1)} \int_{ B_n(0, 2^{j+1}\kappa |x|) \times B_m (y, (2^{j+1}\kappa |x|)^{\mathfrak a}) } |f(\eta)| \, d\eta.
\end{align*}
By Lemma~\ref{le_volume},
\begin{equation*}
    B_n\bigl(0, 2^{j+1}\kappa |x|\bigr) \times B_m \bigl(y, (2^{j+1}\kappa |x|)^{\mathfrak a}\bigr) \subset  B\bigl( (0,y), c_3^{-1} 2^{(j+1)(1-\alpha)} \kappa^{1-\alpha} |x|^{1-\alpha} \bigr).
\end{equation*}
Therefore the $j$th term is bounded by
\begin{equation*}
    2^{-j(1-\alpha)} |x|^{\alpha-1} Mf(\xi),
\end{equation*}
and the $L^p$ estimate for $T_3$ follows by summing over $j\ge 0$ and applying the maximal theorem.

Next, for $T_4$, applying dyadic decomposition twice shows that $|T_4 f(\xi)|$ is bounded by
\begin{equation*}
    |x|^{1-\alpha} \sum_{i\ge 0} \sum_{j\ge 0} \int_{B_n(0,2^{i+1}\kappa |x|) \setminus B_n(0, 2^i \kappa |x|)} \int_{B_m(y,2^{j+1}|x'|^{\mathfrak a}) \setminus B_m(y, 2^j |x'|^{\mathfrak a})} \frac{|f(\eta)|\,d\eta}{|y-y'|^{\frac{(\mathcal{Q}+1)(1-\alpha)}{\mathfrak a}}}.
\end{equation*}
For each $i,j\ge 0$, one has
\begin{equation*}
    B_n\bigl(0,2^{i+1}\kappa |x|\bigr) \times B_m\bigl(y,2^{j+1}|x'|^{\mathfrak a}\bigr)  \subset B \bigl( (0,y), c_3^{-1} 2^{\frac{(j+1)(1-\alpha)}{\mathfrak a} + (i+1)(1-\alpha)} |x|^{1-\alpha} \bigr),
\end{equation*}
and an argument parallel to that for $T_3$ yields the bound
\begin{equation*}
    2^{-i(1-\alpha)} 2^{-j \frac{1-\alpha}{\mathfrak a}} Mf(\xi).
\end{equation*}
The conclusion follows at once.
\end{proof}

By Lemma~\ref{T3T4}, it remains to treat $T_5$. Since $|x|\sim |x'|$, one has $T_5 f(\xi) \lesssim T_2f(\xi)$ with $\varepsilon = 1-\alpha$; see \eqref{def_T2}. Hence, by Lemma~\ref{T1,T2}, $T_5$ is bounded on $L^q$ for $q> \frac{n}{n-1+\alpha} = \Bigl(\frac{n}{1-\alpha}\Bigr)'$. Therefore, by Lemma~\ref{lemma_key2} and Theorem~\ref{hardy_x}, the bilinear form satisfies
\begin{align}\label{eq_final}
    |\mathcal{B}(f,g)|
    &\le C \int_{\mathbb{R}^{n+m}} \frac{|f(\xi)|}{|x|^{1-\alpha}} \, \bigl|(T_3+ T_4 + T_5)g(\xi)\bigr| \, d\xi \nonumber\\
    &\le C \Bigl\| \frac{f(\xi)}{|x|^{1-\alpha}} \Bigr\|_p \sum_{j=3}^5 \|T_j g\|_{p'} \nonumber\\
    &\le C \| \nabla_L f\|_p \|g\|_{p'}
\end{align}
for all $p\in \bigl(1,\frac{n}{1-\alpha}\bigr)$, as required. This completes the proof of Theorem~\ref{prop_RR}.
\end{proof}





\begin{remark}\label{re:harmonic:annihilation}
For later use, we focus on the case $n=1$ and $1/2 \le \alpha <1$. We use the convention
\begin{equation*}
\operatorname{div}_L F = -\partial_x(|x|^\alpha F_x) +\operatorname{div}_y(|x|^\beta F_y),
\end{equation*}
where $F = (F_x, F_y)$. 
Let $u:= L^{-1} \operatorname{div}_L \omega$ with $\omega \in C_c^\infty(\mathbb R^{1+m};\mathbb C^{1+m})$, and let $g\in C_c^\infty((\mathbb R\setminus\{0\})\times\mathbb R^m)$. For $0<t_1 < t_2 < \infty$, define
\begin{equation*}
F_{t_1,t_2}(\xi) := \int_{\mathbb R^{1+m}} \int_{t_1}^{t_2} \vartheta_{\kappa^{-1} |x'|}(x) \nabla_{L,\xi}e^{-tL}(\xi,\eta) g(\eta) \, \frac{dt}{\sqrt{\pi t}} \, d\eta ,
\end{equation*}
where $e^{-tL}(\xi,\eta)$ is the heat kernel of the one-dimensional separated Friedrichs realization. We want to check that the key step \eqref{eq:HA} holds for $u$. The subtle point is that, because of the separation phenomenon, $u$ can have different traces across $x=0$, which may cause trouble in \eqref{eq:HA}.

For $\varepsilon>0$, integration by parts on $\{|x|>\varepsilon\}$ gives
\begin{align*}
\int_{|x|>\varepsilon} \nabla_L u \cdot F_{t_1,t_2} \, d\xi &= \int_{|x|>\varepsilon} u \operatorname{div}_L F_{t_1,t_2} \, d\xi\\
&\qquad + \int_{\mathbb R^m} \Bigl[ u(-\varepsilon,y)\varepsilon^\alpha (F_{t_1,t_2})_x(-\varepsilon,y) - u(\varepsilon,y)\varepsilon^\alpha (F_{t_1,t_2})_x(\varepsilon,y) \Bigr] dy.
\end{align*}
We claim that the above boundary contribution tends to zero as
$\varepsilon\downarrow0$. Indeed, the $x$-component of $F_{t_1,t_2}$ is
\begin{equation*}
(F_{t_1,t_2})_x(\xi) = \int_{\mathbb R^{1+m}} \int_{t_1}^{t_2} \vartheta_{\kappa^{-1} |x'|}(x) |x|^\alpha\partial_xp_t(\xi,\eta) g(\eta) \, \frac{dt}{\sqrt{\pi t}} \, d\eta.
\end{equation*}
Therefore
\begin{align*}
\varepsilon^\alpha (F_{t_1,t_2})_x(\pm\varepsilon,y) &= \int_{\mathbb R^{1+m}} \int_{t_1}^{t_2} \vartheta_{\kappa^{-1} |x'|}(\pm\varepsilon) \varepsilon^{2\alpha} \partial_x e^{-tL}((\pm\varepsilon,y),\eta) g(\eta) \, \frac{dt}{\sqrt{\pi t}} \, d\eta.
\end{align*}
Since $\alpha\ge 1/2$, the Friedrichs realization separates the two half-lines.
Thus $e^{-tL}(\xi,\eta)=0$ when $\xi$ and $\eta$ lie on different half-lines, and on each half-line the heat kernel satisfies the zero-flux endpoint condition
\begin{equation*}
    \lim_{\varepsilon \to 0}
    \varepsilon^{2\alpha}\partial_x e^{-tL}((\pm\varepsilon,y),\eta)=0
\end{equation*}
in the Friedrichs trace sense. Because $g$ is compactly supported away from
$\{x=0\}$ and $t\in[t_1,t_2]$, the heat-kernel estimates give a dominating
function in the $y$-variable. Hence
\begin{align*}
\lim_{\varepsilon\to 0} \int_{\mathbb R^m} \Bigl[ u(-\varepsilon,y)\varepsilon^\alpha (F_{t_1,t_2})_x(-\varepsilon,y) - u(\varepsilon,y)\varepsilon^\alpha (F_{t_1,t_2})_x(\varepsilon,y) \Bigr] dy = 0.
\end{align*}
Consequently,
\begin{equation*}
\int_{\mathbb R^{1+m}} \nabla_Lu \cdot F_{t_1,t_2} \, d\xi = \int_{\mathbb R^{1+m}} u \operatorname{div}_L F_{t_1,t_2} \, d\xi.
\end{equation*}
Letting $t_1 \to 0$ and $t_2 \to \infty$ gives \eqref{eq:HA}.
\end{remark}

\begin{remark}\label{re:applicable:HA}
Without ambiguity, this harmonic annihilation method (Theorem~\ref{prop_RR}) applies to a broader class of functions. In fact, the argument of Theorem~\ref{prop_RR} works verbatim for a function $u$ if
\begin{itemize}
    \item \eqref{eq:HA} holds for $u$,
    \item the Hardy inequality \eqref{eq_hardy_x} holds for $u$.
\end{itemize}
In particular, by Remarks~\ref{re:Hardy} and \ref{re:harmonic:annihilation} above, in the case $n=1$ and $1/2\le \alpha<1$, the reverse Riesz inequality \eqref{RR_p} holds for $u=L^{-1} \operatorname{div}_L \omega$ with $\omega \in C_c^\infty(\mathbb{R}^{1+m}; \mathbb{C}^{1+m})$ provided $u(\xi) \to 0$ as $d(\xi,0)\to \infty$ and $|\nabla_L u|\in L^p$.
\end{remark}










\part{Riesz transform for \texorpdfstring{$n=1$}{n=1} and \texorpdfstring{$\frac{1}{2}\le \alpha<1$}{1/2<alpha<1}}\label{part4}



In this final part we treat the strongly degenerate one-dimensional regime $n=1$ and $1/2\le \alpha<1$. The route from Parts~\ref{part1} and~\ref{part2} cannot be used
directly: the Poincar\'e inequality across the singular line fails, so the abstract reverse-H\"older characterization of the Riesz transform is no longer available (see Remark~\ref{re:H}). At the same time, the Friedrichs realization separates the two half-lines and selects the zero-flux branch on each of them. This separated structure allows one
to prove the required Hodge projector estimate. Combining this Hodge estimate with the reverse Riesz inequality from Part~\ref{part3} then recovers the full $L^p$ range
for the Riesz transform.



\section{Proof of Theorem~\ref{thm:n=1:alpha>1/2}}


\subsection{From reverse H\"older to Hodge projector}


Consider the Hilbert space $L^2(\mathbb R^{1+m};\mathbb C^{1+m})$. We denote by
\begin{equation*}
    \operatorname{div}_L:=\nabla_L^*
\end{equation*}
the $L^2$-adjoint of $\nabla_L$. Thus, for smooth compactly supported
$u$ and vector fields $F=(F_x,F_y)$,
\begin{equation*}
    \int_{\mathbb R^{1+m}}
     \nabla_Lu \cdot F \, d\xi =
    \int_{\mathbb R^{1+m}}
    u \operatorname{div}_L F \, d\xi ,
\end{equation*}
where explicitly
\begin{equation*}
\operatorname{div}_L F = -\partial_x(|x|^\alpha F_x) - \sum_{j=1}^m\partial_{y_j}(|x|^\beta F_{y_j}).
\end{equation*}
With this convention, $L=\operatorname{div}_L\nabla_L$ in the sense of the Friedrichs form. The Hodge projector is defined on $L^2(\mathbb{R}^{1+m}; \mathbb{C}^{1+m})$ by
\begin{equation*}
   \mathcal P:=\nabla_L L^{-1}\operatorname{div}_L =  \nabla_L L^{-1/2} \bigl(\nabla_L L^{-1/2}\bigr)^*.
\end{equation*}
Since $\nabla_L L^{-1/2}$ is a partial isometry on $L^2$, the operator $\mathcal P$ is an orthogonal projection on $L^2$. In particular,
\begin{equation*}
    \mathcal P^2= \mathcal P,\qquad \mathcal P^*= \mathcal P,\qquad \|\mathcal P\|_{2\to2}\le 1,
\end{equation*}
and we say the Hodge projector is bounded on $L^p$ if
\begin{equation}\label{Pi_p}\tag{$\Pi_p$}
    \| \nabla_L L^{-1} \mathrm{div}_L \omega \|_p \le C \| \omega \|_p.
\end{equation}
The aim of this subsection is first to prove \eqref{Pi_p} in the strongly degenerate one-dimensional regime $n=1$ and $1/2 \le \alpha <1$. Once this is obtained, we combine it with the reverse Riesz inequality from
Theorem~\ref{prop_RR} to prove Theorem~\ref{thm:n=1:alpha>1/2}.

The following property is due to Shen \cite{Shen}.

\begin{theorem}[{\cite{Shen}}]\label{thm:Shen}
Let $(E,\mu)$ be a measured metric space satisfying the doubling property (D). Let $T$ be a bounded sublinear operator on $L^2(E,\mu)$. Assume that for some $q\in (2,\infty]$, $1<\lambda_1<\lambda_2$, and $C>0$, we have
\begin{equation}\label{eq:Shen}
    \biggl( \fint_B |Tf|^q \, d\mu \biggr)^{\!\frac{1}{q}} \le C \biggl( \fint_{\lambda_1 B} |Tf|^2 \, d\mu \biggr)^{\!\frac{1}{2}}
\end{equation}
for all balls $B \subset E$ and all $f\in L^2(E,\mu)$ supported in $E\setminus \lambda_2 B$. Then $T$ is bounded on $L^p(E,\mu)$ for $2<p<q$.  
\end{theorem}

To apply Theorem~\ref{thm:Shen} to $\mathcal{P}$, we first recall the following more natural reverse H\"older inequality used in \cite{Shen}:
\begin{equation}\label{eq:RH:Shen}\tag{$\widetilde{\mathrm{RH}}_p$}
     \biggl( \fint_B |\nabla_L u|^p \, d\mu \biggr)^{\!\frac{1}{p}} \le C \biggl( \fint_{\lambda B} |\nabla_L u|^2 \, d\mu \biggr)^{\!\frac{1}{2}}.
\end{equation}

Under \eqref{D} and \eqref{UE}, the Caccioppoli inequality together with the mean value estimate \eqref{meanvalue} gives the implication $\eqref{eq:RH:Shen} \implies \eqref{eq:RH:Jiang}$. However, for the opposite direction, one usually needs \eqref{P2}. 

From our discussion in Part~\ref{part2}, we are able to establish the following.

\begin{theorem}\label{thm:RH:n1:alpha>1/2}
Let $n=1$, $m\ge1$, $\alpha \in [\frac{1}{2},1)$, and $\beta \ge0$. Then \eqref{eq:RH:Jiang} holds for all $1<p\le \infty$.
\end{theorem}



\begin{proof}
We prove the anchored estimate; the remote case follows from Lemma~\ref{RH_remote} and the reduction from arbitrary balls to remote and anchored balls is the same as in Lemma~\ref{reduce_RH}. After rescaling an anchored ball, it is enough to prove the following fixed-cylinder estimate. If $v$ is $L$-harmonic in $Q_4=(-4,4)\times B_m(0,4^{\mathfrak a})$, then
\begin{equation*}
    \sup_{Q_1}|\nabla_L v| \le C\|v\|_{L^\infty(Q_3)}.
\end{equation*}
The $L^p$ reverse H\"older inequality then follows immediately for every
$1<p\le \infty$. Let $\chi\in C_c^\infty(B_m(0,3^{\mathfrak a}))$ be such that $\chi=1$ on $B_m(0,2^{\mathfrak a})$, and set $\omega(x,y):=\chi(y)v(x,y)$. Then $\omega=v$ on $Q_2$, and
\begin{equation*}
L\omega=f:=[L,\chi]v = |x|^{2\beta}v\Delta_y\chi - 2|x|^{2\beta}\nabla_y\chi\cdot\nabla_yv.
\end{equation*}
The function $f$ is supported in the annulus $2^{\mathfrak a}\leq |y|\leq 3^{\mathfrak a}$.

Since $1/2\le \alpha<1$, the Friedrichs realization separates the two half-lines. Thus the restrictions
\begin{equation*}
    \omega_\pm(\rho,y):=\omega(\pm\rho,y),
    \qquad
    f_\pm(\rho,y):=f(\pm\rho,y),
    \qquad
    h_\pm(y):=\omega(\pm2,y)
\end{equation*}
solve independent half-line Friedrichs problems on $(0,2)\times\mathbb R^m$.
By the Friedrichs representation,
\begin{equation*}
    \omega_\pm=\mathcal T h_\pm+ \mathcal S f_\pm,
\end{equation*}
where $\mathcal T$ and $\mathcal S$ are respectively the Poisson and Green operators for the half-line problem
\begin{align*}
\begin{cases}
    Lg=0, & 0<\rho<2,\\
    g(2,y)=h_\pm(y),\\
    \displaystyle \lim_{\rho\to0}\rho^{2\alpha}\partial_\rho g(\rho,y)=0,
\end{cases}\qquad \mathrm{and}\qquad \begin{cases}
    Lg=f_\pm, & 0<\rho<2,\\
    g(2,y)=0,\\
    \displaystyle \lim_{\rho\to 0}\rho^{2\alpha}\partial_\rho g(\rho,y)=0.
\end{cases}   
\end{align*}
We now identify the Fourier--Bessel branch. Taking the Fourier transform in
the $y$-variable gives the ordinary differential equation
\begin{equation*}
-\partial_\rho\left(\rho^{2\alpha}\partial_\rho \hat g(\rho,Y)\right)+\rho^{2\beta}|Y|^2\hat g(\rho,Y)=0,\qquad 0<\rho<2.
\end{equation*}
Equivalently,
\begin{equation*}
\hat g''(\rho,Y)+\frac{2\alpha}{\rho}\hat g'(\rho,Y)-|Y|^2\rho^{2(\beta-\alpha)}\hat g(\rho,Y)=0.
\end{equation*}
Set $\nu:=\frac{\mathfrak b}{\mathfrak a}$ and $\tau_\rho:=\frac{|Y|\rho^\mathfrak a}{\mathfrak a}$. Note that $\mathfrak b\ge 0$, $\nu\ge 0$, and $\mathfrak b=\mathfrak a \nu$. With the substitution $\hat g(\rho,Y)=\rho^{-b}U(\tau_\rho)$, the equation becomes the modified Bessel equation
\begin{equation*}
    \tau^2U''(\tau)+\tau U'(\tau)-(\tau^2+\nu^2)U(\tau)=0.
\end{equation*}
Hence the two local branches are $\rho^{-\mathfrak b}I_\nu(\tau_\rho)$ and $\rho^{-\mathfrak b}K_\nu(\tau_\rho)$.

The Friedrichs finite-energy condition excludes the $K_\nu$ branch, while the regular branch
\begin{equation*}
    \Phi_Y(\rho):=\rho^{-b}I_\nu(\tau_\rho)
\end{equation*}
satisfies the zero-flux endpoint condition. Indeed, as $\rho\to 0$, $\Phi_Y(\rho)=C_Y\left(1+O(\rho^{2a})\right)$, and therefore
\begin{equation*}
\rho^{2\alpha}\Phi_Y'(\rho) = O\left(\rho^{2\alpha+2a-1}\right) = O\left(\rho^{2\beta+1}\right) \to 0.
\end{equation*}
Thus the Poisson multiplier is
\begin{equation*}
M(\rho,Y) := \frac{\Phi_Y(\rho)}{\Phi_Y(2)} = \frac{\mathcal I_\nu(\tau_\rho)}{\mathcal I_\nu(\tau_2)}, \qquad \mathcal I_\nu(t) := t^{-\nu} I_\nu(t).
\end{equation*}
For the Green operator, let
\begin{equation*}
\Psi_Y(\rho) := \rho^{-\mathfrak b} \left[ K_\nu(\tau_\rho) - \frac{K_\nu(\tau_2)} {I_\nu(\tau_2)} I_\nu(\tau_\rho) \right].
\end{equation*}
Then $\Psi_Y(2)=0$, and the half-line Green kernel is
\begin{equation*}
G_Y(\rho,\sigma) = \mathfrak a^{-1}
\begin{cases}
\Phi_Y(\rho)\Psi_Y(\sigma), & 0<\rho<\sigma<2,\\
\Phi_Y(\sigma)\Psi_Y(\rho), & 0<\sigma<\rho<2.
\end{cases}
\end{equation*}
The fixed-order estimates from Part~\ref{part2} now apply with the order $\nu\ge 0$. More precisely, the proof of the zero-flux Poisson estimate in
Proposition~\ref{prop:Poisson:n=1}, together with
Remark~\ref{re:even:Poisson}, gives
\begin{equation*}
\sup_{Q_1} |\nabla_L \mathcal T h_\pm| \le C\|h_\pm\|_{L^\infty(\mathbb R^m)} \le C\|v\|_{L^\infty(Q_3)}.
\end{equation*}
It remains to estimate the Green term. The fixed-order Green estimates obtained
in Proposition~\ref{prop:Green:n=1:even} and Theorem~\ref{thm:Green:n=1:secondtolast}
apply verbatim to the above zero-flux kernel. In particular, the row estimates
for $\nabla_L \mathcal G$ and for $\nabla_L\partial_{y_k'} \mathcal G$ (see Theorem~\ref{lem:Green_kernel}) imply
\begin{equation*}
\sup_{Q_1}|\nabla_L \mathcal S f_\pm| \le C\|v\|_{L^\infty(Q_3)}.
\end{equation*}
Indeed, the part of $f_\pm$ containing $v\Delta_y\chi$ is controlled directly by the row bound for $\nabla_L \mathcal G$. The part containing $\nabla_y\chi\cdot\nabla_yv$ is treated by integrating by parts in the
$y_k'$-variable, exactly as in the proof of Theorem~\ref{thm:RH};
the support of $\nabla_y\chi$ is separated from $Q_1$ in the $y$-variable, and
the row bound for $\nabla_L\partial_{y_k'} \mathcal G$ gives the desired estimate.

Combining the Poisson and Green estimates gives
\begin{equation*}
\sup_{Q_1}|\nabla_L\omega| \le C\|v\|_{L^\infty(Q_3)}.
\end{equation*}
Since $\omega=v$ on $Q_2$, we conclude that
\begin{equation*}
\sup_{Q_1}|\nabla_Lv| \le C\|v\|_{L^\infty(Q_3)}.
\end{equation*}
This completes the proof of Theorem~\ref{thm:RH:n1:alpha>1/2}


\end{proof}












Next, we pass from \eqref{eq:RH:Jiang} to \eqref{eq:RH:Shen}.

\begin{proposition}\label{prop_u_RH_to_Shen_RH_large_alpha}
Let $n=1$, $m\ge1$, $\alpha \in [\frac{1}{2},1)$, and $\beta \ge0$. Then \eqref{eq:RH:Shen} holds for all $1<p\le \infty$.
\end{proposition}

\begin{proof}
The key point is that although the global scale-invariant Poincar\'e inequality across the singular set $\{x=0\}$ fails, a separated $L^2$-Poincar\'e inequality remains valid. More precisely, for $1/2\le \alpha<1$, Robinson--Sikora \cite[Theorem~1.1 III]{RS2} proves that
\begin{equation}\label{eq_separated_Poincare_ball}
\int_B \bigl| f - f_{B^+}\mathbf 1_{B^+} - f_{B^-}\mathbf 1_{B^-}\bigr|^2 \, d\xi \le C r(B)^2 \int_{\lambda B}|\nabla_Lf|^2 \, d\xi
\end{equation}
for some $\lambda \ge 1$, where $B^+:= B \cap \{x>0\}$ and $B^-:=B\cap \{x<0\}$.

We first consider the case where $B$ is remote. Then the implication \eqref{eq:RH:Jiang} $\implies$ \eqref{eq:RH:Shen} is immediate. Indeed, if $B$ is remote, then necessarily $B\cap \{x=0\} = \emptyset$. Moreover, by Lemma~\ref{PE}, the $L^2$-Poincar\'e inequality holds on $B$. Now suppose that $u$ is $L$-harmonic in $\lambda B$; here $\varepsilon_0$ is chosen sufficiently small so that $\lambda B$ remains remote. Then Theorem~\ref{thm:RH:n1:alpha>1/2}, together with H\"older's inequality, yields
\begin{align*}
    \biggl( \fint_B |\nabla_L u|^p \, d\xi \biggr)^{\!\frac{1}{p}}
    &=\biggl( \fint_B |\nabla_L (u-u_{\lambda B})|^p \, d\xi \biggr)^{\!\frac{1}{p}} \\
    &\le \frac{C}{r(B)}  \fint_{\lambda B} |u-u_{\lambda B}| \, d\xi  \\
    &\le \frac{C}{r(B)} \biggl(\fint_{\lambda B} |u-u_{\lambda B}|^2 \, d\xi \biggr)^{\!\frac{1}{2}}\\
    &\le C \biggl(\fint_{\lambda' B} |\nabla_L u|^2 \, d\xi \biggr)^{\!\frac{1}{2}}.
\end{align*}

Next, consider an anchored ball $B=B((0,y_0), r)$, and let $u$ be $L$-harmonic in $\lambda B$. Define the separated constant
\begin{equation*}
\widetilde{c} := u_{(\lambda B)^+}\mathbf 1_{\{x>0\}}(x) + u_{(\lambda B)^-} \mathbf 1_{\{x<0\}}(x).
\end{equation*}
Since $\alpha\in [\frac{1}{2},1)$, the function $\widetilde{c}(x)$ belongs to $\mathcal{F}_{\mathrm{loc}}$ and is locally $L$-harmonic, while
\begin{equation*}
    \nabla_L u=\nabla_L(u-\widetilde{c})
\end{equation*}
(see Remark~\ref{rem_separated_constant_Friedrichs} below). Therefore, the same argument as in the remote case gives
\begin{align*}
    \biggl( \fint_B |\nabla_L u|^p \, d\xi \biggr)^{\!\frac{1}{p}}
    &=\biggl( \fint_B |\nabla_L (u-\widetilde{c})|^p \, d\xi \biggr)^{\!\frac{1}{p}} \\
    &\le \frac{C}{r(B)}  \fint_{\lambda B} |u-\widetilde{c}| \, d\xi  \\
    &\le \frac{C}{r(B)} \biggl(\fint_{\lambda B} |u-\widetilde{c}|^2 \, d\xi \biggr)^{\!\frac{1}{2}}\\
    &= \frac{C}{r(B)} |\lambda B|^{-\frac{1}{2}} \biggl( \sum_{\pm} \int_{(\lambda B)^{\pm}} |u-u_{(\lambda B)^{\pm}}|^2 \, d\xi \biggr)^{\!\frac{1}{2}} \\
    &\le C \biggl(\fint_{\lambda' B} |\nabla_L u|^2 \, d\xi \biggr)^{\!\frac{1}{2}}.
\end{align*}

Finally, if $B$ is neither anchored nor remote, then the same reduction argument as in Lemma~\ref{reduce_RH} applies. We omit the details.
\end{proof}






\begin{remark}\label{rem_separated_constant_Friedrichs}
We justify more explicitly the use of the separated constant
\begin{equation*}
\widetilde c
:=
c_+\mathbf 1_{\{x>0\}}(x)
+
c_-\mathbf 1_{\{x<0\}}(x),
\qquad
c_\pm:=u_{(CB)^\pm}.
\end{equation*}
It may be written as
\begin{equation*}
\widetilde c
=
\frac{c_++c_-}{2}
+
\frac{c_+-c_-}{2}\operatorname{sgn}(x).
\end{equation*}
When $1/2\le \alpha<1$, the function $\operatorname{sgn}(x)$ belongs to the local
Friedrichs form domain. Moreover, $\nabla_L\widetilde c=0$ a.e.\ in the form sense and $\widetilde c$ is locally $L$-harmonic, and
\begin{equation*}
\nabla_L(u-\widetilde c)=\nabla_Lu.
\end{equation*}
Strictly speaking, $\widetilde c$ need not belong to $\mathcal D(L)$ on $\mathbb R^{1+m}$, because it need not be in $L^2$. What is used
in Proposition~\ref{prop_u_RH_to_Shen_RH_large_alpha} is precisely its membership in the local Friedrichs form domain and its local harmonicity. This is sufficient to apply \eqref{eq:RH:Jiang} to
$u-\widetilde c$ on the enlarged ball $\lambda B$.
\end{remark}

Immediately, Proposition~\ref{prop_u_RH_to_Shen_RH_large_alpha} implies the following.






\begin{theorem}\label{thm:Hodge:n1}
Let $n=1$, $m\ge1$, $\alpha \in [\frac{1}{2},1)$, and $\beta \ge0$. Then \eqref{Pi_p} holds for all $1<p<\infty$, i.e.,
\begin{equation*}
    \| \mathcal{P} \omega \|_p = \| \nabla_L L^{-1} \operatorname{div}_L \omega \|_p \le C \| \omega \|_p
\end{equation*}
for all $1<p<\infty$.
\end{theorem}

\begin{proof}
The argument is based on \cite[Section~2.2]{AC}. Note that $|\mathcal{P}|$ is a sublinear operator bounded on $L^2$. Let $2<p<\infty$. Let $B$ be a Grushin ball, and let $\omega \in C_c^\infty(\mathbb{R}^{1+m}; \mathbb{C}^{1+m})$ be a smooth vector field with
\begin{equation*}
    \operatorname{supp} \omega \subset (5B)^c.
\end{equation*}
Let $u$ be a weak solution to
\begin{equation*}
    Lu = \operatorname{div}_L \omega
\end{equation*}
such that $|\nabla_L u|\in L^2$. Clearly,
\begin{equation*}
    |\nabla_L u| = |\mathcal{P} \omega|.
\end{equation*}
By the support condition on $\omega$, the function $u$ is $L$-harmonic in $4B$. Therefore, Proposition~\ref{prop_u_RH_to_Shen_RH_large_alpha} implies that \eqref{eq:Shen} holds for all $q>p$, with $|Tf|$ replaced by $|\mathcal{P}\omega|$. Together with \eqref{D} and the $L^2$-boundedness of $\mathcal{P}$, Theorem~\ref{thm:Shen} applies immediately, and the proof follows.
\end{proof}









\subsection{From Hodge projector to Riesz transform}

We now complete the proof of Theorem~\ref{thm:n=1:alpha>1/2}.

\begin{proof}[Proof of Theorem~\ref{thm:n=1:alpha>1/2}]
Let $p>2$, so that
\begin{equation*}
    1<p'<2<\frac{1}{1-\alpha}.
\end{equation*}
For $\omega\in C_c^\infty(\mathbb R^{1+m};\mathbb C^{1+m})$, define
\begin{equation*}
    u=L^{-1}\operatorname{div}_L \omega.
\end{equation*}
We first verify that Theorem~\ref{prop_RR} applies to $u$. By Remark~\ref{re:applicable:HA} (see also Remarks~\ref{re:Hardy} and \ref{re:harmonic:annihilation}), $u$ must satisfy the following two conditions:
\begin{enumerate}[label=(\roman*)]
    \item $|\nabla_L u|\in L^{p'}$,
    \item $u(\xi) \to 0$ as $d(\xi,0)\to \infty$.
\end{enumerate}
Condition (i) follows from Theorem~\ref{thm:Hodge:n1}. Indeed, \eqref{Pi_p} implies
\begin{equation*}
    \|\nabla_L u\|_{p'} = \| \nabla_L L^{-1} \mathrm{div}_L \omega\|_{p'}\le C \|\omega\|_{p'} <\infty.
\end{equation*}
Condition (ii) follows from the decay of the Green kernel. Since $\omega$ is smooth and compactly supported, the decay of $u$ at infinity is a consequence of the decay of the kernel of $L^{-1}$. Recall that the homogeneous dimension is
\begin{equation*}
    \mathcal{Q}=\frac{1+m(\beta+1-\alpha)}{1-\alpha}.
\end{equation*}
Since $m\ge1$, $\beta\ge 0$, and $1/2\le\alpha<1$, we have
\begin{equation*}
    \mathcal{Q}\ge 1+\frac{1}{1-\alpha}\ge3.
\end{equation*}
Hence, for $\eta \in \operatorname{supp} \omega$, Theorem~\ref{RS_Grushin} yields
\begin{align*}
L^{-1}(\xi,\eta)
&:=\int_0^\infty e^{-tL}(\xi,\eta) \, dt\\
&\le C\int_0^{|x|^{2-2\alpha}} |x|^{-\alpha-m\beta} t^{-\frac{m+1}{2}} \exp\!\Bigl(-c\frac{d(\xi,\eta)^2}{t}\Bigr)\,dt \\
&\quad + C\int_{|x|^{2-2\alpha}}^\infty t^{-\mathcal{Q}/2} \exp\!\Bigl(-c\frac{d(\xi,\eta)^2}{t}\Bigr)\,dt \\
&\le C d(\xi,\eta)^{2-\mathcal{Q}}
    \to d(\xi,0)^{2-\mathcal{Q}}
    \to 0,
\end{align*}
as $d(\xi,0)\to \infty$. This proves condition (ii).

Therefore, $(\mathrm{RR}_{p'})$ applies to $u$, namely
\begin{equation*}
   \|L^{-1/2}\operatorname{div}_L\omega\|_{p'} = \| L^{1/2} u\|_{p'} \le C \| \nabla_L u\|_{p'}.
\end{equation*}
Since $\nabla_L u=\mathcal P\omega$, Theorem~\ref{thm:Hodge:n1} gives
\begin{equation*}
\|L^{-1/2}\operatorname{div}_L\omega\|_{p'} \le C\|\mathcal P\omega\|_{p'} \le
C\|\omega\|_{p'}.
\end{equation*}
By duality, this yields the boundedness of $\nabla_L L^{-1/2}$ on $L^p$.
\end{proof}



\begin{remark}
In fact, the implication: \eqref{RR_p} $+$ \eqref{Pi_p} $\implies$ $(\mathrm{R}_{p'})$ is more subtle than it looks. Indeed, a counterexample is found in \cite[Theorem~1.5]{HeSi} within the framework of manifolds with ends. Particularly, on the connected sum $\mathbb{R}^2 \# \mathbb{ R}^2$, \eqref{R_p} holds if and only if $1<p\le 2$, however, \eqref{RR_p} and \eqref{Pi_p} both remain valid for all $1<p<\infty$.
\end{remark}



We conclude with the following corollary.

\begin{corollary}
Let $n=1$, $m\geq1$, $\alpha\in(0,\frac12)$, and $\beta\geq0$. Then
\eqref{Pi_p} holds for $(1-\alpha)^{-1}<p<\alpha^{-1}$, and fails whenever $1<p<(1-\alpha)^{-1}$ or $\alpha^{-1}<p<\infty$.
\end{corollary}

\begin{proof}
The positive assertion follows directly from
Theorem~\ref{thm:n=1:alpha<1/2}. To prove the failure below the stated
range, suppose that \eqref{Pi_p} holds for some
$p<(1-\alpha)^{-1}$. Since $0<\alpha<1$, we have $\mathcal Q\geq 1+(1-\alpha)^{-1}>2$. The proof of Theorem~\ref{thm:n=1:alpha>1/2}, together with Theorem~\ref{thm_main2}, therefore implies that \eqref{RR_p} is valid for functions of the form $u:=L^{-1}\operatorname{div}_L\omega$, with $\omega\in C_c^\infty(\mathbb{R}^{1+m};\mathbb{C}^{1+m})$. The standard argument combining \eqref{Pi_p} and \eqref{RR_p} then yields $(\mathrm{R}_{p'})$. Since $p<(1-\alpha)^{-1}$ implies $p'>\alpha^{-1}$, this contradicts Theorem~\ref{thm:n=1:alpha<1/2}. The failure for $p>\alpha^{-1}$ follows by duality.
\end{proof}

\begin{remark}
Because \eqref{RR_p} is not known at $p=(1-\alpha)^{-1}$, the preceding
argument does not address either of the endpoint cases $p=(1-\alpha)^{-1}$ and $p=\alpha^{-1}$.
\end{remark}















% {\bf Acknowledgments.} 




\appendix


\section{Telescoping argument and covering lemma}\label{app:tele}



\begin{proof}[Proof of Lemma~\ref{le_tele}]
Set 
\[
g:=|\nabla_L f|\,\mathbf{1}_{2\gamma B}.
\]
We prove \eqref{eq:remote-lemma-312} for Lebesgue points $\xi,\eta$ of $f$.
Let
\[
\rho:=d(\xi,\eta).
\]
We divide the proof into two cases.

\textit{Case 1: $\rho\ge r/8$.}
We first show that for a.e.\ $\xi \in B$,
\begin{equation}\label{eq:point-vs-ball}
|f(\xi)-f_B|\le C r\, M_2(g)(\xi).
\end{equation}
Fix such a $\xi$, and define for $j\ge0$, $B_j:=B(\xi,2^{-j}r)$. Note that
\begin{equation*}
    B_j \subset B_0 = B(\xi,r) \subset B(\xi_0, 2r) = 2B.
\end{equation*}

\textit{Claim:} $B_j$ is remote for every $j\ge 0$.

It is enough to verify that $B_0$ is remote. Since $\xi = (x,y) \in B=B(\xi_0,r)$ and $B$ is remote, we have $|x-x_0|\le |x_0|/2$. Next, because $3B$ is remote, we get $3r \le \varepsilon_0 |x_0|^{1-\alpha}$, and hence $|x_0|\ge \bigl(3r/\varepsilon_0\bigr)^{\frac{1}{1-\alpha}}$. Also, $|x|\ge |x_0|- |x-x_0|\ge \frac{|x_0|}{2}$. Therefore, $|x|\ge 2^{-1} \bigl(3r/\varepsilon_0\bigr)^{\frac{1}{1-\alpha}}$, which implies $r \le \frac{2^{1-\alpha}}{3} \varepsilon_0 |x|^{1-\alpha} \le \varepsilon_0 |x|^{1-\alpha}$. Thus $B_0$ is remote, and hence so is every $B_j$.

Since each $B_j$ is remote, Lemma~\ref{PE} applies on every $B_j$. By Lebesgue differentiation,
\[
f(\xi)=\lim_{j\to\infty} f_{B_j}.
\]
Hence
\[
|f(\xi)-f_{B_0}|
\le \sum_{j=0}^\infty |f_{B_j}-f_{B_{j+1}}|.
\]
Moreover,
\[
|f_{B_j}-f_{B_{j+1}}|
\le \frac{1}{|B_{j+1}|}\int_{B_{j+1}} |f-f_{B_j}|
\le \frac{C}{|B_{j}|}\int_{B_j}|f-f_{B_j}|,
\]
where the last inequality follows from doubling. By Cauchy--Schwarz and Lemma~\ref{PE},
\[
\frac{1}{|B_{j}|}\int_{B_j}|f-f_{B_j}|
\le
\biggl(\frac{1}{|B_{j}|}\int_{B_j}|f-f_{B_j}|^2\biggr)^{\!1/2}
\le
C\, 2^{-j}r\biggl(\frac{1}{|B_{j}|}\int_{\gamma B_j} g^2\biggr)^{\!1/2}.
\]
Therefore
\[
|f(\xi)-f_{B_0}|
\le
C\sum_{j=0}^\infty 2^{-j}r\biggl(\frac{1}{|B_{j}|}\int_{\gamma B_j} g^2\biggr)^{\!1/2}
\le
C r\, M_2(g)(\xi).
\]
Since $B_0=B(\xi,r)\subset 2B$, we further decompose
\[
|f_{B_0}-f_B| \le |f_{B_0} - f_{2B}| + |f_B - f_{2B}|.
\]
For the first term, doubling gives
\begin{align*}
    |f_{B_0} - f_{2B}|
    &\le \frac{1}{|B_0|}\int_{B_0}|f-f_{2B}|\\
    &\le \frac{|B(\xi_0,2r)|}{|B(\xi,r)|} \frac{1}{|2B|}\int_{2B} |f-f_{2B}|\\
    &\le \frac{|B(\xi,3r)|}{|B(\xi,r)|} \biggl(\frac{1}{|2B|}\int_{2B} |f-f_{2B}|^2\biggr)^{\!1/2}\\
    &\le C r \biggl(\frac{1}{|2B|}\int_{2\gamma B} g^2\biggr)^{\!1/2}
    \le C r M_2(g)(\xi).
\end{align*}
The term $|f_B - f_{2B}|$ is estimated in the same way. Thus \eqref{eq:point-vs-ball} follows.

Applying \eqref{eq:point-vs-ball} to both $\xi$ and $\eta$, we obtain
\[
|f(\xi)-f(\eta)|
\le
|f(\xi)-f_B|+|f(\eta)-f_B|
\le
C r\bigl(M_2(g)(\xi)+M_2(g)(\eta)\bigr).
\]
Since in the present case $\rho\ge r/8$, this gives
\[
|f(\xi)-f(\eta)|
\le
C \rho \bigl(M_2(g)(\xi)+M_2(g)(\eta)\bigr).
\]

\textit{Case 2: $\rho< r/8$.}
For $j\ge0$, define
\[
B_j^\xi:=B(\xi,2^{-j+1}\rho),
\qquad
B_j^\eta:=B(\eta,2^{-j+1}\rho).
\]
Because $\rho<r/8$, it is clear that
\[
B(\xi,2\rho)\subset B(\xi_0,r+2\rho)\subset 2B,
\]
and similarly
\[
B(\eta,2\rho)\subset 2B.
\]
By the same argument as in the Claim above, every ball $B_j^\xi$ and $B_j^\eta$ is remote. Again by Lebesgue differentiation,
\[
f(\xi)=\lim_{j\to\infty} f_{B_j^\xi},
\qquad
f(\eta)=\lim_{j\to\infty} f_{B_j^\eta}.
\]
Hence
\[
|f(\xi)-f(\eta)|
\le
\sum_{j=0}^\infty |f_{B_j^\xi}-f_{B_{j+1}^\xi}|
+
|f_{B(\xi,2\rho)}-f_{B(\eta,\rho)}|
+
\sum_{j=1}^\infty |f_{B_j^\eta}-f_{B_{j+1}^\eta}|.
\]
For the telescoping terms, exactly as above,
\[
|f_{B_j^\xi}-f_{B_{j+1}^\xi}|
\le
C\, 2^{-j}\rho \biggl(\frac{1}{|B_j^\xi|}\int_{\gamma B_j^\xi} g^2\biggr)^{\!1/2}
\le
C\, 2^{-j}\rho\, M_2(g)(\xi),
\]
and similarly
\[
|f_{B_j^\eta}-f_{B_{j+1}^\eta}|
\le
C\, 2^{-j}\rho\, M_2(g)(\eta).
\]
For the middle term, since $d(\xi,\eta)=\rho$, we have
\[
B(\eta,\rho)\subset B(\xi,2\rho) \subset B(\eta,3\rho).
\]
Therefore
\[
|f_{B(\xi,2\rho)}-f_{B(\eta,\rho)}|
\le
\frac{1}{|B(\eta,\rho)|}\int_{B(\eta,\rho)} |f-f_{B(\xi,2\rho)}|
\le
\frac{C}{|B(\xi,2\rho)|} \int_{B(\xi,2\rho)} |f-f_{B(\xi,2\rho)}|.
\]
Using Cauchy--Schwarz and Lemma~\ref{PE} on the remote ball $B(\xi,2\rho)$, we get
\[
|f_{B(\xi,2\rho)}-f_{B(\eta,\rho)}|
\le
C\, \rho\biggl( \frac{1}{|B(\xi,2\rho)|}\int_{\gamma B(\xi,2\rho)} g^2\biggr)^{\!1/2}
\le
C\, \rho\, M_2(g)(\xi).
\]
Combining the above estimates and summing the geometric series, we conclude that
\[
|f(\xi)-f(\eta)|
\le
C\, \rho\bigl(M_2(g)(\xi)+M_2(g)(\eta)\bigr),
\]
that is,
\[
|f(\xi)-f(\eta)|
\le
C\, d(\xi,\eta)\bigl(M_2(|\nabla_L f|\mathbf{1}_{2\gamma B})(\xi)+M_2(|\nabla_L f|\mathbf{1}_{2\gamma B})(\eta)\bigr).
\]
This proves \eqref{eq:remote-lemma-312}.
\end{proof}





\begin{proof}[Proof of Lemma~\ref{lemma_remotecover}]
Set $B_n^0 = B_n(0,\varepsilon)$ and define
\[
A_N:= B_n(0, \varepsilon 2^{N}) \setminus B_n(0, \varepsilon 2^{N-1}), \qquad N\ge 1.
\]
Clearly,
\begin{align*}
    \mathbb{R}^n = B_n^0 \cup \Bigl(\bigcup_{N\ge 1}A_N\Bigr),\qquad
    A_N \subset \bigcup_{x\in A_N} B_n \bigl(x, \varepsilon 2^{N-\delta-13}\bigr).
\end{align*}
By Vitali's covering lemma, one can find an index set $I_N$ and a collection of pairwise disjoint balls
\[
\{B_n(x_{N,i},\varepsilon 2^{N-\delta-13})\}_{i\in I_N},
\qquad x_{N,i}\in A_N,
\]
such that
\[
A_N \subset \bigcup_{i\in I_N}B_n\bigl(x_{N,i}, \varepsilon 2^{N-\delta-10}\bigr).
\]
By \eqref{D}, the cardinality $\# I_N$ is finite. More precisely,
\begin{align*}
    (\# I_N) |B_n(0,\varepsilon 2^{N+2})|
    &\le \sum_{i \in I_N} |B_n(x_{N,i}, \varepsilon 2^{N+3})| \\
    &\le C(n,\delta) \sum_{i \in I_N} |B_n(x_{N,i}, \varepsilon 2^{N-\delta-13})|\\
    &\le C(n,\delta) |B_n(0,\varepsilon 2^{N+2})|,
\end{align*}
where the first inequality follows from the inclusion
\[
B_n(0,4\varepsilon 2^{N}) \subset B_n( x_{N,i}, 8 \varepsilon 2^N ),
\qquad i\in I_N,\ N\ge 1.
\]
The second inequality follows directly from \eqref{D}, and the last inequality uses the disjointness of the balls $B_n(x_{N,i},\varepsilon 2^{N-\delta-13})$ together with
\begin{equation*}
    \bigcup_{i\in I_N} B_n\bigl(x_{N,i},\varepsilon 2^{N-\delta-13}\bigr)
    \subset B_n\bigl(0, \varepsilon 2^N + \varepsilon 2^{N-\delta-13}\bigr)
    \subset B_n( 0, \varepsilon 2^{N+2} ).
\end{equation*}
In particular, the bound for $\# I_N$ is independent of both $N$ and $\varepsilon$.

Now set $B_n^j = B_n(x_{j}, r_j)$, where $x_j = x_{N,i}$ and $r_j = \varepsilon 2^{N-\delta-10}$, after relabeling. This yields a sequence of balls satisfying
\[
\bigcup_{j \ge 0}B_n^j = \mathbb{R}^n.
\]
Moreover, for $j \ne 0$, it is easy to check that
\begin{align}\label{eq_|xj|,r_j}
2^{-10-\delta} |x_j| \le r_j \le 2^{-9-\delta} |x_j|,
\end{align}
and for all $x\in B_n^j$,
\begin{align}\label{eq_|x|,r_j}
    2^{8+\delta} r_j \le |x| \le 2^{11+\delta}r_j.
\end{align}
This proves \eqref{(1)} and \eqref{(2)}.

It remains to prove \eqref{(3)}. By construction, $B_n^0 = B_n(0,\varepsilon)$ is the only ball in the collection containing $0$. Now let $x\in \mathbb{R}^n \setminus \{0\}$, and consider the index set
\[
J_x = \{j: x\in B_n^j\}.
\]
If $x\in A_N$ for some $N\ge 1$ and $x\in B_n^j$, then $j$ must belong to $I_N$, $I_{N-1}$, or $I_{N+1}$ if $N\ge 2$, and to $I_N$ or $I_{N+1}$ if $N=1$. Therefore, by \eqref{eq_|x|,r_j}, the doubling condition, and \eqref{eq_|xj|,r_j},
\begin{align*}
    (\# J_x) \bigl|B_n (x,\tfrac{|x|}{2})\bigr|
    &\le \sum_{j \in J_x} \bigl|B_n( x_j, \tfrac{|x|}{2} + \tfrac{|x|}{2^{8+\delta}} )\bigr|\\
    &\le C(n,\delta) \sum_{j\in J_x} \bigl| B_n(x_j, \tfrac{|x|}{2^{14+\delta}} )\bigr|\\
    &\le C(n,\delta) \sum_{j \in J_x} \bigl| B_n(x_j, \tfrac{r_j}{2^{3}} )\bigr|\\
    &= C(n,\delta) \sum_{l\in \{-1,0,1\}} \sum_{j \in J_x \cap I_{N+l}} \bigl| B_n(x_j, \tfrac{r_j}{2^{3}} )\bigr| \\
    &= C(n,\delta) \sum_{l\in \{-1,0,1\}} \Bigl|\bigcup_{j \in J_x\cap I_{N+l}} \tfrac{B_n^j}{8} \Bigr|\\
    &\le C(n,\delta) \bigl|B_n (x,\tfrac{|x|}{2})\bigr|,
\end{align*}
where the second-to-last equality follows from the disjointness of the balls
\[
\frac{B_n( x_{N,i}, \varepsilon 2^{N-\delta-10} )}{8},
\]
and the last inequality uses the fact that, for any $j\in J_x \cap I_{N+l}$,
\begin{align*}
    \frac{B_n^j}{8} = B_n\bigl(x_j, \tfrac{r_j}{8} \bigr)
    \subset B_n \bigl(x,\tfrac{9r_j}{8} \bigr)
    \subset B_n\bigl(x, \tfrac{9 |x|}{2^{11+\delta}}\bigr)
    \subset B_n \bigl(x,\tfrac{|x|}{2}\bigr).
\end{align*}
The case $x\in B_n^0$ is similar. Indeed, in that case, $x\in B_n^j$ implies that $j=0$ or $j\in I_1$, and the same argument applies. This proves \eqref{(3)}, and hence Lemma~\ref{lemma_remotecover}.
\end{proof}






\section{Bessel analysis package}\label{app:B}

% Throughout, $I_\nu$ and $K_\nu$ denote the standard modified Bessel functions. We use several classical facts about these functions, including their power-series expansions, recurrence formulae, Wronskian identity, integral representation, and small-argument asymptotics; see Watson~\cite[Chapter~III]{Watson} and Abramowitz--Stegun~\cite[Chapter~9]{AS}. For $\nu>-1$ we also use the normalized regular branch
% \begin{equation*}
%     \mathcal I_\nu(t):=t^{-\nu}I_\nu(t),
%     \qquad t>0.
% \end{equation*}
% For real $\nu$ we write
% \begin{equation*}
%     \mathcal K_\nu(t):=t^\nu K_\nu(t),
%     \qquad t>0.
% \end{equation*}
% Finally, set
% \begin{equation*}
%     S_\nu(t):=\frac{I_{\nu+1}(t)}{I_\nu(t)}.
% \end{equation*}







% \begin{lemma}\label{lem:Bessel_regular_branch}
% Let $\nu_*>-1$. For every integer $N\geq0$, there exists a constant
% $C_{N,\nu_*}>0$ such that, for all $\nu\ge\nu_*$,
% \begin{equation*}
%     \sup_{0\leq t\leq1}
%     \left|
%         \frac{d^k}{dt^k}\mathcal I_\nu(t)
%     \right|
%     \leq
%     C_{N,\nu_*},
%     \qquad
%     0\leq k\leq N.
% \end{equation*}
% Moreover,
% \begin{equation*}
%     \frac{d}{dt}\mathcal I_\nu(t)=t\mathcal I_{\nu+1}(t).
% \end{equation*}

% \end{lemma}




\begin{proof}[Proof of Lemma~\ref{lem:Bessel_regular_branch}]
The power-series expansion gives
\begin{equation*}
I_\nu(t)=\frac{1}{\Gamma(\nu+1)}\left(\frac{t}{2}\right)^\nu\sum_{q=0}^\infty \frac{1}{q!(\nu+1)_q} \left(\frac{t^2}{4}\right)^q.
\end{equation*}
Hence
\begin{equation*}
\mathcal I_\nu(t) = \frac{2^{-\nu}}{\Gamma(\nu+1)} \sum_{q=0}^\infty \frac{1}{q!(\nu+1)_q} \left(\frac{t^2}{4}\right)^q.
\end{equation*}
Since $\nu\geq\nu_*>-1$, the coefficients and their finitely many $t$-derivatives are uniformly bounded on $0\leq t\leq1$, with constants depending only on $N$ and $\nu_*$. This proves the first assertion. The derivative identity follows from the standard recurrence
\begin{equation*}
\frac{d}{dt}\bigl(t^{-\nu}I_\nu(t)\bigr)=t^{-\nu}I_{\nu+1}(t) = t\,\mathcal I_{\nu+1}(t).
\end{equation*}

\end{proof}





% \begin{lemma}\label{lem_S_derivative_bound}
% Let $\nu_*>-1$. For every integer $N\ge 0$, there exist constants
% $C_{N,\nu_*},A_N>0$ such that, for all $\nu\ge\nu_*$,
% \begin{equation*}
%     \left|S_\nu^{(k)}(t)\right|
%     \le
%     C_{N,\nu_*}(1+\nu+t)^{A_N},
%     \qquad
%     0\le k\leq N,\quad t>0.
% \end{equation*}
% Consequently, for every multi-index $\gamma$ with $|\gamma|\leq N$,
% \begin{equation*}
% \left| \partial_Y^\gamma\bigl(|Y|S_\nu(|Y|)\bigr) \right| \le C_{N,\nu_*}(1+\nu+|Y|)^{A_N}, \qquad Y\in\mathbb R^m.
% \end{equation*}

% \end{lemma}





\begin{proof}[Proof of Lemma~\ref{lem_S_derivative_bound}]
We first treat $0<t\le1$. By the power-series expansion,
\begin{equation*}
    S_\nu(t)
    =
    \frac{t}{2(\nu+1)}
    \frac{E_{\nu+1}(t^2)}{E_\nu(t^2)},\qquad E_\nu(s):=
    \sum_{q=0}^\infty
    \frac{1}{q!(\nu+1)_q}
    \left(\frac{s}{4}\right)^q.
\end{equation*}
For $\nu\ge \nu_*>-1$, we have $E_\nu(s)\ge 1$, and, for every $N\ge 0$,
\begin{equation}\label{eq_A_derivatives_small}
    \sup_{\nu\geq\nu_*}\sup_{0\le s\leq1}
    \left|E_\nu^{(k)}(s)\right|
    \le
    C_{N,\nu_*},
    \qquad
    0\leq k\le N.
\end{equation}
Since $\nu+1\ge \nu_*+1>0$, differentiating the formula for $S_\nu$ gives
\begin{equation*}
    |S_\nu^{(k)}(t)|
    \le
    C_{N,\nu_*},
    \qquad
    0<t\leq1,\quad 0\leq k\le N.
\end{equation*}
We now consider $t\ge 1$. The recurrence identities
\begin{equation*}
    I_\nu'(t)=I_{\nu+1}(t)+\frac{\nu}{t}I_\nu(t),
    \qquad
    I_{\nu+1}'(t)=I_\nu(t)-\frac{\nu+1}{t}I_{\nu+1}(t)
\end{equation*}
give 
\begin{equation*}
    S_\nu'(t)
    =
    1-S_\nu(t)^2-\frac{2\nu+1}{t}S_\nu(t).
\end{equation*}
Since $S_\nu(t)\le C_{\nu_*}$ for all $t>0$ and $\nu\ge \nu_* >-1$, we have, for $t\geq1$,
\begin{equation*}
    |S_\nu'(t)|
    \leq
    C_{\nu_*}(1+\nu).
\end{equation*}
Repeatedly differentiating the equation expresses $S_\nu^{(j+1)}$ as a finite sum of products involving $S_\nu^{(k)}$, $k\leq j$, powers of $t^{-1}$, and the coefficient $2\nu+1$. An induction on $j$ gives
\begin{equation*}
    |S_\nu^{(k)}(t)|
    \leq
    C_{N,\nu_*}(1+\nu)^{A_N}
    \leq
    C_{N,\nu_*}(1+\nu+t)^{A_N},
    \qquad
    t\geq1,\quad 0\leq k\leq N.
\end{equation*}
This proves the first assertion.

For the estimate in the $Y$-variable, set $F_\nu(r):=rS_\nu(r)$. If $|Y|\geq1$, the standard formula for derivatives of radial functions gives
\begin{equation*}
    \left|\partial_Y^\gamma F_\nu(|Y|)\right|
    \leq
    C_\gamma
    \sum_{j=0}^{|\gamma|}
    |Y|^{j-|\gamma|}
    \left|F_\nu^{(j)}(|Y|)\right|
    \leq
    C_{N,\nu_*}(1+\nu+|Y|)^{A_N}.
\end{equation*}
If $|Y|\leq1$, then
\begin{equation*}
    |Y|S_\nu(|Y|)
    =
    \frac{|Y|^2}{2(\nu+1)}
    \frac{E_{\nu+1}(|Y|^2)}{E_\nu(|Y|^2)}
    =
    g_\nu(|Y|^2), \qquad g_\nu(s)
    :=
    \frac{s}{2(\nu+1)}
    \frac{E_{\nu+1}(s)}{E_\nu(s)}.
\end{equation*}
The preceding bounds for $E_\nu$ imply
\begin{equation*}
    |g_\nu^{(k)}(s)|
    \leq
    C_{N,\nu_*},
    \qquad
    0\leq s\leq1,\quad 0\leq k\leq N.
\end{equation*}
Applying the chain rule to $g_\nu(|Y|^2)$ yields
\begin{equation*}
    \left|
        \partial_Y^\gamma\bigl(|Y|S_\nu(|Y|)\bigr)
    \right|
    \leq
    C_{N,\nu_*},
    \qquad
    |Y|\leq1,\quad |\gamma|\leq N.
\end{equation*}
Combining the two cases proves the lemma.
\end{proof}







% \begin{lemma}\label{lem_bessel_inputs_high_frequency}
% Let $\nu_*>-1$. There exist constants $c,C>0$, depending only on $\nu_*$, such
% that, for all $\nu\geq\nu_*$,
% \begin{equation}\label{eq_S_lower_large}
%     \frac{I_{\nu+1}(t)}{I_\nu(t)}
%     \geq
%     c,
%     \qquad
%     t\geq C(1+\nu).
% \end{equation}
% Moreover, $\mathcal K_\nu(t)=t^\nu K_\nu(t)$ is decreasing, and
% \begin{equation}\label{eq_K_normalized_exp_decay}
%     \frac{\mathcal K_\nu(t_2)}{\mathcal K_\nu(t_1)}
%     \leq
%     C\exp\{-c(t_2-t_1)\},
%     \qquad
%     t_2\geq t_1\geq C(1+\nu).
% \end{equation}
% In addition, if $t>0$ and $\nu\ge 0$. Then
% \begin{equation}\label{eq:product:Bessel}
%     I_\nu(t)K_\nu(t)
%     \le
%     \frac{C}{t+\nu}.
% \end{equation}

% \end{lemma}

\begin{proof}[Proof of Lemma~\ref{lem_bessel_inputs_high_frequency}]
The first estimate follows from the standard Amos-type lower bound
\[
    \frac{I_{\nu+1}(t)}{I_\nu(t)}
    \geq
    \frac{t}
    {\nu+\frac12+\sqrt{(\nu+\frac32)^2+t^2}},
    \qquad
    \nu>-1,\quad t>0.
\]
Since $\nu\geq\nu_*>-1$, choosing $C$ sufficiently large gives
\[
    \frac{I_{\nu+1}(t)}{I_\nu(t)}
    \geq c
\]
whenever $t\geq C(1+\nu)$.

For the $K$-branch, the identity
\[
    \frac{d}{dt}\bigl(t^\nu K_\nu(t)\bigr)
    =
    -t^\nu K_{\nu-1}(t)
\]
shows that $\mathcal K_\nu$ is decreasing. Moreover,
\[
    \frac{d}{dt}\log\mathcal K_\nu(t)
    =
    -\frac{K_{\nu-1}(t)}{K_\nu(t)}.
\]
Using
\[
    -\frac{K_\nu'(t)}{K_\nu(t)}
    =
    \frac{\nu}{t}+\frac{K_{\nu-1}(t)}{K_\nu(t)}
\]
and the integral representation
\[
    K_\nu(t)=\int_0^\infty e^{-t\cosh s}\cosh(\nu s)\,ds,
\]
we get
\[
    -\frac{K_\nu'(t)}{K_\nu(t)}
    =
    \frac{\int_0^\infty \cosh s\,e^{-t\cosh s}\cosh(\nu s)\,ds}
    {\int_0^\infty e^{-t\cosh s}\cosh(\nu s)\,ds}
    \geq1.
\]
Hence
\[
    \frac{K_{\nu-1}(t)}{K_\nu(t)}
    \geq
    1-\frac{\nu}{t}.
\]
If $\nu\leq0$, the right-hand side is at least $1$. If $\nu>0$ and
$t\geq2\nu$, it is at least $1/2$. Therefore, after enlarging $C$ if necessary,
\[
    \frac{K_{\nu-1}(t)}{K_\nu(t)}\geq c,
    \qquad
    t\geq C(1+\nu).
\]
Integrating the logarithmic derivative of $\mathcal K_\nu$ over $[t_1,t_2]$
gives
\[
    \log\frac{\mathcal K_\nu(t_2)}{\mathcal K_\nu(t_1)}
    =
    -\int_{t_1}^{t_2}
    \frac{K_{\nu-1}(s)}{K_\nu(s)}\,ds
    \leq
    -c(t_2-t_1).
\]
This proves the exponential estimate.

Finally, Put
\[
    P_\nu(t):=I_\nu(t)K_\nu(t),
    \qquad
    r_\nu(t):=\frac{I_\nu'(t)}{I_\nu(t)},
    \qquad
    q_\nu(t):=-\frac{K_\nu'(t)}{K_\nu(t)}.
\]
The Wronskian identity
\[
    I_\nu'(t)K_\nu(t)-I_\nu(t)K_\nu'(t)=\frac1t
\]
gives
\[
    P_\nu(t)\bigl(r_\nu(t)+q_\nu(t)\bigr)=\frac1t.
\]
By the recurrence formulae
\[
    I_\nu'(t)=\frac{\nu}{t}I_\nu(t)+I_{\nu+1}(t),
    \qquad
    -K_\nu'(t)=\frac{\nu}{t}K_\nu(t)+K_{\nu-1}(t),
\]
and positivity of $I_{\nu+1}$ and $K_{\nu-1}$, we have, for $\nu>0$, $r_\nu(t)+q_\nu(t)\geq\frac{2\nu}{t}$. Thus
\[
    I_\nu(t)K_\nu(t)\leq\frac1{2\nu},
    \qquad
    \nu>0.
\]
On the other hand, the integral representation of $K_\nu$ gives
\[
    q_\nu(t)=-\frac{K_\nu'(t)}{K_\nu(t)}\geq1.
\]
Therefore $r_\nu(t)+q_\nu(t)\geq1$ and hence
\[
    I_\nu(t)K_\nu(t)\leq\frac1t.
\]
Combining the two estimates gives
\[
    I_\nu(t)K_\nu(t)
    \leq
    \min\left\{\frac1t,\frac1{2\nu}\right\}
    \leq
    \frac{C}{t+\nu}
\]
for $\nu>0$. The case $\nu=0$ follows directly from the bound
$I_0(t)K_0(t)\leq t^{-1}$.
\end{proof}





\begin{remark}
The restriction $\nu\geq0$ in the third assertion of  Lemma~\ref{lem_bessel_inputs_high_frequency}
is essential. If $-1<\nu<0$, then
\[
    I_\nu(t)K_\nu(t)\sim c_\nu t^{2\nu},
    \qquad t\to 0,
\]
and the product is unbounded at the origin. In the paper this product estimate is used only for the higher-dimensional angular orders, where $\nu=\nu_\ell\geq0$. 
\end{remark}





\section{Fa\`a di Bruno's Formula}


We record only the form of Faa di Bruno's formula used in the proof. Let
$I\subset\mathbb R$ be an interval, let $g\in C^N(I)$, and let
$F\in C^N(g(I))$. Then, for every integer $N\geq1$,
\[
    \frac{d^N}{dt^N}F(g(t))
    =
    \sum_{k=1}^N
    F^{(k)}(g(t))
    B_{N,k}\bigl(g'(t),g''(t),\ldots,g^{(N-k+1)}(t)\bigr),
\]
where
\[
    B_{N,k}(X_1,\ldots,X_{N-k+1})
    =
    \sum
    \frac{N!}{j_1!\cdots j_{N-k+1}!}
    \prod_{r=1}^{N-k+1}
    \left(\frac{X_r}{r!}\right)^{j_r}.
\]
The sum is taken over all non-negative integers
$j_1,\ldots,j_{N-k+1}$ satisfying
\[
    j_1+\cdots+j_{N-k+1}=k,
    \qquad
    j_1+2j_2+\cdots+(N-k+1)j_{N-k+1}=N.
\]
In particular, if $f>0$ and $f\in C^N(I)$, then applying the formula to
$f=\exp(\log f)$ gives
\[
    f^{(N)}(t)
    =
    f(t)
    \sum
    C_{j_1,\ldots,j_N}
    \prod_{r=1}^N
    \left((\log f)^{(r)}(t)\right)^{j_r},
\]
where the sum is over all non-negative integers $j_1,\ldots,j_N$ such that
\[
    j_1+2j_2+\cdots+Nj_N=N,
\]
and the constants $C_{j_1,\ldots,j_N}$ depend only on $N$.

Consequently, if for some $A\geq1$,
\[
    \left|(\log f)^{(r)}(t)\right|
    \leq
    C_r A^r,
    \qquad
    1\leq r\leq N,
\]
then
\[
    |f^{(N)}(t)|
    \leq
    C_N A^N f(t).
\]
This is the form used when differentiating positive multipliers written as
$f=\exp(\log f)$.





\bibliographystyle{abbrv}
\bibliography{references}



\end{document}
